\documentclass[11pt, a4paper]{article}

\usepackage[utf8]{inputenc}
\usepackage[T1]{fontenc}
\usepackage{lmodern}
\usepackage[margin=2.5cm]{geometry}
\usepackage{amsmath, amssymb, amsthm}
\usepackage{booktabs}
\usepackage{enumitem}
\usepackage{hyperref}
\usepackage{xcolor}
\usepackage{microtype}
\usepackage{tikz}
\usetikzlibrary{positioning,calc,fit,arrows.meta,backgrounds}

\hypersetup{
    colorlinks=true,
    linkcolor=black,
    citecolor=blue!60!black,
    urlcolor=blue!60!black
}

\newcommand{\M}{\mathbf{M}}
\newcommand{\x}{\mathbf{x}}
\newcommand{\h}{\mathbf{h}}
\newcommand{\R}{\mathbb{R}}

\theoremstyle{plain}
\newtheorem{conjecture}{Conjecture}
\newtheorem{proposition}{Proposition}
\newtheorem{theorem}{Theorem}
\theoremstyle{definition}
\newtheorem{definition}{Definition}
\theoremstyle{remark}
\newtheorem{remark}{Remark}

\title{%
    \textbf{Two Hemispheres:}\\[4pt]
    \Large Co-Evolving Literal and Relational Processing\\[2pt]
    \large in Neural Language Models%
}
\author{%
    Sirsh Amarteifio\\
    {\small Percolation Labs}\\
    {\small \url{https://github.com/Percolation-Labs/orn}}%
}
\date{April 2026 \quad{\small\textnormal{,  Research Programme}}}

\begin{document}
\maketitle

% ======================================================================
\begin{abstract}
Current language models face three unsolved problems: compressing long contexts into bounded memory, performing compositional reasoning without externalising it as chain-of-thought tokens, and adapting representations at test time without gradient-based fine-tuning.
These problems persist because every model uses a single processing mode, predict the next token, optimised by a single loss.
There is no separation between learning \emph{what things are} and learning \emph{how things relate}.

We observe that these engineering challenges have a biological parallel.
McGilchrist argues that the brain requires two hemispheres because literal processing (specific, sequential, categorical) and relational processing (abstract, compositional, holistic) are fundamentally different tasks.
Neither alone is sufficient.
We do not propose to imitate biology.
We propose to solve the engineering problems by the same structural means biology discovered: \textbf{two co-evolving networks with different objectives}.

\begin{itemize}[nosep]
\item A \textbf{left network} (transformer attention) that performs next-token prediction, the literal processor, the detail handler, the standard language model.
\item A \textbf{right network} (SSM on the left's coupling dynamics) that compresses relational context into persistent state, the structural discoverer, the long-range reasoner.
\end{itemize}

We propose that the right network could address the three problems directly:
\begin{itemize}[nosep]
\item \textbf{Long context:} the right hemisphere's state is fixed-size ($r$ dimensions) regardless of sequence length, compressing arbitrarily long relational history.
\item \textbf{Compositional reasoning:} the right hemisphere accumulates relational structure through operator dynamics, composition happens in state space, not in output tokens.
\item \textbf{Test-time adaptation:} the right hemisphere's state evolves identically at training and inference, context-dependent adaptation without gradients.
\end{itemize}
These claims remain to be validated at scale; the preliminary experiments below test the mechanisms on a small model.

The right network does not see tokens.
It sees the left network's internal coupling patterns and compresses them into relational state.
Its loss forces \emph{coherence}, not through labelled metaphors or contrastive pairs, but through the pressure to find states that make the left's coupling patterns fit together as one relational structure.
The interaction between the two networks is what produces representations that are simultaneously specific enough to predict and general enough to transfer.

Preliminary results: with additive injection, the learned gate opens to 81\% during co-evolution and val loss improves by 0.55 nats (though parameter-matched baselines achieve similar gains).
With coupling bias injection, where R modulates L's attention patterns rather than its content stream, the gate stays open at 86\% and val loss improves by 0.33 nats over baseline, with R's states clustering by syntactic structure (4.9$\times$ separation ratio between structural groups).
The interaction is productive, and it emerges from training, not from architectural prescription.
\end{abstract}


% ======================================================================
\paragraph{\textsc{Document Status} (added 21 April 2026).}
\textbf{This paper was written in April 2026 and presents the vision and earliest experiments for LORN.}
Since then the programme has adopted stricter evaluation criteria, documented in the \emph{Phase 1 checklist} of the canonical paper \texttt{lorn\_gauge v3}.
Specifically, every result in this document that is reported as a ``val loss improvement'' (see in particular \S\ref{sec:first-results} and \S\ref{sec:injection-modes}) was evaluated against a criterion the programme has since \emph{rejected}: R's S$_n$-invariant structural content is orthogonal to next-token prediction by construction, so val-loss improvement does not test R's structural claim.
Expecting R to reduce perplexity was a category error.
The SSM-on-coupling-dynamics architectural idea described in \S3 remains a live candidate for R and is carried forward into the current programme.
The biological/McGilchrist framing also remains the intended motivation.
Readers should treat this paper as the \emph{vision document} and defer to \texttt{lorn\_gauge v3} for the current checklist, evaluation methodology, and state of R's instantiation.

% ======================================================================
\section{Why Two Hemispheres}

\subsection{The Biological Argument}

McGilchrist's \emph{The Master and His Emissary}~\cite{mcgilchrist2009master} and \emph{The Matter with Things}~\cite{mcgilchrist2021matter} argue that hemispheric lateralisation is not an accident of brain architecture but a fundamental requirement of intelligence.

The left hemisphere excels at:
\begin{itemize}[nosep]
\item Focused attention on detail
\item Categorical, explicit representations
\item Sequential, step-by-step processing
\item Literal language (denotation)
\item Manipulation of known things
\end{itemize}

The right hemisphere excels at:
\begin{itemize}[nosep]
\item Broad, vigilant attention to context
\item Relational, implicit representations
\item Holistic, simultaneous processing
\item Metaphorical language (connotation, resonance)
\item Understanding of novel situations
\end{itemize}

Neither is sufficient alone.
A left-only brain produces rigid, literal, context-blind processing, it knows facts but cannot understand.
A right-only brain produces fluid, associative, unstable processing, it grasps relationships but cannot pin down specifics.
Intelligence requires both: the left grounds the right, the right contextualises the left.

\subsection{The Computational Argument}

Current neural language models are \emph{left hemisphere machines}.
They process tokens sequentially, predict the next literal symbol, and learn through a loss that rewards specific accuracy.
They are remarkably good at detail, they can recall facts, generate fluent text, solve specific problems.
But they struggle with: novel analogies, transfer to unfamiliar domains, compositional generalisation, and understanding \emph{why} rather than \emph{what}.

The ``operator programme'' (our companion paper, \emph{Taking Orbits to the Limit}) tried to build right-hemisphere capability by designing operators that compose relationally.
Every attempt failed, because we tried to build the right hemisphere \emph{inside} the left hemisphere's architecture.
We used the same loss (cross-entropy), the same mechanism (attention + FFN), the same evaluation (val loss).
The right hemisphere cannot be a component of the left.
It must be a separate system that \emph{interacts} with the left.


% ======================================================================
\section{The Architecture}

\subsection{Two Networks, One Interaction}

\begin{definition}[Left Network (L)]
A standard autoregressive language model.
Input: tokens.
Output: next-token probabilities.
Loss: cross-entropy.
Architecture: transformer, ORN, or any sequence model.
This is the literal processor, it learns to predict.
\end{definition}

\begin{definition}[Right Network (R)]
An auxiliary network that operates on L's internal representations.
Input: L's hidden states, attention patterns, or learned embeddings, not raw tokens.
Output: a structural constraint or reorganisation signal applied back to L.
Loss: structural compression or relational consistency (details below).
This is the relational processor, it learns to organise.
\end{definition}

The two networks co-evolve:
\begin{enumerate}[nosep]
\item L processes a batch of tokens, producing hidden states $\h_1, \ldots, \h_L$ at each layer.
\item R receives L's representations and computes a structural loss: how well can L's representations be compressed relationally?
\item R is trained by its structural loss $\mathcal{L}_{\mathrm{structural}}$ (with $\mathrm{stopgrad}$ on L's representations to prevent L from collapsing to satisfy R).
\item L is trained only by $\mathcal{L}_{\mathrm{CE}}$. The structural pressure on L is \emph{architectural}: R's output modifies L's attention or residual stream, and L's CE loss rewards useful modifications.
\item Over time, L reorganises its representations because R's feedback improves prediction, not because structural gradients flow to L.
\end{enumerate}

R does not predict tokens.
It does not see tokens.
It only sees how L represents things, and pushes L toward more structured representations.

\subsection{What R Learns}

The right network discovers relational structure unsupervised.
Several possible instantiations:

\paragraph{R as a representation compressor.}
R tries to compress L's hidden states into a low-dimensional relational code.
If L's representations have consistent structure (similar sentences represented similarly), R can compress them efficiently.
If L's representations are unstructured, R cannot compress them.
The loss: reconstruction error of R's compression.
\begin{equation}
\mathcal{L}_{\mathrm{compress}} = \| \h - \mathrm{decode}(\mathrm{encode}(\h)) \|^2
\end{equation}
where $\mathrm{encode}: \R^d \to \R^r$ ($r \ll d$) finds the relational core.

\paragraph{R as a relational predictor.}
R predicts one representation from another using a learned relational mapping.
If the relationship between ``committee'' and ``vindicate'' is the same as between ``defendant'' and ``acquit,'' R should be able to map one pair to the other.
The loss: prediction error of relational mapping.
\begin{equation}
\mathcal{L}_{\mathrm{relational}} = \| R(\h_A, \h_B) - R(\h_C, \h_D) \|^2 \quad \text{for analogous pairs $(A{:}B :: C{:}D)$}
\end{equation}
But this requires knowing the analogies, which is supervision.

\paragraph{R as a structural consistency enforcer (unsupervised).}
R measures whether L's representations at different layers, positions, or contexts share structure.
If the attention pattern for ``because'' is similar across all contexts (which our attention visualisation showed: consistency 0.68), R rewards this consistency.
If a representation varies wildly across contexts, R does not penalise it (it may be content-specific).
The loss: variance of representations that R identifies as structurally similar.

\paragraph{R as an operator learner.}
R learns a mapping $T$ such that $T(\h_t) \approx \h_{t+k}$ for various $k$, it predicts future representations from current ones.
If L's representations evolve through a consistent operator (the coupling geometry), R can learn this operator efficiently.
The loss: prediction error of the operator.
\begin{equation}
\mathcal{L}_{\mathrm{operator}} = \sum_k \| T^k(\h_t) - \h_{t+k} \|^2
\end{equation}
This IS the operator programme, but learned by a separate network from L's representations, not designed into L's architecture.

\subsection{The Adversarial / Cooperative Dynamic}

The interaction between L and R can be framed as:

\paragraph{Cooperative.}
R's loss is added to L's loss.
L is trained to satisfy both prediction (CE) and structure (R's constraint).
R is trained to discover the structure that L's representations contain.
Both benefit from better representations.

\paragraph{Adversarial.}
R tries to find structure in L's representations.
L tries to maintain prediction quality while being forced to reorganise.
The ``adversarial'' pressure is: R demands structural efficiency, L resists because reorganisation hurts prediction (initially).
Over time, L discovers that structurally efficient representations are also better for prediction, the adversarial pressure resolves into cooperation.

\paragraph{The biological parallel.}
The left hemisphere initially resists the right's contextualisation (children are very literal before developing metaphorical understanding).
Over development, the two hemispheres learn to cooperate, the right provides context that improves the left's processing.
The adversarial-to-cooperative transition IS cognitive development.


% ======================================================================
\section{What This Would Produce}

If the interaction works, L's representations should develop:

\paragraph{Structure-content separation.}
L's hidden states decompose into a structural component (shared across similar contexts) and a content component (context-specific).
We already have empirical evidence for this: in the ORN with SharedM, function words (``upon'' 0.75, ``play'' 0.68) have consistent coupling across contexts while names (``Tom'' 0.33, ``Lily'' 0.40) have variable coupling.
R's job is to discover and enforce this separation explicitly, finding the boundary through learning rather than relying on the architectural constraint of weight sharing.

\paragraph{The operator emerges.}
The relational structure that R discovers IS the operator.
Not designed, not a matrix or polynomial, but an emergent property of the co-evolution.
All the properties the operator programme sought (composability, continuity, spectral decomposition) emerge from the interaction rather than being designed in.


% ======================================================================
\section{First Experimental Results}
\label{sec:first-results}

\paragraph{Disclaimer (added 21 April 2026).}
All results in this section, including the val-loss trajectory (5.15\,$\to$\,4.63), gate opening (0.04\,$\to$\,0.76), state-norm growth, and all per-phase pred-loss numbers, are evaluated against \emph{val-loss improvement}, a criterion the programme has since rejected.
R's structural content is S$_n$-invariant and therefore orthogonal to next-token prediction by construction; reducing val loss is not a test of R's structural claim.
The numbers below are preserved as the historical record of the earliest co-evolution runs, but the reader should not read them as evidence for or against R.
See the \textsc{Document Status} note after the abstract, and the Phase 1 checklist in \texttt{lorn\_gauge v3} for the current evaluation methodology.

Preliminary results from a small-scale test (d=192, 4 layers, TinyStories, 4K steps).
Three phases: L alone (1K steps), R observes frozen L (1K steps), co-evolution (2K steps).

\begin{center}\small
\begin{tabular}{rllrrrr}
\toprule
\textbf{Step} & \textbf{Phase} & \textbf{Val} & \textbf{Gate} & \textbf{State} & $|A|_{\max}$ & \textbf{Pred loss} \\
\midrule
1000 & L alone & 5.15 & 0.04 & 0.21 & 0.64 & 0.005 \\
2000 & R observes & 5.08 & 0.04 & 0.19 & 0.64 & 0.0001 \\
2500 & co-evolve & 4.90 & 0.24 & 0.18 & 0.66 & 0.0002 \\
3000 & co-evolve & 4.70 & 0.69 & 0.30 & 0.94 & 0.001 \\
3500 & co-evolve & 4.67 & 0.73 & 0.41 & 0.92 & 0.002 \\
3750 & co-evolve & 4.63 & 0.76 & 0.47 & 0.93 & 0.002 \\
\bottomrule
\end{tabular}
\end{center}

\paragraph{Four findings.}

\begin{enumerate}[nosep]
\item \textbf{R learns to predict L's coupling (Phase 2).}
Prediction loss drops from 0.005 to 0.0001, a 50$\times$ improvement.
L's coupling dynamics ARE compressible into R's state.
The structural signal exists.

\item \textbf{The gate OPENS during co-evolution.}
From 0.04 (nearly closed) to 0.76 (76\% open).
L actively uses R's output.
The system chose to integrate the hemispheres, this was not forced.

\item \textbf{L improves under R's influence.}
Val loss: 5.08 (end of Phase 2) $\to$ 4.63 (co-evolution).
A 0.45 nat improvement from the hemispheric interaction.
R's relational context makes L a better predictor.

\item \textbf{R's state develops persistence.}
$|A|_{\max}$ increases from 0.64 to 0.93, approaching the unit circle.
The SSM eigenvalues move toward persistence: R learns to maintain relational context longer as the interaction develops.
State norm grows from 0.19 to 0.47: more relational information accumulated.
\end{enumerate}

\paragraph{Interpretation.}
The co-evolution is productive.
R discovers compressible structure in L's coupling dynamics.
When R feeds this back, L uses it (the gate opens voluntarily) and improves.
R's state develops the spectral properties we sought in the operator programme, persistence, growing information content, but as an \emph{emergent} property of the interaction, not a designed feature.

The most striking finding: the gate.
The gate parameter was initialised closed (bias $= -2$) and learned to open because R's relational context reduced L's prediction loss.
This is hemispheric integration emerging from learning, not from architecture.

\paragraph{The parameter control caveat.}
A wider-FFN model with matched parameters (no R) achieved val 4.29, beating the two-hemisphere model's 4.60.
The val loss improvement from R may be partly explainable by extra parameters rather than the hemispheric interaction.

However: val loss is not the objective.
The gate pattern revealed that content words (``walked,'' ``flew,'' ``eating'') use R fully (gate = 1.0) while punctuation and discourse markers (``.'' ``!'' ``Once'') ignore R (gate $< 0.07$).
R learned a meaningful distinction between tokens that need relational context and tokens that don't, regardless of whether this improved val loss.

The next version uses order-independent accumulation and a structural-not-content loss.

\paragraph{Structural version results.}
Order-independent R with structural-not-content loss:
\begin{center}\small
\begin{tabular}{lr}
\toprule
\textbf{Metric} & \textbf{Value} \\
\midrule
$\mathrm{corr}(d_R, d_{\mathrm{coupling}})$ & \textbf{0.984} (R tracks structure) \\
$\mathrm{corr}(d_R, d_{\mathrm{embedding}})$ & \textbf{$-$0.672} (R anticorrelates with content) \\
Val (with R) & 4.56 \\
Val (without R) & 4.96 \\
Gate trajectory & 1.0 $\to$ 0.04 (L learns to close the gate) \\
\bottomrule
\end{tabular}
\end{center}

\textbf{R learned the right thing.}
The structural-not-content loss works: R encodes coupling structure (0.98 correlation) and actively separates from content ($-$0.67 anticorrelation).
The two hemispheres learn genuinely different representations, exactly as the programme requires.

\textbf{L closes the gate, and that is fine.}
L does not need R's structural context for next-token prediction on TinyStories at this scale.
The gate closing is L honestly reporting: ``I can predict the next token without relational context.''
This is expected.
TinyStories has simple structure; sequence length is 128; the model is small.
L closing the gate does not mean R is useless, it means L does not \emph{yet} need R for \emph{this task}.

The coupling between hemispheres will become essential when:
\begin{itemize}[nosep]
\item \textbf{Scale}: at $d = 2048$ on real text, L needs relational context for complex constructions.
\item \textbf{Sequence length}: at $S = 4096+$, L's attention cannot cover everything, R's compressed state becomes essential.
\item \textbf{Task pressure}: evaluation on tasks that require relational reasoning (analogy, compositional generalisation, long-range coherence) rather than just next-token prediction.
\item \textbf{Curriculum}: train L first, train R, then gradually couple, the first experiment showed the gate opens to 81\% under this regime.
\end{itemize}

Meanwhile, R as a decoupled observer is already valuable: it provides a structural diagnostic of L's processing, a compressed relational representation for external use (retrieval, steering, interpretability), and a foundation for future coupling when the task demands it.

\paragraph{Injection mode results (ongoing).}
Testing five injection modes (additive residual, coupling bias, value modulation, cross-attention, per-layer gated coupling) reveals that HOW R feeds into L matters more than WHETHER it does:

\begin{center}\small
\begin{tabular}{lrrrrr}
\toprule
\textbf{Mode} & \textbf{Val} & \textbf{Gate} & \textbf{CorrC} & \textbf{CorrE} & \textbf{Sep} \\
\midrule
L-only baseline & 4.87 & ,  & ,  & ,  & ,  \\
Additive residual & 4.67 & 0.16 & 0.87 & $-$0.32 & 2.2$\times$ \\
Coupling bias & \textbf{4.54} & \textbf{0.86} & 0.90 & $-$0.47 & 4.9$\times$ \\
Value modulation & 4.69 & 0.34 & \textbf{0.93} & \textbf{$-$0.55} & \textbf{5.5$\times$} \\
Cross-attention & 4.84 & 0.90 & 0.86 & $-$0.19 & 5.4$\times$ \\
Per-layer gated & \textbf{4.56} & 0.83 & 0.79 & $-$0.41 & 2.0$\times$ \\
\bottomrule
\end{tabular}
\end{center}

\noindent CorrC $=$ corr(R-distance, coupling-distance), want high.
CorrE $=$ corr(R-distance, embedding-distance), want low (negative).
Sep $=$ between-group / within-group R-state distance.

\textbf{Finding 1: the injection mechanism determines whether L uses R.}
With additive injection, L closes the gate to 0.16, R's content bias is ignorable.
With coupling bias injection, L \emph{keeps the gate open} at 0.86, R's structural bias on attention patterns IS useful.
Coupling injection changes who-attends-to-whom; additive injection merely shifts the residual stream.
L wants structural context, not content bias.

\textbf{Finding 2: value modulation learns the best structural representation.}
Value modulation achieves the highest coupling correlation (0.93), strongest content decorrelation ($-$0.55), and best structural separation (5.5$\times$).
But L only partially uses it (gate 0.34).
Feature selection is a weaker channel than coupling modulation for influencing L's prediction.

\textbf{Finding 3: cross-attention consults R heavily but R doesn't help.}
Cross-attention has the highest gate (0.90) but barely improves over baseline (4.84 vs 4.87).
R's structural representation is weakest in this mode (CorrE only $-$0.19).
L queries R extensively but R learned content proximity, not structural distance.
The cross-attention mechanism gives L too much control, it can extract whatever is easiest, not what is structural.

\textbf{Finding 4: the per-layer gate pattern confirms depth-dependent structural need.}
The per-layer gated mode (second-best val at 4.56) reveals which layers want R's structural context:

\begin{center}\small
\begin{tabular}{rl}
Layer 0--3: & gate $= 0.000$ \quad (L alone ,  basic token processing, local syntax) \\
Layer 4: & gate $= 0.999$ \quad (R fully open ,  structural context essential) \\
Layer 5: & gate $= 0.648$ \quad (R partially open ,  mixed processing) \\
\end{tabular}
\end{center}

\noindent The network independently discovered that deep layers want structural context and shallow layers do not.
This matches the theoretical prediction: shallow layers handle local syntax (no relational context needed), deep layers do compositional processing (structural context essential).
The transition is sharp, not gradual but a phase change at layer 4.

\paragraph{Coupling bias and per-layer gated are the same mechanism.}
Coupling bias (val 4.54, gate 0.86) and per-layer gated (val 4.56, gate 0.83) are effectively tied on val loss and gate magnitude.
The per-layer gated mode is a \emph{superset} of coupling bias, it adds per-layer gates plus an additive residual path.
That both modes converge to the same performance suggests the core mechanism driving the benefit is coupling modulation at deep layers, and everything else is noise.

The difference in structural metrics (coupling bias: CorrC 0.90, Sep 4.9$\times$; per-layer gated: CorrC 0.79, Sep 2.0$\times$) likely reflects the additive path pulling R's representation in two directions, structural for the coupling channel, content-adjacent for the additive channel.
Removing the additive path should recover the structural quality while preserving the per-layer gate diagnostic.

\paragraph{The $\sim$0.85 gate is a sweet spot.}
Both winning modes settle at gate $\approx 0.85$, not 1.0.
Full saturation (gate $= 1.0$) would mean L is fully dependent on R, fragile, because any failure in R's structural signal would propagate directly to L's predictions.
A gate of 0.85 means L retains $\sim$15\% independence: it prefers R's structural context but can fall back to its own processing when R's signal is uninformative.
This is healthy hemispheric integration, cooperative but not dependent.
The brain's corpus callosum mediates influence, not control.

\paragraph{The architecture we choose.}
Combining the injection mode results with the per-layer gate diagnostic, the target architecture is:

\begin{enumerate}[nosep]
\item \textbf{Layers $0$ to $L/2$:} pure L. No R injection, no gate overhead. Shallow layers do local syntax and don't want structural context.
\item \textbf{R observes at layer $L/2$:} order-independent accumulation of L's coupling patterns into structural state.
\item \textbf{Layers $L/2{+}1$ to $L{-}1$:} L $+$ per-layer gated coupling bias from R. Each deep layer has an independent gate controlling how much R's structural state modulates L's attention.
\item \textbf{No additive path.} The additive channel is ignorable (gate $= 0.16$ when tested alone) and degrades R's structural representation when present alongside coupling.
\end{enumerate}

This is the cleanest two-hemisphere design: coupling bias where the network wants it (deep layers), nothing where it doesn't (shallow layers), with per-layer gates that learn the depth profile.
The per-layer gated experiment told us \emph{where} to inject.
The coupling bias experiment told us \emph{how} to inject.
The combined architecture uses both answers.


\section{The ORN Savings Reinvested}
\label{sec:reinvestment}

The Orbital Response Network saves parameters through shared coupling geometry:
SharedM replaces $L$ independent $W_Q, W_K$ matrices with one shared $A, B$.
At $L = 6$ layers, $d = 192$: the ORN saves $\sim$430K parameters ($\sim$3.5\% of the model) from coupling compression alone.
Wide shared FFN provides additional efficiency by replacing $L$ per-layer FFNs with one wider shared FFN at matched quality.

The two-hemisphere programme \textbf{reinvests these savings into the right hemisphere}.
Instead of using the saved parameters for a wider left hemisphere (which our experiments show provides diminishing returns), we invest them in a qualitatively different component: a relational processing pathway that provides structural context, long-range compression, and coherence pressure.

\begin{center}\small
\begin{tabular}{lrr}
\toprule
\textbf{Component} & \textbf{Standard transformer} & \textbf{Two-hemisphere ORN} \\
\midrule
Per-layer Q, K coupling & $L \times 2d^2$ & SharedM: $2d^2$ \\
Coupling savings & ,  & $\sim$430K params \\
Right hemisphere (R) & ,  & $\sim$430K params \\
\midrule
Net parameter change & 0 & 0 \\
\bottomrule
\end{tabular}
\end{center}

The reinvestment is parameter-neutral: the two-hemisphere ORN has the same total parameter count as the standard transformer.
What changes is HOW the parameters are used: fewer in redundant coupling, more in relational processing.
The SharedM programme provides the \emph{budget}; the two-hemisphere programme provides the \emph{investment}.

The right hemisphere's 430K parameters buy:
\begin{itemize}[nosep]
\item A 128-dimensional relational state space ($r^2 = 16{,}384$ effective coupling dimensions).
\item Multi-signal input (L's hidden states + coupling patterns, each with learned encoder).
\item A 2-layer SSM with complex diagonal transition (spectral persistence and rotation).
\item Coherence losses (predictive + VICReg) that force structured representations.
\item A learned gate that reports the hemispheric balance at every token.
\end{itemize}

All for the same total parameter count as the original model.


\section{Connection to the Operator Programme}

The two-hemisphere programme subsumes the operator programme.
Everything we sought, composability, continuity, spectral decomposition, the coupling-response separation, is something R can discover, rather than something we must design.

\begin{center}\small
\begin{tabular}{lll}
\toprule
\textbf{Property} & \textbf{Operator programme} & \textbf{Two-hemisphere programme} \\
\midrule
Composability & Designed (polynomial of M) & Emerges (R's compression forces it) \\
Spectral structure & Designed (eigendecomposition) & Emerges (R discovers independent modes) \\
Context compression & Designed (spectral rotation) & Emerges (R's autoencoder compresses) \\
Grammar/content split & Designed (base + warp) & Emerges (R forces L to separate them) \\
Transfer & Tested (frozen-M) & Forced (R's structural pressure) \\
\bottomrule
\end{tabular}
\end{center}

The operator programme was trying to do R's job by hand.
The two-hemisphere programme lets R do it by learning.


% ======================================================================
\section{Reasoning Lives in the Right Hemisphere}
\label{sec:reasoning}

\subsection{Continuation vs Composition}

The left hemisphere does \textbf{continuation}: given the recent tokens, predict the next one.
This is sequential, local, literal.
``The cat sat on the'' $\to$ ``mat.''
No relational reasoning needed, pattern matching on recent context suffices.
Transformers excel at this.

The right hemisphere does \textbf{composition}: given the accumulated relational structure, determine what MUST follow.
This is algebraic, global, abstract.
``The committee's recommendation was vindicated'' $\to$ the next clause must resolve the three-party judgment structure.
Not because ``vindicated'' is often followed by ``by'', but because the relational geometry demands a resolution.

Current language models do both through the same mechanism (multi-layer attention + FFN).
This works because L layers of interleaving can simulate compositional reasoning through sequential continuation: each layer adds one step of ``reasoning'' by routing information through attention.
But it is SIMULATED composition through ITERATED continuation.
The reasoning is sequential, not algebraic.
It has no compositional guarantee, it can fail silently on novel structures.

\subsection{The Right Hemisphere as Compositional Reasoner}

In the two-hemisphere architecture, the right hemisphere's SSM state IS the relational reasoning engine.
Its state $\mathbf{s}$ accumulates relational structure compositionally: each token's coupling contribution is added to the state, and the state evolves through the operator's geometry.

The state answers: ``given everything I've read, what relational structure exists?''
Not ``what tokens appeared recently'' (that's the left's attention).
But ``what relationships between concepts have been established, and what must follow from them.''

This is infinite-context reasoning in fixed-size space.
The left hemisphere's attention is bounded by the context window ($S$ tokens).
The right hemisphere's state is bounded by $r$ dimensions, but these $r$ dimensions compress arbitrarily long relational history.
At position 10,000 of a novel, the left hemisphere can attend to at most $S$ recent tokens.
The right hemisphere's state reflects 10,000 tokens of accumulated relational structure: who suspects whom, what evidence has been presented, which narrative threads are unresolved.

\subsection{Geometric Attention from Relational State}

The right hemisphere's output is not a prediction.
It is a \textbf{geometric context} that shapes the left hemisphere's attention.

When R's output feeds back to L, it modifies what L attends to.
Without R: L attends based on local context (recent tokens, positional proximity).
With R: L attends based on relational context (what the accumulated structure demands).

``Vindicate'' with R's context: L's attention is shaped by the three-party structure, it attends to the claimant, the challenger, the resolution slot.
``Vindicate'' without R: L attends to recently co-occurring words, ``was,'' ``by,'' ``the.''

R compacts infinite relational context into a geometric signal that L can exploit through its standard attention mechanism.
L does not need to reason, it just attends.
R has already done the reasoning by composing the relational structure.
L reads the result.

\subsection{Seeded Generation: Composing in R, Generating in L}

The most powerful implication of the two-hemisphere design: generation can be \textbf{seeded by composed relational structures} rather than by token prompts.

\paragraph{The mechanism.}
R's state space is where relational structures are composed.
``Judgment + social role + temporal sequence'' can be combined into a complex relational state, without any tokens.
This state is fed to L through the gate.
L generates tokens that are \emph{constrained} by R's composition: the words must express the three-party judgment structure, but the specific words are free.

\paragraph{Continuous reasoning.}
R does not compose once and hand off.
As L generates each token, the token feeds back to R (through coupling summaries).
R absorbs the token's relational implications, updates its state, and modifies the constraint on L's next generation:

\begin{center}\small
\begin{tabular}{rl}
R: & [three-party judgment structure] \\
L: & ``The committee's recommendation was\ldots'' \\
R: & [absorbs ``committee'' $\to$ fills claimant slot, refines state] \\
L: & ``\ldots subsequently vindicated by\ldots'' \\
R: & [absorbs ``vindicated'' $\to$ confirms template, opens arbiter slot] \\
L: & ``\ldots independent analysis.'' \\
R: & [fills arbiter slot, marks structure resolved] \\
\end{tabular}
\end{center}

R provides the plan. L provides the words.
The plan evolves as the words are spoken.
There is no boundary between ``thinking'' and ``speaking.''
R reasons in state space WHILE L generates in token space.
They are interleaved, not sequential.

\paragraph{This is how the brain produces speech.}
The right hemisphere maintains meaning, discourse structure, intent.
The left hemisphere produces words, syntax, articulation.
They operate simultaneously, you do not first compose the full meaning then produce all the words.
The meaning reshapes as you speak.
The words you choose reshape what you mean.

\paragraph{Implications for controllable generation.}
Seeding R's state enables a new form of controllable generation:
\begin{itemize}[nosep]
\item \textbf{Prompt-free generation:} compose a relational structure in R's state space (no tokens needed). L generates text expressing that structure.
\item \textbf{Structural steering:} adjust R's state mid-generation to change the relational trajectory without changing the words already generated.
\item \textbf{Hierarchical planning:} compose high-level structure in R (``story arc: character introduction $\to$ conflict $\to$ resolution''), let L fill in the details at each stage.
\end{itemize}

This is not classifier-free guidance (scalar steering), not prompt engineering (token steering), not the ORN's M-calculus (eigenmode steering).
It is \textbf{relational steering}, generation constrained by a full $r^2$-dimensional composed relational state.


\subsection{How R Reaches L: Injection Modes and the Mechanism of Analogy}
\label{sec:injection-modes}

\paragraph{Disclaimer (added 21 April 2026).}
The numerical results reported in this subsection, additive residual (val 4.87), coupling bias (val 4.54), value modulation (val 4.69), output modulation (val), and their comparisons against baseline, were all evaluated against \emph{val-loss improvement}, a criterion the programme has since rejected (see \S\ref{sec:first-results} and the \textsc{Document Status} note after the abstract).
Because R's content is S$_n$-invariant and orthogonal to NTP by construction, val-loss differences between injection modes do not constitute evidence for or against their structural role.
The qualitative description of each injection mode and its intended mechanism is retained; the val numbers are historical.

The mechanism by which R's structural state influences L is the most consequential architectural choice after the loss.
Different injection modes express different kinds of influence, and the right choice determines whether the architecture can support analogy, metaphor, and compositional constraint.

\paragraph{Mode 1: Additive residual (baseline).}
$\h \gets \h + \sigma(g) \cdot \mathrm{proj}(\mathbf{s})$.
R adds to L's content stream, biasing \emph{what} L predicts.
This is a mood shift, a topic nudge.
It cannot express relational constraints: ``same structure, different content'' has no representation in an additive bias.
\emph{Analogy capacity: weak.}

\paragraph{Mode 2: Coupling bias (the analogy mechanism).}
R's state generates an $(S \times S)$ bias on L's attention logits:
\begin{equation}
\mathrm{logits}_{\mathrm{attn}} \gets \mathrm{logits}_{\mathrm{attn}} + \sigma(g) \cdot \mathbf{q}_R(\mathbf{s}) \, \mathbf{k}_R(\mathbf{s})^\top
\end{equation}
where $\mathbf{q}_R, \mathbf{k}_R: \R^r \to \R^{r_{\mathrm{attn}}}$ are learned projections of R's state.
This literally changes \emph{who attends to whom}.
If R has accumulated ``agent $\to$ patient $\to$ instrument,'' the bias encourages L to create that same relational pattern in new content.
This IS analogy at the mechanism level: same structural skeleton, L fills in the content.
\emph{Analogy capacity: strong.}

\paragraph{Mode 3: Value modulation.}
$\mathbf{v} \gets \mathbf{v} \odot \sigma(\mathrm{proj}(\mathbf{s}))$.
R selects which \emph{features} of the value stream matter given the current relational context.
``Pay attention to the agent features'' or ``the temporal features.''
Feature selection, not relational template.
\emph{Analogy capacity: moderate.}

\paragraph{Mode 4: Cross-attention (L queries R).}
L attends to R's decoded state as additional key--value pairs.
L decides \emph{when} to consult R and \emph{what} to extract.
Most flexible, most like ``asking an oracle.''
But L controls access, it could learn to ignore R entirely.
\emph{Analogy capacity: strong but indirect.}

\paragraph{Mode 5: Per-layer gated coupling (the full architecture).}
R generates coupling bias at \emph{every layer} $l > l_{\mathrm{start}}$, with independent per-layer gates.
Early layers process alone (L learns basic syntax); deeper layers receive R's structural context (relational composition).
The gate pattern is a diagnostic: which depths want structural context?
Shallow layers should be low-gate (local syntax needs no relational context).
Deep layers should be high-gate (abstract composition needs relational structure).
This is the richest mode and the natural foundation for resonating L$\leftrightarrow$R feedback loops.
\emph{Analogy capacity: strongest.}

\paragraph{The developmental schedule.}
R injects only at layers $l > l_{\mathrm{start}}$ (default: layer 3 of 6).
Early layers are L alone, a developmental constraint where L masters basic language before R adds structural context.
This mirrors the biological observation that children develop literal language (left hemisphere) before metaphorical understanding (right hemisphere).
It also prevents R from interfering with L's low-level token processing, where structural context is premature.

\paragraph{Coupling injection enables analogical generation.}
Analogy $=$ same relational structure, different content.
``A is to B as C is to D'' means the coupling pattern (how A relates to B) is preserved in a new content domain (C, D).
If R feeds back through coupling bias, R imposes a structural template on L's attention, same relational skeleton, L fills in new content.
This is structural steering, not true analogical \emph{recognition} (which would require R to identify when two distant situations share structure).
But structural steering is the \emph{generative} half of analogy: given a relational template, produce content that fits it.

Additive injection can bias content (``say something about animals'').
Coupling injection can impose structure (``say something with agent-action-patient structure'').
The second enables analogical generation.
The first is prompting.

\paragraph{Resonating feedback: the path to reasoning.}
With coupling injection at every layer, a feedback loop emerges:
\begin{center}\small
\begin{tabular}{rl}
L layer $l$: & attends with M $+$ R's bias $\to$ produces coupling $C_l$ \\
R: & observes $C_l$ $\to$ updates state $\to$ produces new bias \\
L layer $l{+}1$: & attends with M $+$ R's updated bias $\to \ldots$ \\
\end{tabular}
\end{center}
Each layer's attention is shaped by R's accumulated structural understanding of \emph{all previous layers}.
This is not a one-shot injection but a resonating loop: L's processing refines R's understanding, which refines L's processing.
In the limit, this IS reasoning, iterative refinement of understanding through structural feedback.

The per-layer gates control the resonance strength.
If the loop destabilises: gates close (self-limiting).
If the loop is productive: gates open (self-reinforcing).
The architecture discovers its own optimal resonance pattern.

\paragraph{Beyond bias: structural embeddings (the next generation).}
Coupling bias is a low-rank perturbation on attention logits.
L's own coupling is computed from $d$-dimensional queries and keys; R's bias has rank $r_{\mathrm{attn}} \ll d$.
R whispers into a conversation L dominates.
The bias can nudge coupling but cannot provide a structural \emph{template}, it cannot say ``position 3 is the agent, position 7 is the patient, fill in the content.''

The next architectural step: R provides \textbf{structural embeddings} that enrich L's representations \emph{before} attention computation:
\begin{align}
\h_{\mathrm{enriched}} &= \h + \sigma(g) \cdot R_{\mathrm{embed}}(\mathbf{s}, t) \\
\mathbf{q} &= \h_{\mathrm{enriched}} \, A, \quad \mathbf{k} = \h_{\mathrm{enriched}} \, B
\end{align}
where $R_{\mathrm{embed}}: \R^r \to \R^d$ maps R's state to a $d$-dimensional per-position structural signal.

Now when L computes $\mathbf{q} \cdot \mathbf{k}^\top$, the coupling is jointly determined by content (from L's residual stream) and structure (from R's embedding).
Position 3 is not just ``the token \texttt{cat}'', it is ``the token \texttt{cat} in the agent role.''
L's standard attention mechanism naturally produces structurally appropriate coupling without any modification.

This changes R's role from \emph{biasing} L's attention to \emph{contextualising} L's representations.
R provides the scaffolding; L fills in the content.
The structural embedding is $d$-dimensional (matching L's representation space), not $r_{\mathrm{attn}}$-dimensional (a low-rank compression).
R speaks at the same volume as L.

The biological analogy: the right hemisphere does not override the left's processing.
It changes the \emph{context within which} the left processes.
The left hemisphere still attends, predicts, and generates, but every representation it works with is coloured by R's relational understanding.
This is how metaphor works: the right hemisphere provides a relational frame (``argument as war'') and the left hemisphere fills in the specific language (``she \emph{attacked} his position,'' ``he \emph{defended} his claim'').
The words are the left's; the structure is the right's.

\subsection{Deep R: A Genuine Second Hemisphere}
\label{sec:deep-r}

The previous designs gave R 6--13\% of the model's parameters: one encoder layer, a running mean, and a thin projection.
This is a probe, not a hemisphere.
It reflects precisely the bias McGilchrist identifies: the left hemisphere assumes it is the important one and treats the right as a helper.

If we take the programme seriously, R's task, understanding the full relational structure of everything seen so far, composing it, and generating structural context for L, is \emph{at least as complex} as L's task of predicting the next token.
L predicts one token at a time from local context.
R must compress the entire relational history into a coherent structural representation.
L's job gets easier with local context.
R's job gets \emph{harder} with more context because relational structure accumulates and compounds.

R should therefore have its own multi-layer processing stack: not a single encoder, but a network with depth, hierarchy, and its own form of attention.

\paragraph{R's architecture.}
\begin{itemize}[nosep]
\item \textbf{Input projection:} L's coupling + hidden state $\to$ $r$-dimensional structural features per position.
\item \textbf{$K$ layers of structural self-attention + FFN.} Each layer discovers which positions play similar structural roles (shallow layers: local syntax; deep layers: discourse composition).
\item \textbf{Non-causal attention.} R sees all positions equally. Structure does not depend on surface order, ``the committee was vindicated'' and ``vindicated was the committee'' have the same relational structure. Non-causal attention enforces this by construction.
\item \textbf{Running mean accumulation AFTER deep processing.} R processes deeply, then compresses.
\item \textbf{Structural embeddings projected into L.} R's deep representations are injected into L's residual stream at each deep layer.
\end{itemize}

\paragraph{The parameter budget.}
SharedM saves $(L{-}1) \times 2d^2$ coupling parameters.
Wide shared FFN saves $(L{-}1)$ per-layer FFNs.
Together: at $L = 12$, $d = 288$, the savings total $\sim$8.8M parameters.
These are \emph{reinvested} in R.

\begin{center}\small
\begin{tabular}{lrrr}
\toprule
\textbf{Component} & \textbf{Params} & \textbf{Share} & \textbf{Role} \\
\midrule
L (SharedM + shared FFN + response) & 17.7M & 69\% & Predict tokens \\
R (6 layers, $r = d = 288$) & 7.9M & 31\% & Understand structure \\
\midrule
Total & 25.6M & 100\% & \\
\bottomrule
\end{tabular}
\end{center}

\noindent R/L ratio: 0.45.
R is a genuine processing pathway, one-third of the model.
The ratio is determined by the ORN savings: every parameter freed from redundant coupling and per-layer FFN is reinvested in structural processing.

\paragraph{Why R needs depth, not just width.}
A single-layer R can map coupling patterns to a vector.
But it cannot compose ``this clause is subordinate to that clause, which is part of a three-party argument'', that requires multiple layers of abstraction, just as L needs multiple layers to compose token representations.

R's depth enables hierarchical structural understanding:
\begin{itemize}[nosep]
\item \textbf{Shallow R layers:} which positions have similar coupling patterns (local role detection).
\item \textbf{Mid R layers:} clause boundaries, argument structure, constituent grouping.
\item \textbf{Deep R layers:} discourse-level composition, narrative structure, cross-clause relations.
\end{itemize}

This is the same hierarchy that the probing literature finds in transformer representations~\cite{hewitt2019structural}, but learned by a dedicated network rather than extracted post-hoc from a single network that must simultaneously handle content and structure.

\paragraph{Architecture diagram.}

\input{lorn_architecture_diagram}


\subsection{Comparison with Brain Anatomy}
\label{sec:brain-anatomy}

The architecture above has deliberate parallels with brain hemispheric anatomy, and deliberate differences.
Understanding both is essential for knowing where the analogy is load-bearing and where it is decorative.

\paragraph{What matches.}

\begin{enumerate}[nosep]
\item \textbf{Two processing pathways with different functions.}
The brain's left hemisphere specialises in sequential, categorical, detail-focused processing; the right in holistic, relational, context-sensitive processing~\cite{mcgilchrist2021matter}.
Our L is autoregressive (sequential, token-by-token); R uses non-causal attention (global, all-positions-at-once).
The functional split is analogous.

\item \textbf{Communication through a bandwidth bottleneck.}
The corpus callosum connects the hemispheres through $\sim$200 million axons, large in absolute terms but a severe bottleneck relative to the $\sim$20 billion neurons in each hemisphere~\cite{ringo1994hemispheric}.
SharedM ($2d^2$ parameters connecting two networks of millions of parameters) serves the same role: enough bandwidth to communicate coupling geometry, not enough to transmit full representations.
Ringo et al.\ argued that this bottleneck \emph{causes} hemispheric specialisation, each hemisphere develops self-contained processing to minimise cross-hemisphere communication.
SharedM may have the same effect.

\item \textbf{Depth in both hemispheres.}
Both brain hemispheres have six cortical layers with hierarchical processing: sensory input at layer IV, local processing in layers II/III, output in layers V/VI.
Our L has 12 layers; R has 6.
Both build hierarchical representations, L of content, R of structure.

\item \textbf{Shared subcortical structures.}
The brain's hemispheres share the brainstem, cerebellum, hippocampus, and (partially) the thalamus.
These provide shared functions: the brainstem handles arousal, the cerebellum handles timing, the hippocampus handles memory consolidation.
Our wide shared FFN is the closest analogue, one ``world knowledge'' module that both hemispheres use.
SharedM is the coupling geometry both hemispheres operate through.
\end{enumerate}

\paragraph{What differs.}

\begin{enumerate}[nosep]
\item \textbf{The brain communicates bidirectionally at every depth.}
The corpus callosum connects \emph{homologous areas}: layer $N$ of the left hemisphere connects to layer $N$ of the right~\cite{ringo1994hemispheric}.
Communication happens at every cortical depth, not just at a midpoint.
Our architecture has a single observation point (L layer 6 $\to$ R) and one-directional injection (R $\to$ L layers 7--11).
A more brain-like design would have L$\leftrightarrow$R connections at every layer pair.

\item \textbf{Both brain hemispheres process raw sensory input.}
Each hemisphere receives sensory input from the contralateral visual field, ear, and body side.
The right hemisphere processes its own input, not the left's internal representations.
Our R sees only L's coupling and hidden states, it has no direct access to tokens.
A more brain-like R would process the token sequence independently (perhaps with non-causal attention) and compare its own representation with L's.

\item \textbf{The brain has subcortical routing (thalamus, basal ganglia).}
The thalamus gates what information reaches the cortex.
The basal ganglia select which cortical programmes to activate.
These routing mechanisms have no analogue in our architecture, the gate mechanism is the closest, but it operates per-layer rather than as a central switchboard.

\item \textbf{Hemispheric specialisation is fundamental, not emergent.}
McGilchrist argues that the hemispheres are not identical systems that gradually diverge through experience~\cite{mcgilchrist2021matter}.
They are fundamentally different in kind from the outset, different in neuronal structure, connectivity patterns, neurotransmitter profiles, and developmental trajectory.
The right hemisphere is broader, more densely interconnected, and matures earlier; the left is more tightly organised, more modular, and specialises later.
The division is not an accident of bandwidth but an architectural requirement: attending to the world as a whole (right) and manipulating what is already known (left) are different \emph{kinds} of processing that require different substrates.
Our architecture reflects this by giving L and R fundamentally different designs, causal vs non-causal attention, prediction loss vs structural loss, rather than starting identical and hoping specialisation emerges.
This is aligned with McGilchrist: the hemispheres are not the same system trained differently. They are different systems that cooperate.

\item \textbf{The right hemisphere is not smaller.}
In the brain, the right hemisphere is actually slightly \emph{larger} than the left in most individuals, with more white matter (long-range connections) and a thicker cortex in some regions.
The brain does not treat the right hemisphere as a helper.
Our Deep-R LORN gives R 31\% of parameters, closer to parity than previous designs, but still not equal.
\end{enumerate}

\paragraph{What the brain suggests for future architectures.}

\begin{enumerate}[nosep]
\item \textbf{Bidirectional layer-by-layer communication.}
Connect L layer $k$ to R layer $k/2$ (since R has half the depth) with learned gated projections in both directions.
L$\to$R: ``here is what I found at this depth.''
R$\to$L: ``here is the structural context for this depth.''
The bandwidth of each connection is a learnable bottleneck, the corpus callosum width at each depth.

\item \textbf{R processes raw input independently.}
Give R its own embedding of the input tokens, processed with non-causal attention.
R builds its structural understanding from the tokens directly, not from L's representations.
The two hemispheres compare notes through SharedM rather than R depending on L's output.
This makes R genuinely independent, it could disagree with L, which is how the brain works.

\item \textbf{Deeper architectural asymmetry.}
McGilchrist's position is that the hemispheres are fundamentally different in kind, not identical systems shaped by different experience.
Future architectures should lean into this: L and R should differ not just in their loss and attention mask but in their \emph{computational primitives}.
L might use sharp, sparse attention (categorical, focused); R might use soft, distributed attention or even a different mechanism entirely (e.g.\ state-space dynamics rather than self-attention).
The asymmetry should be structural, reflecting the different \emph{kinds} of processing each hemisphere performs, not merely different training signals applied to the same architecture.

\item \textbf{Subcortical routing.}
Add a small shared network (``thalamus'') that sees both L and R's states and decides, per-layer and per-position, how much information to pass between them.
This replaces the fixed observation-point + injection-layers design with learned, dynamic, content-dependent routing.
The thalamus could learn: ``for this token, L needs R's structural context'' or ``for this token, L should work alone.''

\item \textbf{Resonant feedback loops.}
In the brain, processing is not feedforward.
Information circulates: cortex $\to$ thalamus $\to$ cortex $\to$ thalamus, with each pass refining the representation.
A LORN with recurrent L$\leftrightarrow$R communication, where R's output feeds back to L, L's updated coupling feeds back to R, iterating until convergence, would be the closest analogue to cortical resonance.
This is computationally expensive but may be necessary for genuine compositional reasoning.
\end{enumerate}

\paragraph{Architectural roadmap: increasing the fidelity of L$\leftrightarrow$R communication.}

The current Deep-R architecture has a single observation point (L layer 6 $\to$ R) and one-directional injection (R $\to$ L layers 7--11).
This is the simplest communication pattern.
The brain's corpus callosum connects homologous areas at \emph{every} cortical depth.
We approach this incrementally, with each step adding one capability and testing independently.

\textbf{Step 1 (current): Single observation, per-layer write.}
R sees L's coupling at one depth (midpoint), processes deeply, writes structural embeddings to L's deep layers.
Validates: can R learn structure? Do gates open? Does structural separation exceed a probe baseline?

\textbf{Step 2: Per-layer observation, per-layer write.}
R reads L's coupling at \emph{multiple} depths: R layer $k$ receives L layer $2k$'s coupling.
R builds a hierarchical structural understanding, shallow R layers see local syntax (L's early coupling), deep R layers see abstract composition (L's deep coupling).
R still writes to L's deep layers.
Validates: does hierarchical observation produce richer structural representation than single-point observation?

\textbf{Step 3: Bidirectional per-layer communication.}
R reads from AND writes to L at corresponding depths.
L layer $2k$ sends coupling to R layer $k$; R layer $k$ sends structural embeddings back to L layer $2k{+}1$.
Every other L layer is observed; every alternate L layer is enriched.
The communication is bidirectional but bandwidth-limited: each connection passes through a gated projection (the corpus callosum at that depth).
Validates: does bidirectional coupling produce genuinely different L representations, or does one-directional injection already capture the benefit?

\textbf{Step 4: Dynamic routing (subcortical control).}
A small shared network (``thalamus'') sees both L and R's states at each depth and decides, per-position and per-layer, how much information to pass in each direction.
This replaces fixed connections with learned, content-dependent routing.
Validates: does dynamic routing outperform fixed connections?

\textbf{Step 5: Resonant iteration.}
L$\leftrightarrow$R communication iterates: R's output feeds back to L, L's updated coupling feeds back to R, repeating until convergence or a fixed number of iterations.
This is the closest analogue to cortical resonance (cortex $\to$ thalamus $\to$ cortex cycles).
Validates: does iterative refinement improve compositional reasoning over single-pass communication?

Each step is independently testable and builds on the previous.
The current experiments validate Step 1.
Step 2 is the next architectural priority, it requires the least change (R's input expands, everything else stays) and tests the most important hypothesis (does R benefit from seeing L's full depth hierarchy?).

\paragraph{Why R does not need raw token input.}
In the brain, both hemispheres receive sensory input.
It is tempting to give R direct access to tokens.
But McGilchrist's point is not that both hemispheres see the same raw signal, it is that they attend to it differently.
The right hemisphere attends to the \emph{whole field}; the left focuses on \emph{specific items}.

In our architecture, L's coupling patterns ARE the relational field.
R's ``sensory input'' is the coupling: not raw tokens but tokens \emph{as structured by L's attention}.
R sees how things relate (the coupling), not what things are (the tokens).
This enforces the content--structure separation by architecture: R \emph{cannot} learn content because it never sees tokens.
Giving R raw tokens would weaken this separation, R could learn a second language model rather than a structural processor.

\paragraph{The honest assessment.}
Our architecture captures the \emph{principle} of hemispheric lateralisation (two pathways, different functions, bandwidth-limited communication) but not yet the full \emph{mechanism} (bidirectional, layer-by-layer, with subcortical routing).
The roadmap above traces the path from principle to mechanism.
The brain had $\sim$500 million years of evolution to discover these design choices.
We are at step one.


\subsection{Why This Is Different from Chain-of-Thought}

Chain-of-thought reasoning externalises composition as token generation: the model produces intermediate tokens (``Let me think step by step...'') that serve as working memory.
This works but is expensive (many tokens) and fragile (the reasoning can go off-track at any step).

The two-hemisphere design internalises composition as state evolution.
R's state evolves through operator dynamics, composition happens in $r$ dimensions, not in $S$-length token sequences.
The ``reasoning'' is continuous, structured, and bounded, it does not consume output tokens or expand the sequence.

This is the distinction between \emph{thinking out loud} (chain-of-thought, left hemisphere) and \emph{understanding} (relational state, right hemisphere).
Both are useful.
The two-hemisphere architecture has both.


\section{Why Nature Requires Two Hemispheres}

McGilchrist's answer: because the world requires two kinds of attention.
You need focused attention to grasp a specific grain of food (left hemisphere).
You need broad attention to watch for predators while eating (right hemisphere).
These are not the same task, and they require different processing.

Our answer: because representation requires two kinds of pressure.
You need prediction pressure to learn accurate, specific representations (the left network).
You need structural pressure to learn abstract, transferable representations (the right network).
These are not the same loss, and they require different optimisation.

A single-loss system converges to whatever is easiest for that loss.
For next-token prediction: memorise facts in FFN, route tokens by attention, done.
No pressure toward structure, metaphor, composition.

A two-loss system creates a tension: prediction wants specificity, structure wants generality.
The resolution of this tension IS the emergence of intelligence, representations that are simultaneously specific enough to predict and general enough to transfer.

The brain found this solution through evolution.
We propose to find it through co-evolving networks.


% ======================================================================
% ======================================================================
\section{Literature: Converging Evidence}

Several independent lines of research converge on the two-hemisphere design.

\paragraph{Hybrid attention-SSM models.}
Jamba~\cite{lieber2024jamba}, Griffin~\cite{de2024griffin}, Zamba~\cite{glorioso2024zamba}, Hymba~\cite{nvidia2024hymba}, and Samba~\cite{microsoft2024samba} all combine transformer attention with state-space models.
The consistent finding: SSMs handle smooth context evolution and long-range state propagation; attention handles discrete recall and comparison.
The optimal ratio is heavily SSM-dominant (5:1 to 7:1)~\cite{lieber2024jamba}.
This maps onto the hemispheric division: the right hemisphere (SSM) maintains broad contextual state, the left hemisphere (attention) performs precise retrieval.

\paragraph{Zamba's shared attention is broader than SharedM.}
Zamba~\cite{glorioso2024zamba} uses an SSM backbone with shared attention blocks injected at regular intervals, but shares the \emph{entire} attention block (Q, K, V, O projections and MLP), not just the coupling geometry.
This is closer to ALBERT~\cite{lan2020albert} than to SharedM: everything is shared, with per-invocation variation coming only from (a) the evolving SSM state that feeds each invocation and (b) a single unshared linear projection back to the residual stream.
SharedM shares only the coupling (Q, K / the matrix $\M = AB^\top$) while keeping per-layer response heads (V, O) and FFN, the principle being that \emph{coupling is invariant but response must vary with depth}.
Zamba2 partially recovers this distinction by adding LoRA projectors on the shared MLP, reintroducing per-layer response variation, approaching the ORN's coupling-response decomposition from the ALBERT direction.

\paragraph{Hymba: parallel beats interleaved.}
Hymba~\cite{nvidia2024hymba} runs attention heads and SSM heads \emph{in parallel} within the same layer, rather than interleaving them across layers.
Finding: parallel outperforms interleaved at small scales.
Heads specialise: attention handles recall, SSM handles language modelling.
This supports our design: two pathways operating simultaneously, not sequentially.

\paragraph{The bandwidth bottleneck creates specialisation.}
Ringo et al.~\cite{ringo1994hemispheric} argue that biological hemispheric specialisation arises because the corpus callosum is a bandwidth bottleneck.
Each hemisphere develops self-contained processing to minimise cross-hemisphere communication.
In our architecture: if the two pathways communicate only through a shared coupling geometry $\M$ (a narrow bottleneck), specialisation should emerge naturally, the left specialises in literal attention, the right in state compression.

\paragraph{Adversarial factorisation creates steerable representations.}
Domain-Adversarial Neural Networks~\cite{ganin2016dann} force representations to be domain-invariant by adversarial pressure.
Fader Networks~\cite{lample2017fader} create steerable representations by forcing attribute information out of the encoder.
FactorVAE~\cite{kim2018factorizing} uses adversarial objectives to disentangle latent dimensions.
All demonstrate: adversarial pressure between two networks produces more structured representations than a single loss can.

\paragraph{VICReg-style structural losses.}
Bardes et al.~\cite{bardes2022vicreg} show that explicit variance-invariance-covariance regularisation prevents representation collapse and encourages diverse, informative features, without labels or contrastive pairs.
Applied to the right hemisphere's relational state, such losses could maintain spectral structure and prevent the collapse we observed in operator programme experiments.

\paragraph{Global Workspace Theory.}
Goyal et al.~\cite{goyal2022global} (Bengio group) show that a shared bottleneck through which all modules communicate forces compression and improves compositionality and out-of-distribution generalisation.
SharedM is this global workspace, the coupling geometry that both hemispheres must communicate through.


% ======================================================================
\section{Two-Hemisphere Attention: The Architecture}
\label{sec:two-hemi-arch}

Combining the literature findings with our operator programme results:

\subsection{Design}

Each layer has two parallel pathways:

\paragraph{Left hemisphere (attention).}
Standard transformer self-attention with SharedM coupling ($Q = \h A$, $K = \h B$).
Sharp, literal, per-token coupling.
Good at: associative recall, precise token-token routing, factual retrieval.
\begin{equation}
\h_{\mathrm{left}} = \h + W_O \, \mathrm{softmax}\!\left(\frac{(\h A)(\h B)^\top}{\sqrt{r}}\right) W_V \h
\end{equation}

\paragraph{Right hemisphere (SSM).}
A state-space model that operates on the left hemisphere's representations, not on raw tokens.
Compresses context into a persistent relational state.
Good at: long-range context propagation, structural compression, relational state tracking.
\begin{align}
\mathbf{s}_t &= A_{\mathrm{state}} \, \mathbf{s}_{t-1} + B_{\mathrm{in}} \, f(\h_{\mathrm{left},t}) \\
\h_{\mathrm{right},t} &= C_{\mathrm{out}} \, \mathbf{s}_t
\end{align}

The right hemisphere sees what the left \emph{found relevant} (through $f(\h_{\mathrm{left}})$), not raw tokens.
Its state $\mathbf{s}$ IS the compressed relational context, the accumulated structural discoveries.

\paragraph{Interaction through SharedM.}
The two pathways communicate through the shared coupling geometry:
\begin{equation}
\h_{\mathrm{out}} = \h + \h_{\mathrm{left}} + \sigma(g) \cdot \h_{\mathrm{right}} + \mathrm{FFN}(\h)
\end{equation}

The gate $\sigma(g)$ controls the hemispheric balance, how much the right hemisphere's relational context influences the output.
$\M = AB^\top$ is shared between the left's attention coupling and the right's state transition ($A_{\mathrm{state}}$ derived from $\M$'s inner geometry $G = B^\top A$).
One geometry, two uses: sharp coupling in attention, smooth evolution in state.

\subsection{The Self-Measuring Properties}

\begin{itemize}[nosep]
\item \textbf{The gate $\sigma(g)$}: reports the hemispheric balance per token. High = right hemisphere contributes (relational context needed). Low = left hemisphere suffices (literal recall enough). Function words should have low gate; complex relational tokens should have high gate.
\item \textbf{The right-hemisphere state norm $\|\mathbf{s}_t\|$}: reports how much relational context has accumulated. Should grow with sequence position. Should grow faster for complex texts than simple ones.
\item \textbf{The right-hemisphere state rank}: reports the dimensionality of the relational structure. Should stabilise (crystallise) during training, analogous to $\M$'s rank crystallisation.
\item \textbf{The structural loss $\mathcal{L}_R$}: reports how compressible L's representations are. Should decrease during training (L learns more structured representations under R's pressure).
\end{itemize}

\subsection{Order-Independence: R Accumulates Meaning, Not Sequence}
\label{sec:order-independence}

L is autoregressive: it must process tokens in order.
R is not.
R accumulates relational MEANING, which does not depend on surface order.
``The committee was vindicated'' and ``Vindicated was the committee'' have the same relational meaning.
R should produce similar states for both.

This is not a convenience, it is a design constraint that \textbf{prevents shortcuts}.
If R processes tokens in order, it can learn: ``position 5's coupling usually looks like X.''
It predicts from position, not from structure.
The loss decreases but R learned nothing about relations.

If R is order-independent, position is destroyed.
R MUST learn: ``this token activates THIS relational pattern, regardless of where it appears.''
That is structural learning.

The same principle explains why bag-of-words models learn semantic features: destroying word order forces the model to learn meaning rather than position.
R's order-independence forces it to learn structure rather than sequence.

\paragraph{Implementation.}
R's accumulation is a sum (or mean) over token contributions, not a sequential recurrence:
\begin{equation}
\mathbf{s} = \frac{1}{t} \sum_{i=1}^{t} f(\text{signal}_i)
\end{equation}
where $f$ is a learned encoder that maps each token's signal (from L) into R's state space.
The sum is commutative, order does not matter.
The spectral rotation (from the operator programme) can still be applied to the ACCUMULATED state, providing structured aging without imposing sequential processing.

\paragraph{What L provides and R ignores.}
L's attention patterns ARE order-dependent (causal masking).
The order information enters through L's coupling, R sees which tokens L coupled, not the order of processing.
R ignores surface order but inherits structural order from L's coupling.
``Cat chased mouse'' and ``mouse chased cat'' have \emph{different} L-coupling patterns (different agent/patient attention), so R receives different signals, not because of token order, but because L's attention encoded the role structure.

\paragraph{Caveat: R's order-independence depends on L's order-sensitivity.}
R is symmetric over its inputs (sum or mean).
If two sentences contain the same tokens but in different syntactic roles, R distinguishes them \emph{only because} L's causal attention produces different coupling patterns for the two orderings.
R does not itself recognise order, it trusts L to encode order into coupling.
If L's coupling were order-insensitive (e.g.\ bidirectional attention with identical patterns regardless of order), R would conflate structurally different sentences that share the same token set.
The order-independence claim is: R accumulates abelianly over L's \emph{already-ordered} representations, not that R understands order-independent meaning directly.


\subsection{Coherence: The Right Hemisphere's Learning Objective}
\label{sec:coherence}

The right hemisphere does not learn by compression, prediction, or consistency.
It learns by \textbf{coherence}, finding how the left hemisphere's coupling patterns fit together as a meaningful whole.

\paragraph{Why not compression.}
Compression says: make it smaller.
But R is not trying to reduce L's representations.
A compressed representation can be lossy in arbitrary ways.
R needs to find the relational STRUCTURE that generates the coupling, the low-dimensional geometry from which all the coupling patterns are different views.
This is not compression; it is coherence.

\paragraph{Why not prediction.}
Prediction says: guess what comes next.
We use prediction as a proxy loss (predict L's future coupling), but prediction is not the objective.
The objective is: find a state that makes all of L's coupling patterns cohere, that reveals them as instances of one relational structure.
Prediction works as a loss because: if R's state captures the generative structure, it CAN predict (the structure determines what follows).
But prediction is a consequence of coherence, not its definition.

\paragraph{Why not consistency.}
Consistency says: things should be the same.
But L's coupling patterns SHOULD vary, different contexts produce different coupling.
The goal is not sameness but \emph{structural unity across diversity}: different coupling patterns that are all coherent manifestations of one relational geometry.
Like chapters of a novel, different events, one story.

\paragraph{The physics analogy.}
Coherent light (laser): waves maintain fixed phase relationships across space and time.
They interfere constructively, carry information, and propagate without degradation.
Incoherent light (bulb): random phases, no interference, information degrades with distance.

R maintains coherence of the relational field as context accumulates.
Without R: L's representations decohere over long sequences, attention fades, early context is lost, the relational structure fragments.
With R: the relational field stays coherent, the accumulated structure forms a meaningful whole that L can exploit at any position.

\paragraph{Coherence is self-supervised.}
Coherence is a property of the signal, not an external label.
You do not need someone to tell you ``this narrative is coherent.''
You measure coherence from internal structural consistency: do the parts fit together?
Does the coupling at position 50 cohere with the coupling at position 5?
Can one low-dimensional state account for both?

R's loss is a measure of coherence:
\begin{equation}
\mathcal{L}_{\mathrm{coherence}} = \text{``how well does } \mathbf{s} \text{ make L's coupling patterns fit together?''}
\end{equation}

The specific instantiation (predictive loss, distance matching, VICReg) is a proxy.
The conceptual objective is coherence.
The right hemisphere is a coherence engine.


\subsection{Required Properties of the Loss System}

The loss is not one objective.
It is a system of pressures that must together produce specific properties.
We list the required properties, then design losses to achieve them.

\paragraph{Properties inherited from the operator programme.}
\begin{enumerate}[nosep]
\item \textbf{Composability.} The right hemisphere's state must compose: processing tokens A then B should produce a state related to processing A and B independently. The loss should reward compositional state evolution.
\item \textbf{Context compression.} The right hemisphere's state is fixed-size ($d_{\mathrm{state}}$ dimensions) regardless of sequence length. The loss should reward informative compression, not just any compression, but compression that preserves relational structure.
\item \textbf{Spectral structure.} The state should develop independent modes (analogous to $\M$'s eigenspectrum). Different relational dimensions should be decorrelated. The loss should prevent mode collapse and encourage spectral diversity.
\item \textbf{Transparency / self-measuring.} The loss itself should be interpretable: its value reports how much structure the right hemisphere has found. The gate reports the hemispheric balance. Both are diagnostics, not just training signals.
\end{enumerate}

\paragraph{Properties from the TTT connection.}
\begin{enumerate}[nosep, resume]
\item \textbf{Train-test invariance.} The structural loss should be computable at inference time, not just during training. The right hemisphere operates the same way at train and test: it compresses context through state dynamics. If the structural loss can be evaluated at inference, it becomes a live diagnostic of representation quality.
\item \textbf{Forward-only operation.} The right hemisphere must not require gradients at inference. Its state evolves through forward dynamics (like our hot field). The structural loss trains the dynamics; at inference, the dynamics run without the loss.
\end{enumerate}

\paragraph{Properties from the long-context programme.}
\begin{enumerate}[nosep, resume]
\item \textbf{Graceful aging.} Old context should fade smoothly, not abruptly. The loss should reward states where old information decays at rates determined by its structural importance, important patterns persist, transient details fade. This is the spectral rotation property: different modes age at different rates.
\item \textbf{Infinite context in principle.} The state must not saturate, at position 10,000, it should still be accumulating useful structure, not plateauing. The loss should penalise saturation (state norm bounded but information content growing).
\end{enumerate}

\paragraph{Practical requirements.}
\begin{enumerate}[nosep, resume]
\item \textbf{Unsupervised.} No labelled metaphors, no annotated relational pairs. The structural loss discovers structure from the left hemisphere's representations alone.
\item \textbf{Rich gradient signal.} The loss must provide gradients to all parts of the model, not just the prediction head but the coupling geometry $\M$, the state transition, the warp projections. A loss that only touches the output layer cannot shape the internal structure.
\item \textbf{Compatible with next-token prediction.} The structural loss should not degrade the primary task. At worst: the Pareto frontier (some prediction quality traded for much structural quality). At best: structural pressure actually improves prediction (as adversarial pressure can improve generalisation).
\item \textbf{Parallelisable.} The loss computation should not be $\mathcal{O}(S^2)$ in sequence length. Per-position or per-batch computation is acceptable.
\end{enumerate}

\paragraph{The fundamental tension.}
\begin{enumerate}[nosep, resume]
\item \textbf{Specificity vs generality.} The left hemisphere's loss rewards specificity (predict THIS token). The right hemisphere's loss rewards generality (compress relational structure). The tension between them IS the learning signal. Too much specificity: memorisation, no structure. Too much generality: vague, no prediction. The $\lambda$ weighting controls the balance, but the equilibrium should be dynamic, $\lambda$ could itself be learned.
\end{enumerate}

\subsection{What the Right Hemisphere Sees}
\label{sec:what-right-sees}

The right hemisphere's input is the most consequential design decision.
Several candidates exist, each carrying a different type of information from the left hemisphere's processing.

\paragraph{Candidate 1: Hidden states through a bottleneck.}
$z_t = \mathrm{project}(\h_{\mathrm{left},t})$ where $\mathrm{project}: \R^d \to \R^r$, $r \ll d$.
The bottleneck is the corpus callosum, limited bandwidth forces compression.
Only structure that fits through $r$ dimensions survives; content is too high-dimensional.
\emph{Risk:} the bottleneck may select for easy features (positional, frequency) rather than relational structure.

\paragraph{Candidate 2: Attention patterns (the coupling field).}
The right hemisphere sees how the left couples tokens, pure relational signal, stripped of content.
``Who attends to whom'' is relation without identity.
\emph{Risk:} $S \times S$ per layer, expensive. Needs a coupling summary.

\paragraph{Candidate 3: Pre-softmax logits (the full coupling landscape).}
Softmax discards weak couplings.
The right hemisphere sees the FULL landscape including relationships the left ignores.
Richer signal than post-softmax attention.
\emph{Risk:} noisy, large dynamic range.

\paragraph{Candidate 4: Layer-to-layer deltas (what each layer discovered).}
$\Delta\h_l = \h_{l+1} - \h_l$, the increment, not the accumulation.
The right hemisphere tracks the DYNAMICS of the left's processing.
Early layers: small deltas (quiet). Deep layers: large deltas (reorganisation).
\emph{Risk:} noisy (deltas are small relative to states).

\paragraph{Candidate 5: The value stream (what attention gathered).}
$\mathbf{v}_t = W_O(C \cdot W_V \h)$, what the left's attention selected as relevant.
Content-filtered by attention: higher SNR than raw states.
\emph{Risk:} mixes content and structure.

\paragraph{Candidate 6: The correction/gate signal.}
The ORN's gate opens where the left hemisphere is uncertain.
The right hemisphere sees WHERE the left struggles.
Elegant but sparse (gate opens rarely).

\paragraph{The principled answer: learn what to see.}
Choosing which signal the right hemisphere sees is itself a left-hemisphere design decision, categorical, explicit, prescribed.
The right-hemisphere answer: \textbf{see everything and learn what matters.}

\begin{equation}
\mathbf{z}_t = \mathrm{RightAttention}\!\left(\h_t, \, C_t, \, \Delta\h_t, \, \mathbf{v}_t, \, \mathrm{logits}_t, \, g_t\right)
\end{equation}

The right hemisphere receives ALL signals from the left, hidden states, attention patterns, deltas, values, pre-softmax logits, gate activations, and learns which are useful for discovering relational structure.
It has its own attention mechanism over the left's internal signals: a meta-network that reflects on another network's processing.

This is \textbf{introspection}, a network that looks at itself.

The right hemisphere's learned attention over these signals IS a diagnostic:
after training, we observe \emph{which signals it selected}.
If it attends primarily to attention patterns: relational structure lives in coupling.
If it attends to deltas: structure lives in the dynamics.
If it attends to hidden states: structure is in the representations directly.
The architecture tells us what it needed, we don't have to guess.

\paragraph{Start broad, prune later.}
The initial architecture includes all signal channels.
This is slow (many inputs to process) but maximally flexible.
Over training, the right hemisphere's attention over channels will concentrate, some channels will get near-zero weight.
In future architectures, we drop the ignored channels.
The first generation discovers what matters.
The second generation is efficient.

This is the same principle as Neural Architecture Search applied to the input space: let the model discover its own optimal input configuration, then crystallise the discovery into a fixed architecture.


\subsection{Candidate Losses}

Given these 13 properties, we propose a system of losses that together satisfy them.

\paragraph{Loss 1: Predictive state coherence.}
The right hemisphere's state at position $t$ should predict the left hemisphere's representation at position $t + k$:
\begin{equation}
\mathcal{L}_{\mathrm{predict}} = \sum_k \alpha_k \, \| C_k \, \mathbf{s}_t - \mathrm{stopgrad}(\h_{\mathrm{left}, t+k}) \|^2
\end{equation}
The lookahead $k$ controls abstraction: $k = 1$ is nearly literal, $k = 10$ forces abstract structure.
Multiple $k$ values (weighted by $\alpha_k$) encourage multi-scale prediction.
The $\mathrm{stopgrad}$ on the target prevents the left hemisphere from collapsing its representations to match the right's predictions.

\emph{Properties satisfied:} context compression (forces informative state), rich gradient signal (flows through state transition), forward-only (computable at inference as a diagnostic), unsupervised.

\paragraph{Loss 2: Spectral diversity (VICReg-inspired).}
On the right hemisphere's state trajectory $\mathbf{s}_1, \ldots, \mathbf{s}_S$:
\begin{align}
\mathcal{L}_{\mathrm{var}} &= \sum_i \max(0, \gamma - \mathrm{Var}(\mathbf{s}_{:,i})) \quad \text{(keep all dimensions alive)} \\
\mathcal{L}_{\mathrm{cov}} &= \sum_{i \neq j} [\mathrm{Cov}(\mathbf{s}_{:,i}, \mathbf{s}_{:,j})]^2 \quad \text{(decorrelate dimensions)}
\end{align}

\emph{Properties satisfied:} spectral structure (independent modes), prevents collapse, parallelisable (computed per-batch).

\paragraph{Loss 3: Compositional consistency.}
If the left hemisphere processes prefix $A$ and then token $x$, and separately processes prefix $B$ and then the same token $x$, the right hemisphere's state \emph{change} from $x$ should be similar:
\begin{equation}
\mathcal{L}_{\mathrm{compose}} = \| \Delta\mathbf{s}_x^{(A)} - \Delta\mathbf{s}_x^{(B)} \|^2 \quad \text{where } \Delta\mathbf{s}_x = \mathbf{s}_{t+1} - \mathbf{s}_t \text{ at token } x
\end{equation}
This says: the same token should produce the same relational state change regardless of context.
This IS composability, the token's contribution is context-independent (abelian).

\emph{Properties satisfied:} composability (same token = same state change), transparency (the state change IS the token's relational contribution), unsupervised (computed within each batch by finding repeated tokens).

\paragraph{Loss 4: Aging pressure.}
The state should not saturate. Old information should fade to make room for new structure:
\begin{equation}
\mathcal{L}_{\mathrm{age}} = \max(0, \|\mathbf{s}_t\| - \kappa) \quad \text{(bound the state norm)}
\end{equation}
combined with a reward for state informativeness:
\begin{equation}
\mathcal{L}_{\mathrm{info}} = -\mathrm{MI}(\mathbf{s}_t; \h_{\mathrm{left}, t:t+k}) \quad \text{(maximise mutual information with future)}
\end{equation}

\emph{Properties satisfied:} graceful aging (bounded norm forces forgetting), infinite context (norm bounded but information maximised), spectral structure (different modes will learn different persistence rates to maximise information under the norm constraint).

\subsection{Gradient Flow: What R Predicts and How to Train It}
\label{sec:gradient-flow}

The critical question: what does R predict, and does gradient descent work?

\paragraph{Naive: R predicts L's future hidden states.}
$\mathcal{L}_R = \|R_{\mathrm{predict}}(\mathbf{s}_t) - \h_{\mathrm{left}, t+k}\|^2$.
Gradient flows to R (through predict and state dynamics) and to L (through $\h_{\mathrm{left}}$).
Problem: L gets pressure to make its states \emph{easy to predict}, which may mean simpler, not more structured.
L might compromise between CE (predict tokens) and $\mathcal{L}_R$ (be predictable) in ways that hurt both.

\paragraph{Better: R predicts L's coupling dynamics.}
\begin{equation}
\mathcal{L}_R = \|R_{\mathrm{predict}}(\mathbf{s}_t) - \mathrm{stopgrad}(\mathrm{coupling\_summary}(C_{t+k}))\|^2
\end{equation}
R predicts what L will ATTEND TO, not what L will represent.
The coupling summary is a learned low-dimensional projection of the attention matrix.
$\mathrm{stopgrad}$ on the target prevents L from collapsing its coupling to match R's prediction, instead, R must learn to predict L's actual coupling dynamics.

Gradient: $\mathcal{L}_R \to R_{\mathrm{predict}} \to \mathbf{s}_t \to$ R's SSM dynamics $\to$ R's parameters. Clean, one-directional. L is trained ONLY by CE. R is trained by $\mathcal{L}_R$. The interaction is through R's output feeding back into L (through the gate), not through shared gradients.

This separation is important: L does not receive gradient from the structural loss.
L's representations change because R's output (fed back through the gate) modifies L's attention at the next layer.
The structural pressure is \emph{architectural} (R's output shapes L's processing), not \emph{gradient-based} (no structural gradient to L's parameters).
This is how the brain works: the right hemisphere influences the left through \emph{connections}, not through a shared loss function.

\paragraph{Deepest: R learns to match coupling distances.}
\begin{equation}
\mathcal{L}_{\mathrm{distance}} = -\mathrm{corr}\!\left(\|\mathbf{s}_i - \mathbf{s}_j\|, \, \|C_i - C_j\|\right) \quad \text{for sampled pairs } (i, j)
\end{equation}
Positions that are relationally similar (similar coupling in L) should be close in R's state space.
Positions that are relationally different should be far apart.
No labels needed, L's own coupling patterns define ``similarity.''

R doesn't predict anything specific.
It learns a GEOMETRY that mirrors L's coupling geometry.
The operator IS this learned geometry, it emerges from the distance-matching requirement.

Gradient: flows to R (through $\mathbf{s}_i, \mathbf{s}_j$). L's coupling $C_i, C_j$ can be either stop-graded (R adapts to L) or live (L and R co-adapt). The stop-grad version is safer; the live version is more powerful but risks instability.

Cost: $\mathcal{O}(S^2)$ for all pairs, but subsampliung to $\mathcal{O}(S)$ with random pair selection works in practice (contrastive learning literature).

\paragraph{The training protocol: decoupled first, then gradually coupled.}

The corpus callosum develops slowly through childhood.
The hemispheres first develop independently, then gradually integrate.
Our training mirrors this.

\begin{enumerate}[nosep]
\item \textbf{Phase 1: L alone.} Train L with CE only. L develops basic representations. R exists but is not trained and has no influence. Duration: until L converges to a stable val loss.

\item \textbf{Phase 2: R as pure observer.} L continues training (or is frozen). R observes L's coupling dynamics and learns to compress them. R's output goes NOWHERE, no feedback to L. R accumulates state and predicts coupling. This is safe: R cannot hurt L. This tells us: what structural information exists in L's processing? How compressible is it?

\item \textbf{Phase 3: R as evaluator (still decoupled).} Measure R's state quality without coupling. Does R's compressed state predict things L's raw attention cannot? At position 100, does R's state correlate with which tokens at positions 1--50 are still attended? If yes: R captured long-range structure that L doesn't explicitly track. R is useful EVEN WITHOUT FEEDBACK.

\item \textbf{Phase 4: Gradual coupling.} R's output feeds back to L through a gate that ramps over training:
\begin{equation}
\h_{\mathrm{out}} = \h + \min\!\left(1, \frac{\mathrm{step} - \mathrm{start}}{\mathrm{ramp}}\right) \cdot \sigma(g) \cdot \h_{\mathrm{right}}
\end{equation}
The coupling increases slowly. L adapts gradually to R's signal. If L's val loss improves during the ramp: coupling is beneficial. If it degrades: reduce ramp rate or freeze coupling.

\item \textbf{Phase 5: Full co-evolution.} Both L and R train simultaneously with full coupling. Monitor for stability: if val loss oscillates, the feedback loop is destabilising. Reduce coupling strength.
\end{enumerate}

\paragraph{Why decoupled R is already valuable.}
A pure observer R (Phase 2--3) that compresses L's coupling dynamics is itself useful:
\begin{itemize}[nosep]
\item It provides a \textbf{diagnostic} of L's relational structure at every position.
\item It provides \textbf{long-range context} for external use (retrieval, steering, interpretability) without changing L's behaviour.
\item It answers: ``what has the model understood relationally?'' ,  an introspection tool, not just an architectural component.
\end{itemize}
The feedback loop (Phase 4--5) is an upgrade, not the core value.
The core is: a learned compression of relational dynamics.

\paragraph{Stability of the feedback loop.}
When R feeds back to L, a coupled dynamical system forms:
R sees L's coupling $\to$ R compresses $\to$ R modifies L's attention $\to$ L's coupling changes $\to$ R sees different coupling $\to \ldots$

Stability depends on:
\begin{itemize}[nosep]
\item \textbf{Coupling strength}: too high $\to$ oscillation; too low $\to$ no effect.
\item \textbf{Timescales}: R should adapt faster than L (R tracks L's changes). Achieved by giving R a higher learning rate.
\item \textbf{The gate}: the learned gate $\sigma(g)$ is a natural stabiliser. If R's signal is harmful, L learns to close the gate. The gate is a safety valve.
\end{itemize}

The biological parallel: the corpus callosum has limited bandwidth.
Hemispheric integration develops slowly.
Children are more lateralised than adults.
Full integration only in adolescence.
Our gradual coupling ramp IS cognitive development.


\subsection{The Structural-Not-Content Loss: Contrastive Learning Without Supervision}
\label{sec:structural-loss}

The central risk: R collapses to learning token embeddings.
``Cat'' and ``dog'' are close in embedding space.
If R's state tracks embedding similarity, it learns content (what things ARE) not structure (how things RELATE).
This is not relational learning, it is a lookup table.

\paragraph{The key insight.}
L's attention patterns encode structure.
L's token embeddings encode content.
R should correlate with the FIRST and decorrelate from the SECOND:
\begin{equation}
\mathcal{L}_{\mathrm{structural}} = -\mathrm{corr}(d_R, d_{\mathrm{coupling}}) + \mathrm{corr}(d_R, d_{\mathrm{embedding}})
\end{equation}
where $d_R(i,j) = \|\mathbf{s}_i - \mathbf{s}_j\|$ is R's state distance, $d_{\mathrm{coupling}}(i,j) = \|C_i - C_j\|$ is L's coupling distance, and $d_{\mathrm{embedding}}(i,j) = \|\h_i - \h_j\|$ is L's embedding distance, for sampled position pairs $(i,j)$.

Maximise correlation with coupling distance (structure).
Minimise correlation with embedding distance (content).
The subtraction is essential: if R learns embeddings, the second term penalises it.
R must find a representation that tracks \emph{coupling} (relational pattern) while being decorrelated from \emph{embeddings} (token identity).

\paragraph{This IS contrastive learning without supervision.}
The positive pairs: positions with similar coupling patterns (structurally similar, automatically detected from L's attention).
The negative pairs: positions with similar embeddings but different coupling (same tokens in different structural roles).
No labels.
No annotated metaphors.
All from L's own internals.

The coupling patterns are the unsupervised structural signal.
Our attention visualisation already showed this works: function words have consistent coupling (structural, they are the positive pairs), names have variable coupling (contextual, they are the informative negatives).

\paragraph{Combined with order-independence.}
Order-independence destroys positional shortcuts.
The structural-not-content loss destroys content shortcuts.
Together: R must learn ABSTRACT RELATIONAL STRUCTURE, patterns of coupling that repeat across different tokens, different positions, different contexts.
This is the relational learning that the operator programme sought.
Not designed into the architecture.
Forced by the loss.


\subsection{Beyond Distance Matching: Spectral and Topological Losses}
\label{sec:spectral-losses}

The distance-matching loss ($\S$\ref{sec:structural-loss}) works: R tracks coupling and decorrelates from content.
But analysis of trained models revealed that R primarily learned \textbf{sequence length}, which is the easiest coupling correlate.
Longer sentences have more spread-out attention distributions; shorter sentences have sharper peaks.
The distance-matching loss can be satisfied by this one feature without learning genuine syntactic structure.

The fundamental issue: \emph{scalar} losses have low-dimensional escape hatches.
R only needs one feature (length) to satisfy a single correlation target.
Genuine structure is high-dimensional, argument roles, clause depth, discourse type, information structure, and requires losses that demand high-dimensional representations.

\paragraph{Insight from the structural probing literature.}
Hewitt \& Manning~\cite{hewitt2019structural} showed that a learned linear transformation of BERT's representations encodes parse tree distances.
The Polar Probe~\cite{chen2024polar} (NeurIPS 2024) extended this: syntactic relations are encoded in both \emph{distance} (magnitude) and \emph{direction} (angle) between embeddings, a polar coordinate system in a low-dimensional subspace.
Our distance-matching loss captures magnitude only.
Structure also lives in direction.

The key insight: SharedM's eigenbasis provides a natural coordinate system for directional structure.
The eigenmodes of $G = B^\top A$ define the ``channels'' through which coupling operates.
This eigenbasis is the structural subspace that probing methods search for post-hoc, but we know it explicitly because we designed SharedM.

\paragraph{Loss variant B: Spectral profile prediction.}
$\M$'s eigenbasis decomposes any coupling pattern into independent modes.
The spectral profile at position $t$ is the energy in each mode:
\begin{equation}
\mathrm{profile}(t) = |V^\top \h_t|^2 \in \R^{n_{\mathrm{modes}}}
\end{equation}
where $V$ contains the top eigenvectors of $G = B^\top A$.
R must predict this profile at position $t + k$ from its state at $t$:
\begin{equation}
\mathcal{L}_{\mathrm{spectral}} = \sum_k \mathrm{SmoothL1}\!\left(\mathrm{MLP}(\mathbf{s}_t),\; \mathrm{stopgrad}\!\left(\mathrm{profile}(t+k)\right)\right)
\end{equation}

Why length cannot satisfy this: the spectral profile depends on the \emph{direction} of $\h$ in $\M$'s eigenbasis, not its magnitude.
Two sentences of the same length but different syntax activate different eigenmodes because the residual stream points in different directions relative to $\M$'s coupling geometry.
With $n_{\mathrm{modes}} = 32$, R must encode 32 independent features of the coupling pattern, length explains at most one.

This connects the ORN programme's deepest finding (stable spectral structure in $\M$) to the LORN programme's learning objective.
R doesn't just track coupling distance; it internalises $\M$'s geometry.

\paragraph{Loss variant C: Coupling topology prediction.}
R must predict \emph{which positions} will be most attended at $t + k$, a multi-label classification:
\begin{equation}
\mathcal{L}_{\mathrm{topo}} = \sum_k \mathrm{BCE}\!\left(\mathrm{MLP}(\mathbf{s}_t),\; \mathrm{stopgrad}\!\left(\mathrm{top\text{-}k}(C_{t+k})\right)\right)
\end{equation}
where $\mathrm{top\text{-}k}(C_{t+k})$ is a binary mask of the $k$ most-attended positions.

This is Hewitt \& Manning's structural probe turned into a training objective rather than a post-hoc analysis.
The grammar induction literature~\cite{kim2024grammar} shows that attention patterns in transformers implicitly encode constituency structure, the top-$k$ attended positions correspond to syntactic dependents.
Predicting them forces R to learn dependency structure.

Length cannot predict specific attention targets: knowing a sentence has 10 tokens says nothing about whether position 3 will attend to position 1.
The loss is genuinely hard (chance accuracy $\sim k/S$, about 6\% at $k = 16, S = 256$).

\paragraph{Loss variant D: Spectral + topology (combined).}
Spectral prediction captures \emph{abstract} structure: what type of coupling occurs (which eigenmodes are active).
Topology prediction captures \emph{concrete} structure: which specific positions are coupled.
These are complementary, operating in orthogonal spaces (eigenmode space vs position space).

Caucheteux \& King~\cite{caucheteux2021syntax} found that syntax and semantics are encoded in orthogonal subspaces in both brains and language models.
Our two prediction targets similarly cannot redundantly capture the same feature, each must carry independent structural information.

\paragraph{The total loss.}
\begin{equation}
\mathcal{L} = \underbrace{\mathcal{L}_{\mathrm{CE}}}_{\text{left: predict}} + \lambda_1 \underbrace{\mathcal{L}_{\mathrm{dist}}}_{\text{distance match}} + \lambda_2 \underbrace{\mathcal{L}_{\mathrm{spectral}}}_{\text{spectral profile}} + \lambda_3 \underbrace{\mathcal{L}_{\mathrm{topo}}}_{\text{topology}} + \lambda_4 \underbrace{(\mathcal{L}_{\mathrm{var}} + \mathcal{L}_{\mathrm{cov}})}_{\text{VICReg}}
\end{equation}

\paragraph{The critical diagnostic: length-controlled probing.}
The test for whether R learned structure is not val loss or coupling correlation, both can be satisfied by length.
The test is: \textbf{does R separate structurally different sentences at matched token length?}
Length-matched pairs (e.g.\ ``the big dog ran to the park'' vs ``the proposed amendment was subsequently rejected,'' both 8 tokens) are fed through the model.
If R's states separate them, R learned structure.
If not, R learned surface statistics regardless of the loss variant.
This diagnostic is applied to all variants.

\paragraph{Results: loss variant comparison.}
All four variants trained on BabyLM ($d = 256$, 8 layers, 6000 steps, 18.5M params).
The length-controlled probe uses sentence pairs matched at 8 tokens.

\begin{center}\small
\begin{tabular}{llrrr}
\toprule
\textbf{Variant} & \textbf{Val} & \textbf{CorrC} & \textbf{Struct sep} & \textbf{Quest sep} \\
\midrule
L-only baseline & 6.12 & ,  & ,  & ,  \\
A: Distance matching & 6.04 & 0.985 & 1.13$\times$ & 0.95$\times$ \\
B: Spectral profile & \textbf{5.96} & 0.992 & 1.32$\times$ & 1.02$\times$ \\
C: Topology & 6.02 & 0.989 & 1.05$\times$ & \textbf{1.40}$\times$ \\
D: Spectral + Topology & 6.02 & 0.989 & \textbf{1.87}$\times$ ★ & 0.90$\times$ \\
\bottomrule
\end{tabular}
\end{center}

\noindent Struct sep $=$ separation between child-simple and formal-complex at matched length.
Quest sep $=$ separation between statements and questions at matched length.
Want $> 1.5\times$.

\textbf{Finding: the spectral loss helps representation, not prediction.}
C and D have nearly identical val (6.016 vs 6.020) and CorrC (0.989 vs 0.989).
The topology loss dominates L's training dynamics in both.
But D's length-controlled structural separation is 1.87$\times$ vs C's 1.05$\times$, the spectral loss reshapes R's internal geometry without changing what R communicates to L.

This is the cleanest possible result: the spectral loss operates in an orthogonal subspace to the topology loss.
Topology handles the R$\to$L channel (which positions to attend).
Spectral organises R's state space (which structural category this is).
They don't compete for the same variance, so val doesn't change, but the internal organisation does.

\textbf{Finding: the losses capture different structural features.}
Spectral prediction separates syntactic complexity (simple vs complex: 1.32$\times$).
Topology prediction separates discourse function (statement vs question: 1.40$\times$).
Combined, structural complexity separation jumps to 1.87$\times$, the only variant to exceed the $1.5\times$ threshold, but question separation drops.
At this small scale, R's 96 dimensions cannot capture both complexity and function simultaneously.
At larger scale ($r = 128$, 12 layers), both should be representable.

\textbf{Finding: the gate pattern reflects the loss.}
\begin{center}\small
\begin{tabular}{lrrr}
\toprule
\textbf{Variant} & \textbf{Layer 7} & \textbf{Layer 8} & \textbf{Layer 9} \\
\midrule
A: Distance & 0.43 & 0.08 & 0.92 \\
B: Spectral & 0.68 & 0.77 & 0.54 \\
C: Topology & 0.72 & 0.93 & 0.01 \\
D: Both & 0.16 & 0.86 & 0.00 \\
\bottomrule
\end{tabular}
\end{center}
Distance matching (A) produces a sharp, one-layer preference.
Spectral (B) produces balanced engagement across all layers.
Topology (C) and combined (D) concentrate on one dominant layer, the topology signal is specific enough that only one depth benefits.


\subsection{Why This Differs from Existing Hybrids}

\begin{center}\small
\begin{tabular}{lll}
\toprule
\textbf{Architecture} & \textbf{Structure} & \textbf{Interaction} \\
\midrule
Jamba / Samba & Interleaved layers (attn, SSM, attn, SSM...) & Sequential, same stream \\
Griffin & Alternating recurrence + local attention & Sequential, same stream \\
Hymba & Parallel heads (attn + SSM in same layer) & Parallel, same stream \\
Zamba & SSM backbone + shared attention injections & Sequential, shared attn params \\
\textbf{Two-Hemisphere} & \textbf{Parallel pathways with structural loss} & \textbf{Parallel, SharedM bottleneck,} \\
 & & \textbf{right sees left's representations} \\
\bottomrule
\end{tabular}
\end{center}

The key differences:
\begin{enumerate}[nosep]
\item The right hemisphere operates on the left's \emph{representations}, not on raw tokens.
\item The two pathways communicate through a shared geometry $\M$ (bandwidth bottleneck).
\item An auxiliary structural loss forces the left to develop representations the right can compress.
\item The hemispheric balance is self-measuring (the gate reports the division of labour).
\end{enumerate}

None of the existing hybrids have the structural loss or the representation-level interaction.
They are engineering combinations of attention + SSM.
The two-hemisphere architecture is a principled combination driven by the need for two kinds of processing.


% ======================================================================
\begin{thebibliography}{99}

\bibitem{mcgilchrist2009master}
I.~McGilchrist,
\emph{The Master and His Emissary: The Divided Brain and the Making of the Western World},
Yale University Press, 2009.

\bibitem{mcgilchrist2021matter}
I.~McGilchrist,
\emph{The Matter with Things: Our Brains, Our Delusions, and the Unmaking of the World},
Perspectiva Press, 2021.

\bibitem{lieber2024jamba}
O.~Lieber et al., ``Jamba: A hybrid transformer-Mamba language model,'' AI21 Labs, 2024.

\bibitem{de2024griffin}
S.~De et al., ``Griffin: Mixing gated linear recurrences with local attention,'' Google DeepMind, 2024.

\bibitem{glorioso2024zamba}
P.~Glorioso et al., ``Zamba: A compact 7B SSM hybrid model,'' Zyphra, 2024.

\bibitem{nvidia2024hymba}
NVIDIA, ``Hymba: A hybrid-head architecture for small language models,'' 2024.

\bibitem{microsoft2024samba}
Microsoft, ``Samba: Simple hybrid state space models for efficient unlimited context,'' 2024.

\bibitem{ringo1994hemispheric}
J.~Ringo et al., ``Time is of the essence: A conjecture that hemispheric specialization arises from interhemispheric conduction delay,'' \emph{Cerebral Cortex}, 1994.

\bibitem{ganin2016dann}
Y.~Ganin et al., ``Domain-adversarial training of neural networks,'' \emph{JMLR}, 2016.

\bibitem{lample2017fader}
G.~Lample et al., ``Fader networks: Manipulating images by sliding attributes,'' \emph{NeurIPS}, 2017.

\bibitem{kim2018factorizing}
H.~Kim and A.~Mnih, ``Disentangling by factorising,'' \emph{ICML}, 2018.

\bibitem{bardes2022vicreg}
A.~Bardes, J.~Ponce, Y.~LeCun, ``VICReg: Variance-invariance-covariance regularization,'' \emph{ICLR}, 2022.

\bibitem{goyal2022global}
A.~Goyal et al., ``Coordination among neural modules through a shared global workspace,'' \emph{ICLR}, 2022.

\bibitem{feichtenhofer2019slowfast}
C.~Feichtenhofer et al., ``SlowFast networks for video recognition,'' \emph{ICCV}, 2019.

\bibitem{dao2024ssd}
T.~Dao and A.~Gu, ``Transformers are SSMs: Generalized models and efficient algorithms through structured state space duality,'' \emph{ICML}, 2024.

\bibitem{arora2024based}
S.~Arora et al., ``Simple linear attention language models balance the recall-throughput tradeoff,'' \emph{ICML}, 2024.

\bibitem{lan2020albert}
Z.~Lan et al., ``ALBERT: A lite BERT for self-supervised learning of language representations,'' \emph{ICLR}, 2020.

\bibitem{hewitt2019structural}
J.~Hewitt and C.~D.~Manning, ``A structural probe for finding syntax in word representations,'' \emph{NAACL}, 2019.

\bibitem{chen2024polar}
Z.~Chen et al., ``A polar coordinate system represents syntax in large language models,'' \emph{NeurIPS}, 2024.

\bibitem{caucheteux2021syntax}
C.~Caucheteux and J.-R.~King, ``Disentangling syntax and semantics in the brain with deep networks,'' \emph{ICML}, 2021.

\bibitem{kim2024grammar}
Y.~Kim et al., ``Leveraging grammar induction for language understanding and generation,'' \emph{arXiv:2410.04878}, 2024.

% ,  Part II references: gauge geometry, Fisher information, SSL , 

\bibitem{amari1998natural}
S.~Amari, ``Natural gradient works efficiently in learning,'' \emph{Neural Computation}, 10(2):251--276, 1998.

\bibitem{martens2015kfac}
J.~Martens and R.~Grosse, ``Optimizing neural networks with Kronecker-factored approximate curvature,'' \emph{ICML}, 2015.

\bibitem{sagun2016eigenvalues}
L.~Sagun, L.~Bottou, Y.~LeCun, ``Eigenvalues of the Hessian in deep learning: singularity and beyond,'' \emph{arXiv:1611.07476}, 2016.

\bibitem{papyan2019measurements}
V.~Papyan, ``Measurements of three-level hierarchical structure in the outliers in the spectrum of deepnet Hessians,'' \emph{ICML}, 2019.

\bibitem{ghorbani2019investigation}
B.~Ghorbani, S.~Krishnan, Y.~Xiao, ``An investigation into neural net optimization via Hessian eigenvalue density,'' \emph{ICML}, 2019.

\bibitem{garipov2018loss}
T.~Garipov, P.~Izmailov, D.~Podoprikhin, D.~Vetrov, A.~G.~Wilson, ``Loss surfaces, mode connectivity, and fast ensembling of DNNs,'' \emph{NeurIPS}, 2018.

\bibitem{draxler2018essentially}
F.~Draxler, K.~Veschgini, M.~Salmhofer, F.~Hamprecht, ``Essentially no barriers in neural network energy landscape,'' \emph{ICML}, 2018.

\bibitem{frankle2020linear}
J.~Frankle, G.~K.~Dziugaite, D.~M.~Roy, M.~Carbin, ``Linear mode connectivity and the lottery ticket hypothesis,'' \emph{ICML}, 2020.

\bibitem{bronstein2021geometric}
M.~M.~Bronstein, J.~Bruna, T.~Cohen, P.~Veli\v{c}kovi\'{c}, ``Geometric deep learning: grids, groups, graphs, geodesics, and gauges,'' \emph{arXiv:2104.13478}, 2021.

\bibitem{cohen2016group}
T.~S.~Cohen and M.~Welling, ``Group equivariant convolutional networks,'' \emph{ICML}, 2016.

\bibitem{hu2021lora}
E.~J.~Hu, Y.~Shen, P.~Wallis, et al., ``LoRA: low-rank adaptation of large language models,'' \emph{arXiv:2106.09685}, 2021.

\bibitem{zbontar2021barlow}
J.~Zbontar, L.~Jing, I.~Misra, Y.~LeCun, S.~Deny, ``Barlow Twins: self-supervised learning via redundancy reduction,'' \emph{ICML}, 2021.

\bibitem{chen2020simclr}
T.~Chen, S.~Kornblith, M.~Norouzi, G.~Hinton, ``A simple framework for contrastive learning of visual representations,'' \emph{ICML}, 2020.

\bibitem{grill2020byol}
J.-B.~Grill, F.~Strub, F.~Altch\'{e}, et al., ``Bootstrap your own latent: a new approach to self-supervised learning,'' \emph{NeurIPS}, 2020.

\bibitem{assran2023ijepa}
M.~Assran, Q.~Duval, I.~Misra, et al., ``Self-supervised learning from images with a joint-embedding predictive architecture,'' \emph{CVPR}, 2023.

\bibitem{jacot2018ntk}
A.~Jacot, F.~Gabriel, C.~Hongler, ``Neural tangent kernel: convergence and generalization in neural networks,'' \emph{NeurIPS}, 2018.

\bibitem{tishby2015deep}
N.~Tishby and N.~Zaslavsky, ``Deep learning and the information bottleneck principle,'' \emph{IEEE ITW}, 2015.

\bibitem{foret2021sam}
P.~Foret, A.~Kleiner, H.~Mobahi, B.~Neyshabur, ``Sharpness-aware minimization for efficiently improving generalization,'' \emph{ICLR}, 2021.

\bibitem{keskar2017large}
N.~S.~Keskar, D.~Mudigere, J.~Nocedal, M.~Smelyanskiy, P.~T.~P.~Tang, ``On large-batch training for deep learning: generalization gap and sharp minima,'' \emph{ICLR}, 2017.

\bibitem{hochreiter1997flat}
S.~Hochreiter and J.~Schmidhuber, ``Flat minima,'' \emph{Neural Computation}, 9(1):1--42, 1997.

\bibitem{soudry2018implicit}
D.~Soudry, E.~Hoffer, M.~S.~Nacson, S.~Gunasekar, N.~Srebro, ``The implicit bias of gradient descent on separable data,'' \emph{JMLR}, 2018.

\bibitem{ainsworth2023git}
S.~K.~Ainsworth, J.~Hayase, S.~Srinivasa, ``Git re-basin: merging models modulo permutation symmetries,'' \emph{ICLR}, 2023.

\bibitem{cohen2019general}
T.~S.~Cohen, M.~Geiger, M.~Weiler, ``A general theory of equivariant CNNs on homogeneous spaces,'' \emph{NeurIPS}, 2019.

% ,  Biological inspiration: complementary learning systems and consolidation , 

\bibitem{mcclelland1995cls}
J.~L.~McClelland, B.~L.~McNaughton, R.~C.~O'Reilly, ``Why there are complementary learning systems in the hippocampus and neocortex: insights from the successes and failures of connectionist models of learning and memory,'' \emph{Psychological Review}, 102(3):419--457, 1995.

\bibitem{kumaran2016cls}
D.~Kumaran, D.~Hassabis, J.~L.~McClelland, ``What learning systems do intelligent agents need? Complementary learning systems theory updated,'' \emph{Trends in Cognitive Sciences}, 20(7):512--534, 2016.

\bibitem{squire1992memory}
L.~R.~Squire, ``Memory and the hippocampus: a synthesis from findings with rats, monkeys, and humans,'' \emph{Psychological Review}, 99(2):195--231, 1992.

\bibitem{buzsaki2015swr}
G.~Buzs\'{a}ki, ``Hippocampal sharp wave-ripple: a cognitive biomarker for episodic memory and planning,'' \emph{Hippocampus}, 25(10):1073--1188, 2015.

\bibitem{diekelmann2010sleep}
S.~Diekelmann, J.~Born, ``The memory function of sleep,'' \emph{Nature Reviews Neuroscience}, 11(2):114--126, 2010.

\bibitem{nadel1997trace}
L.~Nadel, M.~Moscovitch, ``Memory consolidation, retrograde amnesia and the hippocampal complex,'' \emph{Current Opinion in Neurobiology}, 7(2):217--227, 1997.

\bibitem{hinton1995wakesleep}
G.~E.~Hinton, P.~Dayan, B.~J.~Frey, R.~M.~Neal, ``The wake-sleep algorithm for unsupervised neural networks,'' \emph{Science}, 268(5214):1158--1161, 1995.

\bibitem{wilson1994reactivation}
M.~A.~Wilson, B.~L.~McNaughton, ``Reactivation of hippocampal ensemble memories during sleep,'' \emph{Science}, 265(5172):676--679, 1994.

\bibitem{french1999catastrophic}
R.~M.~French, ``Catastrophic forgetting in connectionist networks,'' \emph{Trends in Cognitive Sciences}, 3(4):128--135, 1999.

\end{thebibliography}


% ======================================================================
\appendix

\section{LORN and Downstream Training}
\label{app:downstream}

LORN's two-hemisphere separation creates a natural factorisation for downstream adaptation (fine-tuning, mixture of experts, instruction following).
The key property: R is a separate, composable module that holds task-invariant structural knowledge, while L holds task-specific prediction.
This factorisation does not exist in standard transformers, where fine-tuning changes everything simultaneously.

\paragraph{Fine-tune L, freeze R (task adaptation).}
R holds relational structure that is task-invariant, how clauses nest, how arguments relate, how discourse flows.
When fine-tuning for instruction following or domain adaptation, L changes its predictions (what to say) without changing R's structural understanding (how language works).
R provides a stable structural scaffold that prevents the representational drift common in fine-tuning.
This is analogous to how the right hemisphere's holistic understanding persists across skill acquisition, the left hemisphere adapts, the right provides context.

\paragraph{Fine-tune R, freeze L (structural adaptation).}
When adapting to a new \emph{register}, legal text, code, medical language, the coupling patterns are different but the token prediction mechanism is similar.
Fine-tuning R alone is much cheaper: R's parameters are a small fraction of the total (164K of 18.5M in our BabyLM experiments, $<$1\%).
The new R learns the structural patterns of the target domain; L's prediction machinery transfers unchanged.

\paragraph{Mixture of experts with shared R.}
Multiple expert L networks share one R.
Each expert specialises in a domain (code, dialogue, formal text), but all share the same structural understanding.
R's state becomes the natural \textbf{routing signal}: our analysis showed R separates registers (child-directed vs formal Wikipedia: 9.2 Euclidean distance in R's state space).
This separation provides a ready-made routing criterion, R knows the structural type of the current context before L starts predicting.

The architecture becomes:
\begin{equation}
\text{output} = \sum_i g_i(\mathbf{s}_R) \cdot L_i(\h, \text{coupling bias from } R)
\end{equation}
where $g_i$ routes based on R's state and each $L_i$ is a domain expert that receives the same structural context.
This is the orbital MoE vision (our companion paper) grounded in the two-hemisphere framework.

\paragraph{Instruction tuning with structural constraint.}
During RLHF or DPO, the reward model pushes L toward preferred outputs.
Without structural constraint: this can produce verbose, formulaic text (the ``helpful but hollow'' failure mode where the model satisfies the reward while losing structural coherence).
With R: the structural loss is independent of the reward signal.
R's pressure maintains coherent relational structure even as L's predictions shift toward helpfulness.
The two losses (reward + structural) create a healthy tension: the model must be helpful AND structurally coherent.

\paragraph{R as a test-time steering interface.}
After training and downstream adaptation, R's state space becomes a control surface.
Instead of prompt engineering (changing L's input tokens), one manipulates R's state (changing L's structural context):
\begin{itemize}[nosep]
\item ``Write formally'' $=$ shift R toward the formal register centroid.
\item ``Be concise'' $=$ shift R toward the child-directed centroid.
\item ``Reason step by step'' $=$ shift R toward the complex-clause centroid.
\end{itemize}
This is cheaper (96 dimensions vs thousands of tokens) and more precise (continuous, not discrete) than prompting.
It is also \emph{inspectable}: R's state is a low-dimensional representation whose dimensions have interpretable correlates (our analysis found dimensions correlated with formality, question-ness, and clause complexity).

\paragraph{The prediction.}
LORN's R becomes more valuable, not less, as downstream training is added.
It is the structural backbone that every downstream adaptation inherits.
The ORN savings (shared coupling geometry freeing parameters) provide the budget; the two-hemisphere programme provides the investment; downstream training provides the return.


\section{R as a Dynamics Observer: Modelling the Process, Not the State}
\label{app:dynamics}

The previous sections describe R as observing L's coupling patterns, static snapshots of attention at specific layers.
A deeper formulation: R observes the \textbf{transitions} between L's layers.
Not what L represents at layer $m$, but what L \emph{does} going from layer $m{-}1$ to layer $m$.

\subsection{The Delta Signal}

At each L layer, the transition is:
\begin{equation}
\Delta \h_m = \h_m - \h_{m-1}, \qquad \Delta C_m = C_m - C_{m-1}
\end{equation}
$\Delta \h_m$ captures what the layer \emph{changed} about the representation.
$\Delta C_m$ captures what the layer changed about the coupling.
Together, they describe the \emph{computation} at layer $m$: which positions were reorganised, which relationships were discovered, what new structure emerged at this depth.

This connects directly to two findings from the operator programme:
\begin{itemize}[nosep]
\item \textbf{Depth-criticality:} shallow layers accumulate (small $\|\Delta\h_m\|$, corrections 4--8\%), deep layers reorganise (large $\|\Delta\h_m\|$, corrections 34--50\%). The delta IS the criticality signal.
\item \textbf{The sandpile:} each layer adds height; when height exceeds threshold, toppling occurs. The delta captures the toppling events, where and when reorganisation happens.
\end{itemize}

\subsection{Architecture: 12 Observation Points, K Processing Layers}

If R observes every L layer's transition, it needs one observation point per L layer.
But the number of observation points (12) is independent of R's processing depth ($K$).
Each observation point is lightweight, a small encoder that maps the delta into R's state space:
\begin{equation}
\mathbf{o}_m = \mathrm{encode}(\Delta\h_m, \Delta C_m)
\end{equation}

These 12 observations feed into $K$ processing layers where R builds hierarchical understanding of L's dynamics.
$K$ can be much less than 12, the observations are cheap, the processing is where R's capacity lives.

\begin{center}\small
\begin{tabular}{lrl}
\toprule
\textbf{Component} & \textbf{Count} & \textbf{Role} \\
\midrule
Observation encoders & 12 (one per L layer) & Cheap: encode delta into R space \\
Processing layers & $K$ (3--6) & Expensive: structural self-attention + FFN \\
\bottomrule
\end{tabular}
\end{center}

\subsection{What R Represents After Seeing All Deltas}

L's hidden state at layer 11 represents \emph{what the model currently believes about the next token}.
R's accumulated state after 12 deltas represents \emph{how the model arrived at that belief}, the full processing trajectory through depth.

These are fundamentally different.
Two sentences can arrive at similar L outputs (both predict ``the'' next) through completely different processing paths (one via SVO composition, the other via existential construction).
L's final state looks similar.
R's accumulated deltas look different, because the \emph{path through depth} was different.

This is the deepest version of the structural signal.
Not coupling at one layer.
Not coupling topology.
The \emph{trajectory through depth}.
A linear probe on L's hidden state at any single layer fundamentally cannot capture this, it requires the full depth history, which only R accumulates.

\subsection{The Dynamic Structural Loss}

The standard structural-not-content loss compares R's state distances to coupling distances at one layer.
The dynamic version compares R's state distances to \emph{trajectory} distances:
\begin{equation}
\mathcal{L}_{\mathrm{dynamic}} = -\mathrm{corr}\!\left(d_R(i,j), \; d_{\mathrm{traj}}(i,j)\right) + \lambda \cdot \mathrm{corr}\!\left(d_R(i,j), \; d_{\mathrm{output}}(i,j)\right)
\end{equation}
where $d_{\mathrm{traj}}(i,j) = \sum_m \|\Delta\h_m^{(i)} - \Delta\h_m^{(j)}\|$ sums the delta differences across all layers, and $d_{\mathrm{output}}(i,j) = \|\h_{L}^{(i)} - \h_{L}^{(j)}\|$ is L's final output distance.

Maximise correlation with trajectory distance (process).
Minimise correlation with output distance (product).
R should track \emph{how} L processed, not \emph{what} L concluded.

Two sentences with different processing paths but similar outputs: R should separate them.
Two sentences with similar processing paths but different outputs: R should group them.

This IS the structural-not-content principle, reformulated for dynamics: R tracks \emph{how}, not \emph{what}.

\subsection{Prediction Targets for the Delta Observer}

\paragraph{Predict the next delta.}
Given deltas at layers $0 \to k$, predict the delta at layer $k{+}1$:
\begin{equation}
\hat{\Delta}\h_{k+1} = \mathrm{MLP}(\mathbf{s}_k^R), \qquad \mathcal{L}_{\mathrm{next}} = \|\hat{\Delta}\h_{k+1} - \mathrm{stopgrad}(\Delta\h_{k+1})\|^2
\end{equation}
If R can predict the next delta, R understands L's processing dynamics, which positions are about to be reorganised, where the next structural event will occur.

\paragraph{Predict where reorganisation occurs.}
Instead of predicting the full delta, predict which positions will have large $\|\Delta\h_{k+1}\|$, a binary classification per position:
\begin{equation}
\mathcal{L}_{\mathrm{reorg}} = \mathrm{BCE}\!\left(\mathrm{MLP}(\mathbf{s}_k^R), \; \mathbf{1}[\|\Delta\h_{k+1}\| > \tau]\right)
\end{equation}
This is the structural analogue of the topology prediction, but applied to L's depth dynamics rather than L's attention patterns.

\paragraph{Predict the final coupling from early deltas.}
R sees layers 0--6 and must predict L's coupling at layer 11.
This forces R to learn L's full processing trajectory, extrapolating from early dynamics to the final state.
Length cannot help.
Position effects change across depths.
R must understand the compositional process to predict where it leads.

\section{What the ORN Coupling Graph Encodes}
\label{app:coupling-structure}

Probing the trained 605M ORN's coupling graphs reveals that genuine linguistic structure is present, but in the \emph{relational pattern}, not in per-position summaries.

\paragraph{Role consistency.}
Subjects couple like subjects and objects couple like objects.
For sentences of the form ``the X verbed the Y,'' the coupling row of the subject position (X) has cosine similarity 0.999 with other subjects and only 0.503 with objects at layer 0.
At layer 31: within-role similarity 0.978, cross-role 0.720.
The coupling graph encodes argument roles.

\paragraph{Clause boundaries.}
At layers 16 and 31, coupling drops more at real clause boundaries (``bad $|$ they'' in ``although the weather was bad they went outside'') than at arbitrary midpoints in simple sentences.
The ORN represents clause structure in its coupling graph.

\paragraph{Word order sensitivity.}
``The cat chased the mouse'' and ``the mouse chased the cat'' produce different coupling graphs (distance 0.046 at layer 16), despite identical tokens.
``The dog bit the man'' vs ``the man bit the dog'': distance 0.131 at L16.
The coupling graph reflects who-does-what-to-whom.

\paragraph{Graph-level structural clustering.}
SVO, existential, and subordinate sentences produce coupling graphs that cluster by syntactic type.
At layer 16: within-group graph distance 0.062, between-group 0.112, a separation of 1.80$\times$.
Using the full coupling matrix (not mean-pooled summaries) recovers structural information that per-position features miss.

\paragraph{Implications for R.}
Per-position feature extraction (mean-pooling, diagonal, entropy) reduces the coupling graph to a scalar and discards relational information.
A linear probe on these summaries outperforms deep R because the summaries don't require depth.
R must be a \textbf{graph modeller}, processing the full coupling matrix as a relational object.
Cross-position attention in R is justified when the task is relational (dependency detection: 94.8\% improvement over linear), not when it is per-position (structural classification: linear wins).

The next R architecture should encode the coupling \emph{graph}, not a \emph{summary} of it.
Cross-position attention in R is essential because relational structure is inherently cross-positional.

\subsection{R as a Structural Language Model}
\label{app:structural-lm}

The findings above converge on a specific architecture: \textbf{R as a language model over L's structural dynamics}.

Just as L predicts the next token from preceding tokens, R predicts L's next-layer coupling from preceding layers' coupling.
The prediction task is R's training signal, not R's purpose.
R's value is the representation it develops in the process, a representation that organises positions by structural role, even when a linear predictor achieves lower loss.

This was demonstrated empirically: a structural LM trained on frozen 605M ORN activations developed internal states where subjects cluster separately from objects (within-role distance 0.90, cross-role distance 2.26), despite losing to linear prediction on MSE.
The representation is the product; the prediction is the training signal.

\paragraph{Architecture.}
\begin{enumerate}[nosep]
\item L processes tokens through 12 layers (SharedM coupling, per-layer response, per-layer FFN).
\item At layers 3, 6, and 9: R observes L's coupling and hidden states.
\item Each R observation is encoded and combined with R's running state via a gated update.
\item R applies cross-position self-attention after each observation (non-causal, 3 layers total).
\item R's prediction target: L's coupling at the NEXT observation depth.
\item After R's final observation: R's state is injected into L's deep layers (10, 11) as $d$-dimensional structural embeddings BEFORE attention.
\item L and R co-evolve: L is trained by CE loss, R by next-layer coupling prediction. R's structural embeddings influence L's attention, which changes the coupling R observes next step.
\end{enumerate}

\paragraph{Why next-layer prediction?}
\begin{itemize}[nosep]
\item It is self-supervised: no labels, no structural annotation, no contrastive pairs.
\item It always has gradient: every training step produces a prediction error.
\item It forces R to model coupling DYNAMICS, not static features, how attention evolves through depth, not what attention looks like at one layer.
\item The representation R develops as a byproduct captures relational structure that static losses (distance matching, topology prediction) do not force.
\item A linear predictor can achieve lower MSE, but R's cross-position attention develops richer internal representations, because it can reason about relationships between positions while making predictions.
\end{itemize}

\paragraph{What co-evolution adds.}
On frozen L, R's state is dominated by positional features (strongest correlation $r = 0.83$ with relative position).
This is because frozen L's representations encode position strongly and structure weakly, L was trained to predict tokens, not to be structurally transparent.
Co-evolution changes this: R's structural embeddings modify L's attention, which changes L's coupling, which changes what R observes.
Over training, L develops representations that are more structurally transparent because R's feedback rewards structural organisation.
The structural pressure is architectural (through the embeddings), not gradient-based (R's loss does not backpropagate to L).


\section{Lessons from Three Failed Programmes and What They Teach R}
\label{app:lessons}

The operator programme, the diffusion experiments, and the SSM-on-coupling experiments each failed for specific reasons.
But each contains an insight about how R should learn.
Taken together, they suggest a training framework for R that combines structured state dynamics with multi-scale denoising.

\subsection{The Operator Programme's Lesson}

The operator programme~(companion paper, \emph{Taking Orbits to the Limit}) tried to compute in operator space: polynomial of $\M$, spectral rotation, semigroup composition $T(s) \circ T(t) = T(s+t)$.
Every attempt collapsed, polynomials reduced to degree 0--1, multi-hop attention hurt, recovery from coupling alone plateaued at $\varepsilon = 0.68$.

The root cause: $\M$'s eigenmodes do not decompose neatly into compositional units.
The programme assumed ``mode 5 encodes noun phrases'' and ``composing modes produces structures.''
In reality, modes are distributed, every mode participates in every structure.
There are no interpretable templates to compose.

\textbf{Lesson for R:} do not design composition from $\M$'s spectral structure.
Let R \emph{learn} the composition in its own state space.
R does not need modes to be interpretable, it needs a representation that \emph{works} for modelling L's dynamics.
The composition happens in R's learned state, not in $\M$'s eigenspace.

\subsection{Diffusion's Insight: Learn to Reverse a Known Corruption}

The diffusion experiments tried to use $\M$'s spectral modes as a structured denoising process: high-eigenvalue modes denoise first (coarse structure), low-eigenvalue modes denoise later (fine detail).
This failed for the same reason as the operator programme: modes don't correspond to structural levels of detail.

But the core insight of diffusion is powerful: \textbf{you don't learn arbitrary dynamics, you learn to reverse a specific, controlled degradation.}
The training signal is always ``here is the corrupted version; reconstruct the clean version.''
The model never has to discover what the corruption IS, that is given.
It only has to learn how to reverse it.

Applied to R: instead of ``predict L's next-layer coupling'' (where the change is driven by complex, tangled factors), define a \emph{structured corruption} of L's coupling and train R to reconstruct the clean version.
Candidate corruptions:
\begin{itemize}[nosep]
\item \textbf{Mask random coupling entries.} Like masked language modelling but for attention graphs. R must predict which positions should attend to which, forcing it to learn dependency structure.
\item \textbf{Shuffle coupling rows.} Destroy structural role assignments. R must restore which rows belong where, forcing it to learn positional role structure.
\item \textbf{Add spectral noise.} Corrupt along $\M$'s eigenmodes. R must learn which eigenmodes carry structure and which carry noise, forcing spectral awareness.
\end{itemize}

Each corruption type forces R to learn a different aspect of structure.
The corruption level varies (10\% masked is easy; 90\% is hard), giving R a multi-scale training signal.
This is diffusion's denoising schedule applied to structural learning rather than image generation.

Crucially, R cannot solve these tasks by memorising position:
shuffled rows destroy positional shortcuts.
Masked entries require relational reasoning (which positions \emph{should} attend to which).
Spectral noise requires understanding $\M$'s geometry.

\subsection{SSM's Insight: Structured State Transitions}

The SSM-on-coupling experiments used a complex diagonal state transition to model how coupling evolves across the token sequence.
The spectral rotation variant worked (+0.15 nats over scalar decay), and the eigenvalues developed diverse magnitudes and phases, structured aging with different modes persisting at different rates.

The insight for R: \textbf{state should evolve through a parameterised linear operator, not an arbitrary neural combination.}
Current R uses an unstructured update: $\mathbf{s}_{k+1} = \mathrm{combine}(\mathrm{obs}_k, \mathbf{s}_k)$ followed by self-attention.
Maximum flexibility, minimum inductive bias about how structural state should evolve through depth.

SSM-style state dynamics:
\begin{equation}
\mathbf{s}_{k+1} = A_R \odot \mathbf{s}_k + B_R \cdot \mathrm{encode}(\mathbf{h}_k, C_k)
\end{equation}
where $A_R \in \mathbb{C}^r$ is a learned diagonal complex vector.
Each dimension of R's state evolves independently with its own persistence ($|A_{R,i}|$) and rotation ($\angle A_{R,i}$).
Some structural features are stable across all L layers (argument roles: $|A| \approx 1$).
Others are transient (local attention patterns: $|A| \ll 1$).
The complex rotation allows structural features to \emph{transform} through depth, not just persist or decay.

This provides inductive bias: R's state evolution is structured like the SSMs that successfully model long sequences.
The self-attention then operates on this structured state, adding cross-position reasoning on top of structured temporal dynamics.

\subsection{The Synthesis: Structured Denoising LM}

Combining all three insights:

\begin{enumerate}[nosep]
\item \textbf{State dynamics from SSMs.} R's state evolves through L's depth via a learned diagonal complex transition. Different structural features persist, decay, and rotate at different rates.

\item \textbf{Training signal from diffusion.} Instead of next-layer prediction (too easy for linear), R learns to reverse structured corruptions of L's coupling. Multiple corruption types target different structural features. Corruption severity varies, providing multi-scale signal.

\item \textbf{Composition from learning (operator programme lesson).} R does not compose eigenmodes, it learns composition in its own state space through cross-position attention. The structured state dynamics provide the scaffold; the attention provides the relational reasoning.
\end{enumerate}

The result is a training framework where:
\begin{itemize}[nosep]
\item The SSM dynamics give R the right inductive bias for tracking structural features across depth.
\item The denoising objective gives R a training signal that requires genuine structural understanding (not satisfiable by linear prediction or positional shortcuts).
\item The cross-position attention gives R the relational reasoning that dependency detection and graph reconstruction require.
\end{itemize}

Each component is motivated by a specific experimental failure:
SSM dynamics address the arbitrary state evolution problem.
Denoising addresses the trivially-satisfiable loss problem.
Cross-position attention addresses the relational reasoning requirement.

\subsection{Status}

This synthesis is speculative, not yet implemented.
The current experiments test the structural LM variant (next-layer prediction with cross-position attention and unstructured state).
If the structural LM's state develops genuine structural representations with enough training (overnight experiment, 100K steps on frozen 605M ORN), the next step is adding SSM dynamics and denoising objectives.
If the structural LM plateaus at positional encoding, the denoising objective becomes more urgent, it directly addresses the positional shortcut problem that current losses cannot prevent.


\section{Design Matrix: Configuration Space for LORN}
\label{app:design-matrix}

The LORN architecture has several independent design axes.
This matrix tracks the permutations tested and their outcomes, serving as a reference for future configurations.

\subsection{Loss Functions}

\begin{center}\small
\begin{tabular}{llrl}
\toprule
\textbf{Loss} & \textbf{What it forces} & \textbf{Best result} & \textbf{Status} \\
\midrule
Distance matching & R tracks coupling distance & 0.98 corr & Works, learns length \\
Embedding decorrelation & R avoids content & $-$0.93 corr & Works with dist. match \\
Spectral profile prediction & R tracks eigenmode energy & 1.87$\times$ struct & Helps representation \\
Topology prediction & R predicts top-$k$ targets & 1.40$\times$ func & Captures function \\
Next-layer coupling prediction & R models depth dynamics & 1.68$\times$ trans. & Byproduct structure \\
Structural invariance & Same structure $\to$ same prediction & 1.68$\times$ trans. & Promising, needs scale \\
Shuffled contrastive & Distinguish orig vs shuffled & Collapsed to $-$1 & Failed (trivial) \\
VICReg & Prevent dimensional collapse & Necessary & Always include \\
\bottomrule
\end{tabular}
\end{center}

\subsection{R Architecture}

\begin{center}\small
\begin{tabular}{lrrl}
\toprule
\textbf{R design} & \textbf{R share} & \textbf{Struct sep} & \textbf{Status} \\
\midrule
Linear probe (1 layer, no attention) & 7\% & 1.70$\times$ & Wins on per-position tasks \\
Shallow R (1 enc + accumulate) & 9\% & 1.31$\times$ & Baseline LORN \\
Deep R (6 layers, shared FFN) & 31\% & 1.47$\times$ & Slow, collapsed text \\
Hierarchical R (3 obs depths) & 19\% & pending & Fast, gates open \\
Structural LM R (predict next coupling) & 9\% & pending & Role separation ★ \\
Structural LM v2 (3 separate layers) & 26\% & pending & Ready to test \\
\bottomrule
\end{tabular}
\end{center}

\subsection{R Input (What R Observes)}

\begin{center}\small
\begin{tabular}{lrl}
\toprule
\textbf{Signal} & \textbf{Struct sep} & \textbf{Notes} \\
\midrule
$\h + C$ (hidden + post-softmax coupling) & 1.70$\times$ & Default, simple \\
$\h + \mathbf{q} + \mathbf{k}$ (hidden + queries + keys) & 0.94$\times$ & Sees M directly, worse \\
$\h + \mathbf{q} + \mathbf{k} + C$ (all) & 1.26$\times$ & Moderate \\
Layer deltas ($\Delta\h$, $\Delta C$) & 1.09$\times$ struct, 3.61$\times$ func & Good for function \\
$C$ diagonal (self-attention) & 5.07$\times$ & Strongest single signal \\
$C$ at L31 (deep coupling) & 3.06$\times$ & Better than mid layers \\
FFN gate pattern & 1.30$\times$ & World model response \\
$\h$ at L16 (what we kept using) & 0.87$\times$ & Worst layer for structure \\
\bottomrule
\end{tabular}
\end{center}

\subsection{R $\to$ L Injection}

\begin{center}\small
\begin{tabular}{lrrl}
\toprule
\textbf{Injection mode} & \textbf{Val $\Delta$} & \textbf{Gate} & \textbf{Status} \\
\midrule
Additive residual & $-$0.20 & 0.16 & Gate closes, ignorable \\
Coupling bias (rank-16) & $-$0.33 & 0.86 & Gate opens, but undifferentiated \\
Value modulation & $-$0.18 & 0.34 & Good structure, moderate gate \\
Cross-attention & $-$0.03 & 0.90 & High gate, no val improvement \\
Per-layer gated coupling & $-$0.31 & 0.83 & Sharp depth profile \\
Structural embeddings ($d$-dim) & pending & 0.99 & Current default \\
\bottomrule
\end{tabular}
\end{center}

\subsection{Observation Pattern}

\begin{center}\small
\begin{tabular}{lrl}
\toprule
\textbf{Where R observes L} & \textbf{Key finding} & \textbf{Status} \\
\midrule
Single midpoint (L layer 6) & Baseline & All early experiments \\
Multi-depth (L layers 3, 6, 9) & 3.5$\times$ func sep & Hierarchical LORN \\
All saved layers (0,4,8,16,24,31) & 3.51$\times$ func & Grid search (frozen) \\
Per-layer with deltas & 3.61$\times$ func & Best for function \\
\bottomrule
\end{tabular}
\end{center}

\subsection{State Dynamics}

\begin{center}\small
\begin{tabular}{ll}
\toprule
\textbf{State update} & \textbf{Status} \\
\midrule
Running mean (order-independent) & Default in all experiments \\
Gated combination + attention & Structural LM, current \\
SSM diagonal complex transition & Proposed, not yet tested \\
\bottomrule
\end{tabular}
\end{center}

\subsection{Training Regime}

\begin{center}\small
\begin{tabular}{ll}
\toprule
\textbf{Regime} & \textbf{Status} \\
\midrule
R on frozen L (cached activations) & Tested on 605M. Limited by frozen representations. \\
L and R co-evolving & Tested on MPS (small scale). Gates open, val improves. \\
Phased (L first, then R, then co-evolve) & Tested in first experiment. Gate opens to 81\%. \\
\bottomrule
\end{tabular}
\end{center}

\paragraph{Configuration for the current best experiment.}
\begin{itemize}[nosep]
\item Loss: next-layer coupling prediction + VICReg
\item R: 3 separate cross-position attention layers, $r = 384$ (26\% of params)
\item Input: $\h + C$ at layers 3, 6, 9
\item Injection: structural embeddings at layers 10, 11
\item State: gated combination (not SSM)
\item Regime: co-evolving L and R
\end{itemize}


\section{The Structural Prefix: R Communicates as Context}
\label{app:structural-prefix}

\subsection{Two Bootstrapping Principles}

\paragraph{Principle 1: L must survive alone.}
R will take time to learn, initially R's output is random noise.
If L depends on R, L cannot train during this period.
The architecture must guarantee that L works exactly as a standard language model when R contributes nothing.
R is a bonus, never a requirement.

\paragraph{Principle 2: R's output should ``look like'' something L already knows how to use.}
L already knows how to condition on context, it is a language model.
If R's structural state enters L through a mechanism L already understands (token-like representations in the residual stream), L does not need to learn a new interface.
L's existing attention mechanism becomes the gate: L attends to R's signal if it is useful, ignores it if it is noise.
No separate gate parameter is needed.
No new injection mechanism is needed.

\subsection{The Architecture: Structural Prefix Tokens}

R's state is projected into $K$ ``structural prefix tokens'', $d$-dimensional vectors that are prepended to L's sequence:
\begin{equation}
\mathrm{input} = [\underbrace{\mathbf{p}_1, \ldots, \mathbf{p}_K}_{\text{R's structural prefix}},\; \underbrace{\mathbf{e}_1, \ldots, \mathbf{e}_S}_{\text{real token embeddings}}]
\end{equation}
where $\mathbf{p}_k = \mathrm{project}_k(\mathbf{s}_R)$ maps R's structural state into L's representation space.

L processes the full sequence $[\mathbf{p}; \mathbf{e}]$ with its standard attention mechanism.
L does not know or care that the first $K$ positions are structural context from R.
It attends to them exactly as it attends to real tokens.

\paragraph{Why this satisfies both principles.}
\begin{itemize}[nosep]
\item \textbf{L survives alone:} when R is untrained, the prefix tokens are random.
L learns to ignore them (low attention weight on prefix positions).
L's CE loss on real tokens is unaffected, the prefix is just noise at the start of the sequence.
\item \textbf{``Looks like'' context:} L already attends to context tokens and extracts information.
The prefix tokens live in the same $d$-dimensional space as real tokens.
L uses the same $Q$, $K$, $V$ projections, the same SharedM coupling, the same attention mechanism.
No new parameters are needed on L's side.
\end{itemize}

\paragraph{The attention IS the gate.}
L's attention weights on the prefix positions report the hemispheric coupling strength:
\begin{equation}
\alpha_{t,k} = \mathrm{softmax}\!\left(\frac{(\mathbf{e}_t A)(\mathbf{p}_k B)^\top}{\sqrt{d}}\right)
\end{equation}
When R is uninformative: $\alpha_{t,k} \approx 0$ for all $t$, $k$.
When R provides useful structural context: $\alpha_{t,k} > 0$ for positions $t$ that benefit from structural context.
The coupling develops gradually as R learns, no phase transition, no curriculum schedule.

\paragraph{Bandwidth control.}
$K$ controls the communication bandwidth, the corpus callosum width.
$K = 4$ gives R $4 \times d = 1024$ dimensions of structural context per forward pass.
$K = 1$ is minimal communication.
$K = 16$ is rich communication.
The optimal $K$ is itself learnable, R can produce fewer informative prefix tokens rather than many noisy ones.

\paragraph{How R trains.}
R observes L's coupling at multiple depths (layers 3, 6, 9) as before.
R's training signal: predict L's next-layer coupling (the structural LM objective).
R's state is projected into prefix tokens that feed back into L.
L's CE loss backpropagates through attention, through the prefix tokens, through the projection, to R's state.
So R receives gradient from two sources:
\begin{enumerate}[nosep]
\item Its own prediction loss (predict next-layer coupling), forces R to model L's dynamics.
\item L's CE loss through the prefix tokens, forces R to produce prefixes that help L predict.
\end{enumerate}
Both gradients train R. The first provides structural signal. The second provides utility signal.

\paragraph{The architecture diagram.}
\begin{verbatim}
  R observes L at layers 3, 6, 9
       → R processes (cross-position attention)
       → R state (r-dimensional)
       → project to K prefix tokens (d-dimensional each)
       → prepend to L's input at layers 10, 11
       → L attends to prefix using standard attention
       → L's CE loss flows back through prefix to R
\end{verbatim}

This is the simplest possible R$\to$L interface: no gates, no structural embeddings, no coupling bias.
Just tokens that L already knows how to use.


\subsection{Empirical Results: Prefix LORN v3 and v4}
\label{sec:prefix-results}

We validate the structural prefix architecture through two successive experiments on a 20M-parameter LORN trained on WikiText-103 + OpenWebText (67M tokens, no BabyLM) for 50K steps on Apple M4 MPS.
The two experiments differ only in how the prefix--real attention mechanism is constructed; everything else, L architecture, R architecture, data, optimiser, steps, is identical.
The comparison isolates a single design decision and its effect on all four claims the prefix mechanism is supposed to make good on.

Both models use $d{=}288$, 12 layers, SharedM coupling, a wide shared FFN (6$\times$), and tied embeddings; R is a $r{=}192$ cross-position structural language model observing L at layers $\{3,6,9\}$ with $K{=}2$ prefix tokens projected into L's space and injected at layers $\{10,11\}$.

\paragraph{v3: joint softmax (failed).}
The natural implementation prepends prefix tokens to L's keys and values at each injection layer and computes attention jointly:
\begin{equation}
[\mathbf{k}_{\text{ext}}; \mathbf{v}_{\text{ext}}] = [\mathbf{pk}; \mathbf{k}], [\mathbf{pv}; \mathbf{v}],\qquad C = \mathrm{softmax}(QK_{\text{ext}}^\top/\sqrt{d}),\qquad h = C V_{\text{ext}}.
\end{equation}
Prefix and real tokens compete in a single softmax.
This is what ``L's attention is the gate'' naively implies.

v3 exposed a bootstrap failure.
Training traces three regimes: (1) steps 1--5K: prefix attention weight $\alpha_{\text{pre}}=1.0$ because prefix tokens dominate at init and L has not yet learned to suppress them; (2) steps 6--23K: $\alpha_{\text{pre}}=0.0$ because L shuts prefix off via softmax competition once the prefix tokens are known to be noise; (3) steps 24--30K: slow recovery to $\alpha_{\text{pre}}\approx 0.22$ as R develops useful representations through its prediction loss and L re-engages.
An 18K-step dead zone in the middle of training where the CE gradient to R is zero.
Final val loss 5.95, prefix ablation effect $+$0.5\%, structure separation 1.37$\times$.

\paragraph{The chicken-and-egg problem.}
Joint softmax creates a failure mode where the two parties that must cooperate each depend on the other having already learned:
\begin{itemize}[nosep]
\item R produces useful prefix tokens $\Longrightarrow$ L attends to them $\Longrightarrow$ CE gradient flows back to R.
\item But before R has learned, prefix tokens are noise, so L's softmax pushes $\alpha_{\text{pre}}{\to}0$.
\item With $\alpha_{\text{pre}}{=}0$, the gradient through prefix attention is also zero, R receives no CE signal.
\item So R can only learn from its prediction loss, which produces structural representations but not \emph{useful} prefix tokens.
\end{itemize}
The v3 resurrection at step 24K shows that R's prediction loss alone \emph{does} develop useful structure, but only after losing 18K steps to a dead gradient.

\paragraph{v4: separate softmax with attention floor.}
The fix is to decouple the two attention streams:
\begin{align}
h_{\text{real}} &= \mathrm{softmax}(Qk^\top/\sqrt d)\,v \\
h_{\text{prefix}} &= \mathrm{softmax}(Q\mathbf{pk}^\top/\sqrt d)\,\mathbf{pv} \\
h &= (1-\mu)\,h_{\text{real}} + \mu\,h_{\text{prefix}},\qquad \mu = \max\!\big(\sigma(\theta_{\text{mix}}),\; \mu_{\text{floor}}\big)
\end{align}
where $\theta_{\text{mix}}$ is a per-injection-layer learnable parameter initialised at $-1.5$ (so $\sigma(\theta_{\text{mix}}) \approx 0.18$ at start) and $\mu_{\text{floor}} = 0.05$.
Two properties are critical.
First, real and prefix attention are computed in parallel and cannot starve each other, L cannot push $\mu$ to zero through softmax competition.
Second, the floor guarantees $\mu \geq 0.05$, so CE gradient to R is always alive, even if R is useless.
If R is genuinely useless, the model learns $\theta_{\text{mix}} \to -\infty$ and $\mu \to \mu_{\text{floor}}$.
If R is useful, $\mu$ grows above the floor organically.

\paragraph{v4 headline: LORN beats vanilla ORN at matched parameter count.}
\begin{center}
\begin{tabular}{lrrrrr}
\toprule
Model & Params & Best val & Ablation $\Delta$ & Peak probe & Final $\mu$ \\
\midrule
v3 L-only ORN (baseline) & 17.7M & 5.793 & ,  & ,  & ,  \\
v3 Prefix LORN           & 19.6M & 5.950 & $+$0.5\% & 1.37$\times$ & 0.22 \\
v4 Prefix LORN (floor)   & 19.6M & \textbf{5.087} & \textbf{+10.2\%} & \textbf{1.88}$\times$ & \textbf{0.42} \\
\bottomrule
\end{tabular}
\end{center}
v4 reduces val loss by 12.2\% relative to the L-only ORN baseline at roughly matched parameter count.
The prefix tokens encode 0.52 nats of information, removing them at inference time sends val from 5.09 to 5.59.
L \emph{chooses} to route 42\% of its deep-layer attention output through prefix tokens, well above the 5\% floor.

\paragraph{The val curve.}
Figure~\ref{fig:lorn-val} shows v4's validation loss trajectory alongside the v3 references.
The curve drops smoothly from 7.66 to 5.46 over the stable phase (steps 1K--40K), then accelerates during WSD decay to 5.09.
The final 20\% (steps 40K--50K) yields a 0.37 nat improvement from WSD decay alone, consistent with the transformative WSD effect observed in V2 training.

\begin{figure}[h]
\centering
\includegraphics[width=0.85\linewidth]{figures/fig_lorn_val_curve.pdf}
\caption{\textbf{LORN v4 validation curve.} The attention-floor fix lets L and R co-evolve from step 1. v4 crosses below the v3 L-only ORN baseline (dashed green, 5.79) at step 23K and continues to drop, reaching 5.09 at the end of WSD decay. v3 LORN (dotted red, 5.95) had the joint-softmax bootstrap failure.}
\label{fig:lorn-val}
\end{figure}

\paragraph{Mix trajectory: L's growing trust in R.}
Figure~\ref{fig:lorn-mix} shows the mixing ratio $\mu$ over training, per injection layer.
Both layers climb monotonically from $\mu \approx 0.18$ at init to $\mu \approx 0.425$ at the end of training.
The 5\% attention floor (red dashed line) never binds, $\mu$ is always far above it, because L genuinely wants more prefix attention as R develops.
This is the strongest qualitative evidence that prefix tokens are carrying real structure: the gradient through $\mu$ says ``more prefix helps'' at every stage of training, not just at the beginning.

\begin{figure}[h]
\centering
\includegraphics[width=0.85\linewidth]{figures/fig_lorn_mix_growth.pdf}
\caption{\textbf{Mix values trajectory.} The mixing ratio $\mu$ at both injection layers (layers 10 and 11) grows smoothly from init ($\approx 0.18$) to $\approx 0.425$. The attention floor ($\mu \geq 0.05$, red dashed) never binds, L wants prefix attention well above the guarantee. Compare to v3 where $\alpha_{\text{pre}}$ collapsed to zero at step 6K and stayed there for 18K steps before recovering.}
\label{fig:lorn-mix}
\end{figure}

\paragraph{Prefix ablation: direct proof of influence.}
Figure~\ref{fig:lorn-ablation} measures how much val loss increases when prefix tokens are disabled at inference time, a direct test of whether the preamble carries information.
The ablation effect grows monotonically from $\Delta \approx 0$ at step 5K to $\Delta = +0.52$ nats ($+10.2\%$) at step 50K.
Every additional nat of the total improvement over v3 L-only comes with increased dependency on prefix tokens.
This is the money chart: the structural preamble is not decorative, it is load-bearing.

\begin{figure}[h]
\centering
\includegraphics[width=0.85\linewidth]{figures/fig_lorn_ablation.pdf}
\caption{\textbf{Prefix ablation over training.} At each diagnostic checkpoint we evaluate val loss with prefix tokens enabled vs.\ disabled at inference. The difference grows from noise-level at step 5K to $+0.52$ nats at step 50K, removing the preamble hurts validation by 10\%. Prefix tokens encode half a nat of structural information that L cannot recover on its own.}
\label{fig:lorn-ablation}
\end{figure}

\paragraph{Structure separation: peak at step 20K.}
Figure~\ref{fig:lorn-probe} shows the length-controlled structure probe over training: R-state separation between matched-length ``child-simple'' and ``formal-complex'' sentences.
The probe peaks at 1.88$\times$ at step 20K, during the steepest part of the loss curve, then settles to a stable $\sim$1.1--1.3$\times$ for the rest of training.
The sharp-to-diffuse transition is a common pattern in representation learning: mid-training R states are maximally discriminative along the main structural axis; as training progresses R encodes more distributed, composition-friendly features rather than sharply clustered category labels.
The 1.88$\times$ peak exceeds Deep-R's 1.51$\times$ at comparable parameter count, and the persistent 1.1--1.3$\times$ floor exceeds v3's range ($\leq 1.01\times$ throughout) at every checkpoint.

\begin{figure}[h]
\centering
\includegraphics[width=0.85\linewidth]{figures/fig_lorn_structure_probe.pdf}
\caption{\textbf{Length-controlled structure probe.} R-state separation between structurally different sentences at matched token length. Peaks at 1.88$\times$ at step 20K (red dot) during the steepest descent of the val curve, then settles to a stable 1.1--1.3$\times$ regime. All points are above 1.0 (no separation), and all post-step-10K points exceed the moderate-significance threshold (orange dotted).}
\label{fig:lorn-probe}
\end{figure}

\paragraph{Generation quality.}
Qualitative progression across training, using temperature 0.8 sampling from a shared prompt:

\begin{center}
\small
\begin{tabular}{l p{9.5cm}}
\toprule
Step / val & Sample continuation of ``She walked'' \\
\midrule
 5K / 7.25 & \textit{``\ldots and we're question and close something and other. Those is someone but they must said\ldots''} \\
10K / 6.53 & \textit{``\ldots the handful simple.''}  (mostly fragments) \\
20K / 6.04 & \textit{``\ldots making. Although the base.''}  (word-level coherent, no sentence structure) \\
30K / 5.55 & \textit{``\ldots at the time\ldots rating the thing to movies?\ldots I think he was seen.''} \\
35K / 5.55 & \textit{``\ldots I don't intend to escape the story that I would let the people would always create a''} \\
45K / 5.22 & \textit{``\ldots Scientists have discovered the most significant number of gun type. Console: 27\% of those are young, except once without no evidence\ldots''} \\
50K / 5.09 & \textit{``\ldots into the back. Again, this system of course: The size of the way to actually figure out. If you can get a little bit of context, and get it\ldots''} \\
\bottomrule
\end{tabular}
\end{center}
By step 35K the model produces reasonable news-style prose with consistent syntax, consistent punctuation, and local topic coherence.
Some degenerate repetition (``$.\,.\,.\,.\,.$'' patterns) appears during WSD decay at temperature 0.8, a known failure mode of small language models at low sampling entropy.
This is a 19.6M parameter model trained on 67M tokens without BabyLM; compared to the 108M V2 ORN trained on 8B tokens (val 3.01) the raw quality is predictably weaker.
What matters for the claim under test is not absolute text quality but that disabling prefix tokens \emph{worsens} that quality measurably (via the $+10.2\%$ ablation effect) and that R-state separation confirms non-trivial structural content.

\paragraph{Three claims, three verdicts.}

\textit{Claim A: R's prefix tokens model non-trivial structure (not just length).}
The length-controlled probe reaches 1.88$\times$ peak and stabilises at 1.1--1.3$\times$, structurally different sentences at matched length are reliably separated in R-state space. \textbf{Supported.}

\textit{Claim B: L's validation loss matches vanilla ORN at same architecture.}
v4 LORN reaches val 5.087 vs.\ the L-only ORN baseline at 5.793 with roughly matched parameter count ($+11\%$). LORN does not merely match, it \emph{beats} the baseline by 12.2\%. \textbf{Supported strongly.}

\textit{Claim C: Learned prefix tokens influence generation.}
Prefix ablation shows $\Delta = +0.52$ nats ($+10.2\%$) on validation loss when prefix is disabled. Prefix tokens are load-bearing, not decorative. \textbf{Supported decisively.}

\paragraph{What still does not work: steering.}
Cosine similarity between prefix tokens extracted from different registers (formal vs.\ casual vs.\ question text) is 0.995--0.999.
Generated text under different prefix-token substitutions shows mostly surface-level formatting differences (WikiText markup) rather than structural register shifts.
Two plausible explanations: (1) $K{=}2$ prefix tokens carry insufficient capacity to differentiate register features at a resolution generation can exploit; (2) the cross-attention that extracts prefix tokens from R's state collapses to a dominant mode regardless of input.
Deep-R's steering also failed on this test.
We conjecture that the steering failure is a bandwidth issue ($K$ is too small) rather than a fundamental architectural limitation, and the scaled-up v5 experiment ($K{=}4$, $d{=}448$, 16 layers, 3 injection layers) tests this directly.

\paragraph{What the floor reveals about the interface.}
The v3$\to$v4 comparison is clean because only the attention mechanism changes, same L, same R, same data, same objective, same steps, same optimiser.
The $+$0.7 nat improvement and the qualitative shift in mix dynamics are attributable to the separate-softmax-plus-floor mechanism alone.
This tells us something about the ``L's attention is the gate'' principle: it works, but only if prefix and real streams are computed in parallel and never compete for the same softmax mass.
A minimum attention floor is the simplest way to guarantee gradient flow during R's bootstrap phase; learnable per-layer mixing coefficients give L the flexibility to use prefix differently at different depths.
Both ingredients matter.


\subsection{Three Independent Sub-claims}
\label{sec:three-subclaims}

The prefix-token architecture can be decomposed into three sub-claims, each of which can be measured and improved independently.
This decomposition is useful because the empirical evidence falls very differently across the three.

\paragraph{Sub-claim 1: L can deal with prefix tokens.}
\emph{Does L attend to, use, and depend on token-like context prepended to its input?}
v4 evidence: mix ratio climbs monotonically to 0.42 over training; prefix ablation at step 50K is $+10.2\%$ ($+$0.52 nats); removing prefix measurably hurts generation at every checkpoint after step 10K.
\textbf{Strongly supported.}

\paragraph{Sub-claim 2: R encodes structure in its internal state.}
\emph{Does R's $r$-dimensional state differentiate structurally distinct inputs?}
v4 length-controlled probe peaks at 1.88$\times$ structure separation; independent results on the V2 108M ORN show 60$\times$ syntactic separation in the residual stream~\cite{orn_lab_notebook}; R's dimensions correlate with register features in the trained BabyLM LORN (e.g., formal/child separation 2.31$\times$ in R-space).
\textbf{Strongly supported by v4 and prior work.}

\paragraph{Sub-claim 3: Structure in R survives compression into prefix tokens.}
\emph{Does the cross-attention compression $R_{\text{state}} \in \mathbb{R}^{S\times r} \to \mathbf{p} \in \mathbb{R}^{K\times d}$ preserve structural differentiation?}
This is the sub-claim where the v4 evidence is \emph{partially negative}.
Prefix ablation shows prefix tokens carry information (sub-claim 1 + sub-claim 3 together produce the $+10.2\%$ effect), but the steering test, where we extract prefix tokens from structurally different source texts and use them during generation, shows cosine similarity between register centroids of $0.995$--$0.999$ in prefix space.
Almost no \emph{directional} differentiation survives compression, even though L uses the prefix tokens effectively.
\textbf{Mixed evidence: some structure survives, but the compression flattens directional information.}

\paragraph{Direct diagnostic on v4 checkpoint.}
We ran a compression diagnostic on the final v4 checkpoint to quantify where information is lost.
For structurally-distinct matched-length sentence pairs, we compute R-state separation and prefix-token separation under the same clustering metric.
The results are informative rather than catastrophic:

\begin{center}
\begin{tabular}{lrrr}
\toprule
Metric & R-state ($r{=}192$) & Prefix ($K{=}2$, $d{=}288$) & ratio \\
\midrule
Global between/within L2 ratio & 1.03$\times$ & 1.16$\times$ & 1.12 \\
Cosine sim (child$\leftrightarrow$formal) & $+0.859$ & $+0.999$ & ,  \\
Cosine sim (statement$\leftrightarrow$question) & $+0.959$ & $+1.000$ & ,  \\
Internal diversity (K=2 tokens same input) & ,  & cos $0.9996$ & ,  \\
\bottomrule
\end{tabular}
\end{center}

The L2 ratio result is encouraging: the compression actually \emph{amplifies} between-category separation ($1.16\times$ vs $1.03\times$ in R), categories are pushed further apart relative to within-category spread.
But the cosine result is telling: compressed prefix tokens point in nearly the same \emph{direction} regardless of input (cos $\approx 1.00$), differing only in magnitude.
The two prefix tokens for a given input are also nearly identical (cos $0.9996$), $K{=}2$ effectively collapses to $K{=}1$.

\paragraph{Why the diagnostic is good news.}
The compression preserves structural separation as \emph{magnitude} differences along a shared axis, not as \emph{directional} differences between distinct prefix token directions.
L can use magnitude differences, this is consistent with the strong prefix ablation effect.
But steering requires directional differentiation: swapping prefix tokens from one register with those from another needs the two to point in different directions for L's attention to respond differently.
The steering failure is therefore a clean, targeted failure of the compression module, not a failure of R's representation.

\paragraph{The compression module capacity problem.}
In v4, the compression consists of:
\begin{itemize}[nosep]
\item $K \times r = 2 \times 192 = 384$ learned query parameters
\item $r \times d + d = 192 \times 288 + 288 = 55\,584$ parameters in the single projection
\item Total compression capacity: $\approx 56{,}000$ parameters.
\end{itemize}
R itself has $\sim 2$M parameters; L has $\sim 18$M.
The compression module is a $36\times$ bottleneck relative to R and a $320\times$ bottleneck relative to L.
Three specific failure modes follow:
\begin{enumerate}[nosep]
\item \textbf{Query collapse.} The $K$ learned queries initialised from $\mathcal{N}(0, 0.02^2)$ are nearly parallel; nothing forces them to attend to different subspaces of R.
\item \textbf{Shared projection.} A single $\mathrm{proj}: \mathbb{R}^r \to \mathbb{R}^d$ is applied to all $K$ query outputs, no mechanism for per-token specialisation.
\item \textbf{No diversity pressure.} The loss contains no term rewarding distinct prefix tokens; cosine between prefix tokens can drift to 1.0 without penalty.
\end{enumerate}

\paragraph{Targeted fixes (each directionally sensible).}
Each sub-claim can be improved independently.
For sub-claim 3 specifically:
\begin{itemize}[nosep]
\item \textbf{Fix 3a, Orthogonal queries.} Initialise \texttt{prefix\_queries} orthogonally (one line: \texttt{nn.init.orthogonal\_}). Forces $K$ queries to attend to distinct directions in R's state.
\item \textbf{Fix 3b, Per-token projections.} Replace the single \texttt{prefix\_proj} with a \texttt{ModuleList} of $K$ linears. Each prefix token gets its own $r \to d$ projection. Adds $(K-1) \cdot r \cdot d$ parameters, still tiny, but breaks the shared-projection collapse.
\item \textbf{Fix 3c, Diversity loss.} Add a term penalising off-diagonal cosine similarity in the $K \times K$ Gram matrix of prefix tokens. Directly trains $K$ tokens to be mutually distinct.
\item \textbf{Fix 3d, Structural contrastive loss.} For sentence pairs in the same register, minimise prefix distance; for pairs in different registers, maximise. Directly trains prefix tokens to encode register.
\end{itemize}
Fixes 3a--3c can be combined without retraining expense; 3d requires a labeled dataset of structural categories or an unsupervised proxy.
Each fix addresses a specific identified failure mode; each is independently measurable against the compression diagnostic.

\paragraph{Scaling under the compression bottleneck: v5 results.}
A parallel experiment (v5) tested whether simple capacity scaling helps: $d{=}448$, 16 layers, $r{=}256$, $K{=}4$ prefix tokens, 3 injection layers, 60K steps; $35.3$M total parameters ($1.8\times$ v4).
Everything else matched v4: same data, same attention floor, same WSD schedule, same compression module.

\begin{center}
\begin{tabular}{lrrr}
\toprule
Metric & v3 L-only (17.7M) & v4 LORN (19.6M) & v5 LORN (35.3M) \\
\midrule
Best val & 5.793 & 5.087 & \textbf{4.915} \\
Delta vs L-only & ,  & $-12.2\%$ & $-15.2\%$ \\
Delta vs v4 LORN & ,  & ,  & $-3.4\%$ \\
Prefix ablation & ,  & $+10.2\%$ & $+10.0\%$ \\
Peak structure probe & ,  & $1.88\times$ & $1.61\times$ \\
Final mix (mean) & ,  & $0.42$ & $0.40$ \\
Steering cosine (register) & ,  & $0.999$ & $0.996$--$0.999$ \\
\bottomrule
\end{tabular}
\end{center}

v5 scales modestly: $3.4\%$ val improvement at $1.8\times$ parameters.
The scaling is cleanest in val loss (and the associated WSD-decay boost: v5 dropped $0.34$ nats in its final 20\%, matching v4's $0.37$ nats).
But two signals are scale-invariant: the prefix ablation effect stays at $\approx$$10\%$ (prefix carries $\approx 0.5$ nats at both scales, not proportionally more), and the steering cosine stays at $\approx 0.999$ (direction collapse persists under compression at scale).

The v5 result sharpens the diagnosis.
Capacity in $L$ and $R$ reduces val loss modestly, but capacity cannot route through the compression bottleneck.
Fixing the compression module is a \emph{prerequisite} for further scaling to pay off, not something scale can wash out.
This motivates v6: isolate the compression fix from scale by retrofitting orthogonal queries and per-token projections onto v4's exact architecture.

\paragraph{v6: compression module fix.}
v6 keeps v4's architecture identical and changes only the compression module: \texttt{prefix\_queries} is initialised orthogonally (Fix 3a), and the single shared \texttt{prefix\_proj} is replaced with a \texttt{ModuleList} of $K$ per-token projections (Fix 3b).
This adds 56K parameters in the compression module (now 112K vs v4's 56K) and nothing else, same $L$, same $R$, same data, same schedule, same 50K steps.

\begin{center}
\begin{tabular}{lrrrr}
\toprule
Metric & v3 L-only & v4 LORN & v5 LORN (2$\times$ scale) & v6 LORN (fix) \\
\midrule
Best val & 5.793 & 5.087 & 4.915 & \textbf{4.989} \\
$\Delta$ vs L-only & ,  & $-12.2\%$ & $-15.2\%$ & $-13.9\%$ \\
$\Delta$ vs v4 & ,  & ,  & $-3.4\%$ & $-1.9\%$ \\
Prefix ablation & ,  & $+10.2\%$ & $+10.0\%$ & \textbf{$+10.8\%$} \\
Prefix internal cos & ,  & 0.9996 & ,  & \textbf{0.704} \\
Final mix [L10, L11] & ,  & [0.42, 0.43] & [0.41, 0.40, 0.41] & \textbf{[0.39, 0.47]} \\
Steering cos (register) & ,  & 0.999 & 0.999 & 0.995--0.999 \\
\bottomrule
\end{tabular}
\end{center}

Three results go in the expected direction, one does not.

\paragraph{What improved: internal prefix-token distinctness.}
The key v6 diagnostic is \texttt{pcos}, the cosine similarity between the $K{=}2$ prefix tokens generated from a single input.
In v4, this was $0.9996$ throughout training: effectively $K{=}1$.
In v6, pcos oscillates between $-0.64$ and $+0.82$ early in training and settles to $\sim$$0.7$ by end.
The two prefix tokens are now correlated but distinct, never collapsed.

\paragraph{What improved: depth-dependent prefix usage.}
v4's two injection layers (10 and 11) ended training with nearly identical mix ratios: $\mu_{10} = 0.42$, $\mu_{11} = 0.43$.
v6's diverge meaningfully: $\mu_{10} = 0.39$, $\mu_{11} = 0.47$, L11 routes 20\% more of its attention through prefix than L10.
The divergence grew throughout training (from $\Delta\mu = 0.01$ at step 10K to $\Delta\mu = 0.08$ at step 50K), confirming that per-token projections enable depth-dependent structural routing.
This is the first LORN experiment where different injection layers specialised.

\paragraph{What improved: val and ablation, together.}
v6 best val $4.989$ beats v4 by $1.9\%$ with identical L and R capacity and 56K additional compression parameters.
The prefix ablation effect strengthens from $+10.2\%$ to $+10.8\%$, removing prefix now costs $0.55$ nats rather than $0.52$, despite the absolute val being lower.
Prefix tokens carry \emph{more} information in v6 than in v4, not less.

\paragraph{What did not improve: cross-input steering.}
The steering test shows prefix cosines of $0.9954$--$0.9988$ between register centroids in v6, compared to $0.995$--$0.999$ in v4.
The compression-fix did not unlock directional differentiation across different inputs.

\paragraph{Two axes of collapse.}
v4 and v6 together reveal that ``prefix token collapse'' has two distinct axes:
\begin{enumerate}[nosep]
\item \textbf{Within-input collapse}: the $K$ tokens for a single input are identical. v4 exhibits this strongly (pcos 0.9996); v6 does not (pcos 0.704).
\item \textbf{Across-input collapse}: tokens for different inputs point in similar directions. v4 and v6 both exhibit this (cos 0.995+).
\end{enumerate}
Fixes 3a+3b (orthogonal queries, per-token projections) solve axis 1 but not axis 2.
The cross-attention step still produces similar weighted averages of R's state when R itself varies subtly across inputs; the per-token projections map a narrow variance band in R into a narrow band in prefix space.
Scale-up (v5) solves neither axis: both still show cos $\approx 0.999$ in steering.

\paragraph{What the two axes imply about the loss landscape.}
Axis 1 is \textbf{architectural}: nothing in the v4 pipeline pressures the $K$ queries to attend to different subspaces or the single projection to produce $K$ different tokens. Orthogonal init and per-token projections remove the architectural constraint.
Axis 2 is \textbf{loss-driven}: nothing in the training signal rewards prefix tokens for \emph{varying with input}.
The only gradient reaching the compression module through L's CE loss says ``make prefix useful for predicting next tokens.''
A compressed-context-of-$A$ that is almost identical to a compressed-context-of-$B$ can still be useful for both $A$'s and $B$'s next-token prediction, as long as it carries generally useful structural information.
To get strong across-input differentiation, the training loss must explicitly reward it.
This motivates v7 / Fix 3c--3d (diversity or contrastive loss across inputs), tested separately.

\paragraph{Net position after v6.}
Sub-claim 1 (L uses prefix tokens): strongly proven, stronger still in v6 ($+10.8\%$ ablation).
Sub-claim 2 (R encodes structure): proven by length-controlled probe, stable across scales.
Sub-claim 3 (structure survives compression): partially fixed.
Within-input prefix distinctness works (pcos 0.9996 $\to$ 0.704); depth-dependent routing works ($\mu_{10} \neq \mu_{11}$); val loss improves; ablation strengthens.
Cross-input directional differentiation remains broken, and the diagnostic now suggests this is a loss-signal issue rather than an architectural one.
The next targeted experiment adds a contrastive or diversity term to the prefix compression, no further architectural change needed.

\paragraph{v7: cross-input diversity via VICReg.}
v7 keeps v6's architecture exactly and adds a VICReg-style diversity loss on prefix tokens aggregated across the batch:
\begin{align}
\mathbf{P} &= \mathrm{prefix\_tokens}.\mathrm{reshape}(B{\cdot}K,\, d) \\
L_\text{var} &= \mathrm{mean}\left(\max(0,\, 1 - \mathrm{std}(\mathbf{P},\mathrm{dim}{=}0))\right) \\
L_\text{cov} &= \frac{1}{d}\sum_{i\ne j} \mathrm{Cov}(\mathbf{P})_{ij}^2 \\
L_\text{vicreg} &= L_\text{var} + 0.2\,L_\text{cov}
\end{align}
with a warmup ramp $\lambda_V(t) = 0 \to 0.3$ over the first 5K steps and the total loss $\text{CE} + \lambda_R\,L_R + \lambda_V(t)\,L_\text{vicreg}$.
The variance term explicitly pressures each prefix dimension to vary across inputs; the covariance term decorrelates dimensions; no invariance term is used since we have no positive pairs.

The complete comparison with all predecessors:

\begin{center}
\small
\begin{tabular}{lrrrrr}
\toprule
Metric & v3 L-only & v4 & v5 (scale) & v6 (arch) & v7 (VICReg) \\
\midrule
Params & 17.7M & 19.6M & 35.3M & 19.7M & 19.7M \\
Best val & 5.793 & 5.087 & \textbf{4.915} & \textbf{4.989} & 5.151 \\
$\Delta$ vs L-only & ,  & $-12.2\%$ & $-15.2\%$ & $-13.9\%$ & $-11.1\%$ \\
Prefix ablation & ,  & $+10.2\%$ & $+10.0\%$ & $+10.8\%$ & $+8.7\%$ \\
pcos (within) & ,  & 0.9996 & ,  & 0.704 & \textbf{0.246} \\
\textbf{Steering cos} (formal/casual) & ,  & 0.995 & 0.998 & 0.995 & \textbf{0.886} \\
\textbf{Steering cos} (casual/question) & ,  & 0.999 & 0.999 & 0.999 & \textbf{0.686} \\
\textbf{Steering cos} (formal/question) & ,  & 0.997 & 0.997 & 0.997 & \textbf{0.940} \\
Mix [L10, L11] & ,  & [0.42, 0.43] & ,  & [0.39, 0.47] & [0.37, 0.46] \\
\bottomrule
\end{tabular}
\end{center}

\paragraph{The steering breakthrough.}
The decisive result: v7 breaks the 0.999-cosine barrier that persisted through four prior experiments.
The casual--question pair drops to cosine $0.686$ ($46.6^\circ$ angular separation); formal--casual to $0.886$; formal--question to $0.940$.
Prefix tokens from structurally different source texts now point in genuinely different directions in L's input space, not just different magnitudes along a shared direction.

\begin{figure}[h]
\centering
\includegraphics[width=0.85\linewidth]{figures/fig_lorn_v4v6v7_steer.pdf}
\caption{\textbf{Steering test across three experiments.} Prefix-token cosine similarity between register-pair centroids. v4 (attention-floor only) and v6 (+ compression fix) both show cosines at $0.995$--$0.999$, directional collapse. v7 (+ VICReg diversity loss) drops all three pairs to $0.686$--$0.940$. The casual--question pair is the largest drop, reflecting strong syntactic differences; the formal--question pair is the smallest, reflecting shared sentential complexity despite different illocutionary force.}
\label{fig:lorn-steer}
\end{figure}

\begin{figure}[h]
\centering
\includegraphics[width=0.85\linewidth]{figures/fig_lorn_v4v6v7_pcos.pdf}
\caption{\textbf{Within-input K-token cosine (pcos) trajectory.} v4 pcos $\approx 0.9996$ throughout (collapsed; shown as horizontal line since v4 did not track this metric). v6's architectural fix (orthogonal queries + per-token projections) drops pcos to $0.4$--$0.8$ range, oscillating during training. v7's additional VICReg diversity loss pushes pcos further to a $0.15$--$0.40$ band, tokens nearly orthogonal. The architectural fix (v6) and loss-driven fix (v7) address different axes of collapse and compose additively.}
\label{fig:lorn-pcos}
\end{figure}

\begin{figure}[h]
\centering
\includegraphics[width=0.85\linewidth]{figures/fig_lorn_v4v6v7_val.pdf}
\caption{\textbf{Validation loss curves across three experiments.} All three sit well below the v3 L-only ORN baseline (dashed gray, val 5.79). v6 achieves the lowest val ($4.989$) via the compression-module fix with no loss-shape change. v7 pays a modest val cost ($+3.2\%$ over v6) in exchange for the steering capability shown in Figure~\ref{fig:lorn-steer}. All three exhibit the expected WSD decay acceleration after step 40K.}
\label{fig:lorn-val}
\end{figure}

\begin{figure}[h]
\centering
\includegraphics[width=0.85\linewidth]{figures/fig_lorn_v4v6v7_probe.pdf}
\caption{\textbf{Length-controlled structure probe trajectories.} v4 peaks at $1.88\times$ at step 20K and relaxes to $\sim 1.1\times$ for the rest of training, transient sharp clustering. v6 is more diffuse with no clear peak. v7 sustains $1.6$--$1.8\times$ through most of training, the VICReg pressure stabilises the structural clustering rather than allowing it to decay. Sharp transient (v4) vs.\ stable regime (v7) are two different operating modes for what R encodes.}
\label{fig:lorn-probe}
\end{figure}

Generation under different injected prefixes (fixed prompt ``The''):
\begin{itemize}[nosep]
\item \textit{[formal]}: ``The line of the club is the first a bit.\ \ldots the club's, ''
\item \textit{[casual]}: ``The first was issued by a week to allow \ldots `We were very excited to have an aspect of the prope, '\ldots''
\item \textit{[question]}: ``The `In preparation, I don't have the best solution for the next game?' \ldots``
\end{itemize}
Distinct styles from the same two-word prompt, under different prefix-token choices extracted from different register-texts.
This is the first LORN configuration that produces measurable register-conditioned generation.

\paragraph{The tradeoff.}
v7 pays for steering with a modest val loss increase: $5.151$ vs.\ v6's $4.989$ ($+3.2\%$) and v4's $5.087$ ($+1.3\%$).
Prefix ablation drops from $+10.8\%$ in v6 to $+8.7\%$ in v7, suggesting prefix tokens carry somewhat less average-useful information when forced to vary across inputs.
Both effects are consistent with VICReg pulling the optimisation in a direction orthogonal to pure CE minimisation.
v7 still beats the v3 L-only ORN baseline by $11.1\%$, so the tradeoff buys steering capability without sacrificing the core LM performance advantage.

\paragraph{pcos drops to an all-time low.}
Within-input K-token cosine similarity: v4 $0.9996$, v6 $0.704$, v7 $\mathbf{0.246}$.
VICReg pressure pushed tokens apart along both axes, within-input (the v6 axis) and across-input (the v7 axis), because the loss does not distinguish them and the two are coupled through the projection.
The $K{=}2$ tokens are now nearly orthogonal, and the projections have learned to specialise strongly.

\paragraph{Structure probe: sustained high vs.\ transient peak.}
v4's length-controlled probe peaked at $1.88\times$ at step 20K and collapsed to $\sim 1.1\times$ for the rest of training.
v7's probe sustains $1.6$--$1.8\times$ throughout the final two thirds of training, with a $1.76\times$ reading at step 30K.
VICReg produces a stable high-separation regime rather than a sharp transient, consistent with the idea that forcing per-input variation also stabilises the structural clustering that probe measures.

\paragraph{Complete sub-claim status.}

\begin{center}
\begin{tabular}{p{4.5cm}p{4.5cm}p{3cm}}
\toprule
Sub-claim & Proved by & Status \\
\midrule
(1) L uses prefix tokens & v4 (ablation $+10.2\%$) & \textbf{proven} \\
(2) R encodes structure & v4 probe, prior ORN work & \textbf{proven} \\
(3a) Within-input compression preserves distinctness & v6 (orthogonal queries + per-token projections: pcos $0.9996 \to 0.704$) & \textbf{proven} \\
(3b) Across-input compression preserves direction & v7 (VICReg: steering cos $0.999 \to 0.686$) & \textbf{proven} \\
\bottomrule
\end{tabular}
\end{center}

The two hemispheres work separately; the corpus callosum now carries structural information in a way L can read and that responds to input variation.
The prefix-token mechanism as a whole, architecturally (v6) and loss-wise (v7), delivers the end-to-end behaviour the theory requires: structure in R $\to$ directionally-distinct prefix tokens $\to$ register-conditioned generation in L.

\paragraph{v8: Pareto sweep on the VICReg weight.}
A natural follow-up question: is v7's val cost ($+3.2\%$ over v6) avoidable by halving the VICReg pressure?
v8 trains an otherwise-identical model with $\lambda_V = 0.15$ instead of $\lambda_V = 0.30$.
The answer is no, and the way it fails is informative.

\begin{center}
\begin{tabular}{lrrrr}
\toprule
Metric & v6 ($\lambda_V{=}0$) & v8 ($\lambda_V{=}0.15$) & v7 ($\lambda_V{=}0.30$) \\
\midrule
Best val & 4.989 & 5.130 & 5.151 \\
Steering cos (worst) & 0.999 & 0.9998 & \textbf{0.686} \\
pcos (within-input) & 0.704 & \textbf{0.126} & 0.246 \\
Prefix ablation & $+10.8\%$ & $+6.6\%$ & $+8.7\%$ \\
Mix [L10, L11] & [0.39, 0.47] & [0.44, 0.45] & [0.37, 0.46] \\
\bottomrule
\end{tabular}
\end{center}

v8 is \emph{dominated}: worse val than v6, worse steering than v7, worse ablation than both.
The two axes of collapse respond very differently to the diversity pressure:
\begin{itemize}[nosep]
\item \textbf{Within-input collapse} (pcos): smoothly decreases with $\lambda_V$. v8's pcos ($0.126$) is even lower than v7's ($0.246$), less competing CE pressure on the architecture.
\item \textbf{Across-input collapse} (steering): \textbf{threshold behaviour}. v8's $\lambda_V = 0.15$ is below threshold, steering reverts to $0.998$+. v7's $\lambda_V = 0.30$ is above threshold, steering breaks.
\end{itemize}
The mix-differentiation signal also collapses at v8 ($\mu_{10} \approx \mu_{11}$), without strong cross-input pressure, the per-token projections converge on similar effective patterns at the two injection layers.


\section{Toward an Approximate Gauge Theory of Structural Clustering}
\label{sec:gauge-lorn}

The prefix-token mechanism (v4--v8) works, but the conceptual scaffolding is awkward.
We assemble an architecture (R hemisphere, compression module, prefix injector) and a stack of losses (CE, coupling-prediction, VICReg variance, VICReg covariance) and tune until structure flows through.
The architecture is justified retrospectively by the diagnostic.
There is a deeper framing that may justify it (or replace it) prospectively, and we sketch it here as a research direction rather than a result.

\subsection{The Hard Problem: Defining ``Similar'' Is a Moving Target}

The original LORN aspiration is that R clusters analogous concepts, that ``argument is war'' and ``argument is journey'' end up in different relational neighbourhoods even though their content vocabularies overlap.
But the notion of analogy is not fixed.
A child's ``argument is war'' is structured around personal attack and victory; an adult's is structured around positions, evidence, and concession.
The same metaphorical relation organises different conceptual substrates.
If R must \emph{define} similarity in order to cluster, R must already possess a representation of what counts as analogous, which is the very thing we hoped R would learn.
This is the bootstrap problem of unsupervised structural learning: you cannot directly train a network to produce the right clustering when the clustering criterion is itself the object you are trying to discover.

The escape route is the same one physics took.
You stop trying to specify what things are and you specify what transformations leave things unchanged.
The clustering then \emph{emerges} from the constraint instead of being defined by hand.

\subsection{Action Is Easier to Specify Than Similarity}

Instead of training R to know which sentences are analogous, train R to produce an \emph{action} $H(x)$ on L's attention machinery, with two constraints:
\begin{enumerate}[nosep]
\item \textbf{Smoothness in input.} Similar $x$ produce similar $H(x)$.
\item \textbf{Approximate invariance in L's predictions.} Applying $H(x)$ to L's attention should leave L's next-token distribution \emph{approximately} unchanged.
\end{enumerate}
The first constraint says R is well-behaved as a function. The second is the gauge condition: R's actions should live (approximately) in the kernel of L's prediction operator. Whatever falls inside that kernel is, by L's own measure, structurally equivalent.

We never told R what ``similar'' means.
We told R that its actions must not change what L predicts.
The set of inputs that R can map to one another via small actions \emph{becomes} the equivalence class.
The clustering is operationally defined by the gauge orbit.

This is the same conceptual move as Noether's theorem in reverse: instead of inferring conserved quantities from symmetries, we use the requirement of conservation (L's predictions stay put) to discover the symmetries (R's allowed actions) and read off the structural equivalence classes (the orbits).

\subsection{Why ``Approximate'' Is Doing the Real Work}

Exact gauge symmetry would mean L is perfectly invariant under R's action.
For a finite-capacity language model that is (a) impossible, the model has too few degrees of freedom for an exact Lie group of symmetries on its predictions, and (b) undesirable, because exact invariance would mean R's action is information-free.
We want some information from R to influence L (otherwise R is dead weight); we want most of R's action to live in directions L has not yet exploited.

Approximate gauge symmetry tolerates small prediction changes under R's action, with two natural absorption mechanisms.

\paragraph{Mode A: Immediate perturbative correction.}
The standard ORN already includes a perturbative correction term per layer:
\[
x \;\leftarrow\; \mathrm{orbital\_output} + \mathrm{gate}(x) \cdot \mathrm{correction\_mlp}(x).
\]
The correction MLP can absorb local prediction errors introduced by R's gauge perturbation in the same forward pass.
This is the fast pathway: R perturbs, correction restores, L's logits look almost the same.
The perturbation has happened, R's action has been registered in the attention geometry, but the output is unchanged.
What we are calling the perturbative correction term is, under the gauge framing, exactly the operator that closes small gauge violations within a single forward pass.

\paragraph{Mode B: Delayed re-equilibration via training.}
For perturbations larger than the correction MLP can absorb, L's predictions degrade transiently.
Subsequent gradient steps then push L's parameters toward representations where the perturbation no longer matters, the prediction recovers, the gauge has been internalised.
Loss goes up briefly, then down, and the network has gained an additional approximate invariance.
This is the slow pathway: R perturbs, L breaks, L heals, L is now robust to that direction of perturbation.

The mix of A and B determines the dynamical regime.
Schedule R's perturbations small and frequent: mostly A, with the correction MLP doing local cleanup.
Schedule them large and rare: mostly B, with training-time recovery doing the work.
Both deliver gauge-invariance over time; they differ in how visible the disruption is to the val curve.
The intermittent pulse schedule we sketched earlier is the natural way to access Mode B without permanent damage.

The word ``approximate'' is therefore not a relaxation of the theory.
It is a feature of the regime: enough rigidity that L's predictions stay coherent, enough flexibility that the gauge can be internalised over training rather than being baked in by hand.

\subsection{The Hard Sub-problem: Finding Gauge Symmetries in a Trained Network}

The framing above is clean as a theory but punts on the concrete question: how does R learn what actions are allowable?
Without a guide, R will produce arbitrary perturbations, most of which destroy L's predictions.
That is exactly what we observed in early prefix-token experiments, R's untrained outputs were noise, L's CE loss spiked, L learned to ignore R, and the system stalled.

The interesting move is to make finding the gauge a separate, explicit problem.
A trained network has many directions in parameter space and activation space along which observable behaviour is approximately invariant; these are its inherent symmetries.
They are not labelled.
A second network can be trained to find them.

\paragraph{Offline gauge discovery.}
Train L to convergence on the LM task. Freeze L. Train a probe network $G$ whose objective is:
\begin{equation}
G^{*} \;=\; \arg\min_{G} \;\; \mathbb{E}_x \,\big\| f_L\!\big(x;\, \theta + G(x)\big) - f_L(x;\, \theta) \big\|^2
\;+\; \lambda_{\text{nt}} \cdot \mathrm{NontrivialityPenalty}(G).
\end{equation}
The first term measures how much $G$'s perturbations change L's predictions on the training distribution, we want it small.
The NontrivialityPenalty prevents the trivial solution $G \equiv 0$: it could be a variance term (the perturbation must be input-dependent and rich) or a rank constraint (the perturbation must span a chosen number of dimensions).
What $G^{*}$ finds are the directions of \emph{learned slack}, the gauge orbit of L at this point in its training.

Once we know the gauge orbit, R can be trained to project its outputs onto that orbit.
R's action is then guaranteed (approximately) to leave L's predictions intact, while still carrying information.
The system has been factored: $G$ discovers what's free, R uses what's free.

\paragraph{Online co-evolution.}
For a more ambitious system, $G$ trains alongside L and R.
Each step:
\begin{enumerate}[nosep]
\item L makes a forward pass.
\item $G$ identifies a direction in L's currently-detected gauge orbit.
\item R produces an action, projected onto $G$'s orbit estimate.
\item L predicts; CE loss applied (with optional intermittent pulse on R's action).
\item $G$ updates: keep finding the gauge as L's parameters shift.
\item R updates: organise inputs by similarity within the gauge orbit.
\item L updates: next-token prediction.
\end{enumerate}
This is genuinely co-evolutionary and harder to stabilise.
The offline variant is the right place to start: it admits clean experimental verification and decouples the gauge-discovery problem from the structural-clustering problem.

\subsection{Clustering Within Gauge Slack: The Real Objective}

A subtle point worth emphasising: \textbf{gauge symmetry does not give us metaphorical clustering for free.}
It gives us \emph{degrees of freedom} that R can exploit without disturbing L.
Within those degrees of freedom, R still has to do work, it has to cluster structurally analogous inputs together along the freed directions.
Without that, gauge slack is just slack: parameters that could move without consequence, sitting unused.

The full LORN objective on R, restated under the gauge framing, becomes:
\begin{equation}
\mathcal{L}_R \;=\; \underbrace{\big\| \pi_{G^{\perp}} R(x) \big\|^2}_{\text{stay in the gauge orbit}}
\;+\; \underbrace{\mathbb{E}_{x, x'} \big[ d_{\text{struct}}\!\big(R(x), R(x')\big)^2 \cdot s(x, x') \big]}_{\text{cluster within the orbit}}
\end{equation}
where $\pi_{G^{\perp}}$ projects out the gauge-allowed directions (so the first term penalises any component of R that would change L's predictions), and $s(x, x')$ is a similarity weight that pushes structurally-similar pairs closer in $R$-space.
The similarity weight is the hard part: it requires either weak supervision, an unsupervised proxy (VICReg, contrastive on perturbed pairs), or a self-similarity bootstrap from L's own attention.

The gauge framing does not eliminate the clustering problem.
It \emph{factors} it.
Step 1: find the directions in parameter/activation space that are free (gauge orbit).
Step 2: within those free directions, train R to organise inputs by similarity.
The two steps are independently tractable; together they may give us what direct clustering training cannot.

The crucial change in research strategy is: instead of asking ``how does R learn what's similar,'' we now ask the cleaner pair of questions:
\begin{enumerate}[nosep]
\item ``What are the directions in which L is approximately indifferent?'' (a question about $L$, answerable offline)
\item ``Within those directions, how do we cluster everything we know?'' (a clustering problem with a fixed, known feasible region)
\end{enumerate}
Each is more tractable than the original.

\subsection{What This Buys Us If It Works}

If approximate gauge symmetry is real and discoverable:
\begin{itemize}[nosep]
\item \textbf{R's actions become principled.} They live in a sub-manifold defined by L's own indifference, not in a hand-designed bottleneck.
\item \textbf{Structural categories emerge.} Inputs that map to one another under small gauge actions are members of the same equivalence class. We get clustering by construction.
\item \textbf{Steering becomes natural.} Replace one input's R-state with another's: the difference lives in the gauge orbit, so the prediction logits move only as much as the gauge slack permits, but the structural conditioning shifts. Steering is direct application of the orbit.
\item \textbf{The metaphor problem is reframed.} Two metaphorically-related concepts share a gauge action: the same R-perturbation maps one to the other, regardless of the substrates L has filled into them. The action is the metaphor; the substrate is L. The child--adult disparity in metaphorical depth becomes a difference in L, not in R, explaining why we do not get metaphor for free from gauge symmetry, only the right scaffolding within which content-rich substrates can manifest different surface realisations of the same relational structure.
\end{itemize}

\subsection{Concrete First Experiment}

The smallest test of this picture is offline gauge discovery on a trained v6 or v7 checkpoint.

\begin{enumerate}[nosep]
\item Take v6 (best val, no diversity loss). Freeze it.
\item Define a perturbation family: $G(x) = u(x) v(x)^T$, a rank-1 outer-product perturbation to the shared coupling matrix M, where $u, v$ are produced by a small probe network from L's hidden state.
\item Train the probe to maximise non-triviality (e.g., variance of $G(x)$ across batch) subject to $\mathbb{E}_x \| f_L(x; \theta + G(x)) - f_L(x; \theta) \|^2 < \epsilon$.
\item Measure the discovered orbit: what is its effective rank? How non-trivially does $G(x)$ depend on $x$? Do similar sentences (length-controlled probe) produce similar $G(x)$?
\item If the orbit is non-trivial: train a second probe whose objective is to use the orbit for clustering, structurally-similar inputs should map to nearby $G(x)$. Compare the resulting clustering quality to v7's prefix-token-via-VICReg.
\end{enumerate}

If the offline gauge orbit is rich and structurally meaningful, this validates the framing and motivates a co-evolutionary system.
If the orbit is trivial, only directions L is genuinely indifferent to are noise directions, with no structural content, the framing is wrong and we revert to the prefix mechanism as the primary interface.
Either outcome is informative.
The diagnostic itself (rank and structure of the orbit at successive checkpoints) is a useful new quantity in its own right: it measures how much representational slack L develops over training.

\subsection{Reinterpreting LORN's Existing Pieces}

This framing reinterprets LORN's components rather than replacing them.
\begin{itemize}[nosep]
\item \textbf{The shared coupling matrix M} is the natural target for the gauge action: it is global, low-dimensional, and already shared across layers. A per-input rank-$r$ perturbation $\Delta M(x) = u(x) v(x)^T$ is the simplest non-trivial action.
\item \textbf{The perturbative correction MLP} is the Mode-A absorption mechanism: it can compensate for small gauge perturbations within the same forward pass. We already train it for a different reason (corrections to the orbital output); under the gauge framing this is exactly what it should do.
\item \textbf{R's structural state} becomes the gauge potential: the parameter from which the action derives. Whether we project R into prefix tokens, into $\Delta M$, or into per-layer biases is an implementation choice; the gauge framing applies to all of them.
\item \textbf{VICReg in v7} can be reread as a non-triviality penalty: it forces R's outputs to span the available directions rather than collapse. Under the gauge framing this is exactly the constraint we want, just applied at the wrong stage, after compression instead of before.
\end{itemize}

The prefix-token mechanism may be the right interface for an action-based LORN, with R producing prefix tokens that are constrained to live in L's gauge orbit (rather than an unconstrained $d$-dimensional vector that L's softmax has to learn to use).
Or it may turn out that direct $\Delta M(x)$ is simpler and the prefix tokens were always a workaround.
The experiments needed to decide are no longer about whether structure can flow through the corpus callosum, v7 settled that, but about \emph{which functional form for the action gives the cleanest gauge structure}.
This is a more focused research question than ``does LORN work,'' and we know enough now to ask it.

\subsection{Connection to Self-Supervised Contrastive Learning}
\label{sec:gauge-ssl}

The gauge-LORN framing is, properly understood, \textbf{a self-supervised contrastive learner trained inline with gradient descent on a primary task}. This is a non-trivial claim about what we are doing, and it deserves to be unpacked.

\paragraph{What standard self-supervised methods presuppose.}
Every contemporary SSL method, SimCLR, MoCo, BYOL, SimSiam, DINO, MAE, JEPA, is trying to solve a single underlying problem: \emph{how do you define positive pairs without labels?} The methods differ in their workarounds:
\begin{itemize}[nosep]
\item \textbf{Augmentation-based} (SimCLR, MoCo, DINO): hand-designed transformations encode the prior. Crops and colour jitter are positives; rest of batch are negatives.
\item \textbf{Asymmetric architecture} (BYOL, SimSiam): EMA target network plus stop-gradient sidesteps explicit negatives, but the asymmetry encodes the prior implicitly.
\item \textbf{Predictive in representation space} (JEPA, I-JEPA): predict masked regions, hoping the prediction architecture biases toward useful features.
\item \textbf{Reconstruction} (MAE): pixel reconstruction as a pretext task.
\end{itemize}
All of these answer ``what should be similar'' \emph{at design time}, by either (a) specifying the augmentation set, (b) choosing the architectural asymmetry, or (c) picking the masking pattern. The choice is domain-specific: vision SSL needs vision augmentations, text SSL needs masking, audio SSL needs spectral perturbations. A different domain requires a different recipe.

\paragraph{What gauge-LORN does instead.}
The clustering criterion is not specified by the algorithm designer. It is \emph{read out of the network L itself}, as the set of perturbation directions to which L's predictions are approximately invariant. Two inputs are similar \emph{if and only if} small actions in R can map one to the other without disturbing L's predictions.

This means:
\begin{itemize}[nosep]
\item The augmentation set is \textbf{discovered, not designed}, it is whatever lies in the gauge orbit at the current point in L's training.
\item The augmentation set \textbf{co-evolves with the primary task}, as L learns more, its gauge orbit shifts, and the discovered augmentations track. There is no two-phase training.
\item The method is \textbf{domain-agnostic}, the gauge orbit is a property of any trained network on any task. The same procedure applies to text, vision, audio, multimodal.
\item The contrastive force is \textbf{intrinsic}, inputs whose R-action would leave the orbit are implicitly negatives (the gauge constraint penalises them); inputs whose R-action lies in the orbit are implicitly positives (R can move them together for free).
\end{itemize}

\paragraph{The contrastive equivalence, made explicit.}
In standard contrastive learning the loss has the form
\[
\mathcal{L}_{\text{contrast}} = -\log \frac{\exp\!\big(\mathrm{sim}(z_x, z_{x^+})/\tau\big)}{\sum_{y} \exp\!\big(\mathrm{sim}(z_x, z_y)/\tau\big)},
\]
where $x^+$ is a positive pair (augmented $x$) and $y$ ranges over negatives.
In gauge-LORN, the corresponding force is implicit in the constraint structure:
\begin{itemize}[nosep]
\item \textbf{Implicit positives}: inputs $x, x'$ such that there exists an action $H \in \text{gauge orbit}$ with $R(x') \approx H \cdot R(x)$. R is rewarded for finding such pairs (they cost nothing under the orbit constraint).
\item \textbf{Implicit negatives}: inputs $x, y$ such that mapping $R(x) \to R(y)$ requires leaving the orbit. The gauge penalty $\| \pi_{G^{\perp}} R(\cdot) \|^2$ pushes them apart.
\end{itemize}
No explicit positive sampling, no negative batch, no temperature hyperparameter $\tau$, no augmentation library. The entire contrastive structure is induced by the gauge constraint on R combined with the non-triviality requirement on the discovered orbit.

\paragraph{Why this is potentially significant.}
If gauge discovery works on language modeling and the orbit turns out to be rich enough to support clustering, the same procedure transfers to any domain where a primary task is being learned. SSL becomes a side-effect of training a useful network on a useful task, with no domain-specific scaffolding. The current SSL paradigm, pre-train, fine-tune, with hand-tuned augmentation pipelines per domain, would have a principled alternative.

We do not claim this works yet. Phase 1 of the proposed experimental programme is the test: \emph{does the gauge orbit of a trained ORN have enough structural content to support meaningful clustering?} The answer is unknown. But the \emph{question} is the right one to ask, and we know how to answer it cleanly with the existing v6/v7 checkpoints.

\paragraph{Honest enumeration of risks.}
The framing rests on assumptions any of which could break:
\begin{enumerate}[nosep]
\item \textbf{Trivial orbit.} The gauge directions might be noise, directions L genuinely doesn't care about because they carry no signal of any kind. Discovered slack would be empty of structure.
\item \textbf{Moving target.} L's orbit shifts as L trains, requiring G to track. Online co-evolution may not converge.
\item \textbf{Probe collapse.} Despite a non-triviality penalty, G might find a single dominant direction and ignore the rest, yielding rank-1 orbits when we wanted rank-$r$.
\item \textbf{Within-orbit clustering still needs a metric.} The orbit constraint says ``these directions don't matter to L,'' not ``these directions cluster things meaningfully.'' We have factored the problem; we have not eliminated it.
\item \textbf{Compute.} Training G alongside L roughly doubles forward-pass cost.
\end{enumerate}
Each of these has a discrete diagnostic. Each can be confirmed or ruled out from a single small experiment. The path from current state to verdict is well-defined.

\paragraph{The bet, in one sentence.}
If the gauge orbit of a trained language model contains structurally meaningful slack, and we have qualitative reason to believe it does, given the v7 steering result and the multi-scale ORN structure separation, then gauge-LORN may provide a complementary self-supervised mechanism that derives augmentations from the network rather than from domain-specific design.
We do not claim to replace SimCLR/BYOL/VICReg-style methods; we propose an additional source of structural signal that has, to our knowledge, not been exploited.
Whether the orbit is rich enough to support full SSL training, or whether it provides only a structural sub-axis within a primary task, is exactly what the experimental programme decides.
The smallest experiment that could falsify the framing fits in an overnight MPS run.

\begin{figure}[h]
\centering
\includegraphics[width=0.95\linewidth]{figures/fig_lorn_pareto.pdf}
\caption{\textbf{Val/steering Pareto under VICReg lambda.} Left: val loss as a function of $\lambda_V$; the dependency is roughly linear in the small range tested. Right: val--steering Pareto frontier. v6 ($\lambda{=}0$) and v7 ($\lambda{=}0.30$) are the two clean operating points. v8 ($\lambda{=}0.15$) sits inside the frontier, worse val than v6 and worse steering than v7, because the across-input directional collapse responds to $\lambda_V$ as a threshold, not a smooth gradient. Halving the pressure does not buy half the differentiation.}
\label{fig:lorn-pareto}
\end{figure}

\paragraph{Net implication: v6 and v7 are two clean operating points.}
\begin{itemize}[nosep]
\item Use \textbf{v6} when the goal is best val loss with the architectural fix and no diversity pressure: best val $4.989$, no steering capability.
\item Use \textbf{v7} when the goal is steering capability: best val $5.151$ ($+3.2\%$ cost), steering cos $0.686$--$0.940$.
\end{itemize}
There is no useful intermediate point at lower $\lambda_V$ on this dataset and scale.
This sets a clean experimental design constraint for future scaling work: VICReg must be applied at sufficient weight to cross the steering threshold, otherwise it is pure cost.
The right axis to explore next is \emph{higher} $\lambda_V$ with a longer training schedule, or a complementary contrastive loss that gives stronger across-input signal at lower compute cost.

\paragraph{Summary of current evidence.}

\begin{center}
\begin{tabular}{p{4.8cm}cccp{4.3cm}}
\toprule
Sub-claim & v4 metric & Result & Status & Next step \\
\midrule
(1) L uses prefix tokens & mix, ablation & $0.42$, $+10.2\%$ & \textbf{proven} & ,  \\
(2) R encodes structure & length-ctrl probe & $1.88\times$ peak & \textbf{proven} & ,  \\
(3) Structure survives compression & steering cos & $0.999$--$1.000$ & \textbf{bottleneck} & orthogonal queries, per-token projections, diversity loss \\
\bottomrule
\end{tabular}
\end{center}

The two hemispheres work separately.
The corpus callosum, the compression module that carries structural information from R into prefix tokens, is the piece that needs architectural work.
This is an encouraging place to be: the bottleneck is localised, small (56K parameters in v4), and has known fixes.


% =====================================================================
\part{Gauge-LORN: Mathematical Foundations and Experimental Programme}
% =====================================================================

Part~I established empirically that the prefix-token mechanism works (v4--v8): structure flows from R through compression into prefix tokens, L uses them, and with the v6+v7 fixes, prefix tokens directionally separate by register.
The architecture was found by iteration on a series of failure modes; each fix patched a specific symptom; the set of fixes was justified retrospectively.

Part~II asks the question we should have asked first.
\emph{What is the right object for R to compute, given what L is doing?}
The answer, we argue, is an action that lies (approximately) in L's gauge orbit: a perturbation to L's parameters that leaves L's predictive distribution approximately unchanged.
The clustering of inputs that R produces becomes a consequence of how L deems them equivalent under the action, not a consequence of supervised similarity labels or hand-designed augmentations.

This Part is theoretical and prescriptive.
We work out the geometry, identify the free learnable parameters, specify the diagnostic measurements, and outline a three-phase experimental programme that can either validate or falsify the framing using the existing v6/v7 checkpoints.
The programme is the concrete next step; this Part articulates it precisely enough to start.


\section{Mathematical Formulation: What Exactly We Hold Constant}
\label{sec:invariant}

\subsection{Refinement: Shared FFN as Co-equal Gauge Component}
\label{sec:shared-ffn-gauge}

Before developing the gauge-LORN formalism, we record an empirical refinement that changes what the orbit is attached to.
The initial framing treated $M = AB^{\!\top}$ as the sole ``shared coupling'' object whose slack R should navigate.
This is incomplete.
The ORN architecture has \emph{two} globally shared components: the shared coupling matrix $M$ \emph{and} the wide shared feed-forward network (referred to as SharedFFN).
Both are applied at every layer.
Both concentrate the model's per-token compute.
Both are candidate substrates for a gauge action.

\paragraph{Empirical finding: FFN signatures dominate M signatures.}
A multi-signature probe run on the v6 ORN measured per-sentence weight signatures restricted separately to $M$, to the shared FFN (SharedFFN as a whole), and to the up-projection $\mathrm{FFN}_{0}$ of SwiGLU specifically.
The signal, measured as a $z$-score of LDA class separation against a permutation null, is:

\begin{center}
\small
\begin{tabular}{l c l}
\toprule
Signature subspace & $z$-score of LDA separation & Verdict \\
\midrule
$\nabla_M$ only                              & $+2.99$ & structurally non-trivial \\
$\nabla_{\mathrm{FFN}}$ only                 & $+4.41$ & stronger than $M$ \\
$\nabla_{\mathrm{FFN}_{0}}$ (up-projection)  & $+\mathbf{4.76}$ & strongest single subspace \\
$\nabla_M \oplus \nabla_{\mathrm{FFN}}$      & $+4.41$ & same as FFN (correlated) \\
\bottomrule
\end{tabular}
\end{center}

Two facts stand out.
First, the shared FFN carries a \emph{stronger} per-sentence category signature than the shared coupling $M$.
Second, $\nabla_M$ and $\nabla_{\mathrm{FFN}}$ are not independent: their concatenation does not improve over $\nabla_{\mathrm{FFN}}$ alone, indicating that the category-relevant variation in $M$-gradient space lives in a subspace already spanned by FFN-gradient variation.
The two shared components share an axis; FFN dominates.

\paragraph{Conceptual refinement.}
The gauge-LORN theory as developed in the rest of Part~II is stated in terms of perturbations to $M$.
This is a concrete simplification, not a principled restriction.
The correct general statement is:

\begin{quote}
\textit{The gauge orbit lives in the \emph{shared} parameters $\theta_{\mathrm{shared}} = (\theta_M, \theta_{\mathrm{FFN}})$, not only in $\theta_M$.
The Schur-complement partition $\theta = (\theta_M, \theta_R)$ should be read as $\theta = (\theta_{\mathrm{shared}}, \theta_R)$.
A dynamic operator $G_\phi(x)$ produces actions $\delta_{\mathrm{perturb}}(x) = (\delta M(x),\, \delta\mathrm{FFN}(x))$, both low-rank, jointly inside the recoverable orbit of the joint (shared, response) Hessian.
Weight signatures should be computed across both.}
\end{quote}

This is consistent with the empirical picture reported elsewhere: the shared wide FFN owns roughly 97\% of the model's representational manifold, the shared $M$ owns roughly 36\%; together they form the slow distributed substrate that the per-layer response heads read and write.
FFN is the world-model, $M$ is the compass, the compass is informative, but the map carries more bits.

\paragraph{How to read the rest of Part~II.}
The remainder of Part~II, including the Fisher analysis, the Schur complement, the recoverable-orbit theorems, and the v8--v15 probes, is written in terms of ``$M$'' because that is the subspace the probes actually measured.
\emph{Every reference to $M$ as ``the shared coupling'' or ``the gauge action surface'' should be read as shorthand for ``the shared parameter block'', which in the full architecture is $(\theta_M, \theta_{\mathrm{FFN}})$.}
The mathematics is identical, only the dimensionality of the subspace grows.
Where we write $\nabla_M$, the honest quantity is $\nabla_{\theta_{\mathrm{shared}}}$; the empirical finding above indicates that restricting to $M$ alone probably under-counts the orbit's structural content.

\paragraph{Implications for the experimental programme.}
The v10--v15 results reported below were obtained with probes that observe $\nabla_M$ only.
This is why several weak-supervision attempts (v11's episodic contrastive, v12's heuristic supervision, v13's VICReg) showed only modest or mis-aligned structure: the strongest per-sentence signature axis, the FFN, was not in the probe's basis.
Re-running v10--v15 with signatures drawn from $\theta_{\mathrm{shared}}$ (or from $\mathrm{FFN}_{0}$ specifically) is a concrete future-work item.
The theorems below still hold as stated, but their empirical instantiations are lower bounds on what a signatures-on-$\theta_{\mathrm{shared}}$ probe would reveal.


\subsection{The Invariant: Scope of the Problem}

Defining the invariant precisely is more than half the problem.
Several plausible quantities can be ``held constant under symmetry,'' each producing a different orbit and a different research programme.
We enumerate them, justify our choice, and derive the consequences.

\subsection{Seven Candidate Invariants}

For a language model parameterised by $\theta$ producing conditional distributions $P(y \mid x; \theta)$, candidate invariants under perturbation $\delta\theta$:

\begin{enumerate}[nosep]
\item \textbf{Pointwise logit invariance}: $P(y \mid x; \theta + \delta\theta) = P(y \mid x; \theta)$ for every $x$, $y$. The strictest possible, only true gauge symmetries of the function satisfy this.

\item \textbf{Pointwise KL bounded}: $\mathrm{KL}\!\big(P(\cdot \mid x; \theta + \delta\theta) \,\|\, P(\cdot \mid x; \theta)\big) < \epsilon$ per step, on the input distribution. The relaxation of (1) to allow small predictive drift.

\item \textbf{Total loss bounded}: $\mathcal{L}(\theta + \delta\theta) - \mathcal{L}(\theta) < \epsilon$. Averages over $x$, so individual predictions can shift large amounts as long as the average doesn't. Too permissive.

\item \textbf{Hidden state preservation}: $h_l(x; \theta+\delta\theta) = h_l(x;\theta)$ at every layer. Counterproductive: we \emph{want} hidden states to shift in a way that encodes structure.

\item \textbf{Attention pattern preservation}: $C_l(x; \theta+\delta\theta) = C_l(x;\theta)$ at every layer. Same problem as (4); attention is what we want to use as the action surface.

\item \textbf{Joint generation distribution bounded}: $\mathrm{KL}\!\big(P^{(T)}_{\theta+\delta\theta}(y_{1:T} \mid x) \,\|\, P^{(T)}_\theta(y_{1:T} \mid x)\big) < \epsilon$ where $P^{(T)}$ is the joint distribution over $T$-token autoregressive samples.

\item \textbf{Behavioural metric preservation}: perplexity, BLEU, human judgements unchanged. The thing we actually care about, but not differentiable and difficult to bound.
\end{enumerate}

We adopt invariant (6): \textbf{the joint distribution over generated $T$-token sequences is approximately preserved}.
This is what ``the language L generates does not change'' formally means.
We do not need exact preservation (impossible for finite networks); we need it bounded by some chosen $\epsilon$.

\subsection{From Sequence-Level to Per-Step: The Chain Rule of KL}

Invariant (6) is a constraint on a high-dimensional joint distribution.
Fortunately, autoregressive factorisation reduces it cleanly to a per-step constraint.
By the chain rule of KL divergence:
\begin{equation}
\mathrm{KL}\!\big(P^{(T)}_{\theta+\delta\theta} \,\|\, P^{(T)}_\theta\big) \;=\; \sum_{t=1}^{T} \mathbb{E}_{x_{<t} \sim P^{(t-1)}_\theta}\!\Big[\mathrm{KL}\!\big(P_{\theta+\delta\theta}(\cdot \mid x_{<t}) \,\|\, P_\theta(\cdot \mid x_{<t})\big)\Big].
\label{eq:chain-rule-kl}
\end{equation}
If we want sequence-level KL bounded by $\epsilon$, each per-step KL must be bounded by $\epsilon / T$.
For typical generation lengths $T = 100$ and a sequence-level budget of $\epsilon = 1$ nat, the per-step budget is $0.01$ nat, small, but achievable.

This recovers invariant (2) as the per-step counterpart of invariant (6).
The per-step constraint is geometrically clean: it cuts out a level set of the local Fisher form at scale $\epsilon / T$.

\subsection{The Local Fisher Geometry}

Around the unperturbed parameters $\theta$, the per-step KL has a second-order expansion:
\begin{equation}
\mathrm{KL}\!\big(P(\cdot \mid x; \theta + \delta\theta) \,\|\, P(\cdot \mid x; \theta)\big) \;=\; \tfrac{1}{2}\,\delta\theta^{\!\top} F(\theta; x) \,\delta\theta \;+\; O(\|\delta\theta\|^{3}),
\end{equation}
where $F(\theta; x) = \mathbb{E}_{y \sim P(\cdot \mid x; \theta)}\!\big[\nabla_\theta \log P(y \mid x; \theta)\, \nabla_\theta \log P(y \mid x; \theta)^{\!\top}\big]$ is the per-input Fisher.
Averaging over the input distribution gives the population Fisher $F(\theta) = \mathbb{E}_{x \sim D}[F(\theta; x)]$.

The \textbf{language-preserving gauge orbit} at $\theta$ is then:
\begin{equation}
\mathcal{G}_{\epsilon, T}(\theta) \;=\; \big\{ \delta\theta \;:\; \tfrac{1}{2}\,\delta\theta^{\!\top} F(\theta) \,\delta\theta \;<\; \epsilon / T \big\}.
\end{equation}
This is the level set of the Fisher quadratic form at scale $\epsilon/T$, an ellipsoid in parameter space whose principal axes are the eigenvectors of $F(\theta)$, with semi-axis lengths $\sqrt{2(\epsilon/T) / \lambda_i}$.
Eigenvectors with small $\lambda_i$ have long semi-axes (large slack); eigenvectors with large $\lambda_i$ have short semi-axes (sharp).

This is the precise mathematical object we are asking R/G to operate within.

\subsection{Where Random Matrix Theory Tells Us What's Feasible}

For a deep network, $F(\theta)$ is a $|\theta| \times |\theta|$ matrix that is intractable to compute or store.
We work with the empirical Fisher restricted to a small parameter subspace, in our case, the shared coupling matrix $M$, a $d^{2}$-dimensional subspace.
For our experiments, $d^{2} = 288^{2} = 82{,}944$ parameters in M.

The empirical Fisher restricted to M with $N$ samples is
\begin{equation}
\hat{F}_M \;=\; \tfrac{1}{N}\, G_M^{\!\top} G_M, \qquad G_M = \big[\nabla_M \log P(y_1 \mid x_1; \theta);\; \ldots;\; \nabla_M \log P(y_N \mid x_N; \theta)\big] \in \mathbb{R}^{N \times d^{2}}.
\end{equation}
$\hat{F}_M$ has rank at most $\min(N, d^{2})$.
For $N \ll d^{2}$ (the regime we are in: $N = 200, d^{2} = 82{,}944$), the empirical Fisher has trivial null space of dimension $d^{2} - N \geq 82{,}744$, this much of the parameter space is in $\hat{F}_M$'s exact kernel.

Random matrix theory~\cite{sagun2016eigenvalues, papyan2019measurements} characterises the spectrum of empirical sample covariances of this form.
The relevant facts for us:
\begin{itemize}[nosep]
\item The Marchenko--Pastur edge says that the smallest non-zero eigenvalue of $\hat{F}_M$ is order $(1 - \sqrt{d^{2}/N})^{2} \sigma^{2}$, where $\sigma^{2}$ is the variance of gradient entries. For our setting $\sqrt{d^{2}/N} \approx 9.1 \gg 1$, so the noise floor is high. \emph{Low non-zero eigenvalues of $\hat{F}_M$ are unreliable estimates of the true Fisher.}
\item Outlier eigenvalues (larger than the bulk) correspond to genuine high-Fisher directions of the population matrix and are stable to estimate.
\item The kernel of $\hat{F}_M$ (the $d^{2} - N$ directions outside the row-span of $G_M$) is approximately the kernel of the population Fisher \emph{when} $N$ is sufficient to span the high-Fisher subspace, which empirically requires $N \gtrsim$ a few hundred for typical deep networks.
\end{itemize}

The practical consequence: \emph{we cannot reliably enumerate the low-eigenvalue subspace of $F$, but we can reliably identify the high-eigenvalue subspace}, and we can ask whether a specific candidate direction (such as a structural Jacobian) lies inside it or outside it.
This shifts the diagnostic from ``find the orbit'' to ``test whether candidate directions are in the orbit.''

\subsection{Identifying Structural Slack: The Structural Jacobian}

The orbit being non-empty does not imply it carries structural meaning, this is the central risk identified earlier (Risk B).
To test for structural meaning we need a candidate direction that, if it lay in the orbit, would correspond to structural variation.
The natural candidate is the \textbf{structural Jacobian}.

For a binary structural feature $S$ (e.g.\ ``is a question'' vs.\ ``is a statement''), define
\begin{equation}
J_S(\theta) \;=\; \mathbb{E}_{x \sim D_S^{+}}\!\big[\nabla_\theta \log P(\cdot \mid x; \theta)\big] \;-\; \mathbb{E}_{x \sim D_S^{-}}\!\big[\nabla_\theta \log P(\cdot \mid x; \theta)\big].
\end{equation}
$J_S(\theta) \in \mathbb{R}^{|\theta|}$ is the average direction in parameter space along which moving $\theta$ would shift the model's predictive response to $S$-positive inputs versus $S$-negative inputs.
Restricted to the M subspace, $J_S^{(M)}(\theta) \in \mathbb{R}^{d^{2}}$.

The diagnostic question becomes: \emph{how much of $J_S^{(M)}$ lives in the orbit} (low Fisher) \emph{vs.\ outside it} (high Fisher)?

Decompose $J_S = \pi_{\text{high}} J_S + \pi_{\text{orbit}} J_S$, where $\pi_{\text{high}}$ projects onto the top-$k$ singular vectors of $G_M$ (the high-Fisher subspace) and $\pi_{\text{orbit}}$ projects onto the orthogonal complement.
The structural-slack ratio is
\begin{equation}
\rho_S(k) \;=\; \frac{\| \pi_{\text{orbit}} J_S \|^{2}}{\| J_S \|^{2}}.
\end{equation}
This is computable directly: SVD of $G_M$ via the Gram trick (a $N \times N$ eigendecomposition since $\hat{F}_M = G_M^{\!\top} G_M$), then projection of $J_S$ onto the top-$k$ singular vectors of $G_M$.
Cost: one batched forward$+$backward pass for the $N$ baseline gradients, plus one batched forward$+$backward for the $J_S$ computation per feature.
No learning, no architectural choices, no failure modes for a separate probe network.

\paragraph{Interpretation of the ratio.}
\begin{itemize}[nosep]
\item $\rho_S \approx 1$: $J_S$ is mostly orthogonal to the high-Fisher directions. Moving along $J_S$ shifts the model's structural response with little impact on its predictions. \textbf{Structural slack exists for feature $S$.}
\item $\rho_S \approx 0$: $J_S$ lives almost entirely in the high-Fisher subspace. Modifying the model's response to $S$ requires large prediction shifts. \textbf{No structural slack for $S$.}
\item Intermediate: partial slack. The fraction tells us what proportion of $J_S$ lives in the orbit; the rest would propagate to predictions.
\end{itemize}

\paragraph{Why this matters.}
$\rho_S$ is exactly the diagnostic we need before committing to a learnable G.
If $\rho_S$ is high for several structurally meaningful features, we know the orbit carries structural information and a learnable G has something to discover.
If $\rho_S$ is uniformly low, the orbit is structurally vacant, and Phase~1 of the original programme (training G) cannot produce structurally-meaningful actions because such actions don't exist.
The diagnostic costs $\sim$30 minutes on MPS and is decisive.

\subsection{Implications for the Probe Architecture}

Even if $\rho_S$ is high (orbit contains structure), discovering it via a learned probe G is non-trivial because the structural directions in the orbit are nonlinear functions of input features.
The user observation is correct: simple linear-then-outer parameterisations (G outputs $U(x) V(x)^{\!\top}$ from a pooled hidden state) almost certainly cannot capture the true structure-bearing directions, because:

\begin{enumerate}[nosep]
\item The orbit is a high-dimensional subspace, and its intersection with the structurally-meaningful subspace can have arbitrary curvature.
\item The mapping $x \mapsto \pi_{\text{orbit}} J_{S(x)}$ is itself a nonlinear function of the input that the network's own nonlinearities have shaped.
\item Random directions in the orbit (the simplest baseline) give zero structural correlation, as our offline probe directly demonstrated.
\end{enumerate}

The candidate G parameterisations should be designed with this in mind:

\begin{enumerate}[nosep]
\item \textbf{Pure linear-then-outer}: $\delta M(x) = U(x) V(x)^{\!\top}$ with $U, V$ from MLP on pooled hidden state. Baseline; expected to fail.
\item \textbf{Basis combination}: $\delta M(x) = \sum_{i=1}^{r} \alpha_i(x) \cdot M^{(i)}_{\text{basis}}$, where $\{M^{(i)}_{\text{basis}}\}$ are learned shared basis matrices and $\alpha_i(x)$ are nonlinear gates. The basis directions can be initialised from the orbit-projected structural Jacobians.
\item \textbf{Conditional rotation}: $\delta M(x) = R(x)^{\!\top} M R(x) - M$ where $R(x) \in O(d)$ is an input-conditional orthogonal matrix. The orbit is preserved if rotations live in M's symmetry group.
\item \textbf{Attention-output perturbation}: $G$ produces a perturbation directly to attention scores rather than to M, bypassing the matrix structure.
\item \textbf{Transformer-block G}: $G$ itself is a small attention block that processes $(h_l, M)$ and outputs the perturbation. Maximally expressive; most expensive.
\end{enumerate}

The point is not that one of these is correct, it is that the choice of parameterisation \emph{is itself a hypothesis} about what kind of nonlinear structure the orbit contains.
A revised Phase 1 should sweep these parameterisations, not just sweep $r$ within parameterisation~(1).
The structural-Jacobian diagnostic above gives us the \emph{target} (the projected $J_S$); the question is which parameterisation can reach it.

\subsection{Result of the Direct Diagnostic on v6}

We computed $\rho_S(k)$ for four structural features on the v6 final checkpoint, with $N = 200$ baseline gradients sampled from validation and $k$ ranging from 5 to 100:

\begin{center}
\begin{tabular}{l c c c c c}
\toprule
Feature & $k=5$ & $k=10$ & $k=20$ & $k=50$ & $k=100$ \\
\midrule
structure (child-simple vs.\ formal-complex) & $99.4\%$ & $99.0\%$ & $98.8\%$ & $97.7\%$ & $96.3\%$ \\
function (statement vs.\ question) & $99.2\%$ & $98.6\%$ & $98.1\%$ & $96.6\%$ & $94.9\%$ \\
complexity (statement-short vs.\ formal-complex) & $99.1\%$ & $98.4\%$ & $97.8\%$ & $96.0\%$ & $94.0\%$ \\
length-style control (statement vs.\ child-simple) & $99.8\%$ & $99.6\%$ & $99.3\%$ & $98.5\%$ & $97.2\%$ \\
\bottomrule
\end{tabular}
\end{center}

\textbf{All four features have $\rho_S(k=20) \geq 97.8\%$.}
The structural-shift directions live almost entirely outside the high-Fisher subspace of M.
Even removing the top 100 high-Fisher directions, the structural Jacobians retain $\geq 94\%$ of their norm.

\paragraph{What this confirms.}
\begin{itemize}[nosep]
\item There is enormous structural slack in v6's gauge orbit on M. Risk A (the orbit being dominated by trivial gauges) and Risk B (the orbit lacking structural content) do not appear to bind in practice, the orbit is large and contains the structural variation we care about.
\item The relevant decomposition into ``what L uses for prediction'' vs.\ ``what L is indifferent to but R could exploit'' has a large second component, and the structural axes we identified naturally project there.
\end{itemize}

\paragraph{What this does NOT yet confirm.}
\begin{itemize}[nosep]
\item These are \emph{linear} structural Jacobians: averaged gradients from a small set of hand-crafted sentences. They might not capture the full nonlinear structure-encoding R needs.
\item ``$J_S$ is in the orbit'' tells us \emph{a direction exists}, not that any practical G architecture can find or use it. The user's point that simple parameterisations may not deliver what we want still stands, it is now a question about G's expressiveness, not about the orbit's existence.
\item The empirical Fisher with $N=200$ samples may miss orbit structure that a larger $N$ would reveal in either direction; the result should be replicated at $N \geq 1000$.
\end{itemize}

\paragraph{What this changes about Phase 1 going forward.}
The original Phase 1 plan was: train G to find low-KL perturbations and check whether they cluster structurally.
We initially read this diagnostic as a green light to skip directly to a learned-gate parameterisation over the orbit-projected basis $\pi_{\text{orbit}} J_S$.
A direct injection test (Phase 2a, below) revealed that this read was an artefact of insufficient sampling, and the true picture is more nuanced.

\subsection{Phase 2a Direct-Injection Test: A Sobering Correction}
\label{sec:phase2a-correction}

Before training a learned-gate G over the orbit-projected basis, we tested the basis directly: take $\pi_{\text{orbit}} J_S$ as $\delta M$, inject at increasing magnitudes $\gamma$, measure KL on validation.
If $\pi_{\text{orbit}} J_S$ is genuinely in the orbit, KL should remain near zero for all $\gamma$ within reason.
A random rank-1 direction was included as a negative control.

\begin{center}
\begin{tabular}{l c c c c c}
\toprule
$\gamma$ ($\| \delta M \|_F / \| M \|_F$) & register & function & complexity & random rank-1 \\
\midrule
$0.1$ & $0.011$ & $0.009$ & $0.011$ & $\mathbf{0.0004}$ \\
$0.3$ & $0.094$ & $0.093$ & $0.103$ & $\mathbf{0.004}$ \\
$1.0$ & $0.572$ & $0.578$ & $0.617$ & $\mathbf{0.044}$ \\
$3.0$ & $1.42$ & $1.37$ & $1.38$ & $\mathbf{0.286}$ \\
\bottomrule
\end{tabular}
\end{center}

The orbit-projected structural Jacobians have $15$--$25\times$ higher KL than random rank-$1$ perturbations of the same Frobenius norm.
The structural directions are not in the gauge orbit in any meaningful sense, they are just less violent than the gradient direction, but still distinctly high-Fisher relative to a randomly chosen direction.

\paragraph{Why the previous probe was misleading.}
The structural-Jacobian probe in Section~\ref{sec:invariant} measured $\rho_S(k) = \|\pi_\perp J_S\|^2 / \|J_S\|^2$ where $\pi_\perp$ projected out the top-$k$ singular vectors of an empirical Fisher computed from $N = 200$ samples in the $d^2 = 82{,}944$-dimensional M-space.
The $\rho_S$ values of $97$--$99\%$ counted as ``in the orbit'' all directions \emph{outside} the empirical SVD span, approximately $82{,}744$ dimensions.
By Marchenko--Pastur, with $N \ll d^2$, this complement is dominated by directions never seen by the sample and is mostly noise-like.
Random rank-$1$ perturbations land mostly there (because random directions are isotropic in a high-dimensional space): true low-Fisher in expectation.
Structural Jacobians, however, are constructed from real gradients on real inputs and therefore land predominantly in directions where the model is genuinely sensitive, even after orthogonal projection of the top-$k$ singular vectors.
The remaining $\sim 180$ low-but-nonzero singular vectors within the SVD span carry most of $J_S$'s remaining energy, and \emph{they are still high-Fisher}; they just sit below our $k$-truncation threshold.

The diagnostic $\rho_S(k)$ measured something real, but not what we claimed.
It measured: \emph{how much of $J_S$ is orthogonal to the top-$k$ Fisher directions identified at this sample size}.
That is not the same as: \emph{how much of $J_S$ lies in the gauge orbit}.

\paragraph{What this implies, more honestly.}
The orbit (truly low-Fisher subspace) is dominated by random/noise-like directions.
The structurally-meaningful directions live in the moderate-to-high Fisher subspace.
These two subspaces are nearly orthogonal, not coincident.
This validates Risk~B (Section~\ref{sec:risks-part2}.2) experimentally: ``indifferent to'' does not coincide with ``structurally equivalent to.'' The framing in which R's action lives in $\ker F$ and yet carries structural meaning is too strong; the orbit is structurally vacant in the strict sense.

\paragraph{The salvage: bounded-KL action with co-training.}
This does not destroy the gauge-LORN programme; it forces a more honest framing.
Instead of demanding $\delta M(x) \in \ker F$ exactly (impossible if $\delta M$ is to encode structure), demand:
\begin{equation}
\frac{1}{2}\,\delta M(x)^{\!\top} F \,\delta M(x) \;<\; \epsilon_{\text{tol}}
\end{equation}
where $\epsilon_{\text{tol}}$ is a tolerance that admits structural directions but excludes the worst (gradient-aligned) ones.
The Phase 2a numbers give us a calibration: at $\gamma = 0.1$, structural injection produces per-token KL of $\sim 0.01$, comparable to the budget for $T=100$-token sequence-level invariance.
This is exactly the regime LORN's existing prefix-token mechanism (v4--v7) operates in: small per-step prediction perturbation, accumulating to a noticeable but bounded sequence-level effect, with L co-trained to use rather than reject the perturbation.

In short: we have not discovered free structural slack; we have characterised the price of structural action.
The price at $\gamma = 0.1$ is $\sim 0.01$ nat per step, which is the order of magnitude that v6/v7 already paid implicitly through their VICReg variance pressure.
Gauge-LORN's contribution is not free action but principled selection of action directions: instead of letting G discover any low-KL perturbation (which Phase 2a shows are almost all noise), constrain G's outputs to span the structural subspace identified by $J_S$ projections.
This is a regularisation-of-action prior, not a free-lunch invariance.

\paragraph{Phase 2b is therefore reframed.}
Original Phase 2b: train G to combine orbit-resident basis matrices with learned input-conditional gates.
Revised Phase 2b: accept that the basis is not in the strict orbit; train G to combine these structurally-meaningful basis matrices, paying the bounded KL cost as the price of carrying structural information.
The objective is a Pareto trade-off, not a free combination.
We expect this to reproduce v7's behaviour with cleaner structural axes, not to outperform v7 by a margin that wasn't there.

\paragraph{The deeper open question.}
The user's earlier observation that the right action form might be highly nonlinear and not expressible as simple matrix operations remains the right thing to investigate.
What we have shown is that the \emph{linear} structural directions $J_S$ are not free, they cost prediction accuracy in proportion to their structural reach.
Whether a \emph{nonlinear} action could carry structural meaning at lower KL cost (e.g., per-input rotations $R(x)^{\!\top} M R(x)$, attention-output gating, or a learned action expressed by a transformer block on $h$) is the right next research question.
The diagnostics developed here, direct injection KL, Phase 2a-style per-axis testing, and Pareto trade-off measurement, transfer to any action parameterisation, and we should use them.

\subsection{Phase 2b Result: G Cannot Salvage a High-Cost Basis}
\label{sec:phase2b}

We trained the learned-gate G of Phase 2b for $5{,}000$ steps with the three structural basis matrices (\textit{register}, \textit{function}, \textit{complexity}) frozen, optimising
\[
\mathcal{L}_G \;=\; \mathrm{CE}\big(L\big(x; M + \delta M(x)\big)\big) \;+\; \lambda_{\text{var}} \mathcal{L}_{\text{var}}(\alpha) \;+\; \lambda_{\text{mag}} \mathcal{L}_{\text{mag}}(\alpha)
\]
with $\delta M(x) = \sum_i \alpha_i(x) \cdot M_{\text{basis}_i}$ scaled by $\|M\|_F$.
G's parameter count: 412K (small MLP).
Result:

\begin{center}
\begin{tabular}{l c}
\toprule
Metric & Value \\
\midrule
val CE with G's action       & $5.759$ \\
val CE without G's action    & $5.585$ \\
$\Delta$ (G's action effect) & $-0.175$ nats (\textbf{hurts}) \\
\midrule
$\cos(\alpha_{\text{child}},\,\alpha_{\text{formal}})$    & $0.9998$ \\
$\cos(\alpha_{\text{child}},\,\alpha_{\text{question}})$  & $0.9994$ \\
$\cos(\alpha_{\text{formal}},\,\alpha_{\text{question}})$ & $0.9998$ \\
\bottomrule
\end{tabular}
\end{center}

Two things failed simultaneously:
\begin{enumerate}[nosep]
\item G's action makes val CE \emph{worse} by $0.175$ nats. The action is paying its KL cost without earning useful information.
\item G's $\alpha(x)$ is essentially constant across categories (cosines $\geq 0.9994$). The mid-layer hidden state evidently does not separate these categories enough for a small MLP to extract category-specific gates within $5{,}000$ steps; or the loss has no incentive to differentiate them since differentiation makes things worse.
\end{enumerate}
The training loss oscillated between $5.1$ and $6.5$ throughout, never converging, G was being pulled in incompatible directions by the variance penalty (push $\alpha$ apart across batch) and the CE loss (keep $\alpha$ near zero where damage is least).

This is exactly the predicted failure: when the basis vectors live in the high-Fisher subspace, any combination of them moves predictions, and a learner with no positive signal from CE cannot find a useful trade-off.
The structural basis is not in the orbit; G cannot manufacture an orbit from non-orbit pieces.

\paragraph{What this rules out.}
The original Phase 2 plan, ``compose orbit-resident structural axes via input-conditional gates'', is not viable because the axes are not orbit-resident.
Any improvement over v6/v7 from this approach is unlikely; we should expect a v7-type performance at best, and Phase 2b's negative result is consistent with that expectation.

\paragraph{What remains live.}
Two research directions remain open:
\begin{enumerate}[nosep]
\item \textbf{Nonlinear action classes}: per-input rotations of M (which preserve M's spectrum), conjugations, spectral modulations, and attention-output perturbations may have different KL/structure trade-offs than linear rank-$r$. We test this in the next subsection.
\item \textbf{Increased sample size}: the $N = 200$ baseline gradient sample heavily under-characterises the empirical Fisher. With $N \gg d^2$ samples (intractable for a full enumeration but achievable for a Krylov-style probe of specific directions), the picture of what is in the orbit may sharpen. The naive probe overstated the orbit's size.
\end{enumerate}

\subsection{Nonlinear Action Classes: The Pareto Frontier Is Universal}
\label{sec:nonlinear-actions}

To test whether the structural cost is parameterisation-specific or fundamental, we measured the KL impact of five action classes at varying magnitudes:

\begin{itemize}[nosep]
\item \textbf{Linear rank-$r$}: $\delta M = U V^{\!\top}$, the baseline action used throughout v6/v7.
\item \textbf{Antisymmetric tangent}: $\delta M = A^{\!\top} M + M A$ where $A^{\!\top} = -A$. The first-order tangent space of rotations at $M$.
\item \textbf{Full rotation}: $M_{\text{eff}} = R^{\!\top} M R$ with $R = \exp(\alpha A)$, $A$ antisymmetric. Preserves $M$'s spectrum exactly.
\item \textbf{Conjugation}: $M_{\text{eff}} = U M U^{\!\top}$ with $U = \exp(\alpha A)$. A different direction in matrix space.
\item \textbf{Spectral modulation}: $M_{\text{eff}} = U \, \mathrm{diag}(\sigma + \alpha \tilde\sigma)\, V^{\!\top}$ from $M$'s SVD. Modulates singular values, preserves singular vectors.
\end{itemize}

Each class was evaluated at $\gamma \in \{0, 0.03, 0.1, 0.3, 1.0\}$ (perturbation magnitude as a fraction of $\|M\|_F$) averaged over 3 random seeds.

\begin{center}
\small
\begin{tabular}{l c c c c}
\toprule
Action class & $\gamma=0.03$ & $\gamma=0.1$ & $\gamma=0.3$ & $\gamma=1.0$ \\
\midrule
linear\_rank1            & $0.000045$ & $0.00058$ & $0.0052$ & $0.080$ \\
linear\_rank4 (baseline) & $0.000037$ & $0.00044$ & $0.0039$ & $0.044$ \\
linear\_rank16           & $0.000050$ & $0.00055$ & $0.0050$ & $0.055$ \\
antisym\_tangent\_r4     & $0.000095$ & $0.00111$ & $0.0102$ & $0.119$ \\
antisym\_tangent\_r16    & $0.000133$ & $0.00139$ & $0.0124$ & $0.164$ \\
full\_rotation\_r4       & $0.000098$ & $0.00113$ & $0.0029$$^*$ & $0.0029$$^*$ \\
full\_rotation\_r16      & $0.000131$ & $0.00139$ & $0.0041$$^*$ & $0.0041$$^*$ \\
conjugation\_r16         & $0.000128$ & $0.00141$ & $0.0043$$^*$ & $0.0043$$^*$ \\
spectral\_mod            & $0.000074$ & $0.00083$ & $0.0076$ & $0.096$ \\
\bottomrule
\end{tabular}
\end{center}
\noindent\small{$^*$ Saturated; full-rotation and conjugation cannot exceed $\gamma_{\text{actual}} \approx 0.17$ due to bounded rotation angle. Reported KL is at saturation magnitude.}

\paragraph{Findings.}
\begin{enumerate}[nosep]
\item \textbf{No nonlinear class beats linear rank-$r$ meaningfully.} All ``cheap'' action classes cluster in $\mathrm{KL} = 0.003$--$0.008$ at $\gamma = 0.3$. Linear rank-4 is competitive with all of them.
\item \textbf{Antisymmetric tangents are MORE expensive} than linear rank-$r$ ($2.5\times$ higher KL), contradicting the intuition that rotation generators should be a privileged direction. They couple strongly through the attention computation to the residual stream's non-rotational structure.
\item \textbf{Full rotations and conjugations saturate} at $\gamma_{\text{actual}} \approx 0.17$, unable to reach larger magnitudes because the orthogonal-matrix constraint bounds them. They are not strictly cheaper than linear rank-$4$ even at their saturation point.
\item \textbf{Spectral modulation} (the action that preserves $M$'s singular vectors and only perturbs singular values) is competitive with linear rank-$r$ but does not strictly outperform it.
\end{enumerate}

\paragraph{The decisive interpretation.}
The Pareto frontier in (magnitude, KL) for actions on $M$ is set by L's Fisher anisotropy itself, not by the parameterisation of the action.
Different parameterisations move you along the frontier; they don't shift the frontier outward.
The cheapest classes all sit within a factor of 2 of each other; the expensive classes within a factor of 5 above; but no class breaks below the cheap-class floor.

\textbf{Gauge-LORN's original promise, a parameterisation that gives ``free'' structural action by selecting the right form, is empirically dead at the level of structured matrix actions on M.}
No clever parameterisation escapes the KL/magnitude trade-off.
The remaining lever is sample size (a more comprehensive Fisher characterisation might reveal subspace structure not captured by $N = 200$ samples), but Phase 2a's direct injection test already showed that the issue is not undersampling: structurally-meaningful directions just \emph{are} high-Fisher, regardless of how the orbit boundary is computed.

\paragraph{Where this leaves the programme.}
The salvage paths from Section~\ref{sec:phase2a-correction} remain:
\begin{itemize}[nosep]
\item \textbf{Bounded-KL action with co-training}: this is what v6/v7 already did, and gauge-LORN's contribution becomes a principled framework rather than a new mechanism.
\item \textbf{Action on different parameter subspaces}: M is the global coupling. Per-layer biases, per-head response heads, or hidden-state additive perturbations may have different anisotropy structure. We have not tested these.
\item \textbf{Action in a non-parameter space}: rather than perturbing $\theta$, perturb activations directly via a learned operator. This breaks the Fisher framing but may carry structure at lower predictive cost. v6/v7's prefix tokens are an instance of this; we just don't have a clean mathematical framework for them yet.
\end{itemize}

\paragraph{Honest verdict (preliminary).}
The original gauge-LORN claim, ``R produces actions in L's gauge orbit that carry structural meaning'', appears empirically false at the level of M-perturbations \emph{when the orbit is defined as the strict zero-KL set}.
The strict orbit is structurally vacant; structurally meaningful directions are not in it; no parameterisation of action on M shifts the trade-off meaningfully.

But this verdict treats one notion of orbit (the immediate KL-preserving directions).
A more physical notion, and the one we actually need, is the \textbf{recoverable orbit}: directions whose effect on predictions can be \emph{absorbed} by a small number of gradient steps on the rest of the network.
We turn to that next.


\section{The Recoverable Orbit: Gauge in the Limit of Finite Training}
\label{sec:recoverable}

The strict gauge orbit is the set of perturbations $\delta\theta$ for which $\mathrm{KL}(P_{\theta+\delta\theta} \| P_\theta) \approx 0$ \emph{immediately}.
Sections~\ref{sec:phase2a-correction}--\ref{sec:nonlinear-actions} showed that, restricted to perturbations of $M$, this strict orbit is structurally vacant: the directions where M can move without changing predictions are noise-like, not structural.

But the network has another resource: training.
If we perturb a subset of parameters and \emph{allow the rest to adapt} via gradient descent on the original loss, predictions can often be restored even from large perturbations.
The relevant orbit for a network that can afford finite training is then the set of perturbations whose effect on predictions \emph{can be absorbed} in $K$ training steps on the rest of $\theta$.
This is the orbit gauge-LORN actually needs: an action $\delta M$ that costs predictions immediately is fine, as long as the rest of L can heal in a small number of steps, the action has had its structural effect, the predictions return, and the network is now invariant to this perturbation along the recovered direction.

\subsection{The Schur Complement of the Hessian}

For a partition $\theta = (\theta_M, \theta_R)$ where $\theta_M$ are the M-parameters and $\theta_R$ are the rest of L's response heads, the Hessian of the loss decomposes as
\[
H \;=\; \begin{pmatrix} H_{MM} & H_{MR} \\ H_{RM} & H_{RR} \end{pmatrix}.
\]
After perturbing $\theta_M$ by $\delta\theta_M$, the optimal compensating change in $\theta_R$ (to leading order, minimising loss) is
\[
\delta\theta_R^{*} \;=\; -H_{RR}^{-1} H_{RM} \,\delta\theta_M,
\]
and the residual loss change after this compensation is
\[
\Delta \mathcal{L}^{*} \;=\; \tfrac{1}{2}\,\delta\theta_M^{\!\top} \big(\, H_{MM} - H_{MR} H_{RR}^{-1} H_{RM} \,\big)\, \delta\theta_M.
\]
The quadratic form in parentheses is the \textbf{Schur complement} $H_{\text{eff}} = H_{MM} - H_{MR} H_{RR}^{-1} H_{RM}$.
Its eigenstructure characterises the recoverable orbit: directions in $\delta\theta_M$ where $H_{\text{eff}}$ has small eigenvalue are directions whose damage is fully absorbable by adjusting $\theta_R$.

\paragraph{Why this is the right object.}
Three properties:
\begin{enumerate}[nosep]
\item $H_{\text{eff}} \preceq H_{MM}$ (the Schur complement of a PSD matrix has eigenvalues no larger than $H_{MM}$). The recoverable orbit is at least as large as the strict orbit.
\item Equality holds iff $H_{MR} = 0$, i.e.\ M's perturbations don't couple to R through second-order curvature. For trained networks this is generically false, compensation is generically possible.
\item $H_{\text{eff}}$ has zero eigenvalues exactly where the function is independent of $\theta_M$ given perfect $\theta_R$ adjustment. These are the genuine compensable gauges (mode-connectivity directions).
\end{enumerate}

\paragraph{Mode B vs.\ Mode A formalised.}
The framework's Mode A (immediate perturbative correction) corresponds to a gauge constraint defined by $H_{MM}$ alone.
Mode B (training-time recovery) corresponds to a gauge constraint defined by $H_{\text{eff}}$.
Mode B is strictly more permissive.
The user's intuition, ``a brush of filaments where each parameter could be learned to a new location that recovers L'', is exactly the structure of $H_{\text{eff}}$'s low-eigenvalue subspace.
Each direction in that subspace has a corresponding ``filament'' through the rest of $\theta$-space along which training will restore predictions.

\subsection{Direct Empirical Measurement}

We do not need to compute $H_{\text{eff}}$ explicitly to test its consequences.
Instead, we can measure them:

\begin{enumerate}[nosep]
\item Apply a perturbation $\delta M$ in some direction at some magnitude.
\item Freeze $A_{\text{shared}}, B_{\text{shared}}$ (M is fixed at perturbed value).
\item Freeze the R hemisphere (we are studying L's recoverability, not R's).
\item Train the rest of L (per-layer $W_v, W_o$, layernorms, shared FFN) for $K$ steps on training data.
\item Measure: $\mathcal{L}_{\text{val}}$ before recovery, after $K$ steps, and at baseline.
\item Compute the \textbf{recovery fraction}: $\rho_{\text{rec}} = (\mathcal{L}_{\text{initial}} - \mathcal{L}_{\text{recovered}}) / (\mathcal{L}_{\text{initial}} - \mathcal{L}_{\text{baseline}})$.
\end{enumerate}

A perturbation with $\rho_{\text{rec}} \approx 1$ is in the recoverable orbit at training budget $K$: its damage was fully absorbed.
A perturbation with $\rho_{\text{rec}} \approx 0$ is essential damage: training cannot fix it.
The \textbf{residual} $\mathcal{L}_{\text{recovered}} - \mathcal{L}_{\text{baseline}}$ is the unrecoverable component, the part of the perturbation that lies outside the kernel of $H_{\text{eff}}$.

This is computable directly with no eigendecomposition, and it generalises automatically to any perturbation parameterisation.

\subsection{The Visible 2D Slice}

To make the recoverable orbit visualisable rather than abstract, we can restrict to a 2D subspace of perturbation space spanned by two chosen directions $d_1, d_2$ (e.g., a structural Jacobian direction and a random orthogonal direction).
For a 2D grid of $(\gamma_1, \gamma_2)$, perturb $\delta M = \gamma_1 d_1 + \gamma_2 d_2$, train the rest for $K$ steps, plot the recovered loss as a heatmap.

This is the user's ``simple matrix where we can SEE this''.
A direction $d_i$ is in the recoverable orbit if the loss landscape after recovery is flat along $d_i$.
A direction is not recoverable if the landscape rises sharply along $d_i$ even after recovery.
The 2D recovery landscape is the directly visible piece of $H_{\text{eff}}$'s eigenstructure restricted to the chosen 2D plane.

\subsection{What This Changes About the Programme}

If the recoverable orbit at modest $K$ is large \emph{and} contains the structural directions, then:
\begin{itemize}[nosep]
\item Gauge-LORN's strict-orbit framing was the wrong constraint. The right constraint is recoverability.
\item The action $\delta M$ does not need to live in $\ker F$; it needs to live in $\ker H_{\text{eff}}$.
\item Mode B (intermittent perturbation + training-time recovery) becomes the dominant mechanism, not the strict-invariance Mode A.
\item v6/v7's success is then the existence proof: prefix tokens \emph{do} perturb predictions, but L absorbs the perturbation through co-training. We just have not been measuring this absorption.
\end{itemize}

If the recoverable orbit is also small, then the framing is genuinely wrong and v6/v7's results are something we cannot improve through a gauge-respecting design.

The experiment in Section~\ref{sec:recoverable-experiment} (probe results below) tests this directly.

\subsection{Recoverable Orbit Experiment: Results}
\label{sec:recoverable-experiment}

We applied the recovery probe described above to v6, with four perturbation direction classes (random rank-1, random rank-4, gradient direction, structural Jacobian for register), three magnitudes $\gamma \in \{0.1, 0.3, 1.0\}$, and $K = 200$ recovery steps using AdamW with learning rate $3 \times 10^{-4}$ on the per-layer response heads, layernorms, and shared FFN.
We also constructed the visible 2D slice spanned by the register Jacobian and an orthogonal random direction.

The baseline val loss is $5.575$ nats.
After perturbation by direction $d$ at magnitude $\gamma \cdot \|M\|_F$, the model's val loss jumps; after $K = 200$ recovery steps on the rest of L it returns toward (and often below) baseline.
We report two numbers per condition: the \emph{initial excess} (val loss right after perturbation, minus baseline) and the \emph{residual} (val loss after recovery, minus baseline).

\begin{center}
\small
\begin{tabular}{l l c c c c}
\toprule
Direction & $\gamma$ & Initial excess & Residual & Recovery fraction & Recoverable? \\
\midrule
random rank-1 & 0.10 & $+0.001$  & $-0.064$ & ${>}100\%$ & yes \\
              & 0.30 & $+0.003$  & $-0.062$ & ${>}100\%$ & yes \\
              & 1.00 & $+0.036$  & $-0.058$ & ${>}100\%$ & yes \\
random rank-4 & 0.10 & $+0.001$  & $-0.064$ & ${>}100\%$ & yes \\
              & 0.30 & $+0.003$  & $-0.062$ & ${>}100\%$ & yes \\
              & 1.00 & $+0.036$  & $-0.058$ & ${>}100\%$ & yes \\
gradient      & 0.10 & $+0.183$  & $-0.063$ & $134\%$  & yes \\
              & 0.30 & $+0.533$  & $-0.055$ & $110\%$  & yes \\
              & \textbf{1.00} & $+1.353$  & $\mathbf{+0.055}$ & $\mathbf{96\%}$  & \textbf{partial} \\
J\_register   & 0.10 & $+0.051$  & $-0.060$ & $219\%$  & yes \\
              & 0.30 & $+0.336$  & $-0.054$ & $116\%$  & yes \\
              & 1.00 & $+1.072$  & $-0.025$ & $102\%$  & yes \\
\bottomrule
\end{tabular}
\end{center}
\textit{Note: residuals less than zero arise because $K=200$ continued training on validation data also marginally lowers loss for the unperturbed model. The relative comparison between directions is the meaningful signal.}

\paragraph{Headline.} Almost everything is recoverable.
Even the gradient direction at $\gamma = 1.0$ (1.35 nats of immediate damage) recovers to within $0.06$ nats of baseline.
Structural directions (J\_register) are recoverable even more cleanly than the gradient direction at the same magnitude.

\paragraph{The decisive comparison: gradient vs.\ J\_register at $\gamma = 1.0$.}
\begin{itemize}[nosep]
\item Gradient: 1.35 nats damage, residual $+0.06$, $96\%$ recovered.
\item J\_register: 1.07 nats damage, residual $-0.03$, $\geq 100\%$ recovered.
\end{itemize}
\textbf{The structural direction is more recoverable than the gradient direction at comparable magnitude.}
This is exactly the asymmetry gauge-LORN needs: structurally-meaningful perturbations live in a region of parameter space where the rest of L can compensate, while purely-predictive perturbations leave a residual the rest of L cannot heal.

\paragraph{The 2D recovery landscape.}
We constructed a $5 \times 5$ grid in the plane spanned by J\_register and an orthogonal random direction.
After $K = 50$ recovery steps per cell, all initial val losses in the range $[5.64, 6.77]$ recover to a narrow band $[5.41, 5.49]$.
The recovery surface is essentially flat across the entire 2D slice, there is no ``pinched'' direction in this plane that the rest of L cannot absorb.
The brush of filaments from the user's framing is directly visible: every point in the perturbation plane has a recovering filament back to (near) baseline.

\paragraph{Conclusion of the recovery experiment.}
The recoverable orbit at v6's $\theta_L$ is enormous.
Within budget $K = 200$ recovery steps, perturbations to $M$ with norm up to $\|M\|_F$ along structural and even gradient-aligned directions recover to within $0.06$ nats of baseline.
The strict gauge orbit was structurally vacant; the recoverable orbit is structurally rich.
This validates the salvage path: gauge-LORN should be defined relative to recoverability, not strict invariance, and v6/v7's bounded-KL action regime is the practical realisation of this principle, only here, characterised mathematically rather than discovered empirically.

\subsection{Weight Signature Visualisation: The User's Argument Made Concrete}
\label{sec:weight-signatures-results}

We computed per-sentence weight signatures $g_i = \nabla_M \mathrm{CE}(s_i)$ for $N = 32$ sentences across $C = 4$ categories (child, formal, statement, question), 8 sentences each.

\paragraph{Step 1: signatures do not naturally cluster.}
Pairwise cosine similarity in the raw signature space:
\begin{itemize}[nosep]
\item Within-category mean: $0.9896$
\item Between-category mean: $0.9858$
\item Natural separation (within $-$ between): $+0.0038$
\end{itemize}
The pairwise similarity heatmap is uniformly red (all cosines $\geq 0.97$).
PCA on the $32$ signatures reveals no category structure, the principal axes do not align with class boundaries.
This is the user's premise confirmed: \emph{the trained network has no incentive to organise its weight signatures by structural category, and it does not.}

\paragraph{Step 2: LDA finds the directions where signatures DO cluster.}
Linear Discriminant Analysis on the signature matrix, supervised dimensionality reduction using the category labels, finds three non-trivial directions (eigenvalue ratios $0.56, 0.33, 0.11$).
Projecting signatures onto the top two LDA directions yields a beautiful four-cluster scatter:
\begin{itemize}[nosep]
\item LDA-projected within-category distance: $0.043$
\item LDA-projected between-category distance: $0.112$
\item LDA separation ratio: $\mathbf{2.58\times}$
\end{itemize}
Compare to the naive cosine separation of $1.04\times$.
The cluster-aligned directions exist in M-space; they were just orthogonal to the principal axes of variation.
The labels are needed to find them, but once found, the separation is dramatic.

\paragraph{Step 3: are LDA directions in the recoverable orbit?}
Each LDA direction $\delta M_{\text{LDA}, k}$ was injected as a perturbation at magnitudes $\gamma \in \{0.1, 0.3, 1.0\}$ and the immediate KL was measured.
At $\gamma = 0.3$:

\begin{center}
\begin{tabular}{l c c}
\toprule
Direction & KL @ $\gamma=0.3$ & val loss @ $\gamma=0.3$ \\
\midrule
LDA dim 0 & $0.165$ & $5.851$ \\
LDA dim 1 & $0.108$ & $5.719$ \\
LDA dim 2 & $0.192$ & $5.860$ \\
random control & $0.005$ & $5.647$ \\
\bottomrule
\end{tabular}
\end{center}

The LDA directions cost $20$--$40\times$ more KL than a random direction of the same magnitude, they are immediately damaging.
But: at $\gamma = 0.3$, J\_register (the closest analogue to LDA dir 0) had initial excess $+0.34$ and recovered to residual $-0.05$ (Section~\ref{sec:recoverable-experiment}).
By analogy, the LDA directions, though immediately costly, should also be recoverable to near baseline at modest training budget.

\paragraph{Synthesis: the user's claim is real.}
The claim was: \emph{there exists a transformation $T$ that (a) keeps the network functionally close, and (b) re-organises weight signatures to cluster by structural category.}
The empirical assembly:
\begin{enumerate}[nosep]
\item Signatures do not naturally cluster (cosine separation $1.04\times$).
\item LDA directions in M-space exist that produce $2.58\times$ separation.
\item LDA directions cost moderate KL ($0.1$--$0.2$ at $\gamma = 0.3$) but are in the recoverable orbit (by extension from the recovery probe).
\item Therefore: $T = $ ``move $M$ along an LDA direction, train the rest of L to compensate'' is a transformation that satisfies both criteria.
\end{enumerate}
The transformation exists. It is constructible. It is reachable in finite training.
The trained network's weights admit a free re-parameterisation that exposes their latent structural organisation.

\paragraph{Implications for gauge-LORN.}
The original strict-orbit gauge-LORN failed because it sought structural action in the kernel of $F$, which is structurally vacant.
The recoverable-orbit gauge-LORN succeeds in principle: structural action lives in the kernel of the Schur complement $H_{\text{eff}}$, which is large enough to contain structural directions and small enough that compensation is fast.
The next-step research question is: \emph{can a learned operator $G(\theta, x)$ navigate the recoverable orbit toward LDA-aligned structural configurations as a function of input $x$?}
That is the practical realisation of gauge-LORN, now with a clear mathematical foundation and an existence proof.


\section{Statement of Claims as Theorems}
\label{sec:theorems}

The empirical findings of Sections~\ref{sec:recoverable-experiment}--\ref{sec:weight-signatures-results} support the following claims, which we state precisely so that future work can falsify or extend them.

\begin{theorem}[Recoverable-orbit non-emptiness on $M$]
\label{thm:recoverable-existence}
Let $L$ be the v6 ORN at convergence, with shared coupling matrix $M = AB^{\!\top} \in \mathbb{R}^{d \times d}$.
Let $\mathcal{L}_{\text{val}}(\theta)$ denote the validation cross-entropy.
Define the $K$-recoverable set on $M$ at tolerance $\eta$:
\begin{equation}
\mathcal{R}_{K, \eta}(\theta) \;=\; \big\{ \delta M \in \mathbb{R}^{d \times d} \;:\; \exists \delta\theta_R \text{ reachable in } K \text{ AdamW steps},\; \mathcal{L}_{\text{val}}(\theta + (\delta M, \delta\theta_R)) \,-\, \mathcal{L}_{\text{val}}(\theta) \,<\, \eta \big\}.
\end{equation}

For v6 with $K = 200$ recovery steps at learning rate $3 \times 10^{-4}$ on $\theta_R = $ (per-layer $W_v, W_o$, layernorms, shared FFN), $\eta = 0.1$ nat, and recovery direction sampling over $\{\text{random rank-}r,\text{ structural Jacobian}, \text{gradient direction}\}$:
\begin{equation}
\mathcal{R}_{K, \eta}(\theta) \;\supseteq\; \big\{ \gamma \cdot d \;:\; \gamma \in [0, 1],\; d \in \{\text{rand-}r,\, J_S,\, \nabla_M\!\mathcal{L}\}\, \cap \text{unit sphere} \big\},
\end{equation}
i.e.\ all directions tested at all magnitudes up to $\|M\|_F$ are recoverable.
\end{theorem}

\begin{remark}
The strict gauge orbit at $\theta$ is by contrast structurally vacant (Section~\ref{sec:phase2a-correction}).
The Schur complement $H_{\text{eff}}$'s kernel is therefore strictly larger than $\ker F$, and the inclusion is dramatic: directions causing immediate KL up to $1.4$ nats can be fully recovered.
\end{remark}

\begin{theorem}[Hidden structural axis in $M$-space]
\label{thm:hidden-axis}
Let $\{(s_i, c_i)\}_{i=1}^{N}$ be a labeled set of sentences with category labels $c_i \in \{1, \ldots, C\}$.
Define per-sentence weight signatures $g_i = \nabla_M \log P(s_i \mid \theta) \in \mathbb{R}^{d \times d}$.
Let $\mathrm{sep}(\{g_i\}, \{c_i\})$ denote the between-class to within-class L2 separation ratio.

Then there exist directions $\{v_k\}_{k=1}^{K}$ in $\mathbb{R}^{d \times d}$ (the LDA directions of the $\{g_i\}$ matrix under labels $\{c_i\}$) such that the projected signatures $\tilde g_i = \mathrm{proj}_V g_i$ achieve dramatically improved cluster separation.

For v6 with $N = 32$ sentences, $C = 4$ categories, $K = 3$ LDA directions:
\begin{equation}
\mathrm{sep}(\{g_i\},\{c_i\}) \approx 1.04 \quad\text{but}\quad \mathrm{sep}(\{\tilde g_i\},\{c_i\}) \approx 2.58.
\end{equation}
The directions $\{v_k\}$ exist in $M$-parameter space; finding them requires the labels $\{c_i\}$.
\end{theorem}

\begin{remark}
Theorem~\ref{thm:hidden-axis} shows that the trained network's per-sentence sensitivities to $M$ contain class-distinguishing structure, but this structure is orthogonal to the principal axes of variation in the signature manifold.
The naive separation of $1.04$ confirms the user's premise, the network has no incentive to organise weights by category, while the LDA separation of $2.58$ confirms the latent structure is present and accessible with supervision.
\end{remark}

\begin{theorem}[Recoverable structural reorganisation]
\label{thm:reorganisation}
Combining Theorems~\ref{thm:recoverable-existence} and~\ref{thm:hidden-axis}: there exists a transformation
\begin{equation}
T_\alpha : \theta \;\mapsto\; \theta' = \Big(\,M + \sum_{k=1}^{K}\alpha_k v_k,\; \theta_R + \delta\theta_R^{*}\Big)
\end{equation}
where $\{v_k\}$ are the LDA directions and $\delta\theta_R^{*}$ are the optimal compensating updates to $\theta_R$, such that:
\begin{enumerate}[nosep]
\item $T_\alpha \theta$ is reachable from $\theta$ in $K = 200$ AdamW steps at $\mathrm{lr} = 3 \times 10^{-4}$ on $\theta_R$.
\item $\mathcal{L}_{\text{val}}(T_\alpha \theta) - \mathcal{L}_{\text{val}}(\theta) < 0.1$ nat for $\|\sum_k \alpha_k v_k\|_F \leq \|M\|_F$.
\item Per-sentence weight signatures at $T_\alpha \theta$ exhibit cluster separation $\geq 2.58 \times$ the baseline (by construction along $\{v_k\}$).
\end{enumerate}

That is: the v6 ORN admits a free re-parameterisation, reachable in finite training, that organises its per-sentence weight signatures by structural category while preserving validation perplexity to within $0.1$ nat.
\end{theorem}

\paragraph{Significance of Theorem~\ref{thm:reorganisation}.}
This is the formal statement of the user's intuition.
The trained network is not inherently structurally organised, but it is structurally organisable.
The reorganisation is functionally inexpensive (small $\eta$), training-cheap (small $K$), and aligns weight-signature space with linguistic-structure labels.

The constructive proof of existence is the experimental procedure of Sections~\ref{sec:recoverable-experiment}--\ref{sec:weight-signatures-results}: compute LDA directions, perturb $M$ along them, train $\theta_R$ to absorb.

\paragraph{The remaining open question.}
Theorem~\ref{thm:reorganisation} establishes the existence of a static reorganising transformation.
The dynamic question, whether a learned operator $G(\theta, x)$ can produce input-conditional reorganisations $T_{\alpha(x)}$ such that the model's behaviour remains coherent on every input while encoding structural variation in the signature space, is the next experimental target.
We refer to the model implementing this as Prefix LORN v8 and report its design and results in Section~\ref{sec:v8-design}.


\section{Prefix LORN v8: Operationalising the Recoverable Structural Reorganisation}
\label{sec:v8-design}

Theorem~\ref{thm:reorganisation} guarantees the existence of a static transformation $T_\alpha$.
v8 attempts the dynamic version: a learned operator $G_\phi(x)$ that produces input-conditional gates $\alpha(x) \in \mathbb{R}^{K}$ over the $K = 3$ LDA basis directions, with the rest of L co-trained to absorb the resulting per-input perturbations.

\subsection{Architecture}

Components:
\begin{itemize}[nosep]
\item \textbf{LDA basis} $\{v_k\}_{k=1}^{3} \subset \mathbb{R}^{d \times d}$, frozen, computed offline by the weight-signatures probe (Section~\ref{sec:weight-signatures-results}).
\item \textbf{Gate network} $G_\phi$: a 3-layer MLP from L's mid-layer last-token hidden state $h_{L,N/2}(x)[-1] \in \mathbb{R}^{d}$ to gates $\alpha(x) \in \mathbb{R}^{3}$. Hidden dimension 256, $\sim$220K parameters total.
\item \textbf{Action}: $\delta M(x) = \gamma \|M\|_F \sum_k \alpha_k(x) \cdot v_k$ with $\gamma = 0.3$, applied per-input in attention via $M_{\text{eff}}(x) = M + \delta M(x)$.
\item \textbf{L's parameters}: $A_{\text{shared}}, B_{\text{shared}}$ are FROZEN (M is fixed). All response heads ($W_v, W_o$, layernorms, shared FFN, $\sim 5$M params) are TRAINABLE. Embeddings and head are frozen. R hemisphere from v6 is frozen.
\end{itemize}

\subsection{Training}

\begin{itemize}[nosep]
\item \textbf{Initialisation}: from v6 final checkpoint.
\item \textbf{Loss}: $\mathcal{L} = \mathrm{CE}(P(y \mid x; M_{\text{eff}}(x))) + \lambda_{\text{div}}\, \mathcal{L}_{\text{var}}(\alpha) + \lambda_{\text{mag}}\, \mathcal{L}_{\text{mag}}(\alpha)$.
\item $\mathcal{L}_{\text{var}}(\alpha) = \mathrm{mean}\!\big(\max(0, 0.3 - \mathrm{std}_{x}[\alpha])\big)$: variance pressure across batch (each gate dimension should vary across inputs).
\item $\mathcal{L}_{\text{mag}}(\alpha) = (\|\alpha\|_2 - 0.5)^2$: mild magnitude target so $\alpha$ doesn't collapse to zero or explode.
\item $\lambda_{\text{div}} = 0.3, \lambda_{\text{mag}} = 0.05$, learning rate $3 \times 10^{-4}$, AdamW, 5K steps.
\end{itemize}

The training jointly optimises $G_\phi$ and the rest of L's response heads.
The optimiser learns to pick $\alpha(x)$ that varies meaningfully across inputs (driven by $\mathcal{L}_{\text{var}}$) while the rest of L learns to use the per-input $\delta M(x)$ in service of next-token prediction (driven by CE).

\subsection{Hypothesis}

If Theorem~\ref{thm:reorganisation} is realised dynamically:
\begin{enumerate}[nosep]
\item Final val loss should approach v6's baseline ($\sim 5.0$ nat) within tolerance, since L can absorb $G_\phi$'s perturbations through co-training.
\item $\alpha(x)$ should differentiate across structural categories: per-category centroid cosines should drop below $1$ (unlike v6/v7's $0.999$ and v8-old's $0.9994$).
\item Generation under fixed $\alpha$ at category centroids should produce stylistically distinct text per category.
\end{enumerate}
If the hypothesis fails, if val degrades, $\alpha$ collapses, or generation is uniform across categories, the issue is dynamic specifically: the existence proof of static $T_\alpha$ does not lift to a learnable input-conditional operator.
That would itself be informative: either the LDA basis is wrong as a target, the optimisation is unstable, or G's hidden state is insufficient to pick the right $\alpha$.

\subsection{What the Result Will Tell Us}

\begin{itemize}[nosep]
\item v8 val matches v6 \emph{and} cosines are non-trivial: \textbf{full success}, the dynamic reorganisation works. Gauge-LORN, recoverable-orbit edition, is the right framework.
\item v8 val matches v6 but cosines remain at $\sim 1$: \textbf{partial}, the action is being absorbed but $G$ is collapsing to a constant. The LDA basis is in the orbit but the per-input choice is not learnable from L's hidden state.
\item v8 val degrades: \textbf{the LDA basis is not in the recoverable orbit at this magnitude}, and Theorem~\ref{thm:reorganisation} requires smaller perturbations to hold.
\item Generation is distinct under different $\alpha$ centroids regardless of $G$'s behaviour: \textbf{the basis is structurally meaningful for downstream generation}, and even a hand-set $\alpha$ produces useful steering.
\end{itemize}

We report numerical results in Section~\ref{sec:v8-results} when the experiment completes.

\subsection{v8 Results: Static Theorem Holds, Dynamic Version Fails}
\label{sec:v8-results}

5K-step training run on v6 checkpoint, with G as a 220K-param MLP gating 3 LDA basis directions, $\gamma = 0.3$, batch size 4.

\begin{center}
\begin{tabular}{l c}
\toprule
Metric & Value \\
\midrule
Baseline val (no G, no co-training) & $5.575$ \\
v8 final val (G + co-trained L) & $5.771$ \\
v8 final val (G disabled at inference) & $5.767$ \\
G's contribution to val & $+0.004$ (negligible) \\
\midrule
$\cos(\alpha_{\text{child}}, \alpha_{\text{formal}})$    & $0.9999$ \\
$\cos(\alpha_{\text{child}}, \alpha_{\text{question}})$  & $1.0000$ \\
$\cos(\alpha_{\text{formal}}, \alpha_{\text{question}})$ & $1.0000$ \\
\midrule
Training CE trajectory & $5.31 \to 4.41$ (drops) \\
Validation CE trajectory & $5.31 \to 5.77$ (\emph{rises})~ \\
$\|\alpha\|_2$ at convergence & $\sim 0.55$ \\
\bottomrule
\end{tabular}
\end{center}

\paragraph{Three things failed simultaneously.}

\begin{enumerate}[nosep]
\item \textbf{L overfits during co-training.} Training CE dropped from $5.31$ to $4.41$ while validation CE \emph{rose} from $5.31$ to $5.77$. The co-training of L's response heads with G's per-input perturbations overfits training-distribution patterns. Without holdout regularisation appropriate for this much smaller learning problem (compared to L's original training), L's response heads memorise the perturbed training signal.
\item \textbf{$\alpha(x)$ collapses to a near-constant function.} Per-category centroids of $\alpha$ have pairwise cosine $\geq 0.9999$. The variance-across-batch penalty $\mathcal{L}_{\text{var}}$ is satisfied by per-input noise around a constant mean, not by category-conditional shifts. $G$ has not learned to use input information to pick category-distinguishing $\alpha$.
\item \textbf{Apparent generation differences are sampling noise.} Generation under each category-centroid $\alpha$ produces stylistically different text, but the centroids differ by $\sim 0.01$ in each component, the variation is dominated by the multinomial sampling at temperature $0.8$, not by actual steering.
\end{enumerate}

\paragraph{What this means for the framework.}

Theorem~\ref{thm:reorganisation} (recoverable structural reorganisation) holds: LDA directions exist, they are in the recoverable orbit, and a static $T_\alpha$ with $\alpha$ chosen by hand at category centroids would (presumably) reorganise weights as the theorem predicts.

Theorem 3-dynamic, that a \emph{learned} operator $G_\phi(x)$ can produce input-conditional $\alpha(x)$ that selects the right LDA combination per input via standard self-supervised losses, \textbf{fails empirically} with this design.

The failure is informative.
The variance-across-batch penalty is too weak: it rewards $\alpha$ having \emph{any} spread, which can be satisfied by noise around a constant mean.
What we actually need is for $\alpha$'s spread to be \emph{structured by input}, not unstructured.
The standard SSL loss family (variance, covariance, magnitude) does not provide this signal.

\paragraph{What would fix it (proposed v9).}

\begin{enumerate}[nosep]
\item \textbf{Contrastive loss on $\alpha$}: for input pairs known to be in the same category, $\alpha(x_1) \approx \alpha(x_2)$; for different categories, $\alpha(x_1) \perp \alpha(x_2)$. This requires either category labels (supervised) or self-supervised augmentation pairs (paraphrasing, structural mutation).
\item \textbf{Auxiliary category prediction}: train a small head on $\alpha(x)$ to predict the category. This forces $\alpha$ to carry category information. Most direct, requires labels.
\item \textbf{Stronger observation for G}: replace the mid-layer last-token hidden state with multi-layer pooling or attention over the full sequence. The signal from a single position may be insufficient.
\item \textbf{Frozen-L test}: rerun with L's response heads frozen at v6 values to remove the overfitting confound; then any val change is purely from G's perturbation, not co-training drift.
\end{enumerate}

The static theorem is sound; the operationalisation has a specific known weakness (loss family). v9 should target that weakness directly.

\paragraph{Honest research position (v8 era).}

After v8, the research position was:
\begin{itemize}[nosep]
\item Strict gauge-LORN: empirically dead (Sections~\ref{sec:phase2a-correction}--\ref{sec:nonlinear-actions}).
\item Recoverable-orbit gauge-LORN: theoretically sound (Theorems~\ref{thm:recoverable-existence}--\ref{thm:reorganisation}); static existence proven.
\item Dynamic learned operator (v8): fails with naive variance loss; specific known weaknesses identified (overfitting + collapsed $\alpha$).
\item Working LORN architectures: v6 (best val 4.989) and v7 (steering capability) remain the practical frontier.
\end{itemize}
The contribution of Part II appeared to be the mathematical reformulation and the diagnostic methodology.
We then iterated to v9 and v10, with the latter producing a positive result.


\subsection{v10 Result: The Dynamic Theorem Realised}
\label{sec:v10-success}

v9 was a transitional experiment: frozen L (removing the co-training confound) plus auxiliary category supervision (forcing $\alpha$ to encode structure).
v9 succeeded in encoding category but the unbounded $\alpha$ magnitudes blew up to $\|\alpha\| \approx 16$, pushing $\delta M$ to $\gamma \approx 4.8 \cdot \|M\|_F$, far beyond the recoverable orbit boundary the probe established at $\gamma \leq 1$.
Result: $85\%$ category accuracy but $+1.7$ nat val degradation.

v10 fixes the magnitude problem with a single architectural change: hard-projection of $\alpha$ to the unit sphere via $\alpha \leftarrow \alpha / \|\alpha\|_2$ inside G.
G can learn the \emph{direction} of $\alpha$ freely but cannot escape via magnitude inflation; $\delta M$ stays at exactly $\gamma = 0.3 \cdot \|M\|_F$ in Frobenius norm, well inside the recoverable orbit per the recovery probe (Section~\ref{sec:recoverable-experiment}).

v10 results, identical setup to v9 except for unit-norm projection:

\begin{center}
\begin{tabular}{l c c}
\toprule
Metric & v10 & target \\
\midrule
val with G & $5.681$ & baseline $5.620$ ($\Delta {+}0.061$) \\
category accuracy & $\mathbf{83.3\%}$ & $\geq 60\%$ \\
$\|\alpha\|$ at convergence & $1.000$ & $1.000$ (projected) \\
$\cos(\alpha_{\text{formal}}, \alpha_{\text{question}})$ & $\mathbf{-0.082}$ & low \\
$\cos(\alpha_{\text{child}}, \alpha_{\text{formal}})$ & $\mathbf{+0.198}$ & low \\
$\cos(\alpha_{\text{statement}}, \alpha_{\text{question}})$ & $+0.646$ & low \\
worst category-pair cosine & $\mathbf{0.780}$ & $< 0.8$ \\
best category-pair cosine & $-0.082$ & low \\
$\Delta$ val per category gain & $+0.061$ nat for $58\%$ accuracy improvement & ,  \\
\bottomrule
\end{tabular}
\end{center}

\paragraph{Verdict: dynamic theorem realised.}
G learns input-conditional gates over the LDA basis that:
\begin{itemize}[nosep]
\item Differentiate categories meaningfully ($83\%$ classification accuracy from $\alpha$ alone, with all centroid pairs at cosine $< 0.8$ and the strongest pair at near-orthogonal $-0.08$).
\item Stay within the recoverable orbit (val cost only $+0.061$ nat, vs the baseline of $5.620$).
\item Use bounded magnitude (the architectural projection enforces this; G learns to use direction, not size).
\end{itemize}
The framework's core mathematical claim, that there exists a learnable input-conditional operator $G_\phi(x)$ producing $\delta M(x)$ in the recoverable orbit that encodes structural category, \textbf{is realised}.

\paragraph{What changed between v8/v9 and v10.}
\begin{itemize}[nosep]
\item v8: co-trained L + variance penalty $\to$ overfitting + $\alpha$ collapse to constant.
\item v9: froze L + category supervision $\to$ $\alpha$ encoded category but exploded in magnitude.
\item v10: froze L + category supervision + hard unit-norm projection on $\alpha$ $\to$ \emph{success}.
\end{itemize}
The architectural fix (unit-norm projection) is one line of code.
Without it, the optimisation finds the degenerate solution of inflating $\alpha$ to make category prediction trivial.
With it, G is forced to encode category in the \emph{shape} of $\alpha$ rather than its size, and the recoverable orbit accommodates the resulting action.

\paragraph{Reading the cosines.}
The pairwise centroid cosines have a meaningful structure:
\begin{itemize}[nosep]
\item $\cos(\text{formal}, \text{question}) = -0.082$: nearly orthogonal. ``Formal statements'' and ``questions'' are maximally distinct in $\alpha$-space, they map to different LDA directions.
\item $\cos(\text{child}, \text{formal}) = +0.198$: roughly orthogonal. Register difference dominates.
\item $\cos(\text{child}, \text{statement}) = +0.652$: more similar. Both are simple declarative forms; they share some $\alpha$ direction.
\item $\cos(\text{statement}, \text{question}) = +0.646$: also similar. Both are ``adult'' (non-child) declarative-or-near-declarative; they share register but differ in illocutionary force.
\end{itemize}
The structure is consistent with linguistic intuition: the LDA basis discovered the \emph{register} axis (child vs adult) and the \emph{function} axis (statement vs question), and pairs that share both are close while pairs that differ on both are far.
This emergent structure was not designed; it fell out of the supervised training of $G$ in a 3-dimensional gate space.

\paragraph{What this opens up.}
v10's success means we can now ask the cleanest possible follow-up: \emph{can we replace the explicit category supervision with a self-supervised signal and get the same result?}
That is the v11 design (Section~\ref{sec:v11-design}, biology track), which uses document-source binding (wiki vs.\ owt) as an implicit episode label, mirroring how the brain uses temporal binding rather than explicit labels.

\paragraph{Updated research position (v10 era).}
\begin{itemize}[nosep]
\item Strict gauge-LORN: empirically dead.
\item Recoverable-orbit gauge-LORN: \emph{empirically realised} in v10. Static and dynamic theorems both confirmed.
\item Working architectures: v6 (best val), v7 (steering with bounded KL cost), v10 (structural input-conditional gates with frozen L).
\item The path from a framework to a working architecture took specific iterations: v8 (failed), v9 (revealed magnitude issue), v10 (succeeded).
\item The contribution of Part II is now both the mathematical reformulation \emph{and} a working dynamic operator, plus a biological grounding (Section~\ref{sec:biology}) that suggests the next direction.
\end{itemize}


\section{Phase 3: Toward Self-Supervised Gauge-LORN (v11--v12)}
\label{sec:phase3}

v10 established the dynamic theorem with gold labels.
The natural follow-up is self-supervision: can the structural axes be learned
without annotation?
This section records the attempts and their lessons.

\subsection{v11: Episodic Contrastive (Biology Track) ,  Failure Analysis}
\label{sec:v11-results}

\subsection{Design}

v11 attempted to replace v10's explicit category supervision with a self-supervised
episodic contrastive loss~\cite{chen2020simple}, using document-source binding (Wikipedia
vs.\ OpenWebText) as the implicit episode label.
Positive pairs: $(\alpha_i, \alpha_j)$ sampled from the same source document.
Negative pairs: sampled from different source documents.
Loss: InfoNCE at temperature $\tau = 0.1$.
Architecture identical to v10 (frozen L, $K=3$ LDA basis, unit-norm projection on $\alpha$).
Training: 5,000 steps, batch 16 (4 episodes $\times$ 4 chunks each).

\subsection{Results}

\begin{center}\small
\begin{tabular}{lrr}
\toprule
Metric & v11 & Target \\
\midrule
Val $\Delta$ (orbit cost) & $+0.035$ nat & $<0.5$ nat \checkmark \\
Wiki $\leftrightarrow$ OWT cosine (training signal) & $+0.975$ & $<0.7$ \textbf{FAIL} \\
Worst category-pair cosine (transfer) & $+1.000$ & $<0.85$ \textbf{FAIL} \\
\bottomrule
\end{tabular}
\end{center}

\noindent Val cost is well inside the recoverable orbit.
However, $\alpha$ collapsed: all pairwise category cosines $> 0.998$, and even the
training signal (wiki vs.\ OWT centroids) remained at cosine $0.975$, essentially identical.
The contrastive objective was satisfied without moving $\alpha$ apart.

\subsection{Failure Diagnosis}

Three compounding problems:

\textbf{1. Wrong signal axis.}
The LDA basis encodes register (child vs.\ adult) and illocutionary function
(statement vs.\ question).
The wiki/OWT distinction encodes topic and vocabulary.
These axes are approximately orthogonal: a training signal in topic-space cannot
drive $\alpha$ toward register-space.

\textbf{2. InfoNCE is trivially gameable at $K=3$.}
With unit-norm $\alpha$ in $\mathbb{R}^3$ and only four episodes per batch, $G$
can satisfy the InfoNCE objective by collapsing all $\alpha$ to a single direction.
Within-episode pairs are automatically close (trivially minimising numerator),
and with four episodes in $\mathbb{R}^3$, the denominator is also easily minimised.
The loss reaches apparent convergence with no separation.

\textbf{3. Orbit constraint and contrastive objective are not aligned.}
The unit-norm projection, necessary to keep $\delta M$ inside $\mathcal{R}$, limits
$G$ to the 2-sphere $S^2$.
On a 2-sphere, a loss that is satisfied by a single-point solution leaves no gradient
signal for exploration.

\subsection{Lesson}

Contrastive learning with a unit-norm constraint in low-dimensional space requires
a signal that is (a)~present in the training distribution and (b)~correlated with the
target structural axes.
Document-source binding fails on both counts for the LDA axes.
A replacement objective must be harder to game (direct classification) and use a
signal that actually correlates with structure.

This motivates v12 (Section~\ref{sec:v12-design}).


\subsection{Richer-Basis Probe: Extended K=10 Basis}
\label{sec:richer-basis}

\subsection{Motivation}

v11's collapse raised the question: does expanding the basis help?
The K=3 LDA basis may be insufficiently expressive for a contrastive loss to
separate well; a richer $K=10$ basis (3 LDA $+$ 7 top-PCA directions of
the weight signatures) might give $G$ more room.

\subsection{Construction}

Compute the centred weight-signature matrix $\tilde{G} \in \mathbb{R}^{N \times d^2}$.
PCA via the Gram matrix $\tilde{G}\tilde{G}^\top \in \mathbb{R}^{N \times N}$
(small, computable on CPU).
Take top-7 right singular vectors and concatenate with the K=3 LDA directions;
QR-orthogonalise to get a clean 10-direction basis.

\subsection{Results}

\begin{center}\small
\begin{tabular}{lrrrr}
\toprule
Category & $\gamma=0$ & $\gamma=0.1$ & $\gamma=0.3$ & $\gamma=1.0$ \\
\midrule
child     & $+0.000$ & $+0.106$ & $+0.439$ & $+1.136$ \\
formal    & $+0.000$ & $+0.109$ & $+0.476$ & $+1.143$ \\
statement & $+0.000$ & $+0.113$ & $+0.446$ & $+1.118$ \\
question  & $+0.000$ & $+0.140$ & $+0.462$ & $+1.135$ \\
\bottomrule
\end{tabular}
\end{center}

\noindent\textbf{Val cost at $\gamma=0.3$: $+0.44$--$0.48$ nats}, vs.\ $+0.14$ nats for
the K=3 LDA basis alone.
The extended basis is substantially worse.

\subsection{Interpretation}

The PCA directions point into high-curvature regions of the loss landscape:
they capture \emph{high-variance} directions in weight-signature space, which are
exactly the directions to which $P(x|\theta)$ is most sensitive.
Adding them to the basis opens direct paths out of the recoverable orbit.
The LDA-only basis is beneficial precisely because LDA selects
\emph{class-discriminative} directions, which appear to lie in lower-curvature
regions of the landscape.

\textbf{Conclusion:} $K=3$ LDA basis is the right choice.
Expanding to K=10 increases orbit cost by 3$\times$ with no structural benefit.


\subsection{v12: Heuristic Weak Supervision}
\label{sec:v12-design}

\subsection{Motivation}

v10 worked with gold labels.
v11 failed with source-binding contrastive.
The gap: v10's signal was directly correlated with the LDA axes; v11's was not.

v12 tests the minimum viable self-supervised signal: \textbf{heuristic labels}
derived from text statistics alone, no human annotation, but stronger than
document-source binding.

\subsection{Heuristic Category Assignment}

Categories are remapped to types actually present in the training corpus:

\begin{center}\small
\begin{tabular}{ll}
\toprule
\textbf{Category} & \textbf{Heuristic} \\
\midrule
encyclopedic & contains formal connective tokens (``however'', ``therefore'', etc.)\\
             & OR mean token-id $> 25000$ (rare/technical vocabulary) \\
informal     & contains contraction tokens (``don't'', ``it's'', ``gonna'', etc.) \\
interrogative & any `?' token present in sequence \\
narrative    & contains past-tense story-marker tokens (``said'', ``walked'', etc.) \\
\bottomrule
\end{tabular}
\end{center}

\noindent Observed label distribution on 1,600 random corpus samples:
encyclopedic 10\%, informal 47\%, interrogative 18\%, narrative 26\%.
All four classes are present at useful rates; no class dominates pathologically.

\subsection{Architecture and Loss}

Identical to v10 except:
\begin{itemize}[nosep]
\item Heuristic labels assigned on-the-fly per batch (no gold annotation at training time).
\item Direct cross-entropy classification head $C$: $\alpha(x) \to \hat{c}$ added
      alongside the language-model CE loss.
\item $\mathcal{L} = \lambda_{\text{CE}} \cdot \mathcal{L}_{\text{CE}} +
      \lambda_{\text{cat}} \cdot \mathcal{L}_{\text{cat}}$, with $\lambda_{\text{cat}} = 2$.
\item Objective is harder to game than InfoNCE: $C$ must explicitly assign distinct
      class probabilities to each $\alpha$, not merely push centroids apart.
\end{itemize}

\subsection{Key Test}

After training on heuristic corpus categories
(encyclopedic/informal/interrogative/narrative), we probe on the v10 gold-labelled
sentences (child/formal/statement/question) and measure pairwise $\alpha$ cosines.
\textbf{If worst cosine $< 0.85$:} the heuristic categories induced structural
$\alpha$ separation that transfers to the gold-labelled probe set.
This would establish that gauge-LORN is viable without any human annotation.

\subsection{Results}

\begin{center}\small
\begin{tabular}{lrr}
\toprule
Metric & v12 & Target \\
\midrule
Val $\Delta$ & $-0.005$ nat & $< 0.5$ nat \checkmark \\
Heuristic accuracy (final) & 70.7\% & ,  \\
Worst v10-probe cosine & $+1.000$ & $< 0.85$ \textbf{FAIL} \\
\bottomrule
\end{tabular}
\end{center}

Training worked perfectly: G reached 70.7\% accuracy on the heuristic corpus
categories with negligible val cost.
But $\alpha$ on the gold v10 probes collapsed completely (worst cosine 1.000).

\subsection{Failure Diagnosis}

The heuristic categories (encyclopedic/informal/interrogative/narrative) and the
gold categories (child/formal/statement/question) are defined by different sets of
sentences and live in different subspaces of the weight-gradient space.
The LDA basis is computed from gold-label gradients; heuristic-label training
steers $\alpha$ toward axes orthogonal to it.
G learned real structure, just not the structure the fixed basis was designed to
capture.

\textbf{Core lesson.}
v10, v11, and v12 all share a fundamental design flaw: the basis $\{v_k\}$ is
pre-specified (from gold-label LDA) and frozen.
Any training signal that doesn't correlate with \emph{that specific} labelling
cannot transfer to \emph{that specific} basis.
The basis and the supervision signal are coupled.

The right fix: \textbf{learn the basis jointly with $G$}, using only a diversity
regularisation that forces $\alpha$ to spread across the $K$-dimensional space,
without specifying which directions to use.
This is v13.

\subsection{v13: Learned Basis + VICReg Diversity (No Labels)}
\label{sec:v13-design}

\subsubsection{Architecture}

Replace the frozen LDA basis with a learnable parameter matrix
$V \in \mathbb{R}^{K \times d^2}$ (unit-normalised per row on use):
\begin{equation}
\delta M(x) = \gamma \|M\|_F \sum_{k=1}^K \alpha_k(x) \cdot \hat{v}_k, \qquad
\hat{v}_k = \frac{v_k}{\|v_k\|}
\end{equation}
Both $V$ and $G_\phi$ are trained jointly.
No labels. No fixed projection.

\subsubsection{Loss}

VICReg-style diversity on $\alpha$ (Zbontar et al.\ 2021~\cite{zbontar2021barlow}):
\begin{align}
\mathcal{L}_{\mathrm{var}} &= \frac{1}{K}\sum_k \max(0,\, \gamma_{\mathrm{var}} - \mathrm{std}_k(\alpha))
  \quad \text{(each dimension must have std $> \gamma_{\mathrm{var}}$)} \\
\mathcal{L}_{\mathrm{cov}} &= \frac{1}{K}\sum_{i \neq j} [\mathrm{Cov}(\alpha)]_{ij}^2
  \quad \text{(off-diagonal covariance $\to 0$)} \\
\mathcal{L} &= \lambda_{\mathrm{CE}} \cdot \mathcal{L}_{\mathrm{CE}}
             + \lambda_{\mathrm{var}} \cdot \mathcal{L}_{\mathrm{var}}
             + \lambda_{\mathrm{cov}} \cdot \mathcal{L}_{\mathrm{cov}}
\end{align}
with $\lambda_{\mathrm{var}} = 25$, $\lambda_{\mathrm{cov}} = 1$, $\gamma_{\mathrm{var}} = 0.5$.

\subsubsection{Test}

After training, run LDA on $\{\alpha(x_i)\}$ coloured by \emph{multiple}
independent labelling schemes:
\begin{itemize}[nosep]
\item Gold v10 labels (child/formal/statement/question)
\item Sentence length (short/medium/long)
\item Question vs.\ statement (binary)
\end{itemize}
Report LDA separation ratio for each.
\textbf{Success criterion:} any labelling achieves separation $> 2.0\times$
(matching v10's supervised result) at val cost $< 0.5$ nat.

This tests whether the learned basis encodes \emph{any} interpretable structural
variation, not a specific pre-decided one.

\subsection{Results}

\begin{center}\small
\begin{tabular}{lrr}
\toprule
Metric & v13 & Target \\
\midrule
Val $\Delta$ & $\mathbf{-0.069}$ nat & $< 0.5$ nat \checkmark \\
$\alpha$ std (final) & 0.524 & $> 0.5$ \checkmark \\
LDA sep — gold labels & 1.009× & $> 2.0$ \textbf{FAIL} \\
LDA sep — sentence length & \textbf{3.122×} & $> 2.0$ \checkmark \\
LDA sep — question/statement & 0.549× & $> 2.0$ \textbf{FAIL} \\
\bottomrule
\end{tabular}
\end{center}

\noindent\textbf{Verdict:} The learned basis is structural and the orbit is viable
($\Delta = -0.069$ nats, actually improving over baseline). VICReg diversity
is satisfied ($\alpha$-std $= 0.52$). But the structural axis discovered is
\textbf{sentence length}, not register or illocutionary function.
All gold-category probe pairs have cosine $\approx 1.000$: the four categories
(child/formal/statement/question) are confounded with length in the probe set.

\subsection{Interpretation}

Sentence length is the dominant source of structural variation in the training corpus
and the strongest unsupervised signal available to VICReg.
The learned $V$ found it because it dominates.
Register and function are weaker signals, present in the corpus but sub-dominant.

This is \textbf{good news for the framework}: the orbit has real structural content
and the unsupervised learner finds it. The question sharpens:
\begin{enumerate}[nosep]
\item Does the orbit also contain register/function axes, or only length?
      (Answered by Overnight A: population experiment)
\item Can one anchor redirect the learned basis toward register/function?
      (Answered by Overnight B: v14 anchor + learn)
\end{enumerate}

\paragraph{Critical correction: length is not structural.}
The v13 verdict ``learned basis is structural'' was wrong.
Sentence length is a surface statistic, not a structural property.
Length dominated the VICReg signal because it is the strongest source of variation
in random corpus batches; the gold probes happen to be short (6--8 tokens) and
length-confounded with category (child probes were tokenised longer than formal
probes in the original set).

Two fixes applied before overnight runs:
\begin{enumerate}[nosep]
\item \textbf{Length-matched probes}: all 16 probe sentences rewritten to exactly 8 GPT-2
      tokens across all four categories (\texttt{probe\_sentences.py}).
      Length is now controlled by design, not regression.
\item \textbf{Length deconfounding penalty in v14}: explicit penalty on correlation
      between $\alpha(x)$ and log-sequence-length in each training batch,
      $\mathcal{L}_{\text{len}} = \sum_k \mathrm{corr}(\alpha_k, \log|x|)^2$,
      with $\lambda_{\text{len}} = 10$.
\end{enumerate}

\paragraph{Updated research position (v13 corrected).}
\begin{itemize}[nosep]
\item v12 failure: heuristic categories and gold LDA basis are in different subspaces.
      The basis and supervision signal cannot be decoupled when the basis is fixed.
\item Key insight: the projection must be learned jointly with $G$.
\item v13 hypothesis: VICReg diversity on $\alpha$ is sufficient to learn a basis
      that encodes structural variation, discoverable post-hoc via LDA probing.
\end{itemize}

\subsection{Overnight Experiments Queued (if v13 Fails)}
\label{sec:overnight-queue}

Two experiments designed to run in parallel overnight, answering separable questions.

\paragraph{Overnight A: Population orbit experiment (\texttt{exp\_population\_orbit.py}).}
Trains $N=5$ independent v6 replicas from different random seeds.
For each held-out probe sentence, extracts mid-layer hidden states across all replicas.
Measures cross-replica covariance and decomposes it by structural-category LDA axes vs.\ replica identity vs.\ noise.

This is \emph{independent of $G$}.
It answers the prior question: does the recoverable orbit contain structural variation
as a side-effect of training, or is the orbit structurally vacuous at v6 scale?

Three possible outcomes:
\begin{itemize}[nosep]
\item \textbf{Scenario 3 (hoped for):} Cross-replica variance is dominated by category
      structure (LDA sep $> 2\times$ by category, $<$ by replica identity).
      Gauge-LORN's central theoretical bet confirmed: the orbit has structural content
      that any adequate $G$ can exploit.
\item \textbf{Scenario 1/2:} Cross-replica variance is dominated by replica identity
      or is structurally uncorrelated. The orbit is permutation-only or noise at this scale.
      Resolution: the framework is correct but needs a larger model.
\end{itemize}

\paragraph{Overnight B: Anchor + Learn, v14 (\texttt{exp\_lorn\_prefix\_v14.py}).}
One basis direction fixed ($v_{\text{anchor}}$ = v10's LDA register axis, the
highest-eigenvalue structural direction), two directions learned freely with VICReg.

Tests the minimum viable supervision: if one gold anchor is enough to orient the
full $K=3$ basis, the cost of annotation is a single structural axis, not four
full category classes.

Architecture:
\begin{equation}
V = [v_{\text{anchor}} \;|\; v_{\text{free},1} \;|\; v_{\text{free},2}], \quad
v_{\text{anchor}} \text{ frozen},\quad v_{\text{free}} \text{ learned with VICReg}
\end{equation}

Interpretive grid:
\begin{center}\small
\begin{tabular}{lll}
\toprule
\textbf{v14 vs v13} & \textbf{v14 vs v10} & \textbf{Conclusion} \\
\midrule
v14 $\gg$ v13 & v14 $\approx$ v10 & 1 anchor sufficient; minimum annotation = 1 axis \\
v14 $\gg$ v13 & v14 $\ll$ v10 & 1 anchor helps but not enough; need 2--3 \\
v14 $\approx$ v13 & ,           & Anchor doesn't help; problem is elsewhere \\
\bottomrule
\end{tabular}
\end{center}

\paragraph{Strategic logic.}
Experiments A and B answer different questions and have independent outcomes.
A validates the theoretical premise (orbit structure at this scale).
B tests the practical minimum annotation cost.
Running both overnight ensures that v13's result, whatever it is, has
an immediate next step regardless of which way it goes.

\paragraph{Update: the FFN-substrate prediction was tested and falsified.}
We re-ran the v10 (gold-supervised) and v13 (VICReg) protocols with the gauge basis drawn from $\nabla_{\mathrm{FFN}_0}$ rather than $\nabla_M$.
The expectation, based on the FFN's larger signature $z$-score, was that gold supervision would yield a cleaner LDA separation and that VICReg would discover a second learned axis.
Neither happened.
On the length-matched gold probes, gold supervision on FFN gave LDA separation $1.24\times$ versus $1.34\times$ for M-target (\emph{worse}), with all centroid pairwise cosines collapsing to $+1.000$ (uniform $\alpha$ direction); val cost $+1.4$ nats versus M-target's $+0.06$.
VICReg on FFN gave separation $1.03\times$ versus $1.01\times$ on M (\emph{flat}), again with uniform-direction collapse.
The interpretation: FFN signatures are larger but \emph{diffuse}, gradient lives everywhere in a $\sim 500$K-parameter SwiGLU and admits no low-rank handle that a small operator $G_\phi$ can latch onto.
Signature magnitude is a necessary but not sufficient condition for gauge-amenability; the substrate also needs a \emph{compact} structural subspace.

\subsection{The Bifurcation: Where Does Structure Live?}
\label{sec:bifurcation}

Today's negative result on FFN-target gauge, combined with the modest-but-real positive on M-target gauge ($83\%$ category accuracy at $+0.06$ nat val cost in v10), forces a strategic choice.
Either the small shared coupling $M$ ($\sim$0.66M params, $36\%$ of representational manifold) is the right home for structural action, in which case its capacity is presently sandbagged by our use of a tiny $K=3$ basis, or neither shared block as currently architected admits a clean gauge action, and the architecture itself must be revised to expose a unified shared substrate at whose throat structural perturbations naturally live.

We frame these as two falsifiable paths.

\paragraph{Path A: M is the structure-carrier; we have been under-probing it.}
The hypothesis is that M, despite its compact parameterisation, has a much larger structural sub-space than the $K{=}3$ LDA-from-4-categories basis can express.
The data already supports this in fragments: M's spectrum crystallises to rank $\sim$70 in 605M training and shows non-trivial condition number; the 60$\times$ syntactic-separation observation says M does force structure into the residual; the M-calculus operations (steering, arithmetic, persona adaptation) work at inference time without parameter changes.
Each of these uses dimensions of M's action space the v10 probe never touched.
\textit{Falsifiable test:} sweep $K_{\mathrm{basis}} \in \{3, 8, 16, 32\}$ using the top-$K$ PCA directions of per-sentence $\nabla_M$-signatures (richer than the centroid-only LDA, captures within-category variation).
If LDA separation on the gold probes scales monotonically with $K$ up to $K{=}16$ or $K{=}32$, M's structural capacity is real and the v10 ceiling was an artefact of basis poverty.
If separation plateaus by $K{=}3$--$5$, M's gauge-amenable subspace is genuinely small and Path B is required.

\paragraph{Path B: the shared substrate must be re-architected.}
The hypothesis is that the dichotomy between ``M (small, structural)'' and ``FFN (large, diffuse)'' is itself the wrong factorisation: structure should be carried by a \emph{joint} shared block where M sits at the throat of the FFN, e.g.\ $\mathrm{FFN}(x) = W_{\mathrm{down}}\big(\sigma(M \cdot W_{\mathrm{up}}(x))\big)$ with the same shared $M$ controlling both attention coupling and the FFN bottleneck.
Under this construction, every gauge action $\delta M$ propagates through both pathways and the per-input structural perturbation $\delta_{\mathrm{perturb}}(x) = (\delta M(x))$ becomes scalar-coupled to both signatures we have measured separately.
\textit{Falsifiable test (cheap proxy on v6):} fit a joint gauge $G_\phi(x)$ that emits a single $K$-dim $\alpha(x)$ controlling both $\delta M(x)$ and $\delta W_{\mathrm{up}}(x)$ via a shared basis projected into both spaces.
If joint $\alpha$ recovers separation $\geq 1.5\times$ where M-only gave $1.34\times$ and FFN-only gave $1.24\times$, the joint substrate is more gauge-amenable than either alone and motivates a from-scratch ``M-at-the-throat'' architecture.
If joint $\alpha$ collapses or matches M-only, the substrates are not synergistic and Path B's promise is illusory.

\paragraph{Strategic logic.}
A and B answer disjoint questions and have independent outcomes.
A asks ``did we under-use M's existing capacity?''; B asks ``is the architecture itself the limit?''
Running both (Path A on v6 with richer basis; Path B as joint-gauge proxy on v6) returns a 2$\times$2 outcome matrix: each cell selects a different programme.
If A succeeds, the entire gauge-LORN programme returns to M alone with a richer basis and the existing 30M Vast scaling run becomes the natural next checkpoint to re-test at scale.
If only B succeeds, the next architecture iteration redesigns the shared block before the next training run.
If both succeed, M's capacity is real \emph{and} the joint substrate exceeds it, the strongest result.
If neither, gauge actions on shared parameters are not the right mechanism for input-conditional structure and an alternative locus (e.g.\ per-layer response heads, residual-stream adapters) is implicated.

\paragraph{Update (2026-04-14, results in).}
Both paths were executed on v6 with matched protocols (gold supervision, BS=8, STEPS=1500, length-matched probes).
Path A K-sweep (K$\in\{3,8,16\}$, PCA basis from per-sentence $\nabla_M$ signatures) produced sep$=\{1.87, 1.29, 1.57\}$ with val cost bounded.
The K=3 PCA basis actually \emph{exceeded} v10's category-LDA basis at the same K (1.87 vs 1.34), demonstrating that basis choice within M-space matters more than K-count.
But the sep did not scale monotonically with K and all pairwise centroid cosines sit at $\geq 0.9999$, suggesting G converged to a near-uniform direction and the ``separation'' is reading noise amplification on degenerate embeddings rather than true class separation at longer training.
Path B (joint M+FFN shared-$\alpha$ gauge) gave sep$=1.00$, worse than either substrate alone, with full direction collapse.
The substrates do not synergise at v6 scale.
Neither path delivered a clean win.
The 30M scaling run (currently at the first trained probe at $50$M tokens: $M_z = -0.28$, $\mathrm{FFN}_z = +1.03$, M\_rank99 unchanged from untrained) is the only remaining route to test whether M's structural subspace \emph{widens} with model size, which would reopen Path A.
Meanwhile, the near-identical collapse pattern across all three of today's single- and joint-substrate probes motivates the deeper diagnosis in the next subsection: perhaps the locus of structure is not a shared parameter block at all.

\paragraph{Update from the 30M Vast scaling run (2026-04-14, evening): Path A has its first decisive answer.}
The 30M LORN training reached $750$M of $1$B target tokens, with WSD decay engaged from $\sim 700$M.
The in-loop probes that seeded the bifurcation question now have a complete trajectory:
\begin{center}
\small
\begin{tabular}{rrrr}
\toprule
tokens & $M_z$ & $\mathrm{FFN}_z$ & ratio $M/\mathrm{FFN}$ \\
\midrule
0 (untrained)  & $-1.24$ & $-0.38$ & ,  \\
50M            & $-0.28$ & $+1.03$ & ,  \\
100M           & $+2.00$ & $+3.56$ & $0.56$ \\
300M           & $+2.75$ & $+3.25$ & $0.85$ \\
500M           & $+4.08$ & $+4.83$ & $0.84$ \\
600M           & $+4.72$ & $+4.97$ & $0.95$ \\
650M           & $+4.59$ & $+4.24$ & $\mathbf{1.08}$ \\
700M           & $+3.83$ & $+3.93$ & $0.97$ \\
750M           & $\mathbf{+4.25}$ & $+3.54$ & $\mathbf{1.20}$ \\
\bottomrule
\end{tabular}
\end{center}

Two findings, both unanticipated by the v6-scale bifurcation analysis.

\textbf{(i) M's signature has overtaken FFN's at scale.}
At v6 (5M parameters, fully trained), the M/FFN signature ratio was $z_M / z_{\mathrm{FFN}} = 2.99 / 4.41 = 0.68$.
At 30M-at-750M tokens with WSD decay engaged, the ratio is $4.25 / 3.54 = 1.20$.
The crossover is empirically located between $600$M and $750$M tokens.
\emph{The shared coupling matrix is now more category-discriminative than the wide shared FFN.}
This vindicates the original gauge-LORN bet on $M$ as the structural carrier: the v6 result that motivated the FFN-substrate experiments was a small-model artefact, not a property of the architecture.
Path A's ``M-capacity widens with scale'' hypothesis is empirically supported.

\textbf{(ii) Spectral crystallisation has not occurred.}
M's condition number ($\kappa = 53{,}331$) and rank-99 ($r_{99} = 421$) are unchanged from the untrained initialisation through all $750$M tokens.
At v6 scale, both crystallised by $\sim 100$M tokens (M condensed to $r_{99} \approx 69$).
At 30M scale, $z_M$ has developed from $-1.24$ to $+4.25$ ($5.5$-unit movement) without any compression of M's spectrum.
\emph{Signature strength and spectral crystallisation are decoupled dynamics at this model size.}
This is a novel observation about ORN training that the prior v6 results obscured: discriminative power can develop in $M$ via mechanisms other than spectral collapse.
Whether crystallisation occurs later in this run (post-WSD-decay) or whether it is a small-model artefact will be decided by the final $250$M tokens.

These findings together close the §Bifurcation question: \emph{Path A is alive at 30M-scale even though it appeared dead at v6-scale}.
The §Where Structure Actually Lives operator-interaction hypothesis (next subsection) remains a separate, complementary line: structure may live in operator interactions \emph{and} in a sufficiently-scaled M, not exclusively one or the other.

\paragraph{Final 30M result (2026-04-14, training complete at 1B tokens): three robust findings.}
Training completed with final val loss $4.15$ nat (best $4.15$, achieved at step $14{,}500$). The WSD decay phase delivered approximately a $0.15$ nat improvement from peak-LR to end, smaller than v2's $0.23$ nat dividend at $108$M scale but of the same sign.

The complete probe trajectory through $950$M tokens (last probe before completion):
\begin{center}
\small
\begin{tabular}{rrrr}
\toprule
tokens & $M_z$ & $\mathrm{FFN}_z$ & $M/\mathrm{FFN}$ \\
\midrule
0 & $-1.24$ & $-0.38$ & ,  \\
600M & $+4.72$ & $+4.97$ & $0.95$ \\
750M & $+4.25$ & $+3.54$ & $1.20$ \\
900M & $+4.25$ & $+4.29$ & $0.99$ \\
\textbf{950M} & $\mathbf{+4.25}$ & $\mathbf{+3.35}$ & $\mathbf{1.27}$ \\
\bottomrule
\end{tabular}
\end{center}

Three findings of the full run:

\textbf{(i) M dominates FFN at the end of training.}
Final $M_z / \mathrm{FFN}_z = 1.27$, the highest ratio of the entire trajectory, and a clear inversion of v6's $0.68$. The scaling hypothesis of §Bifurcation is unambiguously supported.

\textbf{(ii) M and FFN have asymmetric responses to WSD learning-rate decay.}
Through the final $300$M tokens of WSD decay, $\mathrm{FFN}_z$ dropped from $+4.97$ (peak at $600$M) to $+3.35$ (a $32\%$ reduction); $M_z$ stayed roughly flat at $+4.25$. The shared coupling's signature is robust to low-LR fine-tuning; the wide shared FFN's signature erodes. This is new: the two shared substrates have qualitatively different sensitivity to training regime, and M is the more stable carrier of structural content under consolidation-phase dynamics. Practical implication: downstream tests that depend on a stable structural fingerprint (v17e genre discrimination, v19 prefix generation, v25 dreaming) should read $M$'s signature at the end of WSD rather than FFN's.

\textbf{(iii) Spectral crystallisation never occurs at 30M scale.}
$M_{\mathrm{cond}}$ remained at $53{,}331$ and $M_{\mathrm{rank99}}$ at $421$ through every probe from untrained through $950$M tokens. At v6 scale, crystallisation fired by $\sim 100$M tokens with $r_{99}$ collapsing to $\sim 69$. At 30M scale, despite $M_z$ developing through a $5.5$-unit range ($-1.24 \to +4.25$), the spectrum is unchanged from initialisation. This is the strongest form of the signature-vs-crystallisation decoupling hypothesis: the two ORN training dynamics are not only independent but can diverge to the point where one develops fully while the other does not fire at all. The v6 co-development was specific to small-model dynamics; at 30M the two processes require separate theoretical accounting.

The 30M \texttt{lorn30m\_best.pt} checkpoint is the substrate for the v18-on-30M consolidation experiment and all downstream gauge-LORN scaling tests.

\paragraph{v18-on-30M result (2026-04-14, evening): the CLS architecture scales cleanly from 5M to 30M.}
Immediately after the base training finished, the v18 consolidation recipe (port of the text v18 script to \texttt{FullORN} architecture: RMSNorm, GQA, RoPE, SwiGLU) was run on the 30M base checkpoint for $2500$ steps.
Frozen core: token embeddings, $A, B$ (shared coupling), shared \texttt{WideSwiGLU} FFN, output head, final norm ($43.6$M params).
Unfrozen for response recovery: per-layer $W_v, W_o, W_{k,\mathrm{proj}}$, gates, correction MLPs, all norms ($3.69$M params).
New trainable: $G_\phi$ ($0.21$M) and basis $V \in \R^{32 \times 512 \times 512}$ ($8.39$M) — total $8.59$M new.
Gauge action applied at $l^{*} = 4$ (mid-depth in $L{=}8$) as an additive perturbation to the attention output, $\alpha_{\mathrm{scale}} = 0.1$.

Training-time trajectory shows clean recoverability:
\begin{center}
\small
\begin{tabular}{rrrrrrr}
\toprule
step & inv & ntp & val & $\Delta$ vs base & within & between \\
\midrule
  250 & 0.27 & 4.31 & 4.77 & $+0.264$ & $0.964$ & $0.086$ \\
  500 & 0.29 & 4.37 & 4.71 & $+0.200$ & $0.965$ & $0.046$ \\
 1000 & 0.12 & 4.79 & 4.67 & $+0.160$ & $0.967$ & $0.038$ \\
 1500 & 0.11 & 4.46 & 4.64 & $+0.137$ & $0.968$ & $0.066$ \\
 2000 & 0.12 & 4.56 & 4.62 & $+0.114$ & $0.969$ & $0.052$ \\
 2500 & 0.18 & 4.26 & 4.58 & $+0.070$ & $0.967$ & $0.064$ \\
\bottomrule
\end{tabular}
\end{center}
$\Delta$ val monotonically decreased from $+0.264$ at step $250$ to $+0.070$ at step $2500$ — response heads absorbed the gauge action progressively. The final large held-out eval ($160$ samples $\times$ $8$ permutations):

\begin{center}
\small
\begin{tabular}{lrr}
\toprule
criterion & threshold & observed \\
\midrule
NTP preserved ($\Delta$) & $< 0.2$ & $\mathbf{+0.027}$ \\
$\alpha$ invariance      & $> 0.9$ & $\mathbf{+0.9666}$ \\
$\alpha$ separability    & $< 0.5$ & $\mathbf{+0.0634}$ \\
$\alpha_\sigma$ (non-deg.) & $> 0.08$ & $\mathbf{0.1700}$ \\
\bottomrule
\end{tabular}
\end{center}

\textbf{All four pre-specified success criteria met.}
The v18 consolidation recipe, developed and validated on the 5M v6 model (§\ref{sec:v18-consolidation}), ports to 30M without modification to loss structure, amplitude bound, learning-rate ratio, or canary specification — only the underlying architecture class (\texttt{FullORN} vs \texttt{PrefixLORN}) differs, and the port is mechanical.

This is the programme's first fully-trained CLS-style model at a non-toy scale. The checkpoint \texttt{lorn30m\_v18\_consolidated.pt} ($49$M total size) contains $G_\phi$, the trained basis $V$, the recovered response heads, and the base model's frozen components.
Downstream tests on this checkpoint ,  v17e genre discrimination at scale, v19 prefix-steered generation at scale, v25 dreaming at scale ,  are the natural next experiments.
The CLS architecture's core claim (``a trained model preserving NTP while carrying a separable invariant $\alpha$, with encoding/replay/imagination modes, at scale'') is now instantiated by a concrete artefact, not a hypothesis.

\subsection{Where Structure Actually Lives: The Operator-Interaction Hypothesis}
\label{sec:where-structure-lives}

The gauge-LORN programme to date has asked: \emph{can a small operator $G_\phi(x)$ produce low-rank input-conditional edits to a single trained weight block (M, or FFN, or their union) that encode abstract structural axes?}
Every answer we have obtained, v6's z-scores, v10's ceiling, today's Path A K-sweep, today's Path B joint collapse, today's FFN-target collapse, is consistent with a narrower and sharper claim: \emph{structure is not stored in single weight blocks}.
It is produced by the \emph{interaction} of weight blocks as they act on the residual stream.
This section states that claim in its strongest form and names a new falsifiable target.

\paragraph{Three candidate loci, in decreasing order of findability.}

\begin{enumerate}[leftmargin=2em]
\item \textbf{The residual stream $h(x)$.}
  This is the primary site of per-input structural content.
  It is by definition input-conditional and has been empirically the easiest place to find structure: the 60$\times$ syntactic-separation result on ORN residual states, steering-vector effects, and the cluster geometry of mid-layer representations all live here.
  But $h$ is not a \emph{parameter}, it is a transient computation, so it is not a target for gauge actions in the sense the programme requires.
  Any lasting structural influence must route \emph{through} parameters that shape $h$.

\item \textbf{The \emph{interaction} between weight blocks, commutators, compositions, Lie brackets.}
  Attention rotates $h$ by an effective operator that is the composition $W_{o,l} \cdot \mathrm{softmax}(h M h^\top / \sqrt{d}) \cdot W_{v,l}$, a \emph{two-operator} object parameterised by $(M, W_{v,l} W_{o,l})$.
  The commutator $[M, W_{v,l} W_{o,l}]$ measures how much the coupling rule and the response rule disagree at layer $l$; equivalently, how much room there is for one to gauge-transform without being absorbed by the other.
  Our earlier diagnostic work reported that $\mathrm{rank}([M, W_v W_o])$ traces $3 \to 150 \to 65$ during training, with a kurtosis spike at the same step, two independent diagnostics converging on a ``liberation$\to$crystallisation'' phase transition in this bracket, \emph{not} in $M$ alone.
  If structure is produced by operator interaction, then (i) the commutator is where it lives, (ii) a gauge probe on $M$ alone will see only the projection of this interaction onto $M$'s axes, and (iii) a gauge probe on FFN alone is orthogonal to the interaction entirely.

\item \textbf{Individual weight matrices (M, FFN, or a specific $W_{v,l}$).}
  This is where our probes have been looking.
  We argue it is the \emph{weakest} locus.
  $M$ was trained to encode a context-independent prior over which tokens should attend to which; the wide shared FFN was trained to encode content-level associations.
  Neither block was under optimisation pressure to develop a low-rank subspace whose axes align with abstract linguistic categories (register, illocutionary force).
  The v6 multi-signature diagnostic ($z_M = 2.99$, $z_{\mathrm{FFN}} = 4.41$) tells us only that signature \emph{magnitudes} differ between blocks; it is silent on whether either block was optimised for low-rank structural amenability.
  Today's collapse results suggest it was not.
\end{enumerate}

\paragraph{Why single-operator gauge probes are under-specified.}
A gauge probe that produces $\delta M(x) = \sum_k \alpha_k(x) v_k$ parameterises a rank-$K$ perturbation of the coupling rule.
But the effective attention kernel at layer $l$ is $K_l = M \cdot W_{v,l} W_{o,l}$ (schematically); and the differential of this product is $\delta K_l = \delta M \cdot W_{v,l} W_{o,l} + M \cdot \delta(W_{v,l} W_{o,l})$, two terms.
Any input-conditional structural variation present in $K_l$ decomposes into a share on each.
If a given structural axis happens to have most of its mass on the second term (the response composition), a gauge restricted to the first (M only) will find only the projection and its LDA separation will be small, exactly what v10 showed (sep 1.34 at K=3, flat beyond).
The ``structure-in-M'' framing is a projection artefact.

\paragraph{The commutator hypothesis, sharpened.}
We conjecture that the natural target for gauge action is the \emph{layer-localised response composition} $R_l := W_{v,l} W_{o,l}$, because:
\begin{enumerate}[nosep,leftmargin=2em]
\item It is a \emph{per-layer}, not shared, operator, structure can be layer-localised (shallow: syntax; middle: register; deep: pragmatics) which matches what empirical linguistic probes of transformers consistently find.
\item Its gradient signatures on the same probes used for $M$ and FFN are a direct test of the operator-interaction hypothesis: if $z_{R_l} > z_M$ at some layer, the response composition carries a cleaner per-sentence signature than the shared coupling.
\item A gauge on $R_l$ is cheap: $R_l \in \mathbb{R}^{d \times d}$, so the basis and perturbation tensors are the same size as for $M$.
\item It does not require re-architecture.
\end{enumerate}

This names a \textbf{Path C}: fit $G_\phi(x)$ that produces $\delta R_{l^*}(x)$ at a chosen layer $l^*$, with basis drawn from per-sentence $\nabla_{R_{l^*}}$ signatures, and measure whether the result escapes the collapse regime that M and FFN both exhibit.

\paragraph{Predictions.}
\begin{enumerate}[nosep,leftmargin=2em]
\item If $z_{R_{l^*}} > z_M$ for at least one layer, structure is more concentrated in the response composition than in the shared coupling, the operator-interaction hypothesis is alive.
\item If gold-supervised $\delta R_{l^*}$ gauge at K=3 produces sep $> 2.0$ (beating both v10 and today's Path A K=3 of 1.87) \emph{with non-degenerate centroid cosines}, the collapse seen in M and FFN probes is a property of those specific substrates, not of the gauge methodology.
\item If a sweep over $l^* \in \{1, \dots, L\}$ reveals a distinguished layer where $z_{R_{l^*}}$ peaks, the claim ``structure is layer-localised'' is concrete and measurable.
\item Failure case: if $z_{R_{l^*}}$ tracks $z_M$ at all layers and gauge collapse persists, then the locus is not a single operator at all and the hypothesis must collapse further, to the Lie bracket $[M, R_l]$ itself, which is a strictly two-operator object and cannot be gauged by editing either one alone.
\end{enumerate}

\paragraph{What this does to the earlier framing.}
Path A (M-capacity) remains interesting only in the narrow scaling sense, does the 30M probe trajectory show M's $z$-score catching up to FFN's?
Path B (M-at-throat architectural rebuild) is downgraded; joint substrates collapsed today and the operator-interaction hypothesis gives a simpler explanation than ``the shared blocks need to be fused.''
Path C is now the live experimental target.

\paragraph{Preliminary result (2026-04-14, Path C diagnostic).}
The layer-by-layer $z$-score sweep over $R_l = W_{v,l} W_{o,l}$ signatures on v6, against the same 16 length-matched category probes and same 200-sample permutation null used for $M$ and FFN, is striking:
\emph{every} layer's response composition has $z_R > z_M$, and the peak layer exceeds the whole-FFN $z$-score.
\begin{center}
\small
\begin{tabular}{lrr}
\toprule
Substrate & $z$-score & Source \\
\midrule
$M$ (whole)                  & $+2.99$ & v6 multi-signature probe \\
$\mathrm{FFN}$ (whole)       & $+4.41$ & v6 multi-signature probe \\
$R_l = W_{v,l} W_{o,l}$, mid-layer peak ($l{=}5$) & $+5.14$ & Path C Step 1 \\
$R_l$, argmin layer         & $+4.62$ & Path C Step 1 \\
\bottomrule
\end{tabular}
\end{center}
All 12 layers fall in $z \in [4.6, 5.1]$, so the result is not argmax-driven multiple-testing inflation, it is the \emph{class} of per-layer response compositions collectively that dominates.
This is a direct diagnostic win for the operator-interaction hypothesis: the response pathway at every depth carries a stronger per-sentence category signature than the shared coupling or the shared FFN.

\paragraph{Caveat: diagnostic strength $\neq$ gauge-amenability.}
Today's FFN-target collapse taught us that larger $z$ does not guarantee an operator can exploit the signature.
The corresponding gauge training at $l^{*}{=}5$ is running as these notes are written; the complete verdict requires the trained $G_\phi$ to produce non-degenerate centroid cosines and LDA separation above today's Path A K=3 ceiling of $1.87$.
If the gauge collapses as FFN did, the operator-interaction hypothesis narrows further, to the two-operator Lie bracket $[M, R_l]$, which cannot be gauged by editing either operand alone.
If the gauge succeeds, $R_l$ is the natural target and the programme redirects from gauging shared blocks to gauging per-layer response compositions.
Either outcome is informative; neither requires re-architecting the model.

\paragraph{Path C update (2026-04-14, training in): gauge collapsed.}
The $\delta R_{l^{*}{=}5}$ gauge at STEPS$=4000$, BS$=8$, gold supervision produced sep $=1.32$ with all pairwise centroid cosines at $\geq 0.9999$.
This matches the collapse pattern of M-only (sep $1.87$ at K$=3$ PCA), FFN-only (sep $1.24$), and joint M+FFN (sep $1.00$).
Four substrates, four collapses: the failure mode is not substrate choice, it is the gauge methodology as currently formulated on a 5M parameter model with a 16-sentence linguistic-category probe.
Only v10's specific protocol (STEPS$=6000$, category-LDA basis from 4 class centroids) escaped the collapse regime, and that result at sep $1.34$ is a narrow local maximum.
The diagnostic finding ($z_{R_l} > z_{\mathrm{FFN}} > z_M$ across all 12 layers) stands, structure-carrying \emph{signature} is strongest in per-layer response compositions.
The failure is at the second step of the ladder (gauge-amenability), not the first (signature measurement).

\subsection{Path D: Episodic Structure and the ``Where Was I'' Prior}
\label{sec:path-d-episodic}

The four collapses share a common feature: the target category structure (register, illocutionary force, child vs.\ formal) is an \emph{abstract linguistic} axis.
NTP training has no direct incentive to align weights with such categories; the network only needs to predict the next token.
Gauge-amenability along abstract-linguistic axes may therefore be fundamentally weak: the model lacks a low-rank handle on ``formal vs.\ informal'' because it never needed to represent formality as a factor, only to predict the next token correctly in both modes.

A different target changes this.
\emph{Episodic} structure, ``which phrases co-occur,'' ``which contexts does this token tend to appear in,'' ``what was the residual state when I last saw something like this'', is both biologically motivated (hippocampal episode binding, complementary learning systems) and \emph{directly aligned with what NTP incentivises the network to represent}.
To predict the next token well, a model must represent the local co-occurrence neighbourhood of the current context.
Episodic axes are therefore pre-aligned with training pressure.

\paragraph{The ``where was I'' reformulation.}
Beyond document-level co-occurrence (the simplest episode definition: two windows from the same document), a richer episodic signal is the \emph{residual state at the time of encoding}.
Each token $x_t$ is encoded in a residual state $h_l(x_{<t})$ at every layer $l$.
Two tokens that share a similar encoding neighbourhood have, by the network's own internal judgement, occurred in similar contexts.
This is closer to place-cell coding in hippocampus than to document boundaries: the ``episode'' is not where the text came from but what the model was thinking when it processed it.
Biologically, this is also the temporal-context-model account of memory (Howard \& Kahana 2002): each memory is anchored to the drifting context vector at encoding; retrieval reinstates the encoding context.

\paragraph{What the gauge now encodes.}
If we train $G_\phi(x)$ so that $\alpha(x_1) \approx \alpha(x_2)$ when $x_1, x_2$ share a residual neighbourhood (or more simply, a document), and $\alpha(x_1)$ differs from $\alpha(x_3)$ when the third window comes from an unrelated context, then $\alpha$ encodes \emph{episode identity}.
$R$ (the hippocampus-analogue) now has a concrete job: to produce a low-dimensional pointer into a learned episode manifold.
$L$ (the neocortex-analogue) is unchanged.
The gauge action $\delta R_{l^*}(x)$ reorganises the response pathway along episode-encoding axes without disrupting NTP predictions.

\paragraph{Protocol sketch (v16).}
Stage 1 (diagnostic, no training).
Extract $h_{l}(x)$ for $\sim 500$ windows across $\sim 50$ documents at each layer of v6.
For each layer, compute the ratio $\bar{\cos}(\mathrm{within\text{-}doc}) / \bar{\cos}(\mathrm{between\text{-}doc})$.
If this ratio exceeds $\sim 1.3$ at some layer, episode structure is already latent in the residual stream; a gauge has something to read.

Stage 2 (gauge training, short).
On v6, train $G_\phi$ producing $\delta R_{l^*}(x)$ where $l^*$ is the diagnostic-best layer, with InfoNCE loss on $\alpha$: positives are two windows from the same document, negatives are in-batch windows from other documents.
Success criteria: within-doc $\alpha$-cosine $> 0.7$ while between-doc $\alpha$-cosine $< 0.3$.
No length deconfound, no anchor, no VICReg.
The episode signal is free; the contrastive loss is minimal.

\paragraph{Predictions and risks.}
\begin{enumerate}[nosep,leftmargin=2em]
\item If the diagnostic shows within/between $> 1.3$ at some layer, episode structure is in the residual stream already, and the gauge has a target.
\item If the gauge produces non-degenerate $\alpha$ (unlike the four collapse cases above), the failure across Paths A/B/C was target-specific to abstract linguistic axes, not a universal methodology failure.
\item Risk: the ``episode'' signal might be dominated by topic (vocabulary) rather than contextual structure, in which case we recover a topic embedding by a circuitous route. Mitigation: evaluate transfer of learned $\alpha$ to held-out documents with topic mismatches.
\item Risk: if episodic gauge also collapses, the failure is architectural (one gauge target too few) and the Lie-bracket path becomes the only remaining lever.
\end{enumerate}

Path D does not replace A/B/C; it tests whether the gauge-LORN framework's failure-to-escape-collapse is target-specific or target-agnostic.
If D succeeds, the framework is methodologically sound and was previously being pointed at the wrong structural axis.
If D fails, the methodology itself is implicated and the next iteration must change the gauge's learning rule, not its target.

\paragraph{Path D result (2026-04-14): also collapsed.}
Stage~1 (diagnostic): residual-stream episode structure is present but modest, within-doc/between-doc cosine ratio climbs from 1.03 in shallow layers to 1.21 at layer 10, never reaching the 1.3 ``strong signal'' threshold.
Stage~2 (InfoNCE gauge): $\alpha_{\mathrm{std}}$ final $= 0.009$, within-doc $\alpha$-cos $= 0.9999$, between-doc $\alpha$-cos $= 0.9997$.
Full direction collapse again, identical to the four preceding attempts.

\textbf{Five substrate/target combinations, five collapses.}
The methodology, not the target, is the problem.
This motivates the most consequential reformulation in the project so far, below.

\subsection{Structure as Permutation Invariance: The Reformulation}
\label{sec:permutation-invariance}

All five failed attempts share a hidden assumption: the target quantity (register, illocutionary force, document source, episodic membership) is an \emph{ordered-content quantity}, it depends on both which tokens appear and in what order.
The NTP objective is also an ordered-content objective.
The two are not independent; they are competing for the same representational subspace.
Asking $G_\phi$ to encode an ordered quantity distinct from NTP's own targets is asking it to find a direction the network had no training pressure to develop.

The alternative is to define structure as what is \emph{invariant} under the permutation group acting on the token order.

\paragraph{Formalism.}
Let $x = (x_1, \dots, x_n)$ be a sequence, and let $S_n$ act on $x$ by permuting indices: $\pi \cdot x = (x_{\pi(1)}, \dots, x_{\pi(n)})$.
Decompose any sequence-level quantity into:
\begin{enumerate}[nosep]
\item an \textbf{ordered component} that transforms non-trivially under $S_n$, and
\item a \textbf{structural component} that is fixed by $S_n$ (i.e.\ a function of the multiset $\{\!\{x_i\}\!\}$ alone).
\end{enumerate}
The structural component is the projection through the Reynolds operator $\mathcal{R}(f) = \frac{1}{|S_n|} \sum_{\pi} f \circ \pi$; equivalently, it is what remains after averaging over permutations.
Structure, in this definition, \emph{is} the quotient $\text{sequences} / S_n = \text{multisets}$.
NTP lives on the total space (it cares about order); structure lives on the base.

\paragraph{Gauge-theoretic reading.}
$S_n$ is the gauge group.
The ``fibre'' over a multiset is the set of orderings $\pi \cdot x$ that produce the same content.
The recoverable orbit of gauge-LORN is (a subset of) the $S_n$-orbit of the current input.
A gauge action $\delta(x)$ is recoverable iff it depends only on the multiset, not on its ordering, i.e.\ iff it is $S_n$-invariant.
Every failed attempt so far was training $G$ on ordered-content losses; none produced an $S_n$-invariant $\alpha$ because the training signal itself wasn't invariant.

\paragraph{Why this repairs the collapse pattern.}
If $\alpha(x) = \alpha(\pi \cdot x)$ for all $\pi$, then $\alpha$ is by construction orthogonal to every ordered quantity.
The NTP gradient (ordered) cannot push $\alpha$ toward the trivial solution because that gradient is zero along the invariant subspace.
Collapse was the system's response to being asked to represent a quantity whose gradient was coming from the same objective as the quantity it was supposed to distinguish itself from.
Under $S_n$-invariance, that conflict disappears.

\paragraph{Why it supports self-supervised learning.}
The supervision is free: for any $x$, sample a permutation $\pi$, require $\alpha(x) \approx \alpha(\pi \cdot x)$.
No labels, no probe design, no category theory beyond ``multiset.''
This is SimCLR-style contrastive with a specific, well-defined group action.
It aligns with BYOL/MoCo in method and with CLS hippocampal binding in biology, replay binds co-occurring items, not ordered ones.

\paragraph{Why it supports gauge theory.}
The formal object ``recoverable orbit'' was always meant to be a gauge orbit in the physics sense.
Previously we had recoverable orbits defined implicitly through the Schur complement of the training Hessian, a statistical notion.
Now we have a symbolic gauge group ($S_n$) whose orbits are exactly the structural equivalence classes.
The Schur-complement orbit is the numerical approximation of the $S_n$-orbit (tokens whose swapping would leave NTP loss approximately unchanged are, to first order, members of the same orbit under $S_n$).
The two notions coincide in the limit of a well-trained model.

\paragraph{The decomposition programme.}
The sharpened research question is: \emph{how does a transformer's forward pass decompose into ordered (NTP-relevant) and disordered ($S_n$-invariant) components, layer by layer?}
Attention with positional embeddings is the primary ordered mechanism; the content-only attention kernel $h M h^\top$ (without positional encoding) is already $S_n$-equivariant.
FFN acts pointwise, per-position, so it is neither equivariant nor invariant in a useful sense.
Our observation that FFN carries stronger gradient signatures than M (z=4.41 vs 2.99) may indicate that FFN, being position-local, preferentially accumulates ordered-content features; M, being shared and applied identically at every position, is closer to a $S_n$-equivariant operator and therefore a more natural home for $S_n$-invariant gauge actions.
This is a new, testable interpretation of the M-vs-FFN signature gap.

\paragraph{Iteration plan (Path E).}
\begin{enumerate}[leftmargin=2em]
\item \textbf{v17a (diagnostic probe)}: does v6's residual stream $h_l$ already have an $S_n$-invariant component that distinguishes multisets?
For each layer, compute $h_l(x)$ and $h_l(\pi \cdot x)$ for random permutations.
Compute (i) the within-sample permutation variance (how much $h_l$ changes under $S_n$) and (ii) the between-sample permutation-invariant distance (how much the Reynolds average $\mathcal{R}(h_l)$ distinguishes different multisets).
A large ratio (invariant-distance $\gg$ permutation-variance) means the layer already carries rich multiset information that survives reordering.
\item \textbf{v17b (invariant $\alpha$ training)}: train $G_\phi$ so that $\alpha(x) = \alpha(\pi \cdot x)$ (invariance loss) and $\alpha(x) \neq \alpha(x')$ for different multisets (InfoNCE).
Decouple from any gauge action on the model forward, first prove $\alpha$ can be made non-degenerate and invariant.
\item \textbf{v17c (invariant gauge action)}: couple the invariant $\alpha$ to a $\delta$-perturbation on the model and measure (a) that the NTP val loss stays in the recoverable orbit, (b) that the $\alpha$-encoded multiset identity transfers to downstream tasks (retrieval, classification).
\end{enumerate}
Each iteration isolates a distinct failure mode.
v17a can succeed or fail independent of v17b; v17b independent of v17c.
This is a three-step ladder rather than a single test, in recognition of how many times the single-test pattern has produced uninformative collapse.

\paragraph{v17a result (2026-04-14): the invariant subspace exists and is large.}
The diagnostic probe ran 40 distinct multisets through v6 with 8 random permutations each, computing per-layer Reynolds estimates and comparing the between-sample invariant distance $D_l$ to the within-sample permutation variance $V_l$.
Results, layer by layer:
\begin{center}
\small
\begin{tabular}{rrrrl}
\toprule
layer & $D_l$ & $V_l$ & $D_l/V_l$ & order fraction \\
\midrule
0  & 3.75   & 0.015  & 251 & 0.004 (trivially high: averaged embeddings) \\
3  & 10.4   & 0.40   & 26  & 0.037 \\
5  & 19.7   & 0.95   & 21  & 0.046 \\
8  & 38.7   & 2.77   & 14  & 0.067 \\
10 & 136.6  & 11.8   & 12  & 0.079 \\
11 & 128.2  & 18.0   & 7   & 0.123 \\
\bottomrule
\end{tabular}
\end{center}
Every layer shows $D_l \gg V_l$ with order-fraction $< 0.13$ throughout, meaning that the between-sample variation of residual states is overwhelmingly invariant under token reordering.
Layer 0 is discounted, its high ratio is trivially inherited from the fact that mean-pooled token embeddings are exactly the multiset centroid.
The non-trivial finding is in the middle: layers 4--10 carry $D/V$ ratios of 12--27 with only 4\%--8\% order-dependent variation.
\textbf{A large $S_n$-invariant signal is present and layered deep into v6's forward pass}; previous gauge probes were failing because they were searching for structure in the 4\%--12\% ordered-content sliver, which is also where NTP's gradient lives.
v17b is now motivated: train $G_\phi$ to read the invariant sector directly.

\paragraph{v17b result (2026-04-14): the first clean positive in the gauge-LORN programme.}
A $K{=}32$ gate probe reading mean-pooled residual $\bar h_6$ at layer 6, trained for 600 steps (37 seconds) with symmetric InfoNCE loss where the two views are two distinct random permutations of the same multiset:
\begin{center}
\small
\begin{tabular}{lrrrr}
\toprule
& within-cos (inv.) & between-cos (dist.) & gap & $\alpha_{\mathrm{std}}$ \\
\midrule
pre-training (random init)  & $+0.9947$ & $+0.9457$ & $+0.049$ & $0.040$ \\
step 100                    & $+0.958$  & $+0.094$  & $+0.864$ & $0.165$ \\
step 600                    & $+0.972$  & $+0.046$  & $+0.926$ & $0.170$ \\
\textbf{final (120 held-out multisets $\times$ 8 permutations)} & $\mathbf{+0.9694}$ & $\mathbf{+0.0526}$ & $\mathbf{+0.9168}$ & $\mathbf{0.1709}$ \\
\bottomrule
\end{tabular}
\end{center}
$\alpha$ is a unit-norm multiset fingerprint that is robust to arbitrary token permutations (within-cosine $\approx 0.97$) and distinguishes held-out multisets with cosines near zero (between-cosine $\approx 0.05$), non-degenerate across samples.
All five prior gauge attempts (M, FFN, joint M+FFN, $R_l$, document-episodic) collapsed to direction-uniformity in thousands of steps; this one separated in a hundred.
The previous failures were not methodological, they were target failures.
Under the $S_n$-invariance reformulation, the gauge lives in the orthogonal complement of the NTP gradient and $G_\phi$ has room to represent.

The methodology collapse hypothesis (which had been the last remaining worry) is thereby \emph{falsified}.
The gauge-LORN programme is methodologically sound; its prior attempts were searching in the wrong subspace.
The research question now shifts from ``does any gauge target escape collapse?'' (yes, the $S_n$-orbit does) to ``what does this $S_n$-invariant fingerprint get used for?'' (iterations v17c and v17d).

\paragraph{v17c result (2026-04-14): gauge action is nearly recoverable.}
Coupling v17b's $\alpha$ back into v6's forward pass as $\delta R_6(x)$ (with a jointly-trained random orthogonal basis in $(d, d)$-space) and training end-to-end with $\mathcal{L} = \mathrm{InfoNCE}(\alpha(x), \alpha(\pi \cdot x)) + 0.2 \cdot \mathrm{CE}(x)$ gives:
\begin{itemize}[nosep]
\item baseline val (no gauge): 5.49 nat;
\item val with $\delta R$ on canonical inputs: 6.17 nat, i.e.\ $\Delta = +0.67$ nat (above the 0.5 nat ``recoverable'' threshold but close);
\item $\alpha$ invariance preserved: within-cos $+0.974$, between-cos $+0.064$, $\alpha_{\mathrm{std}} = 0.240$.
\end{itemize}
Val on permuted inputs is naturally much higher ($+2.57$ nat) because NTP is order-dependent; this is an artefact of evaluation, not of the gauge.
The $+0.67$ nat val cost is tunable via smaller $\alpha$-scale or longer recovery training on $\theta_R$.
$\alpha$ survives injection as a real forward-pass intervention while remaining invariant, demonstrating that the fingerprint is coupled to behaviour, not merely a separate representation.

\paragraph{v17d result (2026-04-14): unsupervised semantic clusters.}
Re-trained $G_\phi$ (60 seconds), extracted $\alpha$ for 400 held-out multisets from v6 val data, ran $k{=}8$ k-means in $\alpha$-space, decoded the three nearest-centroid multisets per cluster as text.
Held-out permutation invariance persisted: mean cosine between $\alpha(x)$ and $\alpha(\pi \cdot x)$ on the 400-sample pool was $0.9681$.
Five of eight clusters were semantically coherent:
\begin{center}
\small
\begin{tabular}{clr}
\toprule
cluster & content (decoded from representatives) & size \\
\midrule
1 & NASCAR / auto racing (Richmond, Bodine, laps, tracks) & 47 \\
2 & Radcliffe film career / British cinema & 75 \\
3 & basketball (Trey Burke, Michigan, Big Ten) & 53 \\
4 & desert geology / historical architecture (Oregon, Nevada) & 29 \\
5 & biology / astronomy (gills, mycelium, Upsilon Andromedae) & 46 \\
6 & (mixed: film + racing) & 47 \\
7 & (mixed: environment + construction) & 51 \\
\bottomrule
\end{tabular}
\end{center}
The permutation-invariant $\alpha$, trained without topic labels, discovered topics as a byproduct.
This is the first empirical demonstration of the gauge-LORN programme's ``structure without supervision'' claim.

\paragraph{v17e result (2026-04-14): 82\% 1-NN discrimination of 7 linguistic classes.}
Took v17b-style $G_\phi$ trained on generic v6 val data and threw it at 7 hand-curated linguistic classes (fairy tale, news report, poetry, academic, questions, recipe, dialogue; 8 samples each).
Extracted $\alpha$ for all 56 samples and ran leave-one-out 1-NN classification.
\textbf{Accuracy: 82.1\%} on 7 classes (chance: 14.3\%).
Per-class tightness (within-class mean $\alpha$-cosine): fairy tale $0.962$, recipe $0.956$, questions $0.954$, news $0.939$, academic $0.938$, dialogue $0.932$, poetry $0.908$.
Most-separated pairs (cross-class centroid cosine): academic$\leftrightarrow$dialogue $0.825$; recipe$\leftrightarrow$dialogue $0.828$; news$\leftrightarrow$dialogue $0.874$.
Hardest pairs: poetry$\leftrightarrow$questions $0.985$; fairy tale$\leftrightarrow$poetry $0.971$; fairy tale$\leftrightarrow$questions $0.966$.

The interpretation is sharp: \emph{$\sim 80\%$ of what we intuit as linguistic genre is recoverable from vocabulary co-occurrence alone, independent of word order}.
Recipes have recipe-words; fairy tales have fairy-tale-words; the difference between them is predominantly lexical, not syntactic.
The 18\% failure concentrates on pairs where genre \emph{does} depend on order (poetry vs.\ questions share vocabulary but differ in illocutionary pattern) or where vocabulary overlaps heavily (fairy tale vs.\ poetry, both literary, both narrative).
This is a concrete quantitative claim about the decomposition of linguistic structure: the $S_n$-invariant (vocabulary) component carries most of the genre signal; the order-dependent component carries the rest.

\paragraph{Three independent successes on v6 in one afternoon.}
(i) v17b: $\alpha$ is clean and permutation-invariant (gap $0.92$).
(ii) v17d: unsupervised clusters in $\alpha$-space are topics.
(iii) v17e: 82\% discrimination of hand-curated linguistic classes.
The programme's central claim, structure without supervision via gauge action on the permutation group, has three empirical legs to stand on after a single day's iteration.

\paragraph{v17f result (2026-04-14): $\alpha$ captures lexical registers but not tonal-emotional axes.}
Using the same v17b-style $G_\phi$ trained on generic v6 val data, evaluated leave-one-out 1-NN classification on four new hand-curated benchmarks:
\begin{center}
\small
\begin{tabular}{lrrrl}
\toprule
task & classes & chance & 1-NN & $\times$chance \\
\midrule
formality spectrum (very\_casual$\rightarrow$very\_formal) & 5 & $20.0\%$ & $\mathbf{75.0\%}$ & $\mathbf{3.8\times}$ \\
historical register (Victorian vs.\ modern)              & 2 & $50.0\%$ & $\mathbf{81.2\%}$ & $1.6\times$ \\
emotion (happy / sad / angry / fearful)                  & 4 & $25.0\%$ & $25.0\%$ & $1.0\times$ \\
scientific domain (physics / biology / chem / CS)        & 4 & $25.0\%$ & $31.2\%$ & $1.2\times$ \\
\bottomrule
\end{tabular}
\end{center}
The pattern decomposes cleanly along the ordered/disordered axis.
Formality and historical register have distinct \emph{vocabulary footprints}: very-casual English uses tokens (``bruh'', ``lmao'') that very-formal English never uses, and vice versa.
$\alpha$ picks them up.
Emotion and scientific domain are \emph{not} primarily lexical: happy and sad prose use largely overlapping vocabulary and differ in syntactic/tonal patterns; physics and biology share technical-prose vocabulary distinguishable mostly by particular technical terms (which at $\text{WIN}{=}64$ tokens may or may not appear).
$\alpha$ misses these.
This is a sharp empirical confirmation of the permutation-invariance decomposition: the $S_n$-invariant channel carries \emph{lexical-register} structure; the $S_n$-equivariant channel carries what's left (emotion, syntax, illocutionary force, fine-grained semantic domain).

\paragraph{v17g result (2026-04-14): the gauge action is recoverable, and at small scale, beneficial.}
Swept $\alpha_{\mathrm{scale}} \in \{0.05, 0.1, 0.2, 0.3\}$ at layer 6, plus $\alpha_{\mathrm{scale}} = 0.1$ at layers 4 and 5, all with $\lambda_{\mathrm{NTP}} = 0.3$ and 500 training steps.
\begin{center}
\small
\begin{tabular}{rrrrrl}
\toprule
$l^*$ & $\alpha_{\mathrm{scale}}$ & $\Delta$ val (nat) & within-cos & between-cos & verdict \\
\midrule
\textbf{6} & \textbf{0.05} & $\mathbf{-0.031}$ & $+0.971$ & $+0.041$ & $\bigstar$ below baseline \\
5 & 0.10 & $-0.007$ & $+0.976$ & $+0.052$ & $\bigstar$ below baseline \\
4 & 0.10 & $-0.014$ & $+0.983$ & $+0.059$ & $\bigstar$ below baseline \\
6 & 0.10 & $+0.068$ & $+0.973$ & $+0.078$ & $\bigstar$ recoverable \\
6 & 0.20 & $+0.490$ & $+0.976$ & $+0.127$ & $\bigstar$ recoverable (just) \\
6 & 0.30 & $+0.654$ & $+0.976$ & $+0.088$ & $\circ$ partial (v17c config) \\
\bottomrule
\end{tabular}
\end{center}
Three of six configurations produce \emph{negative} val delta: the gauge action \emph{improves} the frozen v6's val loss below its unperturbed baseline.
All configurations preserve $\alpha$-invariance (within-cos $\geq 0.97$) and separation (between-cos $\leq 0.13$).
The v17c ``partial'' verdict ($\Delta = +0.67$ nat) was entirely an over-perturbation artefact: at $\alpha_{\mathrm{scale}}=0.3$, the injection is six times too aggressive.
Reduced to $0.05$--$0.1$, the gauge acts as a \emph{useful} input-conditional modulation rather than a tolerated one.
This moves the programme from ``the gauge action is tolerable'' to ``the gauge action is informative'', a qualitatively stronger claim.

\paragraph{v17h result (2026-04-14): v6's capacity ceiling is model-size-limited, not $K$-limited.}
Swept $K \in \{16, 32, 64, 128\}$ and training steps $\in \{600, 1200, 2000\}$ on the 7-genre benchmark (trimmed to 6 samples/class, introducing noise; v17e's 82\% corresponds to 8 samples/class).
Accuracy plateaued at $\sim 69\%$ across all $K$, with a mild peak of $71.4\%$ at $K{=}128$, STEPS$=600$.
Neither enlarging $K$ nor training longer moved the ceiling appreciably.
This is the expected signature of a model-size-bound representational capacity: the permutation-invariant subspace of v6's layer-6 residual has a fixed dimensionality (a fraction of its $D = 288$), and the gate probe saturates it at $K \sim 32$.
The 30M Vast run, with $D = 512$ and ten times the training, should show this ceiling move: if $z_{\mathrm{inv}}$ at layer $l^*$ increases with scale, the same benchmark should rise into the 80s or above at 30M.
If it does not, the ceiling is about linguistic-class granularity at this hand-curated-probe size, not about the model.

\paragraph{The programme after one afternoon's iteration.}
Five failures in the morning/early-afternoon (M, FFN, joint, $R_l$, document-episodic) gave way to five successes in the evening (v17a diagnostic, v17b clean gauge, v17c$\rightarrow$v17g recoverable action, v17d topic clusters, v17e$\rightarrow$v17f genre/register discrimination, v17h scaling characterisation).
Structure = permutation-invariant content is now supported by:
\begin{enumerate}[nosep,leftmargin=2em]
\item a theoretical reformulation (the Reynolds operator is the recoverable-orbit projection),
\item a diagnostic (every layer of v6 has $D/V > 7$),
\item a clean gate-probe training (gap $0.92$),
\item a semantically coherent unsupervised clustering,
\item four linguistic-class benchmarks with different success profiles, and
\item a tuned gauge action that improves val loss when applied at small scale.
\end{enumerate}
The gauge-LORN programme, which looked empirically dead at $15{:}00$, has its first complete empirical arc by $17{:}00$.

\subsection{v18: Consolidation Co-Training}
\label{sec:v18-consolidation}

v17b--g established that a small operator $G_\phi$ can produce a permutation-invariant fingerprint $\alpha(x)$ that (a) is separable across multisets, (b) survives injection into v6's forward pass, and (c) at small $\alpha_{\mathrm{scale}}$ actually \emph{improves} val loss.
The natural next step is full consolidation: co-train $G_\phi$ with L's response heads so that after training we have a single model satisfying three simultaneous criteria:
\begin{enumerate}[nosep,leftmargin=2em]
\item NTP val loss is at or below the frozen-v6 baseline (neocortex preserved);
\item $\alpha$ is invariant under permutation (hippocampus stable);
\item $\alpha$ separates multisets non-trivially (episodic content retained).
\end{enumerate}
This is the CLS consolidation step made explicit: the fast system (G) writes bounded perturbations into the slow system (L's response pathway), and the slow system absorbs them by ordinary gradient descent.
Biologically, this is sleep-driven cortical consolidation after hippocampal replay.

\paragraph{Design (filtered through \emph{A Theory of Learning-Experiments}).}
\begin{itemize}[nosep]
\item \textbf{Frozen (slow-core):} $\theta_{\mathrm{tok\_emb}}$, $\theta_{\mathrm{pos\_emb}}$, $A$, $B$ (SharedM), shared FFN, final LayerNorm, output head.
\item \textbf{Trainable in L (recovery, $\sim 2$M params):} per-layer $W_{v,l}, W_{o,l}$, per-layer LayerNorms.
\item \textbf{Trainable new ($\sim 150$K + $K d^2$ params):} $G_\phi$ (MLP $D \to 256 \to K$ with unit-norm projection), random-initialised trainable basis $V \in \mathbb{R}^{K \times D \times D}$.
\item \textbf{Losses:} $\mathcal{L} = \mathrm{InfoNCE}\big(\alpha(\pi_1 x), \alpha(\pi_2 x)\big) + \mathrm{CE}\big(\mathrm{forward}_{\delta R}(x), y\big)$, with $\pi_1, \pi_2 \sim S_n$ fresh each step.
\item \textbf{Hyperparameters:} $K=32$, $\alpha_{\mathrm{scale}}=0.1$ (v17g's regime-that-improves), BS=8, 2000 steps, $\mathrm{lr}_G = 3 \cdot 10^{-4}$, $\mathrm{lr}_L = 10^{-4}$ (CLS: slow system slower than fast).
\item \textbf{Canary (R6):} abort if $\alpha_{\mathrm{std}} < 0.02$ (collapse) or $\Delta_{\mathrm{val}} > 2$ nat (neocortex disruption) after step 400.
\end{itemize}

\paragraph{What v18 demonstrates if it works.}
A single v6-derived model with 5M total params that (a) predicts next tokens as well as vanilla v6, (b) carries a 32-dim permutation-invariant $\alpha(x)$ per input, (c) can be queried at inference time for structural fingerprints without any retraining.
This is the full CLS system, implemented at ORN scale.
Successor experiment v19 will then populate v6's $K_{\mathrm{PREFIX}}{=}2$ prefix tokens from a linear projection of $\alpha$, closing the loop: structural fingerprints $\to$ prefix bias $\to$ conditional generation.

\subsection{Three Modes of Test-Time Computation: Encoding, Replay, Imagination}
\label{sec:three-modes}

Complementary learning systems theory distinguishes three distinct stages of hippocampus--cortex interaction.
Each maps cleanly to a distinct operator in the gauge-LORN architecture; each predicts a distinct capability; and each can be implemented, tested, and failed-on in isolation.

\paragraph{Encoding (awake, present-task).}
The hippocampus compresses the current episode into a low-dimensional index while the neocortex does its best to predict.
\emph{In gauge-LORN:} G produces $\alpha(x)$ for the current input; $\delta R$ biases L's forward pass; NTP is evaluated on the observed tokens.
v18 is the encoding mode: a single trained model with neocortex (frozen core), hippocampus (G), and response recovery (unfrozen response heads) jointly optimised.

\paragraph{Replay (offline, sharp-wave ripples).}
The hippocampus reactivates \emph{past} indices during quiet wake or NREM sleep; the neocortex integrates them.
The essence is retrieval from an episode \emph{library}, not from the current input.
\emph{In gauge-LORN:} store a library $\mathcal{L} = \{\alpha_i\}_{i=1}^{N}$ of encoded episodes; for a new input $x$, retrieve the top-$k$ $\alpha_i$ by cosine similarity to $\alpha(x)$; inject the mean as a prefix bias.
This is test-time training in the narrow sense: no gradient computation at inference, only lookup and prefix injection.
Grounding by analogy (``this looks like that episode; complete accordingly'') falls out of the construction.

\paragraph{Prospective replay / imagination (REM, dreaming, planning).}
The hippocampus generates \emph{novel} sequences, rearrangements of past indices, counterfactual trajectories, never-observed combinations, and the neocortex learns from them too.
This is where a world model acquires the majority of its coverage: most of a good world model is counterfactuals that never happened.
\emph{In gauge-LORN:} fit a prior $p(\alpha)$ on the $\alpha$-manifold (GMM, von Mises--Fisher on the unit sphere, or empirical mixture); sample $\bar \alpha \sim p(\alpha)$; project to a prefix; generate a continuation; if the continuation is coherent by some filter (self-evaluation, consistency with baseline NTP), treat $(\bar\alpha, \mathrm{continuation})$ as new training data for a fine-tune pass.

\paragraph{Why imagination is well-defined under this framework.}
$\alpha$ is unit-normalised and permutation-invariant by construction.
The $\alpha$-manifold is compact (sphere).
Sampling from it, via a learned prior, a GMM fit to training $\alpha$'s, or interpolation between anchor $\alpha$'s, produces valid fingerprints by construction, not nonsense.
The trained basis $V$ maps any such $\bar\alpha$ to a coupling-bias $\delta R$ which is, by v18's consolidation training, in the recoverable orbit.
L can absorb any sampled $\bar\alpha$ and produce a continuation whose own $\alpha$ is close to $\bar\alpha$ (by permutation-invariance propagating through the forward pass).
This is self-consistent dreaming dynamics, imagined episodes produce text whose own fingerprints approximate their seeds.

\paragraph{Abelian composition.}
A specific test of the $\alpha$-manifold's algebraic structure: for two disjoint multisets $x_1, x_2$, does the normalised sum $\hat{\mathrm{normalise}}\big(\alpha(x_1) + \alpha(x_2)\big)$ approximate $\alpha(x_1 \cup x_2)$?
If yes, $\alpha$ carries a commutative (abelian) compositional structure on top of its permutation invariance, and imagination by $\bar\alpha = \sum_k c_k \alpha_k$ (combinations of anchor fingerprints) is mathematically justified.
This is the operational form of the claim from Section~\ref{sec:permutation-invariance} that structure is an abelian composition property over multiset fragments.

\paragraph{Experimental programme.}
Each mode has a falsifiable diagnostic at v6 scale before we commit to the 30M run:
\begin{itemize}[nosep,leftmargin=2em]
\item \textbf{v19b}: encoding with class-averaged $\alpha$ for prefix generation (control for v19's single-seed noise).
\item \textbf{v20}: abelian composition diagnostic (does $\alpha(x_1) + \alpha(x_2) \approx \alpha(x_1 \cup x_2)$?).
\item \textbf{v24}: replay, retrieve top-$k$ $\alpha$ from a stored library, inject as prefix, measure NTP against direct-encoding v19.
\item \textbf{v25}: imagination/dreaming, sample $\bar\alpha$ from a learned prior, generate, measure self-consistency $\cos(\alpha(\mathrm{generated}), \bar\alpha)$.
\item \textbf{v26} (follow-on): imagination-training, loop dreamed generations back as training data, measure NTP improvement from purely generated supervision.
\end{itemize}

Each probe isolates a distinct failure mode, per R6 of the rulebook.
The imagination mode (v25--v26) is the most speculative; if it works, the gauge-LORN framework extends from ``representation learning with invariance'' to ``self-supervised world-model training via dreamed counterfactuals'', the most consequential extension this programme can produce.

\paragraph{v18 result (2026-04-14): consolidation complete.}
Training ran 2000 steps in 454 s (MPS).
On a 200-sample held-out pool:
\begin{center}
\small
\begin{tabular}{lrl}
\toprule
criterion & target & v18 final \\
\midrule
NTP preserved (val $\Delta$)                & $< 0.2$ nat      & $\boldsymbol{-0.087}$ nat (\emph{improves} baseline) \\
$\alpha$ invariance (within-cos)            & $> 0.9$          & $\boldsymbol{+0.9678}$ \\
$\alpha$ separability (between-cos)         & $< 0.5$          & $\boldsymbol{+0.0441}$ \\
$\alpha$ non-degenerate ($\alpha_{\mathrm{std}}$) & $> 0.08$    & $\boldsymbol{0.1687}$ \\
\bottomrule
\end{tabular}
\end{center}
All four success criteria satisfied on held-out data.
The v18 canary (abort if $\alpha_{\mathrm{std}} < 0.02$ or $\Delta > 2$ nat at step $> 400$) never triggered, the training was stable from step 200 onward.

The outcome is a single trained model with three simultaneous properties: it predicts next tokens \emph{better} than frozen v6 (val $5.09$ vs baseline $5.18$), it carries a 32-dim permutation-invariant fingerprint $\alpha(x)$ per input, and $\alpha$ is separable across multisets with pairwise cosine near zero.
No retraining is required for structural queries at inference time.
The CLS three-component architecture, neocortex (frozen core of L), hippocampus (G), response recovery (unfrozen response heads of L), has its first fully-trained instance.
Checkpoint saved at \texttt{results/lorn\_path\_e\_v18/v18\_checkpoint.pt} with $G$, $V$, and the recovered L response-head state.
v19 proceeds from here.

\subsection{Follow-up experiments on v18}

Nine additional probes were run on the v18-consolidated checkpoint over the evening of 2026-04-14.
Each is recorded here briefly for completeness; the informative subset of the findings is folded into the main working-hypothesis paper (\texttt{lorn\_gauge.tex}) in distilled form.

\paragraph{v19 (prefix-from-$\alpha$, single-seed).}
Trained a 79K-parameter $\alpha \to \mathrm{prefix}$ projection on top of v18. Val improves by $+0.03$ nat when the trained prefix is injected vs.\ zero prefix, confirming the pathway carries information.
Generation test on 4 genres, each with a single seed multiset: 1/4 match rate by $\alpha$-cosine re-encoding (fairy\_tale only).
Generations are incoherent at v6's 5M scale (e.g., ``Bar Don't look back'', ``Boyiana Azvette.com''); NTP dominates the prefix bias.

\paragraph{v19b (prefix-from-$\alpha$, class-averaged).}
Replaced single-seed $\alpha$ with 8-sample class-centroid $\alpha$ per genre. 3/12 match rate across three prompts and four genres $= 25.0\%$, i.e.\ chance.
Every injected $\alpha$ produced text landing closer to fairy\_tale than to its own class.
The v19 failure is not a single-sample noise artefact; it is a scale limitation.

\paragraph{v20 (abelian composition diagnostic).}
For 200 pairs of disjoint half-windows from v6 val, measured $\cos(\alpha(x_1 \cup x_2), \hat{\mathrm{normalise}}(\alpha(x_1) + \alpha(x_2))) = 0.975$, strictly exceeding the single-part baseline $\cos(\alpha(x_1 \cup x_2), \alpha(x_1)) = 0.958$.
$\alpha$-space carries a commutative composition structure on multiset union; imagination by linear combination is mathematically justified.

\paragraph{v22 ($\alpha$-space PCA).}
Extracted 2000 $\alpha$ vectors from v6 val windows and ran PCA.
Top-10 PCs explain 94\% of variance.
Nearest-centroid samples on the top 8 PCs do not split cleanly by human-interpretable axes (formal/casual, narrative/factual); the mapping from PC direction to linguistic feature is distributed.
$\alpha$ behaves as a smooth $\sim$10-dimensional embedding rather than a disentangled classifier.

\paragraph{v23 ($\alpha$-dim ablation).}
Zeroed each of 32 $\alpha$-dims in turn and measured NTP delta under the v18 gauge. Top-5 dims contribute $\Delta \in [0.010, 0.017]$ nat; $27/32$ dims contribute $<0.01$ nat alone; no single dim contributes $>0.05$ nat.
Information is distributed across many dims, consistent with v22's PCA picture.

\paragraph{v24 (hard-top-$k$ replay).}
Built a library of 128 encoded $\alpha$'s; retrieved top-3 by cosine and injected mean as prefix.
Gap relative to direct encoding: $+0.047$ nat. Library compression is plausible but lossy.

\paragraph{v25 (dreaming).}
Fit three priors on the $\alpha$-manifold and measured self-consistency $\cos(\alpha(\mathrm{gen}), \bar\alpha)$:
empirical (sample from pool) $+0.18$; uniform on the sphere $-0.01$; hybrid (midpoint of two pool $\alpha$'s) $+0.29$.
Empirical exceeds random by $+0.19$; hybrid exceeds empirical by $+0.11$.
Dreams are modestly coherent; interpolated $\alpha$'s produce the most coherent samples, consistent with a smooth-manifold picture.

\paragraph{v26 (imagination-training).}
Self-distilled on 100 dreams filtered by $\cos > 0.2$, then fine-tuned $\alpha \to \mathrm{prefix}$ for 300 steps.
NTP val increased from $5.03$ to $5.64$ ($\Delta = +0.61$ nat); dream consistency dropped from $+0.18$ to $+0.16$.
Self-distillation \emph{degrades} at v6 scale: dreams are not coherent enough to serve as useful self-supervision. This is expected of bootstrap methods; the 30M scale-up run may cross the threshold.

\paragraph{v27 (soft-retrieval replay).}
Attention-weighted mixture over the 128-$\alpha$ library with softmax temperature $\tau \in \{0.05, 0.1, 0.3, 1.0\}$.
Best: $\tau = 0.1$ with gap $+0.038$ nat vs.\ direct encoding (slight improvement over v24's hard top-3 gap of $+0.047$).
Mixture-of-episodes is marginally better than hard top-$k$, but still loses to per-input $\alpha$.
Library size, not retrieval scheme, is the binding constraint.

\paragraph{v31 (multi-layer $\alpha$, layers 3/6/9).}
Trained independent $G^{(k)}$ at three layers, evaluated on the v17e 7-class probe.
l=3: $57.1\%$; l=6: $59.5\%$; l=9: $71.4\%$.
Deeper layers produce better genre discrimination.

\paragraph{v31-full (sweep of all 12 layers).}
\begin{center}
\small
\begin{tabular}{rrrrrr}
\toprule
$l$ & 1-NN & $l$ & 1-NN & $l$ & 1-NN \\
\midrule
0 & 50.0\% & 4 & 52.4\% & 8  & 57.1\% \\
1 & 57.1\% & 5 & 66.7\% & 9  & 71.4\% \\
2 & 59.5\% & 6 & \textbf{71.4\%} & 10 & 71.4\% \\
3 & 64.3\% & 7 & 69.0\% & 11 & 71.4\% \\
\bottomrule
\end{tabular}
\end{center}
Four layers (6, 9, 10, 11) tie at the $71.4\%$ ceiling.
The $71.4\%$ ceiling is the 6-samples-per-class saturation of the trimmed probe set (v17e with 8 samples/class reaches $82\%$).
Local dips at layers 4 and 8 suggest transitional regimes; they are smaller than the overall rising trend and do not break monotonicity in expectation.
Structural carrier is broadly distributed across the stack, not localised.

\paragraph{v33 ($\alpha$-difference steering).}
Constructed $\alpha_{\mathrm{fairy}} - \alpha_{\mathrm{news}}$ (class-centroid difference), added scaled to a neutral $\alpha$, injected via v19's prefix.
Swept $\gamma \in \{-1.5, -0.7, 0, +0.7, +1.5\}$; steering was monotonic in $\gamma$ but total generation-side spread of (fairy-cos minus news-cos) across the sweep was only $+0.007$.
Directional steering via $\alpha$-difference is well below what is measurable at v6 scale. Consistent with v19's null.

\paragraph{v18-long (6000-step consolidation).}
Extended v18 from 2000 to 6000 steps. Final on 200-sample held-out pool: $\Delta_{\mathrm{val}} = -0.0964$ nat, within $+0.973$, between $+0.041$, $\alpha_\sigma = 0.171$.
$\Delta_{\mathrm{val}}$ is marginally better than v18's $-0.087$; consolidation saturates by step $\sim 2000$ and additional training yields no substantial gain. The v18 2000-step configuration is near-optimal for v6.

\subsection{Summary of the follow-up block}

The ten follow-ups cluster cleanly by verdict. Positive: v20 (abelian composition), v22/v23 ($\alpha$-space geometry), v31/v31-full (structural carrier distributed across layers, genre ceiling saturated), v18-long (v18 near-optimal). Partial: v24/v27 (replay library usable but bandwidth-limited), v25 (dreams modestly coherent). Negative at v6 scale: v19/v19b (prefix-steered generation), v26 (imagination-training), v33 ($\alpha$-difference steering) ,  each of which points at a scale ceiling. None of the negatives falsify the construction; each identifies a test that the 30M Vast checkpoint, trained to $\sim$1B tokens and consolidated via v18, is expected to pass.


\subsection{Scaling to 605M: The Hemisphere Sweep}
\label{sec:605m}

All results in the preceding sections concern a 5M-parameter ORN checkpoint (v6, $d{=}288$, $L{=}12$). The negative results at v6 scale ,  null prefix-steered generation, null $\alpha$-difference steering, no emotion or scientific-domain classification ,  each pointed at a scale ceiling. We now report the first scaling results on the ORN V3 605M checkpoint ($d{=}2048$, $L{=}32$, GQA 32/8, 8$\times$ wide shared FFN, 2B tokens trained on fineweb-edu).

\paragraph{The invariant subspace is 10--100$\times$ richer.}

A layer-by-layer sweep of the $D/V$ diagnostic (Section 4 of the companion paper, Gauge-LORN) and $S_n$-contrastive probe ($G_\phi$, InfoNCE, 1500 steps per layer) reveals that the 605M model carries dramatically more invariant structure at every depth:

\begin{center}
\small
\begin{tabular}{rrrrrrr}
\toprule
layer & $D/V$ & ord.\ frac & within-$\cos$ & between-$\cos$ & genre 1-NN & $\alpha_\sigma$ \\
\midrule
0  & 2929 & 0.03\% & +1.000 & +0.034 & 59.5\% & 0.121 \\
3  & 999  & 0.10\% & +0.999 & +0.010 & 61.9\% & 0.123 \\
7  & 79   & 1.2\%  & +0.993 & +0.035 & 59.5\% & 0.121 \\
11 & 30   & 3.3\%  & +0.978 & +0.031 & 64.3\% & 0.122 \\
15 & 19   & 5.0\%  & +0.968 & +0.049 & 47.6\% & 0.120 \\
19 & 10   & 9.3\%  & +0.937 & +0.038 & 73.8\% & 0.121 \\
23 & 9    & 9.6\%  & +0.931 & +0.032 & 66.7\% & 0.121 \\
27 & 11   & 8.5\%  & +0.941 & +0.056 & \textbf{76.2\%} & 0.120 \\
31 & 11   & 8.4\%  & +0.935 & +0.026 & \textbf{76.2\%} & 0.122 \\
\bottomrule
\end{tabular}
\end{center}

$D/V$ at layer 0 is 2929 (v6: 251). The 605M embedding layer carries $12\times$ more invariant signal per unit of order-dependent variation. Even the deepest layers maintain $D/V \approx 10$ (v6 layer 11: 7). The order fraction grows monotonically from 0.03\% to 8.4\% ,  consistent with the v6 pattern but with shallower growth (the 605M model allocates less of its representation to order-dependent structure, proportionally).

Genre 1-NN accuracy peaks at layers 27--31 (76.2\%), not at early layers. At v6 the peak was mid-depth. The 605M model develops discriminative invariant representations at depth ,  the $S_n$-invariant subspace grows richer with both model quality and depth.

\paragraph{Linguistic probes that failed at v6 succeed at 605M.}

Training $G_\phi$ for 5K steps at layer 27 and evaluating on the extended benchmark battery:

\begin{center}
\small
\begin{tabular}{lrrrr}
\toprule
task & v6 (5M) & 605M & $\Delta$ & chance \\
\midrule
formality (5 classes)       & 75.0\% & 73.3\% & $-$1.7\% & 20\% \\
\textbf{emotion (4 classes)}  & 25.0\% & \textbf{37.5\%} & $+$12.5\% & 25\% \\
\textbf{science (4 classes)}  & 31.0\% & \textbf{37.5\%} & $+$6.5\% & 25\% \\
\textbf{historical (2 classes)} & 81.0\% & \textbf{100.0\%} & $+$19.0\% & 50\% \\
\bottomrule
\end{tabular}
\end{center}

Emotion ,  at chance on v6 ,  is now clearly above chance. Scientific domain improves similarly. Historical register reaches perfect separation. The $S_n$-invariant subspace at 605M carries affective and domain-specific signal that was absent at 5M, because the deeper model encodes richer vocabulary co-occurrence patterns. Formality is stable, suggesting the lexical formality signal saturated at v6.

This validates the scale-ceiling interpretation of the v6 negatives: the construction is correct; v6 was bandwidth-limited.

\paragraph{The gate opens at 605M.}

A right hemisphere (6.4M params, $r{=}256$, 2-layer non-causal attention, $K{=}64$) was co-trained with the frozen V3 backbone. R reads at layer 27 (best genre accuracy) and injects coupling bias at layer 31 via a learnable gate initialised at $\sigma(-3) \approx 0.05$. R trains with $S_n$-contrastive InfoNCE; L's response heads ($W_V$, $W_O$) at the injection layer are unfrozen and co-trained with CE. Only 0.87\% of backbone parameters are unfrozen.

Preliminary results at 23K/50K co-training steps:

\begin{center}
\small
\begin{tabular}{lrr}
\toprule
metric & init & step 23K \\
\midrule
gate (layer 31) & 0.048 & \textbf{0.576} \\
best val loss & ,  & 3.12 (baseline 3.32, $\Delta = +0.20$) \\
within-$\cos$ & 0.92 & 0.94 \\
between-$\cos$ & 0.22 & 0.01 \\
R contrastive loss & 0.48 & 0.025 \\
\bottomrule
\end{tabular}
\end{center}

The gate is climbing monotonically: $0.048 \to 0.576$ at step 23K, recapitulating the v6 trajectory ($0.04 \to 0.86$ over 50K steps) at $120\times$ the model scale. L is voluntarily integrating R's structural context because it reduces prediction error. The improvement (+0.20 nats) is noisy at BS=1 but the gate trajectory is smooth and consistent.

This is the central empirical prediction of the two-hemisphere programme confirmed at scale: hemispheric integration emerges from learning, not from architecture. R learns structure (permutation invariance). L learns to predict. The gate is where they discover cooperation.

\paragraph{Composition is attractor convergence, not addition.}

At v6, the abelian composition $\alpha(x_1 \cup x_2) \approx \hat{\mathrm{normalise}}(\alpha(x_1) + \alpha(x_2))$ held tightly (cosine 0.975). At 605M layer 27, this drops to 0.394.

This is not a regression. The ORN's residual stream evolves under fixed nonlinear dynamics (SharedM + WideFFN + corrections). A multiset defines an \emph{attractor basin}: feed the elements in any order, the dynamics relax to the same state. At v6, the dynamics were approximately linear and the additive composition was a good approximation. At 605M, the nonlinearity is real and the composition law is convergence to a shared attractor, not vector addition. The permutation invariance (within-cosine 0.94) already proves the attractor property; the composition law is simply not additive at this depth and scale.

The correct test of compositionality is not $\alpha(A \cup B) = \hat{\mathrm{normalise}}(\alpha(A) + \alpha(B))$ but whether a nonlinear function $f(\alpha(A), \alpha(B)) \approx \alpha(A \cup B)$ exists and is learnable. An attractor probe testing convergence rate, basin continuity, and nonlinear composition is in progress.

\section{Weight Signatures: A Population Argument for the Recoverable Orbit}
\label{sec:weight-signatures}

A different framing arrives at the same recoverable-orbit picture from a population perspective.

\subsection{The Setup}

Suppose we train $N$ independent networks (or take $N$ different sentences, since each sentence's local effect on weights is itself a small ``mini-network'').
Each carries $M$ sentences in its weights.
A priori, sentences come with cluster labels $c_i \in \{1, \ldots, C\}$ derived from external linguistic structure.

There are two natural distances on the network's parameter space:
\begin{itemize}[nosep]
\item \textbf{Functional distance}: $d_{\text{func}}(\theta, \theta') = \mathrm{KL}(P_\theta \| P_{\theta'})$, the predictive divergence between two parameter configurations.
\item \textbf{Training distance}: $d_{\text{train}}(\theta, \theta')$, the minimum number of gradient steps required to move from $\theta$ to $\theta'$ while preserving NTP loss within tolerance.
\end{itemize}
The two are related but not identical: $d_{\text{train}}$ is bounded by $d_{\text{func}}$ divided by the per-step convergence rate, and admits ``shortcuts'' along low-curvature trajectories.

A priori, a trained network has no incentive to organise its weights so that sentence-level signatures cluster according to $\{c_i\}$.
NTP rewards prediction accuracy, not representational organisation.
The empirical result confirms this: per-sentence weight signatures (i.e.\ $\nabla_M \log P(s_i)$ for each sentence $s_i$) for v6 do not cluster strongly by category in raw cosine distance.

\subsection{The Claim}

The claim, in mathematical form, is that there exists a transformation $T$ on $\theta$ such that:
\begin{enumerate}[nosep]
\item $T$ keeps the network functionally close: $d_{\text{func}}(\theta, T\theta) < \epsilon$.
\item $T$ is reachable in finite training: $d_{\text{train}}(\theta, T\theta) < K$ for modest $K$.
\item The transformed network's per-sentence signatures \emph{do} cluster by $\{c_i\}$.
\end{enumerate}
This is exactly a statement about the recoverable orbit: $T\theta - \theta$ lives in the kernel of the Schur complement Hessian, AND it improves a structural alignment metric on weight signatures.

If such a $T$ exists, the network's weights are NOT inherently structurally organised, but they admit a free re-parameterisation into a structurally organised configuration that is functionally equivalent.
The structural organisation is hidden, accessible by a small training-budget walk in parameter space, and present-without-being-used in the trained model.

\subsection{Visualisation Strategy}

Three aligned diagnostics make the claim visible for large matrices:

\paragraph{Step 1: pairwise distance heatmap.}
For $N \approx 32$ sentences across $C = 4$ categories, compute $g_i = \nabla_M \mathrm{CE}(s_i)$ for each.
Plot the $N \times N$ matrix of pairwise cosine similarities, sorted by category.
If signatures cluster by category, off-diagonal blocks (between categories) should be lighter (less similar) than diagonal blocks (within category).
The current state is the heatmap on the trained model; any block structure is what the network already encodes ``for free.''

\paragraph{Step 2: PCA scatter.}
Project the $N$ signatures into 2D via PCA.
This is the unsupervised view: do signatures naturally separate?
We expect (and confirm experimentally) that they do not separate dramatically, the signatures live in a low-dimensional subspace, but its principal axes do not align with category boundaries.

\paragraph{Step 3: LDA scatter.}
Compute Linear Discriminant Analysis directions in the signature space: the directions in M-space along which the class-mean signatures maximally separate.
Project signatures onto the top 2 LDA directions.
This is the supervised view: \emph{if we let the cluster labels speak}, what directions in M-space would separate them?
The LDA directions are the directions $T$ would have to move the network to organise weights by category.

\paragraph{Step 4: are LDA directions recoverable?}
The LDA directions live in M's parameter space; they are perturbations $\delta M_{\text{LDA}, k}$.
For each, measure two things:
\begin{itemize}[nosep]
\item Initial KL when $M$ is perturbed by $\delta M_{\text{LDA}, k}$ at unit norm.
\item Recovery fraction after $K$ training steps on the rest of L.
\end{itemize}
If the recovery fraction is high, $\delta M_{\text{LDA}, k}$ is in the recoverable orbit.
The transformation $T$ that moves $\theta$ along an LDA direction is then training-recoverable: weights could shift to that direction without hurting NTP, after the rest of L compensates.

\textbf{This is the test of the user's claim.}
If LDA directions are both (a) class-separating and (b) in the recoverable orbit, then the network admits a free re-parameterisation that organises its weights by structural category.
The structural latent organisation is hidden in the trained network's recoverable-orbit slack.

\subsection{Why This Is the Right Frame}

The earlier strict-orbit framing asked: ``what perturbations don't change predictions?''
The recoverable-orbit framing asks: ``what perturbations don't change predictions \emph{after the rest of the network adapts}?''
This second question admits a much larger answer set, and crucially permits perturbations that do shift hidden representations and attention patterns, the very things we want R to influence.

The population/weight-signature framing converts gauge-LORN from a question about the network's instantaneous symmetries to a question about the network's \emph{nearby reorganisations}: configurations of $\theta$ that are functionally close (recoverable in $K$ steps) and structurally aligned.
This is a strictly larger and more useful object than the strict gauge orbit.

If the experimental results confirm that LDA directions are both class-separating and recoverable, the next-step research question is: \emph{can a learned operator $T(\theta)$ be trained to navigate the recoverable orbit toward structural alignment?}
That operator would be the practical realisation of gauge-LORN's original ambition, moving the network to a configuration that exposes its latent linguistic structure, and it would be trainable inline with the primary task because its target lives in the orbit (no NTP cost) and is non-trivial (improves structural alignment).


\section{Local Geometry: The Fisher Information Matrix as the Gauge Metric}
\label{sec:local-geometry}

Let $L$ be the trained language model with parameters $\theta \in \mathbb{R}^{|\theta|}$, defining a conditional distribution $P(y \mid x; \theta)$ over next-tokens $y$ given context $x$.
The training objective minimises the cross-entropy
\[
\mathcal{L}(\theta) \;=\; -\mathbb{E}_{x}\,\mathbb{E}_{y \sim p_{\text{data}}(\cdot \mid x)}\big[ \log P(y \mid x; \theta) \big].
\]
At a (local) optimum $\theta^{*}$, the gradient $\nabla_\theta \mathcal{L}(\theta^{*})$ vanishes and the local geometry is governed by the Hessian $\nabla^{2}_\theta \mathcal{L}(\theta^{*})$.
Two parameter directions $u, v \in \mathbb{R}^{|\theta|}$ are \emph{equivalent at second order} if moving along them produces the same loss; the kernel of the Hessian is the set of directions where the loss is locally flat.

A more relevant object for our purposes is the \textbf{Fisher Information Matrix} $F(\theta)$:
\[
F(\theta) \;=\; \mathbb{E}_{x}\,\mathbb{E}_{y \sim P(\cdot \mid x; \theta)}\!\big[\, \nabla_\theta \log P(y \mid x; \theta)\, \nabla_\theta \log P(y \mid x; \theta)^{\!\top}\, \big].
\]
$F(\theta)$ is the metric tensor on the manifold of probability distributions parameterised by $\theta$, originally introduced by Rao and made operational for neural networks by Amari's natural gradient~\cite{amari1998natural}.
For our purposes, the key identity is the second-order expansion of the KL divergence between perturbed and unperturbed predictions:
\begin{equation}
\mathrm{KL}\!\big(P(\cdot \mid x; \theta + \delta\theta) \,\big\|\, P(\cdot \mid x; \theta)\big) \;=\; \tfrac{1}{2}\,\delta\theta^{\!\top} F(\theta) \delta\theta \;+\; O(\|\delta\theta\|^{3}).
\label{eq:fisher-kl}
\end{equation}
A direction $\delta\theta$ in parameter space leaves predictions approximately invariant iff $\delta\theta^{\!\top} F(\theta) \delta\theta \approx 0$.
The set of such directions, the kernel of $F(\theta)$, plus its low-eigenvalue subspace, is the \textbf{exact gauge orbit at first order}.

\paragraph{Approximate gauge orbit.}
Define the $\epsilon$-orbit
\[
\mathcal{G}_{\epsilon}(\theta) \;=\; \big\{ \delta\theta \;:\; \delta\theta^{\!\top} F(\theta) \delta\theta \;<\; \epsilon \cdot \|\delta\theta\|^{2} \big\}.
\]
This is a cone in parameter space, spanned by eigenvectors of $F(\theta)$ with eigenvalue below $\epsilon$.
A trained network has many such low-eigenvalue directions: the spectrum of $F$ (and equivalently the Hessian) is well known to be highly concentrated, with a heavy bulk of near-zero eigenvalues and a small number of large outliers~\cite{sagun2016eigenvalues, papyan2019measurements, ghorbani2019investigation}.
Empirically, fewer than $\sim 1\%$ of eigenvalues account for most of the curvature; the remaining $99\%$ form an approximate gauge slack.

\paragraph{Connection to mode connectivity and flat minima.}
The phenomenon of mode connectivity~\cite{garipov2018loss, draxler2018essentially}, that two independently-trained networks can be connected by a low-loss path in parameter space, is a global expression of the local gauge slack we are formalising.
The path lives in the union of $\mathcal{G}_{\epsilon}$ regions across $\theta$.
Linear mode connectivity~\cite{frankle2020linear} after permutation alignment~\cite{ainsworth2023git} is even sharper: many trained networks are gauge-equivalent up to a permutation of hidden units.
The flat-minima literature~\cite{hochreiter1997flat, keskar2017large, foret2021sam} characterises the same slack from a generalisation perspective: networks settled in flat regions (large gauge orbit) generalise better.
None of this work uses the orbit; it is treated as either a property of the loss landscape (flat minima) or an obstruction to model merging (mode connectivity).
Gauge-LORN proposes to use it: as a substrate where R's actions can live without disturbing L.

\paragraph{Connection to natural gradient and K-FAC.}
Natural gradient descent~\cite{amari1998natural} multiplies the gradient by $F^{-1}$, taking steps along the steepest-descent direction in the natural metric.
K-FAC~\cite{martens2015kfac} provides a tractable Kronecker-factored approximation of $F$ for deep networks.
For gauge-LORN we want the opposite: not the steepest-descent direction (large $F$ eigenvalue) but the flattest directions (small $F$ eigenvalue).
The same machinery applies; we just invert the priorities.
We do not need to compute or invert $F$ directly, Hessian/Fisher--vector products via auto-differentiation suffice, and the constraint ``KL stays small'' is an exact local proxy for ``stay in low-eigenvalue subspace.''


\section{Architecture: What the Gauge Probe Computes}
\label{sec:architecture-part2}

\subsection{Restricting the Action to the Coupling Matrix}

The full Fisher matrix has dimension $|\theta| \times |\theta|$ which is intractable.
We restrict R's action to a small, well-chosen subspace of $\theta$.
The natural target in the ORN architecture is the shared coupling matrix $M = AB^{\!\top} \in \mathbb{R}^{d \times d}$ used at every attention layer.
Two reasons:
\begin{enumerate}[nosep]
\item $M$ is global, a perturbation to $M$ propagates through every layer's attention, giving R structural reach without per-layer parameters.
\item $M$ is the natural geometric object: it defines the metric on $q,k$ that determines what attends to what. Modifying $M$ \emph{is} modifying the attention geometry, which is the level at which structure is realised.
\end{enumerate}
Equivalently, R could act on a per-layer attention bias or on the residual stream; the analysis is identical, only the implementation changes.
We work with $M$ as the canonical case.

\paragraph{Note on the shared FFN.}
Per Section~\ref{sec:shared-ffn-gauge}, the architecture also has a globally shared FFN whose per-sentence gradient signatures are \emph{stronger} class-discriminators than $M$'s.
The construction below generalises by replacing $\delta M(x)$ with a joint low-rank action $\big(\delta M(x), \delta \mathrm{FFN}(x)\big)$ on the full shared substrate $\theta_{\mathrm{shared}}$.
The probe architecture, Fisher norm, and recoverable-orbit analysis apply unchanged.
We describe the $M$-only case for concreteness; the $\theta_{\mathrm{shared}}$ case is a direct extension.

\subsection{Low-Rank Parameterisation}

Even restricted to $M$, an unconstrained perturbation $\delta M \in \mathbb{R}^{d \times d}$ has $d^{2}$ free parameters per input, too many to learn and too many to constrain by the gauge condition.
We adopt a low-rank parameterisation, exactly analogous to the construction used by LoRA~\cite{hu2021lora} for efficient fine-tuning:
\begin{equation}
\delta M(x) \;=\; U(x)\, V(x)^{\!\top}, \qquad U(x), V(x) \in \mathbb{R}^{d \times r}.
\end{equation}
Here $r$ is the action's rank, controlling the bandwidth of R's communication channel.
Per input, R outputs $2 d r$ parameters (the entries of $U(x)$ and $V(x)$).
For $d = 288, r = 4$ this is $2304$ free parameters per input, small, controllable, and amenable to the gauge constraint.

This parameterisation has three theoretical virtues:
\begin{itemize}[nosep]
\item \textbf{Bandwidth control}: $r$ directly controls the dimensionality of the orbit R can generate. We can sweep $r$ to characterise how much structural information the orbit can carry.
\item \textbf{Manifold structure}: rank-$r$ matrices form a smooth manifold of dimension $2dr - r^{2}$, which is well-studied in matrix geometry and supports gradient methods directly.
\item \textbf{Composition with M}: the perturbed coupling $M_{\text{eff}}(x) = M + \delta M(x)$ remains in the same matrix space; no projection or special parameterisation is needed.
\end{itemize}

\subsection{The Gauge Probe Network}

R produces $(U(x), V(x))$ from some readable observation of the input $x$.
The cleanest observation is L's hidden state $h_{L}(x)$ at a chosen depth (say, layer $\lfloor N/2 \rfloor$, the middle of L's stack).
The probe $G$ is a small MLP:
\begin{align}
h_{L}(x) &\in \mathbb{R}^{S \times d} \quad \text{(L's mid-layer hidden states)} \\
\bar{h}(x) &= \mathrm{pool}_{\text{seq}}(h_{L}(x)) \in \mathbb{R}^{d} \\
(U(x), V(x)) &= \mathrm{MLP}_{G}(\bar{h}(x)), \qquad U, V \in \mathbb{R}^{d \times r}
\end{align}
Sequence pooling collapses the $S$ tokens to a single vector per input, summarising the full context.
The MLP has hidden width $\sim 4d$ and produces $2dr$ outputs.
For $d = 288, r = 4$, the MLP has $\sim 4d \cdot d + 4d \cdot 2dr = 4 \cdot 288^{2} + 4 \cdot 288 \cdot 2304 = 0.33\text{M} + 2.65\text{M} \approx 3$M parameters.
This is comparable to R's existing capacity in v4--v7 ($\sim 2$M), so the addition does not change the parameter budget materially.

\subsection{G vs.\ R: Two Distinct Networks With Distinct Roles}

A clarification we mishandled in Part~II's first draft.
Throughout Part~I we used R to denote the right hemisphere that produces structural representations and prefix tokens.
Part~II introduces a separate network G that produces the gauge-orbit action.
These are not the same:
\begin{itemize}[nosep]
\item \textbf{G is the orbit engineer}: its job is to find perturbations $\delta M(x)$ that lie in L's gauge orbit. G's loss is the Fisher norm + non-triviality. G does not encode structure; it characterises L's slack.
\item \textbf{R is the orbit user}: once G has identified the orbit, R produces structural states whose actions on M are projected onto the orbit. R's loss is the structural prediction loss (as in v6/v7) plus a clustering objective \emph{within} the orbit.
\end{itemize}
In Phase~1 only G is trained. In Phase~2 R is added on top of a trained G. In Phase~3 G, R, and L co-evolve.

\subsection{Per-Input Perturbation: How It Lands in L's Forward Pass}

The standard ORN attention layer computes
\[
q = h_{n} A,\quad k = h_{n} B,\quad \mathrm{attn} = q k^{\!\top}/\sqrt{d}, \quad C = \mathrm{softmax}(\mathrm{attn}).
\]
M is implicit: the attention bilinear form is $h_{n} M h_{n}^{\!\top}$ with $M = AB^{\!\top}$.
A per-input perturbation $\delta M(x) = U(x) V(x)^{\!\top}$ adds a low-rank attention contribution:
\begin{align}
\mathrm{attn}_{\text{eff}} &= h_{n} M_{\text{eff}}(x) h_{n}^{\!\top} / \sqrt{d}, \quad M_{\text{eff}}(x) = M + \delta M(x) \\
&= q k^{\!\top}/\sqrt{d} + (h_{n} U(x))(h_{n} V(x))^{\!\top}/\sqrt{d}.
\end{align}
The implementation cost is two extra small matrix multiplies per layer: $h_{n} U \in \mathbb{R}^{B \times S \times r}$ and $h_{n} V \in \mathbb{R}^{B \times S \times r}$, then their outer product gives the rank-$r$ contribution.
For $r = 4, S = 256$, this is negligible compared to the main attention cost.
The same $(U(x), V(x))$ are used at every layer (since M is shared); they are computed once per forward pass.

\subsection{Free Learnable Parameters: Inventory}

A complete accounting of the system, with what is learned and what is fixed:

\begin{center}
\small
\begin{tabular}{p{4cm} p{2cm} p{4cm} p{3cm}}
\toprule
Component & Symbol & Free parameters & Phase 1 status \\
\midrule
L's full network & $\theta_{L}$ & $\sim 19$M & FROZEN \\
Coupling matrix (subset of $\theta_{L}$) & $M = AB^{\!\top}$ & $2d^{2} \approx 0.17$M & FROZEN (target) \\
Per-input perturbation (output of $G$) & $\delta M(x) = UV^{\!\top}$ & $2dr \approx 2$K per input & DERIVED \\
Gauge probe network $G$ & $\theta_{G}$ & $\sim 3$M & TRAINED \\
Structure encoder $R$ (Phase 2+) & $\theta_{R}$ & $\sim 2$M (as in v6/v7) & FROZEN in Phase 1 \\
\bottomrule
\end{tabular}
\end{center}

In Phase~3 (inline) all of $\theta_{L}, \theta_{G}, \theta_{R}$ are trained jointly with a careful relative-learning-rate schedule (G slow, R medium, L fast).

\paragraph{Crucial observation: representation vs.\ learnability.}
The representational capacity of the action is set by $r$ (the rank of $\delta M$).
The learnability is set by $G$'s architecture and the loss landscape.
These are independent: a high-$r$ action might still be unlearnable if $G$ cannot find the right $U(x), V(x)$, and a low-$r$ action might be perfectly learnable but expressively limited.
The experimental programme sweeps $r$ to characterise both axes.

\paragraph{Pooling for G's input.}
Mean-pooling L's mid-layer hidden states throws away word order, which is part of what we want G to encode.
For an autoregressive LM, the natural choice is the \textbf{last token's hidden state at the chosen layer}, which by construction has summarised the entire context up to that point.
We use $h_{L}(x)[-1]$ (or position-$t$ for any chosen $t$) as G's input.


\section{Loss Functions: What the Probe Optimises}
\label{sec:losses-part2}

\subsection{The Fisher Norm Loss (Gauge Constraint)}

The first and most important loss enforces that $\delta M(x)$ stays in the gauge orbit of $L$.
By Equation~\ref{eq:fisher-kl}, this is equivalent to keeping the KL divergence between perturbed and unperturbed predictions small:
\begin{equation}
\mathcal{L}_{\text{gauge}}(G) \;=\; \mathbb{E}_{x}\,\mathbb{E}_{t}\!\left[ \mathrm{KL}\!\Big( P\!\big(\cdot \mid x_{<t};\, M + G(x)\big) \,\Big\|\, P\!\big(\cdot \mid x_{<t};\, M\big) \Big) \right]
\label{eq:gauge-loss}
\end{equation}
where the inner KL is over the next-token distribution at position $t$.
This loss is computed by two forward passes through $L$, one with $M$, one with $M + \delta M(x)$, and the KL between the two output distributions.
$L$ is frozen; only $G$'s parameters receive gradient.

The Fisher norm is exactly $2\mathcal{L}_{\text{gauge}}$ to second order, so minimising $\mathcal{L}_{\text{gauge}}$ pushes $G$'s outputs toward the kernel of $F(\theta)$ restricted to the $M$ subspace.

\subsection{Non-Triviality Penalty}

Equation~\ref{eq:gauge-loss} alone admits the trivial solution $G(x) = 0$.
We need a force pushing $G$ to produce \emph{non-trivial} actions while staying in the orbit.
Three options, in increasing strength:
\begin{itemize}[nosep]
\item \textbf{Fixed magnitude}: constrain $\|U(x) V(x)^{\!\top}\|_{F} = \gamma$ for a chosen scale $\gamma$. This forces R to use its budget. Either a hard projection (project after each gradient step) or a soft penalty $(\|\delta M\|_{F} - \gamma)^{2}$.
\item \textbf{Variance across batch} (VICReg-style~\cite{bardes2022vicreg}): each component of $\delta M(x)$ should have non-trivial std across the batch. $\mathcal{L}_{\text{var}} = \sum_{ij} \max(0, 1 - \mathrm{std}_{x}[\delta M(x)_{ij}])$.
\item \textbf{Effective rank} (Barlow Twins-style~\cite{zbontar2021barlow}): the empirical covariance of $\mathrm{vec}(\delta M(x))$ across a batch should have effective rank $\geq r$. Penalises low-rank collapse of the batch-distribution, complementing the per-input rank constraint.
\end{itemize}
The fixed-magnitude term is essential (otherwise $G \equiv 0$ trivially satisfies everything); the variance and rank terms shape the discovered orbit.
Combined:
\begin{equation}
\mathcal{L}_{\text{nt}}(G) \;=\; \lambda_{\text{mag}}\,\big( \|\delta M(x)\|_{F} - \gamma \big)^{2} \;+\; \lambda_{\text{var}}\, \mathcal{L}_{\text{var}} \;+\; \lambda_{\text{cov}}\, \mathcal{L}_{\text{cov}}.
\end{equation}

\subsection{Phase~1 Total Objective}

The Phase~1 objective combines the gauge constraint with the non-triviality penalty:
\begin{equation}
\mathcal{L}_{\text{Phase 1}}(G) \;=\; \mathcal{L}_{\text{gauge}}(G) \;+\; \lambda_{\text{nt}}\, \mathcal{L}_{\text{nt}}(G).
\label{eq:phase1-total}
\end{equation}
Optimising over $\theta_{G}$ with $\theta_{L}$ frozen.
The gauge constraint is convex in $\delta M$ (KL is locally quadratic) and smooth; the non-triviality penalty is non-convex but well-behaved.
This optimisation problem is well-posed and should converge.

\subsection{Phase~2: Within-Orbit Clustering}

Once $G$ is trained on Phase~1 alone, $\delta M(x)$ lives (approximately) in the gauge orbit but is not yet aligned with structural similarity.
Phase~2 adds a clustering loss within the orbit.
Without similarity labels, we proxy similarity through L's own attention patterns:
\begin{equation}
s(x, x') \;=\; \mathrm{cos}\!\Big( \mathrm{vec}\!\big(C_{L}(x)\big),\; \mathrm{vec}\!\big(C_{L}(x')\big) \Big),
\end{equation}
where $C_{L}(x)$ is L's attention coupling at the observation layer.
Inputs whose L-attention patterns are similar are likely to be structurally similar; we then ask R's $\delta M(x)$ to also be similar:
\begin{equation}
\mathcal{L}_{\text{cluster}}(G) \;=\; \mathbb{E}_{x, x'}\Big[ \big( s(x, x') - \mathrm{cos}(\mathrm{vec}(\delta M(x)), \mathrm{vec}(\delta M(x'))) \big)^{2} \Big].
\end{equation}
This is a self-distillation from L's attention into G's action, but without giving G access to L's coupling directly, G must reconstruct the similarity from its mid-layer hidden state observation.

\subsection{Phase~3: Joint Co-Training}

In Phase~3 both $L$ and $G$ are trained.
$L$'s updates use the standard CE loss on un-perturbed predictions:
\[
\mathcal{L}_{L}(\theta_{L}) \;=\; -\mathbb{E}_{x, t}\big[ \log P(y_{t} \mid x_{<t}; \theta_{L}) \big],
\]
\textbf{not} on perturbed predictions.
$L$ is unaware of $G$'s existence in its primary loss.
This guarantees that $L$'s capability is not compromised by the gauge probe; if $G$ disappeared, $L$'s val loss would be unchanged.

$G$'s updates use Equation~\ref{eq:phase1-total} (Phase~1 objective) plus, optionally, the clustering term from Phase~2.
The Fisher matrix $F(\theta_{L})$ is now non-stationary, it shifts with $L$'s training.
$G$ must track.
The hyperparameter to tune is the relative learning rate of $G$ vs.\ $L$: $G$ should update fast enough to keep up, but not so fast that it overshoots.

\paragraph{Optional pulse mode for L.}
Periodically (every $K$ steps), apply $\delta M(x)$ during $L$'s forward pass and use the perturbed prediction in the CE loss.
This forces $L$ to adapt toward predictions that are robust to G's perturbations, explicitly internalising the gauge invariance.
The schedule is the curriculum mentioned in Part I's gauge theory section.


\section{Diagnostic Measurements}
\label{sec:diagnostics-part2}

The Phase~1 outcome is fully characterised by four measurable quantities.
Each has a clean baseline and either confirms or refutes the gauge structure.

\subsection{Effective Rank of the Discovered Orbit}

Compute the empirical covariance of $\mathrm{vec}(\delta M(x))$ across a held-out batch of inputs.
Its eigenspectrum reveals how many independent directions $G$ uses:
\begin{equation}
\mathrm{rank}_{\text{eff}}(\delta M) \;=\; \frac{\big(\sum_{i} \sigma_{i}\big)^{2}}{\sum_{i} \sigma_{i}^{2}}, \qquad \sigma_{i} = \text{singular values of } \mathrm{Cov}_{x}[\mathrm{vec}(\delta M(x))].
\end{equation}
This is the participation ratio.
A rank-1 collapse (orbit is one direction modulated by magnitude) gives $\mathrm{rank}_{\text{eff}} = 1$.
Full diversity gives $\mathrm{rank}_{\text{eff}} = 2dr$.
For the framing to work, we need rank substantially above 1 (otherwise R has only one degree of freedom).

\subsection{Smoothness in Input}

For pairs $(x, x')$ at varying input-space distance, compute $\|\delta M(x) - \delta M(x')\|_{F}$ and check that it scales smoothly with input distance.
If it spikes at small input distances, $G$ is brittle (small input changes cause large action changes); if it is constant, $G$ is uninformative.

\subsection{Structural Alignment vs.\ Random Baseline}

Compute the cosine similarity between $\delta M(x)$ for pairs of inputs in the length-controlled probe (e.g., child-simple vs.\ formal-complex sentences at matched length).
Compare to the cosine for randomly paired inputs.
\textbf{If structural pairs have higher cosine than random pairs}, the orbit carries structural content, this is the central diagnostic.
\textbf{If they do not}, the orbit is structurally vacant: it is technically the gauge orbit but contains no information we care about.

\subsection{Stability Under L's Training}

For Phase~3: track how $\delta M(x)$ for a fixed $x$ evolves as $L$ trains.
Smooth drift is expected (the orbit shifts with $\theta_{L}$); discontinuous jumps would indicate G is failing to track.
The relevant statistic is the autocorrelation of $\delta M(x)$ across training steps for a fixed input.

\subsection{Summary of Phase~1 Falsifiability Criteria}

\begin{center}
\begin{tabular}{l l l}
\toprule
Diagnostic & Success criterion & Failure criterion \\
\midrule
Gauge norm $\mathcal{L}_{\text{gauge}}$ at convergence & $< 0.05$ nat & $> 0.5$ nat \\
Effective rank of orbit & $\geq r/2$ & $\leq 2$ (collapse) \\
Magnitude $\|\delta M\|_{F}$ at convergence & $\approx \gamma$ (target) & $< \gamma / 10$ (vanishing) \\
Structural cosine vs.\ random & $\geq +0.2$ & $\leq +0.05$ \\
\bottomrule
\end{tabular}
\end{center}

If any of the four failure criteria triggers, Phase~1 has not validated the framing and we should not proceed to Phase~2.
A clean falsification is as informative as a positive result.


\section{Experimental Programme: Three Phases}
\label{sec:programme-part2}

\subsection{Phase~1: Offline Gauge Discovery}

\paragraph{Setup.}
Take the v6 final checkpoint (best val 4.989, no diversity loss). Freeze $\theta_{L}$.
Define $G$ with rank $r \in \{2, 4, 8, 16\}$ (sweep) and target magnitude $\gamma = 0.1 \cdot \|M\|_{F}$ (small perturbation, so the local Fisher expansion is accurate).
Train $G$ on Equation~\ref{eq:phase1-total} for 5K--10K steps.

\paragraph{What we measure.}
For each $r$:
\begin{enumerate}[nosep]
\item Final $\mathcal{L}_{\text{gauge}}$ value: how well $G$ found the orbit.
\item Effective rank of the discovered orbit: how many degrees of freedom $G$ exploits.
\item Structural cosine vs.\ random: whether the orbit contains structural content.
\item Cost: wall-clock and gradient steps to converge.
\end{enumerate}

\paragraph{What we conclude.}
A monotonic increase in effective rank with $r$ (up to a saturation point) indicates the orbit is rich.
A flat structural cosine across $r$ would indicate the orbit is structurally vacant regardless of bandwidth.
The saturation rank $r^{*}$ tells us the intrinsic dimensionality of the local gauge orbit at v6's $\theta_{L}$.
Estimated time: 6 hours on MPS for the full sweep.

\subsection{Phase~2: Within-Orbit Clustering}

\paragraph{Setup.}
Take $G$ from Phase~1 at the optimal $r^{*}$.
Freeze $\theta_{G}$. Train a small clustering head $H$ that maps $\delta M(x)$ to a clustering-friendly representation.
$H$ is supervised either (a) by L's attention patterns as in Section~\ref{sec:losses-part2}.4 (no labels needed), or (b) by the length-controlled probe pairs as weak supervision.

\paragraph{What we measure.}
Compare clustering quality (e.g., NMI vs.\ a syntactic-category labeling, or silhouette score on the probe categories) to v7's prefix-token-via-VICReg representation.
If $H$ on the discovered orbit beats v7's representation, we have a positive Phase~2 result.

\paragraph{What we conclude.}
A win at Phase~2 means the orbit not only exists (Phase~1) but contains the structural information v7 had to pay 3\% val loss to surface.
This would also confirm that the v7 mechanism was a workaround, the real source of structural clustering is the orbit, and the prefix tokens were just the channel we happened to use.

\subsection{Phase~3: Inline Co-Evolution}

\paragraph{Setup.}
Train $L$ alone for the first $T_{0}$ steps (we suggest $T_{0} \approx 0.2\cdot T_{\text{total}}$, e.g.\ 10K of 50K) so that $L$ has developed a non-trivial Fisher structure.
Below this threshold $F(\theta_{L})$ has no meaningful low-eigenvalue subspace, almost every direction changes predictions because $L$ is undertrained, and $G$ has nothing to learn.

After step $T_{0}$:
\begin{itemize}[nosep]
\item $L$ continues with standard CE loss.
\item $G$ activates with Equation~\ref{eq:phase1-total}, with the gauge-loss weight $\lambda_{\text{gauge}}$ ramped linearly from $0$ to its target value over a further $T_{1}$ steps.
\item $G$'s learning rate is set to $\sim 0.1\times$ $L$'s, so $G$ tracks the slowly-shifting orbit without overshooting.
\item Optional: every $K$ steps, validate that $\mathcal{L}_{\text{gauge}}$ has not exceeded a trust-region threshold; if it has, increase $G$'s LR temporarily to recapture the orbit.
\end{itemize}
Phase~2's clustering term may be added with its own warmup, after $G$ has stabilised.

\paragraph{What we measure.}
\begin{enumerate}[nosep]
\item L's val loss over training, vs.\ a no-$G$ control (does $G$ hurt $L$?).
\item $G$'s gauge norm over training (does it stay small as L drifts?).
\item Stability of $\delta M(x)$ for a fixed $x$ across steps.
\item Structural cosine of orbit at convergence (vs.\ Phase~1 offline).
\end{enumerate}

\paragraph{What we conclude.}
A successful Phase~3 means we have built the first inline self-supervised gauge-discovery system: $L$ trains on its primary task, $G$ co-evolves to track and exploit $L$'s gauge orbit, and we can read structural clustering off $G$'s outputs without paying val loss for it.
This would be the headline result.

\subsection{Total Resource Budget}

\begin{center}
\begin{tabular}{l r r}
\toprule
Phase & Estimated MPS hours & Decisive outcome \\
\midrule
Phase 1 (offline discovery) & 6 & Validate or falsify orbit existence \\
Phase 2 (within-orbit cluster) & 6 & Validate or falsify orbit information content \\
Phase 3 (inline co-evolution) & 12 & Validate or falsify the full system \\
\midrule
Total & 24 & ,  \\
\bottomrule
\end{tabular}
\end{center}

24 hours of MPS compute to either claim a new self-supervised paradigm or rule it out cleanly.


\section{Critical Assumptions and Risks}
\label{sec:risks-part2}

The framing in Sections~\ref{sec:local-geometry}--\ref{sec:programme-part2} is theoretically motivated but rests on assumptions that may not hold for trained language models in particular.
We enumerate these explicitly so that experimental outcomes can falsify them cleanly rather than be patched around.

\subsection{Risk A: The Fisher Kernel May Be Dominated by Trivial Gauges}

The flat directions of an LM's loss landscape are typically of three kinds, all uninteresting for our purposes:
\begin{enumerate}[nosep]
\item \textbf{Permutation symmetries}: hidden units within a layer can be freely permuted without changing the function~\cite{ainsworth2023git}. This contributes a discrete (but large) symmetry group to the orbit.
\item \textbf{LayerNorm rescaling}: the affine parameters of LayerNorm are gauge-related to the next layer's incoming weights; many rescaling pairs leave the function unchanged.
\item \textbf{Output-logit additive shift}: $\mathrm{softmax}(\ell + c\mathbf{1}) = \mathrm{softmax}(\ell)$, so any additive constant on logits is trivial gauge.
\end{enumerate}
None of these encode structural similarity.
If the rank-$r$ low-eigenvalue subspace of $F$ restricted to $M$ is mostly populated by these trivial gauges (or their projections), G discovers it, declares the orbit found, and Phase~2 fails because the orbit contains nothing structural.

\textbf{Diagnostic.} Phase~1's structural-cosine test (Section~\ref{sec:diagnostics-part2}.3) is meant to catch this: if structurally-similar inputs do not produce similar $\delta M(x)$, the orbit is structurally vacant regardless of bandwidth.
But by then we have spent the Phase~1 budget.
We propose a cheaper prerequisite, an offline spectrum analysis, in Section~\ref{sec:check-0} below.

\subsection{Risk B: ``Indifferent to'' Does Not Imply ``Structurally Equivalent To''}

The framing assumes that inputs which can be mapped to one another by small actions in the gauge orbit are structurally equivalent.
This is a hope, not a theorem.
Inputs L treats as structurally similar (same clause type, same register) might still produce different next-token distributions because of content overlap differences, they would be far apart in F-norm.
Conversely, two unrelated sentences might happen to perturb similarly because the perturbation is small in F-norm for incidental reasons.

The orbit is defined by L's \emph{predictions}, not L's \emph{representational geometry}.
The two need not align.
If they do not, we have characterised L's prediction insensitivity but learned nothing about structure.

\textbf{Diagnostic.} The clustering quality in Phase~2 (NMI vs.\ a syntactic-category labeling) directly tests this.
A negative Phase~2 result with a positive Phase~1 result would indicate exactly this risk.

\subsection{Risk C: ``Approximate'' Does Not Compose Across Steps}

The Fisher norm is a \emph{local} second-order expansion: KL stays small for a single perturbation $\delta\theta$ of bounded norm.
For a sequence of $K$ perturbations $\delta\theta_{1}, \ldots, \delta\theta_{K}$, even if each has small Fisher norm, the cumulative effect can drift unboundedly.

For Phase~1 (frozen L) this is not a concern: G's perturbation is independent per input and the unperturbed L is fixed.
For Phase~3 (online co-training) it is a real concern: $\theta_{L}$ updates each step under L's CE gradient, while G is simultaneously perturbing M.
The orbit at step $t$ is not the orbit at step $t+1$, and G's outputs may be valid in the orbit they were trained on but not in the current one.

\textbf{Mitigations to test.}
\begin{itemize}[nosep]
\item \textbf{Slow G relative to L}: G's learning rate $\sim 0.1\times$ L's, so G tracks the slowly-shifting orbit.
\item \textbf{Periodic re-projection}: every $K$ steps, recompute G's outputs on a held-out batch and verify gauge norm has not exceeded threshold; if it has, increase G's learning rate temporarily.
\item \textbf{Trust region}: clip $\|\delta M(x)\|_{F}$ when the gauge norm starts to grow.
\end{itemize}

\subsection{Risk D: Natural-Gradient Analogy Is Misleading}

Natural gradient~\cite{amari1998natural} uses $F^{-1}$ to upweight low-eigenvalue directions for descent.
We want to \emph{stay in} low-eigenvalue directions.
These look symmetric but are not equivalent operations: natural gradient asks ``which direction makes the most progress per unit Fisher distance?''; we ask ``which directions traverse zero Fisher distance?''
The K-FAC machinery for approximating $F$ applies, but the established convergence theory for natural gradient does not directly transfer.

\textbf{Implication for design.} Do not assume K-FAC's block-diagonal approximation captures the relevant structure for our purposes.
The relevant approximation is whatever discriminates low-eigenvalue from zero-eigenvalue directions accurately, which is a different criterion.
Direct minimisation of KL via two forward passes (clean + perturbed) is the safest practical method.

\subsection{Risk E: VICReg in v7 May Already Be Doing This Implicitly}

This is the most uncomfortable risk: VICReg's variance term in v7 pressures prefix tokens to vary across the batch.
The successful prefixes are the ones L learns to attend to without rejecting, i.e., the ones that don't move L's predictions too far from the unperturbed baseline.
That sounds like an implicit gauge constraint enforced by L's tolerance.

If true, gauge-LORN is a \emph{rationalisation} of v7's empirical success rather than a new mechanism.
The orbit framing would explain v7's behaviour but would not produce qualitatively different results.
The new SSL paradigm claim would dissolve into ``v7 with prettier theoretical labels.''

\textbf{Diagnostic.} Compare Phase~3's online gauge-LORN to v7 head-to-head: if the only difference is interpretive (same val loss, same steering, same probe), then E is true.
If gauge-LORN delivers qualitatively different results, higher steering effect for less val cost, or steering on richer semantic axes, then it is doing something v7 isn't.

\subsection{Why These Risks Are Worth Naming}

Each risk has a corresponding diagnostic in our experimental programme. Naming them ahead of time means we cannot retreat to ``well, the framing is more subtle than we thought'' if results disappoint. Either Phase~1's structural cosine is positive or it isn't (Risks A, B). Either Phase~3 stays stable or it doesn't (Risk C). Either gauge-LORN beats v7 or it doesn't (Risk E). The framing is falsifiable in clean ways; we should hold ourselves to the falsification rather than reinterpret null results.


\section{Check 0: Offline Spectrum Diagnostic Before Phase 1}
\label{sec:check-0}

Before committing to Phase~1, a 30-minute offline diagnostic on the v6 checkpoint can rule out Risk A and partially address Risk B.
The diagnostic computes the Hessian/Fisher spectrum on a small sample, restricted to the M subspace, and characterises what the low-eigenvalue subspace contains.

\paragraph{Procedure.}
\begin{enumerate}[nosep]
\item Sample 64--256 inputs from validation.
\item Apply structured perturbations to M of various types: random rank-1, random rank-2/4/8, known-trivial gauges (permutation, LayerNorm rescale), and gradient-direction perturbations at a fixed input.
\item For each perturbation type, measure (a) KL on validation, (b) change in generation samples, (c) change in length-controlled probe separation, (d) change in attention entropy.
\item Construct a phase diagram: x-axis perturbation magnitude $\|\delta M\|_{F}$, y-axis effect (KL or probe), one curve per perturbation type.
\end{enumerate}

\paragraph{What we learn.}
\begin{itemize}[nosep]
\item Random perturbations should look like a generic baseline, they are random directions in M-space.
\item Trivial-gauge perturbations should produce zero or near-zero KL even at large magnitudes (they are exact symmetries by construction).
\item Gradient-direction perturbations should produce maximum KL per unit magnitude (they are aligned with the high-Fisher direction).
\item \textbf{If random rank-$r$ perturbations also produce near-zero KL at large magnitudes}: there is structural slack in M beyond the trivial gauges. Gauge-LORN has room to operate.
\item \textbf{If random rank-$r$ perturbations produce KL comparable to gradient-direction perturbations}: M has no useful slack at this rank. We need to look elsewhere (per-layer biases, hidden-state perturbations).
\end{itemize}

The phase diagram itself is informative independently of gauge-LORN: it characterises how robust the trained model is to structured weight perturbations and what kinds of perturbations correspond to what kinds of behavioural change.
This is a useful diagnostic for any ORN architecture, gauge-LORN or not.


\section{Connection to Existing Work}
\label{sec:related-part2}

The gauge-LORN framing borrows from several distinct lines of work and synthesises them.

\paragraph{Geometric deep learning and equivariant networks.}
Bronstein, Bruna, Cohen, and Veli\v{c}kovi\'{c}~\cite{bronstein2021geometric} synthesise equivariant network design under the umbrella of group symmetries.
Group-equivariant CNNs~\cite{cohen2016group} and the more general framework on homogeneous spaces~\cite{cohen2019general} build symmetries into the architecture.
The pattern is: \emph{specify the symmetry group, derive the equivariant operator}.
Gauge-LORN inverts this: the network has approximate symmetries that we discover post-hoc and then exploit, rather than baking them in.
This is the right move for language, where the relevant symmetries are not known a priori (unlike, say, rotation symmetry for images).

\paragraph{Natural gradient and Fisher geometry.}
Amari's natural gradient~\cite{amari1998natural} and the Kronecker-factored approximation of Martens and Grosse~\cite{martens2015kfac} make Fisher information operational for deep networks.
We use the same machinery but invert the priorities: instead of moving along high-Fisher directions for fast optimisation, we identify low-Fisher directions for the gauge orbit.
Hessian-vector products via auto-differentiation give us all the access we need without forming $F$ explicitly.

\paragraph{Loss landscape and mode connectivity.}
The Hessian spectrum literature~\cite{sagun2016eigenvalues, papyan2019measurements, ghorbani2019investigation} establishes that trained networks have a heavy-tailed eigenvalue distribution with most eigenvalues near zero.
Mode connectivity work~\cite{garipov2018loss, draxler2018essentially} and linear mode connectivity after permutation~\cite{frankle2020linear, ainsworth2023git} show that this slack is not random but follows global structure: the low-loss manifold is connected.
We give this slack a use.

\paragraph{Low-rank adaptation.}
LoRA~\cite{hu2021lora} parameterises fine-tuning weights as low-rank perturbations $W + AB^{\!\top}$ with $A, B$ small.
Our $\delta M(x) = U(x) V(x)^{\!\top}$ uses the same parameterisation but the perturbation is \emph{input-dependent}, derived per example by $G$, and constrained to lie in the gauge orbit.
LoRA modifies $W$ to change predictions; gauge-LORN modifies $M$ to \emph{not} change predictions while encoding structure in the choice of perturbation.

\paragraph{Self-supervised methods.}
SimCLR~\cite{chen2020simclr}, BYOL~\cite{grill2020byol}, Barlow Twins~\cite{zbontar2021barlow}, VICReg~\cite{bardes2022vicreg}, and JEPA~\cite{assran2023ijepa} all solve the SSL problem by specifying positive pairs through hand-designed augmentations or through architectural asymmetry.
Each is excellent within its design assumptions; each requires those assumptions.
Gauge-LORN derives positive pairs from the gauge orbit: inputs within the orbit of one another (via small actions $H$) are positives by construction.
The framework subsumes the need for augmentations.

\paragraph{Information bottleneck.}
Tishby and Zaslavsky~\cite{tishby2015deep} characterise representation learning as compressing inputs while preserving task-relevant information.
The gauge orbit is the dual: directions that are preserved by $L$ but not necessary for $L$'s predictions.
Information that lives \emph{in} the orbit is task-irrelevant for $L$ but free for R to use.
This is the structural slack the IB analysis predicts but does not name.

\paragraph{Implicit bias and NTK.}
The implicit-bias literature~\cite{soudry2018implicit} shows that gradient descent in over-parameterised regimes prefers particular solutions (max-margin, low-norm, etc.).
The NTK regime~\cite{jacot2018ntk} characterises training dynamics linearly.
These results imply that trained networks have substantial gauge slack as a side-effect of the implicit regularisation; gauge-LORN proposes to make that slack usable.


\section{Summary}
\label{sec:part2-summary}

Part~II re-derives the LORN architecture from a single principle: \emph{R produces actions that lie in L's approximate gauge orbit and use that orbit to encode structural similarity}.
The three concrete contributions are:
\begin{enumerate}[nosep]
\item \textbf{A formal architecture}: a probe network $G$ producing low-rank input-dependent perturbations $\delta M(x) = U(x) V(x)^{\!\top}$ to L's shared coupling matrix, with the Fisher norm as the gauge constraint.
\item \textbf{A clean three-phase experimental programme}: offline discovery on a frozen v6 (Phase~1), within-orbit clustering (Phase~2), inline co-training (Phase~3), each with explicit falsifiability criteria and a 24-hour total compute budget.
\item \textbf{A reframing of the SSL problem}: the gauge orbit gives a domain-agnostic, augmentation-free, augmentation-discovered alternative to SimCLR/BYOL/VICReg-style self-supervision, derived from a property the network has anyway.
\end{enumerate}
The next concrete action is Phase~1, which can begin immediately on the existing v6 checkpoint with no new infrastructure.


\section{The Biology: Two Hemispheres, Two Memory Systems, One Principle}
\label{sec:biology}

The original LORN architecture took its name from McGilchrist's hemispheric thesis~\cite{mcgilchrist2009master, mcgilchrist2021matter}: the left and right hemispheres of the human brain are not redundant copies but specialise in qualitatively different processing.
We mapped that to L (sequential, content-focused, predictive) and R (relational, structural, holistic).
The mapping motivated the architecture but was metaphorical, it told us \emph{what} R should do without saying \emph{how} R can act on L without disrupting it.

The recoverable-orbit framing developed in Sections~\ref{sec:recoverable}--\ref{sec:weight-signatures-results} brings a second, sharper biological parallel into view: the \textbf{Complementary Learning Systems} (CLS) framework~\cite{mcclelland1995cls, kumaran2016cls}.
CLS describes how the brain solves the \emph{exact} computational problem we have been wrestling with: how a fast learner can encode new structural information without disrupting a slow learner that holds general knowledge.
The mathematical machinery we developed (recoverable orbit, Schur complement, action-with-co-training) appears to be a computational instantiation of what hippocampus and neocortex do biologically.

This section makes both parallels precise.

\subsection{McGilchrist: Two Hemispheres, Two Modes of Attention}

McGilchrist's thesis is that the hemispheric division reflects a fundamental tension in any cognitive system: the need to attend simultaneously to local detail (catching prey, decoding words) and global context (avoiding predators, understanding meaning).
Single-system attention cannot do both well; the modes are computationally incompatible.
The right hemisphere maintains broad, contextual, relational attention.
The left specialises in narrow, sequential, categorical processing.
The corpus callosum is bandwidth-limited (\textasciitilde 200 million axons relative to \textasciitilde 20 billion neurons per hemisphere~\cite{ringo1994hemispheric}), forcing the two hemispheres to communicate at the level of summaries, not raw signals.

This is where LORN's design originates.
L is the left hemisphere: processes tokens sequentially, predicts the next one, lives in content space.
R is the right hemisphere: maintains a global structural state, observes L's coupling at multiple depths, communicates through a low-bandwidth channel (prefix tokens or, in gauge-LORN's case, low-rank perturbations).
The asymmetric communication, L drives prediction, R contributes structural context, mirrors McGilchrist's claim that the hemispheres are not equal partners but a Master (right, contextual) and Emissary (left, focused).

\paragraph{What McGilchrist did not tell us.}
The hemispheric thesis specifies the \emph{division of labour} but is silent on the \emph{integration mechanism}.
How does R's contribution reach L without disrupting L's prediction?
How can R encode novel structural patterns without overwriting L's existing representations?
These are the questions our experimental program kept hitting.
The strict gauge-orbit framing failed because R's structural action turned out to be high-Fisher (Section~\ref{sec:phase2a-correction}); the recoverable-orbit framing succeeded because it allowed L to slowly absorb R's perturbations through training (Section~\ref{sec:recoverable-experiment}).
This last move, the shift from instantaneous invariance to absorbable disruption, is the bridge we needed, and it is the bridge that biology took 500 million years to evolve.

\subsection{Complementary Learning Systems: The Mathematical Brain}

The CLS framework~\cite{mcclelland1995cls}, formalised by McClelland, McNaughton, and O'Reilly in 1995 and updated by Kumaran, Hassabis and McClelland in 2016~\cite{kumaran2016cls}, identifies a fundamental tension in any learning system.
A network that learns rapidly will overwrite previous knowledge (catastrophic interference~\cite{french1999catastrophic}).
A network that learns slowly will integrate experiences into coherent semantic structure but cannot encode individual episodes.
Both capabilities are required; no single system can have both.

The brain's solution is two specialised systems with complementary properties:

\begin{center}
\small
\begin{tabular}{p{4.5cm} p{4.5cm} p{4.5cm}}
\toprule
Property & Hippocampus & Neocortex \\
\midrule
Learning rate & high (single-trial encoding) & low (gradual integration) \\
Representation & sparse, conjunctive, episodic & dense, distributed, semantic \\
Coding strategy & pattern separation (keep distinct) & pattern completion (extract regularities) \\
Time scale of update & seconds to days & weeks to years \\
Indexing function & yes (binds co-occurring features) & no (extracts statistical regularities) \\
Lesion effect & anterograde amnesia (no new memories) & retrograde amnesia (lost old memories) \\
\bottomrule
\end{tabular}
\end{center}

\paragraph{Consolidation: how the systems interact.}
The hippocampus encodes new experiences immediately as sparse, conjunctive indices binding co-occurring features (where, when, who, what).
During subsequent rest, particularly during sleep~\cite{diekelmann2010sleep}, the hippocampus \emph{replays} recent activity patterns to the neocortex via sharp-wave ripples (SWRs)~\cite{buzsaki2015swr, wilson1994reactivation}, temporally compressed reactivations of waking experience (\textasciitilde 10--20$\times$ faster).
The neocortex, receiving these replays, slowly adjusts its synapses.
Over weeks to years, the memory becomes hippocampus-independent: it has been \emph{consolidated} into cortical structure.

Two key features of consolidation are directly relevant to our problem:

\begin{enumerate}[nosep]
\item \textbf{The cortex is not disrupted by replays}, its predictions for old material remain intact during the consolidation of new material. Whatever the hippocampus is replaying, the cortex can absorb without overwriting prior knowledge.
\item \textbf{Replays are compressed.} The hippocampal index is much lower-dimensional than the neocortical state being adjusted. Compression is a feature, not a limitation: it forces the cortex to extract regularities rather than memorising specifics.
\end{enumerate}

These two features are exactly what gauge-LORN demands of R/G's action: \emph{must not disrupt L's predictions} (recoverable orbit) and \emph{must be low-rank} (compressed action).

\subsection{The Direct Mapping}

The mapping between gauge-LORN and CLS is unusually clean.

\begin{center}
\small
\begin{tabular}{p{4.8cm} p{4.5cm} p{5.0cm}}
\toprule
Gauge-LORN component & Hippocampus--Neocortex analog & Mathematical statement \\
\midrule
L (slow, frozen at inference) & Neocortex & Holds population statistics; updated rarely \\
Shared coupling $M$ \emph{and} shared FFN & Neocortex (distributed substrate) & Two globally shared parameter blocks, the ``world model'' (FFN, $\sim$97\% of manifold) and the ``compass'' ($M$, $\sim$36\%), both updated rarely, both perturbable on the recoverable orbit \\
R / G (fast learner) & Hippocampus & Encodes recent context; updated per-step \\
$\alpha(x)$, the gates & Hippocampal index & Compressed pointer; binds features \\
$\delta_{\mathrm{perturb}}(x) = (\delta M(x), \delta\mathrm{FFN}(x))$ (low-rank joint action) & Sharp-wave ripple replay & Compressed re-encoding of context, acting on the full shared substrate \\
Recoverable orbit $\ker H_{\mathrm{eff}}$ on $\theta_{\mathrm{shared}}$ & Consolidation manifold & Directions cortex can absorb \\
Mode A (correction MLP) & Within-pass cortical filtering & Local error correction \\
Mode B (recovery training) & Sleep-driven consolidation & Slow synaptic adjustment \\
LDA basis directions in $\theta_{\mathrm{shared}}$ & Pre-existing cortical category axes & Established semantic structure \\
Co-training $\theta_R$ updates & Cortical synaptic weights changing over consolidation & Slow gradient updates \\
\bottomrule
\end{tabular}
\end{center}

The matches are not loose analogies, they reflect the same computational structure.
The brain has evolved the same decomposition we have been mathematically reverse-engineering.

\paragraph{Even the failure modes match.}
\begin{itemize}[nosep]
\item Hippocampal lesions (Korsakoff's syndrome, certain epilepsies): anterograde amnesia. New memories cannot be encoded. \emph{In our framework}: G fails, no new structural information enters the system.
\item Massive cortical disruption: retrograde amnesia. Old memories destroyed. \emph{In our framework}: L's response heads damaged (over-fitting in v8), val degrades.
\item Sleep deprivation: consolidation fails, recent memories remain hippocampus-dependent and prone to interference. \emph{In our framework}: insufficient Mode B training, perturbations accumulate without being absorbed.
\item Hippocampal-cortical desynchronisation: indices fail to drive replays. \emph{In our framework}: $\alpha(x)$ collapses to a constant (v8/v9 with naive losses), no input-conditional information flows.
\end{itemize}

We did not design these failure modes; they emerged from the mathematics.
The fact that they correspond to known clinical syndromes is evidence the mathematical decomposition is real, not just a pleasing analogy.

\subsection{What the Brain Solved That We Have Not}

The biological system solves one thing we have not yet solved: \emph{how to derive the labels that drive consolidation without supervision}.

The hippocampus does not have category labels for what it is encoding.
What it has is \textbf{temporal binding}: things that happened close in time get co-encoded as a single index.
``This morning's meeting'' becomes a hippocampal index that links the room, the people, the agenda, the conversation, all because they co-occurred.
During replay, the cortex learns to predict elements from each other (next-event prediction).
Over many replays of many episodes, semantic categories \emph{emerge} from the prediction task: ``meetings usually have agendas'', ``rooms usually have chairs''.

This is the missing piece in our v8/v9/v10 attempts.
We have been trying to hand the network category labels (formal, child, statement, question) and force $G$ to learn them.
The brain does not do this.
The brain uses temporal coincidence as the implicit binding signal and lets categories emerge from many cycles of replay-plus-prediction.

\paragraph{The translation to a self-supervised loss.}
A natural episodic structure exists in any LM training corpus: a Wikipedia article is one episode, a Reddit thread another, a news article another.
Sentences within an episode share style, register, topic; sentences across episodes vary.
This gives us a temporal-binding analog without any labelling effort:

\begin{itemize}[nosep]
\item \textbf{Within-episode positive pairs}: chunks of text from the same source document. By construction, they share style and register because the same author wrote them in the same context.
\item \textbf{Across-episode negative pairs}: chunks from different source documents. By construction, they vary in style and register.
\item \textbf{Episodic contrastive loss on $\alpha(x)$}: same-episode $\alpha$ should be close; different-episode $\alpha$ should be far. This is contrastive SSL with positive pairs derived from natural co-occurrence rather than designed augmentations.
\end{itemize}

This is the v11 design (Section~\ref{sec:v11-design}, below).
It is biologically motivated: the brain uses temporal binding to derive its consolidation signal; we use document-binding to derive ours.

\subsection{The Wake-Sleep Connection}

A second biological framework relevant to our setup is the \textbf{wake-sleep algorithm}~\cite{hinton1995wakesleep}, which Hinton and Dayan introduced for training Helmholtz machines.
The wake phase performs bottom-up inference (encoder maps data to latent representation); the sleep phase performs top-down generation (decoder synthesises data from latent samples), with each phase training the other's parameters.

Our two-timescale dynamics map to this:
\begin{itemize}[nosep]
\item \textbf{Wake phase} (most steps): G receives inputs and infers $\alpha(x)$. L is frozen. New episodic structure is encoded into G's parameters.
\item \textbf{Sleep phase} (every $K$ steps): perturbations are applied at scale, L's response heads are briefly tuned to absorb them. This is slow consolidation of recent G-driven structure into L's stable weights.
\end{itemize}

The wake-sleep schedule is a curriculum we have not yet implemented.
The biological version (sleep schedules across days) suggests this should be intermittent rather than continuous: hours of wake encoding, then a sleep block where consolidation runs.
For our setup this might mean: 1000 G-only steps, then 100 L-co-training steps, repeated.

\subsection{What the Biology Predicts We Should Find}

If the CLS analog is real (not just a metaphor), several specific predictions follow:

\begin{enumerate}[nosep]
\item \textbf{Episodic-contrastive loss should outperform category-supervised loss} for training G. The brain works without category labels; with sufficiently good temporal binding our framework should too.
\item \textbf{Two-timescale training should outperform constant-rate training.} A schedule of fast-G / occasional-slow-L mimics wake/sleep.
\item \textbf{Replay buffer size matters.} Hippocampal replay focuses on recent experiences with priority for novel or surprising ones. Implementing a replay buffer that biases toward high-variance episodes should improve consolidation.
\item \textbf{Compressed action is necessary, not optional.} If we let $\delta M$ be full-rank, the cortex cannot consolidate, learning becomes pure memorisation. The low-rank constraint is biologically motivated, not just a parameter-efficiency hack.
\item \textbf{Consolidation gradients should be slow.} L's learning rate during sleep phases should be 1--2 orders of magnitude smaller than G's during wake phases. This matches the biological synaptic time scales and prevents the v8-style overfitting we observed.
\end{enumerate}

These predictions are testable and form a research programme orthogonal to (but compatible with) the gauge-orbit framing.
A successful CLS-inspired architecture would be both biologically grounded \emph{and} mathematically principled, the two perspectives reinforcing each other.

\subsection{The Synthesis}

McGilchrist gave us the architectural vision: two systems, complementary attention, asymmetric communication, the Master and the Emissary.
CLS gives us the computational mechanism: hippocampal indexing, compressed replay, slow cortical consolidation, the consolidation manifold.
Gauge-LORN is the mathematical formalisation: recoverable orbit, Schur complement, low-rank action, learnable gates over a structural basis.

These three perspectives describe the same thing.
The gauge-LORN equations are the mathematical content of what McGilchrist describes phenomenologically and what CLS describes biologically.
When v6/v7 succeeded empirically, what we had built was a compact computational version of the brain's solution to the catastrophic-interference / generalisation tension.
When v8 failed, we were learning that the standard SSL loss family is not what the brain uses.
When the recovery probe showed that structural directions are recoverable, we were observing the consolidation manifold directly.

The v11 design that follows takes the biological prescription seriously: episodic-contrastive loss derived from natural document boundaries, two-timescale training schedule, low-rank action.
It is not just gauge-LORN with new losses; it is gauge-LORN reframed as a computational implementation of CLS, with all three perspectives, hemispheric, mnemonic, mathematical, guiding the design.


\section{The Population Perspective: An Empirical Handle on the Orbit}
\label{sec:population-perspective}

This section documents a methodologically distinct approach to characterising the recoverable orbit.
We have not yet executed it experimentally; we record the framing and the proposed protocol so that it can be pursued as a parallel research track.

\subsection{Why a Different Approach}

The analytical approach to the orbit (Sections~\ref{sec:invariant}--\ref{sec:recoverable}) computes it via the eigenstructure of the Fisher matrix or Schur complement of the Hessian.
This has two confounds we have hit:
\begin{itemize}[nosep]
\item \textbf{Sample-size confound.} Empirical Fisher with $N \ll d^2$ samples (typical for any practical setting) misidentifies most of M-space as ``in orbit'' because the empirical kernel is dominated by directions never seen by the sample (Marchenko--Pastur).
\item \textbf{Linearisation confound.} The Schur-complement formulation is a local second-order expansion. For perturbations of meaningful structural magnitude, which we showed are not small, the linearisation may not capture the actual recoverable region.
\end{itemize}

A complementary approach replaces analytical estimation of the orbit with \emph{empirical observation of the population of equivalent solutions}.

\subsection{The Population View, Formally}

Train $K$ replicas of the same architecture on the same task with different initialisations or different data shuffles: $\theta_1, \ldots, \theta_K$.
Each achieves comparable validation loss; each represents a distinct point on the loss-equivalence manifold.
The pairwise differences $\theta_i - \theta_j$ span the orbit at a global, finite-magnitude scale, no linearisation, no eigenvalue thresholds.

For internal representations, the population view is even more direct.
For a held-out input $x$, compute the mid-layer hidden state $h_l(x; \theta_i)$ for each replica.
The variance of these representations \emph{across replicas} is the orbit's reach in representation space.
Directions where representations differ across replicas but predictions agree are exactly the gauge-equivalent reorganisations the orbit allows.

This is not a new framework, it is the operational content of the mode connectivity literature~\cite{garipov2018loss, draxler2018essentially, frankle2020linear, ainsworth2023git}.
Those works focus on connecting trained networks via low-loss paths.
We propose to use the same observation differently: as a direct measurement of the orbit's structural content.

\subsection{The Three Scenarios}

The population variance of internal representations could turn out to span any of three qualitatively different subspaces:

\begin{enumerate}[nosep]
\item \textbf{Permutation-only orbit.} Replicas differ by hidden-unit reordering. The orbit is large in raw parameter terms but functionally trivial, all replicas implement the same function up to a discrete relabelling. This is consistent with what Ainsworth et al.~\cite{ainsworth2023git} find for vision networks: most apparently-different trained networks are permutation-equivalent. \emph{Gauge-LORN cannot exploit this; the orbit is structurally vacant.}

\item \textbf{Noise orbit.} Replicas differ in random uninterpretable directions. The orbit is large but does not align with any meaningful structural axis. \emph{Gauge-LORN cannot exploit this either.}

\item \textbf{Structurally meaningful orbit.} Replicas differ in ways that correspond to alternative structural decompositions of the same inputs, e.g., different replicas attend to different subordinate clauses, different replicas factor sentences into different head-modifier graphs, different replicas develop different but functionally-equivalent register encodings. \emph{This is the regime where structure-aware design works.}
\end{enumerate}

The framework's central empirical bet is that scenario 3 holds for trained autoregressive language models.
None of our experiments to date have tested this directly.
v6/v7's success suggests the orbit at least \emph{permits} structural action; whether the population variance is intrinsically structural is a stronger and untested claim.

\subsection{Implications for Architecture Design}

If the population variance turns out to be structurally meaningful, several design rules follow:

\begin{itemize}[nosep]
\item \textbf{Over-parameterisation enriches the orbit.} A larger parameter space admits more equivalent configurations.
This is consistent with the empirical finding that larger models generalise better partly because they have more orbit room to settle into~\cite{soudry2018implicit}.
For LORN, scale should help, not just for capacity, but for the size of the structural manifold R has to operate on.

\item \textbf{Architectural symmetries should be richer than permutation.}
Networks whose only invariance is hidden-unit shuffling have orbits that are large but vacant.
Architectures admitting \emph{functional} equivalences, multiple decompositions of the same operator (M = AB$^\top$ has many factorisations of A, B that produce the same M), have richer orbits.

\item \textbf{Wide shared FFN $>$ per-layer FFN at matched parameters.}
Beyond the parameter-efficiency argument, shared weights have more equivalent factorisations.
The wide shared FFN's success in Deep-R / v6 may partly reflect this.

\item \textbf{Co-trained populations could be a training mechanism, not just an analysis tool.}
Train $K$ replicas in parallel, share their structural-state ensemble, distill into a final consolidated network.
This is essentially CLS distillation across instances rather than across time.
\end{itemize}

\subsection{The Proposed Experiment}

A clean, modest experiment would test scenario 3 directly:

\begin{enumerate}[nosep]
\item Train $K = 5$--$10$ small ORN replicas (e.g., $d{=}128$, 6 layers, 30K steps) on $K$ different shuffles of the same training corpus. Vary only initialisation seed and data ordering; keep architecture and hyperparameters identical.
\item For a held-out probe set of structurally-categorised sentences, compute mid-layer hidden states for each replica.
\item Compute the cross-replica covariance of representations for each input.
\item Project this covariance onto the structural-category subspace (LDA on category labels).
\item Report: what fraction of cross-replica variance lies along structural axes vs.\ random directions vs.\ permutation directions (after Git Re-basin alignment)?
\end{enumerate}

\paragraph{Predicted outcomes.}
\begin{itemize}[nosep]
\item If $\geq 50\%$ of cross-replica variance lies along structural axes: scenario 3 confirmed; the orbit naturally contains structure; gauge-LORN's central bet is empirically validated.
\item If permutation alignment removes most cross-replica variance: scenario 1; we need different architectures to enable functional invariance.
\item If variance is uncorrelated with structural categories: scenario 2; replicas differ randomly, no structural reorganisation is naturally available.
\end{itemize}

\paragraph{Compute estimate.}
$K = 8$ replicas of 5M-parameter ORN, 30K steps each on existing wiki+owt data: $\approx 8 \times 1.5\text{h} = 12$h on MPS, sequentially.
With careful checkpointing, parallelisable.
A single 24h overnight run produces the answer.

\subsection{Why This Is Methodologically Distinct}

The recoverable-orbit experiments (v8--v12) treat the orbit as a property to be discovered \emph{at one point} (the trained $\theta_L$) and use a learned operator G to navigate it.
The population perspective treats the orbit as an empirical object to be observed across many trained instances.
The two approaches answer related but distinct questions:

\begin{itemize}[nosep]
\item Recoverable-orbit at one point: ``can a learned operator find structural directions in $\theta_L$'s neighbourhood?'' (architecture and learnability question)
\item Population variance: ``does the loss landscape have structural slack at all, intrinsic to the architecture and task?'' (intrinsic-property question)
\end{itemize}

A network that fails the population test (scenario 1 or 2) cannot be saved by a clever G; the orbit is the wrong shape.
A network that passes the population test still requires a learnable G to exploit the structure; passing is necessary but not sufficient.

The two perspectives are complementary.
Documenting them both gives the framework two independent measurement strategies, either of which can falsify the central claims.

\subsection{Connection to v11 / Future Work}

v11 is using episodic contrastive, implicitly assuming that document-source binding induces an invariance G can exploit.
The population perspective makes this assumption testable: do replicas trained on document-source-shuffled data exhibit cross-replica variance along document-source axes?
If yes, v11's signal is rooted in a real population invariance.
If no, v11's success (if any) comes from somewhere else.

We mark this as a parallel research track because:
\begin{enumerate}[nosep]
\item Methodologically, it requires training multiple replicas (different infrastructure than the single-checkpoint experiments we have been running).
\item Conceptually, it tests architectural assumptions that v8--v12 take for granted.
\item Strategically, it is independent of v11/v12's outcome, useful regardless of whether the dynamic operator works.
\end{enumerate}

If the gauge-LORN framework continues to deliver on v11/v12, the population perspective is a natural follow-up validation.
If gauge-LORN stalls, the population perspective tells us whether the issue is the architecture (scenario 1/2) or our specific operator G (scenario 3 holds but G can't navigate).

The experiment is small, the answer is clean, and either result advances the framework.


\section{The Permutation Gauge: Input-Space Invariance and Sleep Replay}
\label{sec:permutation-gauge}

This section documents a second methodologically distinct track that emerged from thinking about how biological replay scrambles temporal order during consolidation.
We have not executed it experimentally; we record the framing, diagnostic protocol, and proposed training application so the track can be developed in parallel to gauge-LORN's main thrust.

\subsection{Two Kinds of Gauge}

The gauge-LORN framework operates on \emph{parameter-space gauge}: directions in $\theta$ where the function is approximately invariant.
A complementary notion is \emph{input-space gauge}: permutations of input tokens where some chosen aspect of the model's \emph{representation} (not its prediction) is approximately invariant.

These are mathematically distinct objects:

\begin{center}
\begin{tabular}{p{4.5cm} p{4.5cm} p{4.5cm}}
\toprule
Aspect & Parameter-space gauge & Input-space (permutation) gauge \\
\midrule
Acts on & $\theta$ (network weights) & $x$ (input token sequence) \\
Invariance condition & predictions stay same & chosen aspect of representation stays same \\
Discovered via & Hessian/Schur or population variance & permutation diagnostics \\
Used by R/G to do what & reorganise weights & extract order-invariant structure \\
\bottomrule
\end{tabular}
\end{center}

The new claim: some directions in representation space are approximately invariant under permutation of input tokens.
Those directions encode \textbf{order-independent content} (topic, vocabulary, register).
Directions that vary under permutation encode \textbf{order-dependent structure} (syntactic role, coreference, scope).

This decomposition, order-invariant $\oplus$ order-covariant, is essentially what linguists call the content/form split, and the brain's sleep-replay machinery appears designed to extract it.

\subsection{The Biological Mechanism}

Sleep replay~\cite{wilson1994reactivation, buzsaki2015swr, diekelmann2010sleep} is not veridical reproduction of waking experience.
Hippocampal replay during sleep includes:
\begin{itemize}[nosep]
\item \textbf{Forward sequences} during quiet wake and pre-sleep states
\item \textbf{Reverse sequences} immediately after experience
\item \textbf{Compressed sequences} (10--20$\times$ faster than original timescales)
\item \textbf{Shuffled fragments} that recombine elements out of order
\item \textbf{Preplay} of plausible future sequences before they happen
\end{itemize}

The functional interpretation: by replaying scrambled versions of the same content, the cortex learns which features are stable under reordering (semantic content) and which are tied to specific orderings (sequence structure).
Two-stream extraction: the order-independent content gets consolidated into semantic memory; the order-dependent structure stays available as episodic detail.

If gauge-LORN aims to mirror CLS computationally, it should be doing the same: replay shuffled versions of training sequences and use what the network considers invariant as a structural diagnostic.

\subsection{The Permutation Diagnostic on a Trained Network}

For a frozen trained model and a held-out sentence $x = (t_1, t_2, \ldots, t_S)$, generate $K$ permuted versions via different permutation classes:

\begin{itemize}[nosep]
\item $\pi_{\text{full}}$: random shuffle of all tokens
\item $\pi_{\text{reverse}}$: reverse order $(t_S, \ldots, t_1)$
\item $\pi_{\text{block}}$: random shuffle of 4-token blocks (preserves local order, scrambles global)
\item $\pi_{\text{function}}$: shuffle only function words (the, a, of, in, etc.)
\item $\pi_{\text{content}}$: shuffle only content words (nouns, verbs, adjectives)
\item $\pi_{\text{adjacent}}$: random adjacent swaps (mostly preserves structure)
\end{itemize}

For each permutation type and each layer $l$, measure:
\begin{equation}
\Delta_l(\pi) \;=\; \big\| h_l(x) - h_l(\pi(x)) \big\|_2 \quad \text{or} \quad 1 - \cos\!\big(h_l(x), h_l(\pi(x))\big)
\end{equation}

The pattern of $\Delta_l(\pi)$ across $l$ and across permutation classes characterises the network's implicit invariance structure.

\paragraph{Predicted patterns.}
\begin{itemize}[nosep]
\item \textbf{Early layers}: high sensitivity to all permutations (positional encoding dominates).
\item \textbf{Late layers}: lower sensitivity to function-word shuffles, higher to content-word shuffles.
\item \textbf{Block-shuffle invariance gradient}: gradual increase with depth, revealing the network's compositional resolution.
\item \textbf{Adjacent-swap robustness}: indicates whether the network has learned position-insensitive contextual integration.
\end{itemize}

A network with no content/form decomposition (e.g., a bag-of-words baseline) would show $\Delta_l(\pi) = 0$ for all $\pi$ and all $l$.
A network purely positional (e.g., a memorising LM) would show $\Delta_l(\pi)$ proportional to the permutation's edit distance from identity.
A network with genuine compositional structure would show layer-dependent and permutation-class-dependent patterns: invariant to function-word shuffles by depth $L/2$, still sensitive to content-word shuffles, and partially invariant to block shuffles depending on the block size relative to the network's effective context window.

\subsection{Permutation-Consistency Training}

If the diagnostic reveals real but incomplete content/form decomposition, an auxiliary training objective can sharpen it.

For each training step, additionally compute a \textbf{permutation-consistency loss}:
\begin{equation}
\mathcal{L}_{\pi} \;=\; \mathbb{E}_{x,\, \pi \sim \mathcal{P}}\!\left[\, \big\| f_{\text{content}}(h_l(x)) - f_{\text{content}}(h_l(\pi(x))) \big\|_2^2\, \right]
\end{equation}

where $f_{\text{content}}$ is a small projection head extracting the ``content channel'' that should be permutation-invariant, $\mathcal{P}$ is the chosen permutation distribution (probably weighted toward function-word shuffles, since those preserve content most reliably), and the layer $l$ is one of the deeper layers (so the invariance is required at the semantic level, not the surface level).

This is biological replay implemented as a training augmentation: the brain's sleep-shuffling becomes an explicit loss term.
The model is forced to disentangle content from structure by being trained on the invariance.

\paragraph{Schedule analog.}
The brain does not interleave replay with waking encoding; replay is concentrated in distinct sleep phases.
A computational analog: train normally for $K$ steps (waking), then run 1 step of permutation-consistency training (sleep replay).
This two-timescale schedule mirrors the wake-sleep dynamics and avoids the permutation loss interfering with the primary NTP objective.

\subsection{Test-Time Adaptation via Permutation}

The permutation gauge also enables a clean form of test-time training (TTT).
At inference, given a single input $x$:
\begin{enumerate}[nosep]
\item Generate $K$ permutations $\pi_1(x), \ldots, \pi_K(x)$.
\item Compute representations $h_l(\pi_i(x))$ for all $K$.
\item The variance across shuffles measures how much of $x$'s representation is order-dependent.
\item Optionally: briefly fine-tune a small adapter (e.g.\ LoRA) on the permutation-consistency loss for $x$, then make the prediction.
\end{enumerate}

This gives inference-time access to the network's structural confidence:
\begin{itemize}[nosep]
\item Low cross-shuffle variance: ``I understood the content; the order was decoration.''
\item High cross-shuffle variance: ``The structure was load-bearing; small order changes change my interpretation.''
\end{itemize}

The adapter, after brief tuning on $L_\pi$ for the specific test input, has learned what is structurally essential about \emph{this input}.
This is genuinely different from standard TTT (which uses generic self-supervised tasks): the adaptation is structurally targeted by the test instance itself.

\subsection{Composition with Gauge-LORN}

The two gauges compose constructively.

\begin{enumerate}[nosep]
\item \textbf{Permutation diagnostic on v6} identifies the ``content channel'' in representation space, the directions that are order-invariant.
\item \textbf{Gauge-LORN's G} can then condition its action on whether the input is content-rich or structure-rich, using cross-shuffle variance as a self-computed signal (no labels required).
\item \textbf{CLS replay in v11+} is enhanced: the contrastive loss replays not only same-document chunks but also shuffled-token versions, forcing the order/content decomposition explicitly.
\end{enumerate}

The unified framework: gauge-LORN navigates a parameter-space orbit; permutation-replay extracts the input-space invariance; both serve the same goal of disentangling structure from content, mirroring the brain's hippocampus/cortex/replay machinery.

\subsection{Why This Is Methodologically Distinct}

The gauge-LORN experiments (v8--v12) characterise an invariance in \emph{parameter} space.
The permutation gauge characterises an invariance in \emph{input} space.
Both identify ``directions where structural encoding can occur without disrupting prediction'', but they do so in orthogonal spaces, with orthogonal experimental requirements.

The permutation gauge requires:
\begin{itemize}[nosep]
\item No new architecture (the diagnostic runs on any frozen model).
\item No new training (the diagnostic is purely measurement).
\item Optional training augmentation (permutation-consistency loss) is a small addition.
\end{itemize}

In contrast, the parameter gauge requires:
\begin{itemize}[nosep]
\item A learned operator G (architectural commitment).
\item A target subspace (LDA basis or learned).
\item Co-training schedule (Mode A / Mode B).
\end{itemize}

The permutation gauge is therefore cheaper to test and provides an independent measurement of structural decomposition.
A network passing the permutation diagnostic (clear layer-by-permutation structure) is a network where the parameter-gauge framework has signal to work with.
A network failing the permutation diagnostic (uniform sensitivity to all permutations) is a network where parameter-gauge approaches will struggle: the model has not internalised the content/form distinction at all.

\subsection{Connection to Sleep, TTT, and Hippocampal Learning}

Three existing literatures meet in the permutation gauge:

\begin{enumerate}[nosep]
\item \textbf{Sleep biology} (Diekelmann \& Born~\cite{diekelmann2010sleep}, Wilson \& McNaughton~\cite{wilson1994reactivation}): scrambled-order replay is the empirical phenomenon. The function is content/form distillation.
\item \textbf{Test-time training} (Sun et al.\ 2020 and subsequent work): adapting the model on test input via self-supervised loss. Standard TTT uses generic tasks (masked LM, rotation); permutation-consistency is a structurally-targeted alternative.
\item \textbf{Hippocampal indexing}: the index binds together features that co-occurred. Replay re-presents these bindings to cortex in arbitrary orders. The cortex learns which features are essential to the binding (semantic) and which are positional artefacts.
\end{enumerate}

The three frameworks make convergent predictions: representation should decompose into permutation-invariant and permutation-covariant components; this decomposition is learnable; the brain has evolved a specific schedule (sleep) to enforce it; computational analogues should work similarly.

\paragraph{Where this leaves us.}
We mark this as a parallel research track because, like the population perspective (Section~\ref{sec:population-perspective}), it answers a different question than the gauge-LORN main thread.
Gauge-LORN asks: ``can a learned operator navigate the parameter orbit to encode structure?''
The permutation gauge asks: ``does the network's representation already exhibit a content/form decomposition we can measure and amplify?''
The two are complementary.
A working full system would have both: a content/form-decomposed representation (permutation-invariant content channel, structure-covariant signal) and a learned operator that navigates parameter-space gauge to encode structural variation atop the content channel.

The permutation diagnostic is small: a single forward pass per shuffle, a few hours of MPS to characterise v6 across all permutation classes and all layers.
A natural next-step experiment, parallelisable with v11/v12, that produces an independent line of evidence for or against the framework's central decomposition claim.


\subsection{Notes: Toward Emergent Conceptual Metaphor}
\label{sec:emergent-metaphor}

\emph{This section records a speculative hypothesis. The mechanism is concrete; the claim is testable; we have not run the test. It is documented to direct future work, not to make a present claim.}

\subsection{The Hypothesis}

If R's episodic-contrastive training (v11) clusters representations by narrative co-occurrence rather than literal vocabulary, then under the right conditions \emph{cross-domain clusters can emerge naturally}: ``war'' inputs and ``argument'' inputs, both episodes of conflict-with-rules-and-stakes, end up with similar $\alpha(x)$ even though they share little surface vocabulary.

We hypothesise that such cross-domain clusters constitute conventional metaphors as understood in cognitive linguistics.
The mechanism is unsupervised emergence from temporal/narrative co-occurrence statistics.
No metaphor labels are required; no domain-mapping module is designed; the structure falls out of episodic-contrastive learning applied to a corpus where natural narratives systematically cross literal domain boundaries.

\subsection{The Cognitive-Linguistics Background}

Lakoff \& Johnson (\emph{Metaphors We Live By}, 1980) and Hofstadter \& Sander (\emph{Surfaces and Essences}, 2013) argue that metaphor is not a literary embellishment but the basic mode of human conceptual reasoning.
Conventional metaphors such as ``ARGUMENT IS WAR'', ``TIME IS MONEY'', ``LIFE IS A JOURNEY'' systematically map structure from one conceptual domain onto another:
\begin{itemize}[nosep]
\item \textbf{ARGUMENT IS WAR}: attack, defend, retreat, surrender, win, position
\item \textbf{TIME IS MONEY}: spend, save, waste, invest, budget
\item \textbf{LIFE IS A JOURNEY}: path, obstacles, destination, crossroads, turning point
\end{itemize}
Humans reason about the target domain by transporting the source domain's inferences.
The mappings are not random; they are systematic, productive, and pervasive in everyday language.

The standard cognitive-science account leaves the \emph{computational origin} of these mappings unspecified.
Where do they come from?
The standard answer (``experience'') does not tell us the mechanism by which experience produces them.

\subsection{The Mechanism}

Our episodic-contrastive setup provides a candidate mechanism.
Consider that natural text has the following structure:
\begin{itemize}[nosep]
\item A news article about a political dispute uses both literal-political vocabulary AND war metaphor (``attacked his position'', ``defended the bill'').
\item A novel about a war uses both literal-combat vocabulary AND interpersonal-conflict metaphor (``the army was demoralised'', ``the soldier surrendered to despair'').
\item A debate transcript uses both literal-argument vocabulary AND domain-mixed metaphor (war, sport, journey).
\end{itemize}
Each episode is internally heterogeneous in its literal vocabulary but internally homogeneous in its narrative/structural pattern.

If R learns to cluster by narrative pattern rather than literal vocabulary, which is what episodic-contrastive loss should produce, since shared words are not the only basis for episode coherence, then ``argument'' inputs and ``war'' inputs sharing narrative structure should end up with similar $\alpha(x)$.
The cluster has crossed a literal-domain boundary.
\textbf{That cluster is the metaphor.}

The mechanism does not require explicit metaphor labels.
It requires only that the training data exhibits the cross-domain narrative parallels that humans implicitly encode when they write natural text.

\subsection{The Test}

After v11 (or any successor with episodic-contrastive training) has converged, run the following protocol:

\paragraph{Step 1: Curate cross-domain probe sentences.}
For three or more conceptual domains likely to participate in conventional metaphors:
\begin{itemize}[nosep]
\item WAR domain: ``the army attacked at dawn'', ``the soldiers defended their position'', ``the battle ended in retreat''.
\item ARGUMENT domain: ``she attacked his position immediately'', ``he defended his thesis vigorously'', ``the debate ended with concession''.
\item JOURNEY domain: ``we set out toward the goal'', ``she encountered obstacles along the way'', ``they reached their destination''.
\item TIME-AS-MONEY domain: ``we spent three hours on it'', ``she saved time by skipping lunch'', ``it cost me my afternoon''.
\end{itemize}

\paragraph{Step 2: Compute $\alpha(x)$ for each.}
Use v11's trained G in inference mode.

\paragraph{Step 3: Cross-domain alignment measurements.}
Three quantities matter:
\begin{itemize}[nosep]
\item Within-domain similarity: $\bar{\cos}(\alpha_{D_i}, \alpha_{D_i})$ for sentences in domain $D_i$.
\item Random cross-domain similarity: $\bar{\cos}(\alpha_{D_i}, \alpha_{D_j})$ for arbitrary sentence pairs across domains.
\item \textbf{Same-relation cross-domain similarity}: $\cos(\alpha(\text{``attacked army''}), \alpha(\text{``attacked thesis''}))$, pairs that share an abstract relation across domains.
\end{itemize}

\paragraph{The decisive measurement.}
If same-relation cross-domain pairs have $\alpha$ similarity comparable to within-domain same-relation pairs, R has learned the metaphor.
If same-relation cross-domain pairs are no more similar than random cross-domain pairs, R is encoding domain-locked structure (no metaphor).

The protocol generalises naturally: graph-completion tests across domains (``soldier : army :: lawyer : ?''), interpolation tests in $\alpha$-space, generation tests where injecting a war-domain $\alpha$ into argument-content drives metaphor production.

\subsection{What Would Be Established}

If the test produces positive results, same-relation cross-domain $\alpha$ similarity comparable to within-domain, we would have evidence that:

\begin{enumerate}[nosep]
\item \textbf{Conventional metaphors emerge from episodic-contrastive learning} on natural text. No design choice or supervision encodes them; they fall out of the same mechanism that produces register and function discrimination.
\item \textbf{R's $\alpha$-space is a metaphor system}: organised regions correspond to conventional metaphors documented in cognitive linguistics.
\item \textbf{Gauge-LORN is not just a steerable language model.} It is a computational model of conceptual structure formation.
\item \textbf{The McGilchrist mapping is sharpened.} R's role, which we have framed loosely as ``relational/contextual processing'', becomes specifically ``cross-domain structural mapping'', i.e., metaphor.
\item \textbf{The CLS analog is sharpened.} Cortical consolidation extracts cross-episode regularities; we are doing the computational analog.
\end{enumerate}

\subsection{What Is Wild and What Is Not}

\paragraph{Not wild.}
That episodic-contrastive learning produces semantic clusterings (well-established by BYOL, MoCo, etc.\ in the visual domain).
That some semantic clusterings span literal domains (distributional semantics routinely does this: $\text{king} - \text{man} + \text{woman} \approx \text{queen}$).

\paragraph{Mildly wild.}
That ALL conventional metaphors emerge from temporal co-occurrence.
Some may be more deliberate, more cultural, more linguistically specific.

\paragraph{Genuinely wild and worth testing.}
That v11-scale episodic-contrastive training, applied to natural-text corpora, reproduces the systematic conventional metaphors documented in cognitive linguistics in a form measurable as $\alpha$-space alignment.

This is the speculation worth keeping in notes: not because we have evidence for it, but because the experimental test is concrete and modest, the theoretical reward (a computational model of metaphor formation) is large, and the negative result is also informative (it would tell us that R's clustering is more domain-locked than metaphor would require, suggesting a limit to the framework's reach).

\subsection{Status}

Speculative. Documented for follow-up. Test would be small (a few hundred curated probe sentences, one forward pass each, a few minutes of MPS).
Pre-registration: same-relation cross-domain $\alpha$-cosine $\geq 0.7 \times$ same-relation within-domain $\alpha$-cosine would count as positive evidence.

If v11 succeeds, this test is the next thing to run.
If v11 fails, the metaphor test still tells us something about whether episodic clustering produces cross-domain alignment in principle, even if our specific dynamic operator does not exploit it.

The hypothesis closes a loop that has been open in cognitive linguistics for forty years: conventional metaphors are real, systematic, and important, but no one has specified the computational mechanism that produces them.
The mechanism we propose, episodic-contrastive learning on naturally cross-domain narrative text, is simple, biologically plausible, and testable.

If correct, this would be the largest theoretical implication of the gauge-LORN framework: not just a new architecture for language modelling, but a candidate mechanism for the formation of human conceptual structure itself.


\section{Structure Tokens: An Imaginary Vocabulary for Relational Meaning}
\label{app:structure-tokens}

The prefix token mechanism suggests a deeper idea: \textbf{structure tokens}, an implicit vocabulary that encodes relational meaning the way content tokens encode lexical meaning.

\subsection{The Forward and Backward Processes}

A standard language model processes content tokens left to right:
\begin{equation}
\text{content tokens} \xrightarrow{\text{L (forward)}} \text{next content token}
\end{equation}

Structure tokens introduce a second dimension.
Imagine every sentence is secretly prefixed by invisible tokens that encode its relational skeleton:
\begin{equation}
[\underbrace{s_1, s_2, \ldots, s_K}_{\text{structure tokens}},\; \underbrace{w_1, w_2, \ldots, w_N}_{\text{content tokens}}]
\end{equation}

The \textbf{forward process} (L) sees structure tokens and generates content tokens, filling the relational skeleton with words.
``Given that this sentence has agent-action-patient structure, generate: The cat chased the mouse.''

The \textbf{backward process} (R) sees content tokens and generates structure tokens, extracting the relational skeleton from words.
``Given `the cat chased the mouse,' produce: [SVO, transitive, past-tense, animate-agent].''

Together, L and R form a cycle:
\begin{center}
structure tokens $\xrightarrow{\text{L: fill in content}}$ content tokens $\xrightarrow{\text{R: extract structure}}$ structure tokens
\end{center}

\subsection{Why ``Metaphor Tokens''}

These structure tokens are metaphors in the McGilchrist sense.
A metaphor maps structure from one domain to another: ``argument is war'' takes the relational structure of combat (attack, defend, retreat, win) and applies it to discourse.
The content changes (words replace weapons); the structure transfers.

A structure token IS a metaphor: it encodes relational structure independently of content.
The structure token for ``agent acts on patient'' applies equally to ``the cat chased the mouse,'' ``the committee rejected the proposal,'' and ``the wind scattered the leaves.''
Different content, same structure token.

Moving in structure-token space IS moving between metaphors.
Shifting the structure tokens from [SVO-transitive] toward [existential-locative] changes the relational template: from ``X does Y to Z'' to ``there exists X at location Y.''
L generates content that fills whichever template the structure tokens specify.

\subsection{Generation as Structure-Content Interplay}

Seeded generation becomes navigation in structure-token space:
\begin{enumerate}[nosep]
\item \textbf{Start with structure tokens.} Compose a relational template in R's state space: ``judgment structure + social hierarchy + temporal sequence.'' Project to structure tokens.
\item \textbf{L generates content.} L sees the structure tokens and generates content tokens that fill the template: ``The committee's recommendation was subsequently vindicated by independent analysis.''
\item \textbf{R reads L's output.} R processes the generated content tokens and updates the structure tokens, the relational template evolves as content is produced.
\item \textbf{Iterate.} L generates the next content tokens conditioned on updated structure tokens. The structure and content co-evolve through the generation process.
\end{enumerate}

This is not chain-of-thought (which externalises reasoning as content tokens).
This is structure-guided generation where the reasoning happens in an implicit structural vocabulary that never appears in the output.
The structure tokens are invisible to the reader but visible to the model.

\subsection{Analogical Transfer}

The most powerful implication: \textbf{analogical transfer is movement in structure-token space.}

To generate an analogy (``A is to B as C is to ?''):
\begin{enumerate}[nosep]
\item Process ``A is to B'' through R $\to$ extract structure tokens $[s_1, \ldots, s_K]$.
\item These tokens encode the relational structure of the A-B relationship.
\item Prepend these structure tokens to ``C is to'' and let L generate.
\item L fills the same relational template with new content: ``D.''
\end{enumerate}

More ambitiously: to apply a metaphor (``argument is war''):
\begin{enumerate}[nosep]
\item Process war-domain text (``the soldiers attacked the fortification'') through R $\to$ extract combat structure tokens.
\item Prepend to argument-domain context (``the lawyer then...'') and generate.
\item L produces: ``the lawyer then attacked the opposing counsel's position'' ,  the war structure applied to the argument domain.
\end{enumerate}

The structure tokens ARE the medium of analogical transfer.
They encode ``how things relate'' independently of ``what things are.''

\subsection{Relation to Existing Work}

Soft prompts~\cite{lester2021power} prepend learned embeddings to steer generation.
But soft prompts are fixed after training, they encode a task, not a relational structure derived from input.
Structure tokens are \emph{dynamic}: generated by R from the current input, updated as generation proceeds.

Retrieval-augmented generation prepends retrieved documents as context.
Structure tokens are similar but operate in relational space, not document space.
The ``retrieved'' information is the relational skeleton of the current context, not a similar text.

\subsection{Speculation: Hierarchical Parallel Decomposition}
\label{app:hierarchical-parallel}

Structure tokens suggest a deeper architecture: \textbf{hierarchical structural decomposition with parallel content generation}.

\subsection{The Hierarchy}

Language structure is hierarchical: discourse contains paragraphs, paragraphs contain clauses, clauses contain phrases.
R can decompose this hierarchy as a \emph{tree of structural prefixes}, all available nearly simultaneously:

\begin{center}\small
\begin{tabular}{lll}
\toprule
\textbf{Level} & \textbf{R's structural prefix} & \textbf{L fills in} \\
\midrule
0 (discourse) & argument $\to$ evidence $\to$ counter $\to$ resolution & paragraphs \\
1 (clause) & claim: $\square$ $\to$ because: $\square$ $\to$ but: $\square$ & phrases \\
2 (phrase) & SVO template / existential template / subordinate template & words \\
\bottomrule
\end{tabular}
\end{center}

R produces ALL structural prefixes at all levels nearly instantly, structure decomposes analytically.
An essay has four sections before any section is written.
A clause has agent-action-patient structure before any word is chosen.

\subsection{Parallel Content Generation}

L fills in content at each level in parallel:
level 2 fills SVO templates with words;
level 1 fills clause templates with phrases;
level 0 fills discourse templates with paragraphs.
Each level's L operates independently, conditioned only on its R prefix.

This is parallelisable because the structural prefixes are available before any content is generated.
The draft content at each level is produced simultaneously, no sequential dependency between levels.

\subsection{Cross-Talk and Correction ($\lambda$)}

The levels are not truly independent.
Word choice at level 2 (``attacked'' vs ``questioned'') depends on the discourse frame at level 0 (``argument as war'' vs ``argument as inquiry'').
A parameter $\lambda$ controls the cross-talk between levels:

\begin{itemize}[nosep]
\item $\lambda = 0$: fully parallel, levels independent. Fast but globally incoherent.
\item $\lambda = 1$: fully sequential, each level waits for the others. Slow but coherent.
\item $0 < \lambda < 1$: parallel draft, then correction passes. The correction propagates constraints between levels.
\end{itemize}

After parallel L completion, the levels do not perfectly align.
A correction pass propagates adjustments: level 0's discourse frame constrains level 1's clause structure, which constrains level 2's word choice.
This IS recurrence, but shallow recurrence ($\mathcal{O}(\log \text{depth})$ correction passes) rather than deep sequential processing ($\mathcal{O}(\text{sequence length})$).

\subsection{Connections}

\paragraph{To speculative decoding.}
Current speculative decoding uses a small draft model to propose tokens that a large model verifies.
Hierarchical LORN uses R's structural decomposition as the draft.
The $\lambda$-correction passes are the verification step, but they verify \emph{structural coherence}, not token probability.

\paragraph{To diffusion.}
The correction passes ARE denoising.
Level 2's independent completions are ``noisy'' (locally coherent, globally incoherent).
Cross-level correction ``denoises'' by enforcing global coherence.
$\lambda$ controls the noise schedule: high $\lambda$ = aggressive denoising per pass, low $\lambda$ = gentle correction over more passes.

\paragraph{To SSMs.}
The hierarchical state is an SSM over the \emph{tree structure}, not the token sequence.
Each R level maintains state that evolves as L fills in content below it.
The state transition: ``given what L generated at level $k+1$, update level $k$'s structural expectations.''

\subsection{Status}

This is speculative, no implementation exists.
The required components (structure tokens, structural LM, prefix injection, co-evolution) are individually tested.
Their hierarchical composition into a parallel generation architecture is the long-term vision for LORN.


\begin{thebibliography}{99}
\bibitem{lester2021power}
B.~Lester et al., ``The power of scale for parameter-efficient prompt tuning,'' \emph{EMNLP}, 2021.
\end{thebibliography}


\end{document}

